% Allegato tecnico. Si compila con: xelatex, bibtex, xelatex, xelatex.
% Preambolo comune dell'allegato tecnico. Si compila con XeLaTeX + BibTeX.
\documentclass[10pt,a4paper]{article}
\usepackage[a4paper,margin=2.2cm]{geometry}
\usepackage{fontspec}
\setmainfont{Fira Sans}[Scale=0.95,BoldFont={Fira Sans SemiBold},ItalicFont={Fira Sans Italic}]
\setmonofont{DejaVu Sans Mono}[Scale=0.82]
\usepackage{polyglossia}\setmainlanguage{italian}
\usepackage{amsmath,amssymb,amsthm,mathtools}
\usepackage{xcolor,tikz,booktabs,tabularx,array,enumitem,microtype,ragged2e,caption,float}
\usetikzlibrary{positioning,arrows.meta,calc,angles,quotes,patterns,decorations.pathreplacing}
\usepackage[numbers,sort&compress]{natbib}
\usepackage[hidelinks]{hyperref}
\definecolor{accent}{HTML}{0B5C8A}
\definecolor{muted}{HTML}{55606C}
\definecolor{boxbg}{HTML}{EEF3F7}
\setlength{\parindent}{0pt}\setlength{\parskip}{4pt}
\setlist{leftmargin=1.4em,itemsep=1.5pt,topsep=2pt}
\captionsetup{font=small,labelfont=bf}
\newcolumntype{L}{>{\RaggedRight\arraybackslash}X}
% Ambienti numerati per capitolo (sezione)
\theoremstyle{plain}
\newtheorem{proposizione}{Proposizione}[section]
\newtheorem{lemma}[proposizione]{Lemma}
\newtheorem{corollario}[proposizione]{Corollario}
\theoremstyle{definition}
\newtheorem{definizione}[proposizione]{Definizione}
\newtheorem{ipotesi}[proposizione]{Ipotesi}
\newtheorem{controesempio}[proposizione]{Controesempio}
\theoremstyle{remark}
\newtheorem{osservazione}[proposizione]{Osservazione}
% Livello di prova di un risultato: \prova{D} dimostrato, {C} calcolo, {S} simulazione, {P} dati pubblici, {N} non verificato sperimentalmente
\newcommand{\prova}[1]{\textsf{\footnotesize\color{accent}[#1]}}
% Riquadro per condizioni di validità e limiti
\usepackage{framed}\colorlet{shadecolor}{boxbg}
\newenvironment{riquadro}[1]{\begin{shaded}\textbf{#1}\par\smallskip}{\end{shaded}}
% Notazione comune
\newcommand{\Reali}{\mathbb{R}}
\newcommand{\vmax}{v_{\max}}
\newcommand{\tocc}{\tau_{\mathrm{occ}}}
\newcommand{\Oss}{\mathcal{O}}      % insieme osservato libero
\newcommand{\Omb}{\Omega}           % ombra
\newcommand{\Pres}{\Pi}             % possibile presenza
\newcommand{\Reg}{\mathcal{R}}      % regione di previsione
\newcommand{\unita}[1]{\,\mathrm{#1}}
\DeclareMathOperator{\dilata}{dil}

\newcommand{\autori}{Ignazio Crocchiolo, Federico Dellini, Danilo Amalfitano, Emanuele Ciotola, Pierfrancesco Amendola (gruppo SDMs)}
\begin{document}
\begin{titlepage}
\centering
\vspace*{4cm}
{\LARGE\bfseries\color{accent} Percezione con incertezza dichiarata:\\[4pt] continuità visiva durante le occlusioni\\[4pt] e sicurezza delle indicazioni,\\[4pt] dal corto raggio al volo ad alta velocità\par}
\vspace{1.2cm}
{\Large Documento tecnico, versione 3\par}
\vspace{2.5cm}
{\large Autori: \autori\par}
\vspace{0.8cm}
{\large novembre 2026\par}
\vfill
\end{titlepage}

\section*{Sommario}
% Sommario (abstract) dell'allegato: solo il testo, da racchiudere in \begin{abstract}...\end{abstract} o in un paragrafo «Sommario».
Si descrive un sistema di percezione che dichiara ciò che non ha osservato. Per ogni oggetto tiene distinte tre risposte: lo vede; non lo vede, ma indica una regione dove può essere; non sa. Il non osservato non è mai dichiarato «libero»; l'avviso di possibile presenza si spegne solo dopo aver guardato tutta la zona; la perdita di affidabilità di un sensore è dichiarata. Il caso centrale è il volo ad alta velocità di aerei e droni, da $20$ a $210\unita{m/s}$, con collisioni, detriti ed eventi improvvisi fortemente impulsivi. Tutto è simulazione dalle schede tecniche, calcolo, o traiettorie pubbliche con un sensore simulato: nessuna prova usa un sensore fisico o hardware dedicato, il classificatore è supposto esistente e le etichette sono assegnate a mano, e il sistema fa solo osservazione, monitoraggio e avviso. Il prototipo è stato costruito e provato in un giorno.

Nei sei scenari di volo simulati, dall'avaria di un aereo alla collisione di due droni veloci, il sistema non ha mai dichiarato «libero» dove c'era l'oggetto o un frammento ($0/200$ $[0;\,0{,}019]$ nel controllo, $0/60$ $[0;\,0{,}060]$ nelle altre celle; intervalli di Wilson al $95\,\%$ per passaggio). Dopo la collisione veloce il punto vero e i frammenti sono rimasti nella regione dichiarata in $60/60$ passaggi $[0{,}940;\,1]$, grazie a un insieme di presenza con contenimento dimostrato per velocità fino a $210\unita{m/s}$, se l'errore iniziale è entro $r_0$. La collisione veloce, invece, non è ancora dichiarata come evento: una regola di contatto in prova la riconosce entro $2\unita{s}$ in 57 passaggi su 60 $[0{,}863;\,0{,}983]$ con un radar a scansione elettronica, ma apre ancora falsi eventi in un quasi incrocio (2 su 60) e non è adottata. A bordo di un drone in moto, durante una collisione a $30$--$100\unita{m}$ con fiammata, fumo e detriti, la regione del drone seguito non si è persa in $60/60$ passaggi per geometria con un radar a $360^\circ$; con il radar accecato fino a $2\unita{s}$ il «libero» resta assente, ma l'avviso dei detriti al portatore non è affidabile.

La perdita di confidenza dopo l'ultima misura è data in forma chiusa, $C(\tau)=1-\exp(-r^2/2\sigma^2(\tau))$ con $\sigma^2(\tau)$ polinomio cubico in $\tau$, dimostrata nel modello a velocità quasi costante ed estesa a una miscela di regimi di moto. Dopo una taratura, su traiettorie di volo pubbliche con un sensore simulato, la regione al $90\,\%$ ha contenuto l'aereo da 85 a 89 volte su 100 fra $1$ e $30\unita{s}$; i casi vengono per nove decimi da quattro voli di una sola rotta, e in virata la copertura scende a $0{,}73$ $[0{,}67;\,0{,}78]$, un limite dichiarato. Sensori allo stato dell'arte e calcolo dedicato sono stimati per analogia: il calcolo dedicato porterebbe l'elaborazione a $0{,}08$--$1{,}2\unita{ms}$ (ipotesi), ma la leva vera è la rivisita del sensore. Ogni limite trovato ha la sua miglioria e la prova che la deciderebbe.

Gli altri usi condividono lo stesso nucleo. A corto raggio, un prototipo rileva elementi mobili e ostacoli con una camera di profondità a stereo attivo, frontale e ferma, e continua a dare un'indicazione dichiarata quando un ostacolo nasconde il soggetto. Il sistema tiene distinte tre informazioni: l'osservato; il non osservato ma previsto, con una regione e un livello; il non noto. Lo stato «libero» è scritto solo per il suolo osservato nel fotogramma corrente; l'avviso di possibile presenza non termina senza osservazione diretta. Nelle ipotesi dichiarate sono dimostrati che ciò che non è osservato non è mai «libero», il contenimento del soggetto nascosto nell'insieme di possibile presenza (per una lettura esatta del suolo, senza il margine che il sistema usa) e la condizione sulla fascia di guardia. In simulazione, con modello ideale del sensore e un soggetto alla volta, la regione di previsione è emessa a ogni ciclo dell'occlusione una volta riconosciuto l'ingresso nell'ombra (0 interruzioni dopo l'ingresso su $23\,402$ passaggi nascosti; in 928 di essi la traccia si chiude prima dell'ombra e non c'è regione, e con una regola aggiuntiva, attiva solo dopo una prova d'ingresso, 24 restano senza avviso); con un ostacolo fermo largo $2\unita{m}$, il bordo dell'ombra noto e il soggetto sul proprio percorso, la regione contiene il soggetto in almeno 9 casi su 10 fino a $2\unita{s}$ e in almeno 85 su 100 fino a $3\unita{s}$; con un ostacolo largo $1\unita{m}$ la soglia di 9 su 10 non è raggiunta, e su profondità simulata dalle schede tecniche questa regola non è stata eseguita. Dietro un ostacolo fermo l'avviso resta acceso, e nella realizzazione provata resta acceso anche dopo la rimozione dell'ostacolo. Il classificatore degli oggetti è un componente richiesto dal sistema: in questo lavoro non è realizzato né valutato, ed è simulato con etichette assegnate a mano. Nessuna prova è stata eseguita su un sensore fisico: i risultati a corto raggio sono ottenuti a camera ferma, al chiuso, fino a $3\unita{m}$, e non sono stati verificati fuori da queste condizioni.

Su un sito aerospaziale sorvegliato, con droni e aerei in rullaggio, i sensori a lunga portata sono simulati dalle schede tecniche (radar per droni, radar del movimento a terra, LiDAR, camere, ADS-B e identificazione remota). Con $3\unita{m}$ di portata nessuno di questi scenari è servito: la distanza d'avviso richiesta va da $24\unita{m}$ a qualche chilometro. Il «libero» si dichiara per classe, con la probabilità di rilevamento cumulata e le passate trattate come correlate; l'avviso usa l'inviluppo dello scenario e si spegne solo guardando tutta la zona in cui il mezzo può essere. In simulazione dalle schede, un drone a 15--$21\unita{m/s}$ nascosto da un edificio è annunciato in tempo all'operatore in 200 passaggi su 200, e un aereo in rullaggio nascosto da un hangar in 200 su 200; una regola di spegnimento più semplice perde l'avviso in 564 passaggi su 2400. Per gli aerei in volo l'avviso non è tempestivo, e il sistema lo dichiara.
Le simulazioni mostrano anche dove il sistema cede. Il «libero» dipende dalla mappa degli ostacoli: con la mappa del sito dichiarata completa (un'ipotesi del simulatore) nessuno di 171 passaggi ha un «libero» sul drone, ma con una mappa dichiarata e sbagliata il «libero» dietro un edificio che non vi compare è falso; per questo, per impostazione, la mappa non è dichiarata completa. Quando il drone rallenta o si ferma l'avviso lo perde per qualche ciclo in 12 passaggi su 200. Dopo una collisione fra droni lenti l'evento è dedotto entro $2\unita{s}$ in 48 passaggi su 60 con l'integratore stabile dei frammenti (51 con il primo, instabile), con 3 eventi aperti prima dell'urto, e la zona dei detriti è di alcuni chilometri e non si riduce guardando. Le probabilità di rilevamento e le latenze dei sensori sono ipotesi.


\tableofcontents
\clearpage
\section*{Simboli e termini}\label{sec:simboli}
\addcontentsline{toc}{section}{Simboli e termini}

{\small
\begin{tabularx}{\textwidth}{@{}p{2.6cm}L@{}}
\toprule
Simbolo & Significato\\
\midrule
$x$, $y$, $z$ & riferimento della camera riportato al suolo: $x$ a destra, $z$ in avanti lungo il suolo, quota $y$ verso l'alto\\
$h$, $\theta$ & altezza e inclinazione verso il basso della camera\\
$\rho$, $\varphi$ & distanza e azimut di un punto del suolo\\
$d$ & profondità misurata lungo l'asse ottico\\
$f$, $b$ & focale in pixel e base stereo\\
$\sigma_d$, $\sigma_z(d)$ & rumore di disparità e rumore di profondità\\
$a$ & altezza di un ostacolo\\
$L$, $w_c$ & larghezza del pannello e larghezza del corpo\\
$\Delta$ & periodo di campionamento\\
$\mathbf x=[p_x,p_z,v_x,v_z]^{\mathsf T}$ & stato di una traccia: posizione e velocità al suolo\\
$F$, $Q$, $q$ & matrice di transizione, covarianza del rumore di processo e sua densità\\
$H$, $R$ & matrice di misura e covarianza del rumore di misura\\
$\nu$, $S$, $\gamma$ & innovazione, sua covarianza, soglia della porta di validazione\\
$P_D$ & probabilità di rilevamento per fotogramma\\
$P_{\mathrm{mob}}$, $P_{fa}$ & probabilità che la regola «mobile» dichiari in movimento un oggetto in moto e un oggetto fermo\\
$\vmax$ & velocità massima ammessa per un oggetto\\
$\tau$ & latenza dall'evento all'uscita\\
$\tocc$ & tempo trascorso dall'ultima osservazione\\
$t_a$ & tempo entro cui un soggetto fuori dall'ombra, in campo e in portata, è riportato come osservato\\
$t_{\mathrm{conf}}$ & tempo di conferma di una traccia\\
$\Oss(t)$ & insieme del suolo osservato libero\\
$\Omb(t)$, $\hat\Omb(t)$ & ombra di un ostacolo sul suolo, vera e stimata\\
$\varepsilon_\Omb$, $e_\Omega$ & limite dell'errore del bordo dell'ombra; errore misurato, con segno\\
$\Pres(t)$ & insieme di possibile presenza\\
$W(t)$, $w$ & area d'avviso e larghezza della fascia di guardia\\
$\Reg_\alpha(t)$ & regione di previsione al livello dichiarato $\alpha$ (qui $\alpha=0{,}90$)\\
$\dilata(A,r)$ & dilatazione dell'insieme $A$ di raggio $r$\\
$T_r$, $v_e$, $a_f$ & tempo di reazione, velocità propria e decelerazione di frenata di chi porta il sensore\\
$N$ & numerosità di una prova\\
$\ell^-$ & estremo inferiore dell'intervallo di Wilson al 95\,\%\\
\bottomrule
\end{tabularx}

\medskip
\begin{tabularx}{\textwidth}{@{}p{3.6cm}L@{}}
\toprule
Termine & Significato\\
\midrule
passaggio & un soggetto attraversa il campo, scompare dietro un ostacolo fermo e di norma ricompare; è l'unità statistica delle prove\\
traccia & la stima dello stato di un oggetto seguito, tenuta dal filtro da un ciclo al successivo\\
regola «mobile» & regola che dichiara «in movimento» un gruppo in base allo spostamento su una finestra di tempo\\
insieme di possibile presenza & parte del suolo in cui un oggetto non osservato può trovarsi; si riduce solo per osservazione\\
area d'avviso & l'ombra stimata dell'ostacolo allargata della fascia di guardia, su cui esce l'avviso di possibile presenza\\
fascia di guardia & allargamento dell'ombra che copre il tempo $t_a$, la latenza, la tenuta di un ciclo e l'errore del bordo\\
regione di previsione & regione emessa a ogni ciclo per una traccia entrata nell'ombra, con il suo livello\\
«previsto ampio» & stato della regione quando il livello non è dichiarabile: la regione resta, senza punto e senza livello\\
identità presunta & identità riassegnata alla ricomparsa: è un'ipotesi, e l'avviso non ne dipende\\
affidabilità persa & segnale emesso quando il sensore non è in salute; la mappa è dichiarata «non osservato»\\
\bottomrule
\end{tabularx}}

Il livello di prova di un risultato è indicato da una lettera: \prova{D} dimostrato nel modello dichiarato; \prova{C} verificato per calcolo; \prova{S} verificato in simulazione; \prova{P} misurato su dati pubblici; \prova{N} non verificato sperimentalmente.

\clearpage
\section{Introduzione: problema, requisiti e sintesi dei risultati}\label{sec:intro}

\subsection{Il problema}\label{sec:intro-problema}

Il testo della sfida a cui il progetto risponde osserva che «riconoscere un oggetto o un soggetto in movimento non basta: un ostacolo può nasconderlo temporaneamente, mentre la scena continua a cambiare». Chiede un prototipo capace di «rilevare elementi mobili e ostacoli, gestire interruzioni della visibilità e distinguere le informazioni osservate da quelle incerte», che dia «feedback spaziali chiari e tempestivi quando le condizioni lo consentono, segnalando invece la perdita di affidabilità quando necessario», con «continuità percettiva, bassa latenza ed efficienza elaborativa, senza rinunciare alla sicurezza delle indicazioni».

Il riconoscimento non basta perché un rilevatore non produce nulla nell'istante in cui il soggetto esce dalla vista. Da quell'istante un sistema può tacere, e il silenzio è letto come assenza; può continuare a mostrare l'ultima posizione, o una posizione estrapolata, come se fosse osservata; oppure può dichiarare che cosa sa e con quale livello. Le prime due scelte sono indicazioni non sicure: un filtro privo di misure accumula errore per tutta l'occlusione \cite{cao2023}, e in simulazione, con modello ideale del sensore, un punto lasciato fermo all'ultima misura dista dal soggetto $1{,}48$--$1{,}54\unita{m}$ in mediana dopo $2\unita{s}$, contro $0{,}37$--$0{,}39\unita{m}$ del punto previsto (una disposizione, miscela di comportamenti, $N$ fra 2077 e 2115 passaggi) \prova{S}. Questo allegato descrive un sistema che adotta la terza scelta.

Il sistema usa una camera di profondità a stereo attivo, frontale e ferma, ad altezza $h=1{,}2\unita{m}$ e inclinazione $\theta=15^\circ$ verso il basso, con portata utile di $3\unita{m}$. La catena di elaborazione è: mappa polare del suolo a quattro stati (libero, occupato osservato, in ombra, non osservato); gruppi e regola «mobile»; associazione e filtro di Kalman \cite{kalman1960} a velocità quasi costante; durante un'occlusione, una regione di previsione $\Reg_\alpha(t)$ con livello dichiarato e un avviso di «possibile presenza» sull'ombra dell'ostacolo; regole scritte, non apprese; uscita e segnale di «affidabilità persa». Il sistema osserva e avvisa: nessuna uscita è un comando, e non vi sono funzioni d'ingaggio o di puntamento.

\subsection{Requisiti misurabili}\label{sec:intro-requisiti}

La Tabella~\ref{tab:intro-requisiti} traduce ogni requisito della sfida in una grandezza, con il modo di misurarla e la soglia di riferimento. L'unità statistica è il \emph{passaggio}: un soggetto attraversa il campo, scompare dietro un ostacolo fermo e di norma ricompare. Le definizioni complete e gli intervalli stanno in \S\ref{sec:val}; le soglie sono criteri di progetto e il loro esito è nelle Tabelle~\ref{tab:intro-sintesi} e~\ref{tab:intro-sintesi-due}. Ai requisiti della sfida il progetto aggiunge un componente di classificazione su più classi di oggetti in movimento (persone; veicoli di tipo diverso, per esempio auto, treni, aerei a terra; altri ostacoli), ciascuna con un profilo cinematico, e un profilo deterministico, da configurazione iniziale, di chi porta il sensore.

\begin{table}[!t]
\centering\footnotesize
\caption{Requisiti della sfida e grandezze che li misurano. $\tocc$: tempo trascorso dall'ultima osservazione; $\ell^-$: estremo inferiore dell'intervallo di Wilson al 95\,\% \cite{wilson1927}; $\Delta=1/15\unita{s}$: periodo di campionamento.}\label{tab:intro-requisiti}
\begin{tabularx}{\textwidth}{@{}>{\RaggedRight\arraybackslash}p{2.5cm}LL>{\RaggedRight\arraybackslash}p{3.7cm}@{}}
\toprule
Requisito & Grandezza & Come si misura & Soglia di riferimento\\
\midrule
Rilevare elementi mobili & quota dei soggetti a $0{,}5\unita{m/s}$ dichiarati «in movimento» entro $1\unita{s}$; quota dei passaggi in cui un oggetto fermo è dichiarato «in movimento» & conteggio per passaggio, con intervallo & almeno 9 su 10 fino a $2\unita{m}$; la seconda quota è riportata con il suo intervallo\\
Rilevare ostacoli & ostacolo fermo localizzato entro $15\unita{cm}$ dalla posizione vera & conteggio per prova & 5 prove su 5 fra $1{,}1$ e $3\unita{m}$\\
Gestire le interruzioni della visibilità: continuità & quota dei cicli senza osservazione in cui è emessa una regione $\Reg_\alpha(t)\neq\emptyset$; passaggi con almeno un ciclo senza regione & sui passaggi con almeno 4 misure prima della scomparsa & 0 passaggi con un'interruzione\\
\quad copertura & quota dei passaggi ancora nascosti al tempo $\tocc$ con il soggetto in $\Reg_\alpha$, al livello dichiarato $\alpha=0{,}90$; area di $\Reg_\alpha$ rispetto all'ombra & per comportamento e per $\tocc\in\{0{,}5;1;2;3\}\unita{s}$ & $\ell^->0{,}90$ fino a $2\unita{s}$ e $\ell^->0{,}85$ a $3\unita{s}$, in ogni combinazione\\
\quad identità & quota dei passaggi con la stessa identità alla ricomparsa & sui passaggi con ricomparsa & almeno $0{,}95$ per chi resta sul proprio percorso\\
Distinguere osservato e incerto & uscite «osservato» emesse mentre il soggetto è nascosto & conteggio sui cicli con il soggetto nascosto & 0\\
Feedback chiari e tempestivi & tempo dal primo fotogramma alla chiamata dell'uscita sonora; presenza di stato, età e livello in ogni uscita & misura di tempo; controllo sullo schema dell'uscita & al più $300\unita{ms}$\\
Segnalare la perdita di affidabilità & ritardo del segnale dedicato; segnali falsi & prova con flusso interrotto, sensore coperto o spostato & entro $1\unita{s}$; 0 falsi in $60\unita{s}$\\
Bassa latenza ed efficienza & latenza $\tau$ dall'evento all'uscita; tempo d'orologio per ciclo su un filo di CPU & 95º percentile; numero di cicli oltre $\Delta$ & 95º percentile di $\tau$ entro $150\unita{ms}$; ciclo entro $\Delta=66{,}7\unita{ms}$\\
Sicurezza delle indicazioni & celle «libero» contenute nell'ombra vera; istanti con il soggetto nascosto fuori dall'area d'avviso; avvisi tolti senza osservazione diretta & conteggi su celle, istanti e passaggi & 0 per ciascuna\\
Classificazione e profili (requisito di progetto) & differenze fra due esecuzioni sugli stessi fotogrammi che differiscono solo nelle etichette di classe & confronto di celle «libero», area d'avviso e livelli & 0 celle «libero» in più; 0 avvisi ridotti; nessun livello più basso\\
\bottomrule
\end{tabularx}
\end{table}

\subsection{L'idea}\label{sec:intro-idea}

Il sistema tiene distinte tre informazioni e non le mescola in nessuna uscita (Figura~\ref{fig:intro-catena}). La prima è ciò che è \emph{osservato} nel fotogramma corrente: il suolo visto libero $\Oss(t)$, gli ostacoli, i soggetti con una misura nel ciclo. La seconda è ciò che \emph{non è osservato ma è previsto}: per un soggetto seguito ed entrato nell'ombra $\Omb(t)$ di un ostacolo, una regione $\Reg_\alpha(t)$ emessa a ogni ciclo con il tempo trascorso $\tocc$ e, quando è dichiarabile, il livello; quando il livello non è dichiarabile la regione resta, con lo stato «previsto ampio». La terza è ciò che \emph{non è noto}: fuori campo, oltre la portata, pixel senza dato, ombra in cui non è entrato alcun soggetto seguito.

Il principio di sicurezza è uno: l'assenza di un rilevamento non è mai un'indicazione di via libera. Lo stato «libero» è scritto solo per il suolo osservato nel fotogramma corrente, e nessun contatore di tempo, modello appreso, etichetta di classe, identità o previsione lo produce. L'insieme di possibile presenza $\Pres(t)$ si riduce solo dove il suolo è rivisto libero; previsione e classe possono allargare un avviso, mai restringerlo. Il principio è analogo a quello dichiarato da un costruttore di sistemi industriali di sorveglianza tridimensionale, che trattano come occupato ciò che nessun sensore vede \cite{therobotreport-veo}, e a quello dei metodi che sovrastimano lo stato degli ostacoli nascosti \cite{orzechowski2018}; l'analogia riguarda il principio, non il livello di integrità né la conformità a una norma. Quando il sensore non è in salute il sistema emette un segnale dedicato, dichiara tutta la mappa «non osservato» e continua a uscire a ogni ciclo.

\begin{figure}[H]
\centering
\begin{tikzpicture}[font=\scriptsize,>=Stealth,node distance=0.42cm,
  b/.style={draw,rounded corners=1pt,align=center,text width=2.05cm,minimum height=0.95cm,inner sep=2pt}]
\node[b] (cam) {camera di profondità ferma};
\node[b,right=of cam] (map) {mappa polare del suolo a quattro stati};
\node[b,right=of map] (gru) {gruppi e regola «mobile»};
\node[b,right=of gru] (fil) {associazione e filtro di Kalman};
\node[b,right=of fil] (occ) {occlusione: $\Reg_\alpha(t)$ e $W(t)$};
\node[b,right=of occ] (reg) {regole e uscita};
\draw[->] (cam)--(map); \draw[->] (map)--(gru); \draw[->] (gru)--(fil); \draw[->] (fil)--(occ); \draw[->] (occ)--(reg);
\node[b,below=0.5cm of cam] (sal) {salute del sensore};
\node[b,dashed,below=0.5cm of fil] (cla) {classe: richiesta, qui simulata};
\node[b,text width=4.55cm,below=0.5cm of $(occ)!0.5!(reg)$] (usc) {mai mescolati: osservato; previsto, con regione, livello e $\tocc$; non noto};
\draw[->] (cam)--(sal);
\draw[->,dashed] (cla)--(fil);
\draw[->] (reg.south)--(reg.south|-usc.north);
\draw[->] (sal.south)--++(0,-0.3)-|node[pos=0.25,below]{«affidabilità persa»: scavalca le regole}(usc.south);
\end{tikzpicture}
\caption{La catena di elaborazione e le tre informazioni in uscita. Il tratteggio indica il componente di classificazione, richiesto dal sistema e in questo lavoro simulato con etichette assegnate a mano.}\label{fig:intro-catena}
\end{figure}

\subsection{Contributi}\label{sec:intro-contributi}

\begin{enumerate}[label=\arabic*.]
\item La geometria d'osservazione in forma chiusa: dove il suolo è visibile, l'ombra di un ostacolo, il rumore di profondità e l'effetto di un errore di posa (\S\ref{sec:geo}).
\item Una mappa polare del suolo a quattro stati in cui «libero» è scritto solo dal fotogramma corrente, con la proposizione di sicurezza, i controesempi alle varianti con memoria e il confronto con le mappe persistenti e a voxel (\S\ref{sec:map}).
\item La catena di rilevamento, associazione e stima, con la prova che la massa non è identificabile da misure di posizione e che l'accelerazione non è stimabile con precisione utile; il componente di classificazione su più classi, con i profili cinematici di classe e i requisiti che il sistema gli pone (\S\ref{sec:sti}).
\item Il trattamento dell'occlusione: insieme di possibile presenza con contenimento e monotonia, fascia di guardia, regione di previsione che, per costruzione, non si svuota finché la traccia è nell'ombra (se l'ingresso nell'ombra è stato riconosciuto e l'ostacolo è fermo), con un livello dichiarato solo nel dominio in cui è stato misurato in simulazione, e l'enunciato che previsione ed etichetta non riducono l'avviso (\S\ref{sec:occ}).
\item Un metodo di valutazione definito sul passaggio: misure, intervalli, numerosità, limiti della taratura per quantili, lettura di un conteggio «dentro $k$ su $n$» (\S\ref{sec:val}).
\item Il costo di calcolo e le alternative escluse, ciascuna con il controesempio che la esclude (\S\ref{sec:cos}).
\item Un'uscita che porta sempre stato, età e livello, il profilo deterministico di chi porta il sensore con la distanza d'avviso, il segnale di affidabilità persa e il dominio di validità (\S\ref{sec:lim}).
\end{enumerate}

\subsection{Sintesi dei risultati e livello di prova}\label{sec:intro-sintesi}

Le Tabelle~\ref{tab:intro-sintesi} e~\ref{tab:intro-sintesi-due} riportano i risultati principali. Ogni risultato porta il livello di prova: \prova{D} dimostrato nel modello dichiarato; \prova{C} verificato per calcolo; \prova{S} verificato in simulazione; \prova{P} misurato su dati pubblici; \prova{N} non verificato sperimentalmente. Le simulazioni sono di due tipi: \emph{ideale}, cioè con modello ideale del sensore (posizioni al suolo con rumore gaussiano di $5\unita{cm}$, probabilità di rilevamento per fotogramma $P_D$, sensore e ostacolo fermi, un soggetto per ombra), e \emph{da scheda}, cioè su profondità simulata dalle schede tecniche di tre camere a stereo attivo a $424\times240$ pixel e, dove indicato, a $848\times480$, a posa nota. Gli intervalli fra parentesi quadre sono di Wilson al 95\,\%; gli estremi superiori degli intervalli per zero eventi sono arrotondati per eccesso ($2{,}85\,\%$ è scritto $2{,}9\,\%$ e $0{,}0155\,\%$ è scritto $0{,}016\,\%$); gli altri estremi al valore più vicino, salvo l'estremo inferiore $0{,}875$ del pannello largo $1\unita{m}$, scritto $0{,}87$. Le cifre vengono da catene di prova distinte, che non sono state eseguite tutte insieme in un solo programma: la simulazione d'insieme con modello ideale del sensore (regione persistente e area d'avviso); la catena su profondità simulata nella configurazione di riferimento (gruppi per componenti connesse, senza regione persistente); la stessa catena con i gruppi a celle forti, le condizioni aggiuntive del riaggancio e l'avviso che resta dopo il riaggancio; l'elaborazione di traiettorie pubbliche con il solo filtro. Ogni riga delle tabelle dice da quale viene. I gruppi a celle forti, le condizioni aggiuntive del riaggancio e l'avviso per le tracce chiuse vicino a un'ombra si attivano all'avvio solo se le prove d'ingresso del \S\ref{sec:lim-salute} sono superate.

\begin{table}[p]
\centering\footnotesize
\setlength{\tabcolsep}{4pt}
\caption{Risultati principali, con condizioni, numerosità e livello di prova.}\label{tab:intro-sintesi}
\begin{tabularx}{\textwidth}{@{}>{\RaggedRight\arraybackslash}p{2.8cm}>{\RaggedRight\arraybackslash}p{4.1cm}L>{\RaggedRight\arraybackslash}p{1.2cm}>{\RaggedRight\arraybackslash}p{1.05cm}@{}}
\toprule
Risultato & Condizioni & Valore & Prova & Dove\\
\midrule
Nessuna cella non osservata è dichiarata «libero» & regola di scrittura unica; per l'inverso, limitato ai raggi dei pixel che leggono il suolo e a meno del margine di lettura: profondità valida corretta, posa esatta, suolo piano & proposizione; fuori dalle ipotesi (vetro, errore di posa di $1^\circ$) cade & \prova{D} & \S\ref{sec:map}\\
«Libero» della cella & ideale, errore d'inclinazione entro $0{,}5^\circ$; 300 condizioni per 30 fotogrammi, due lanci & 0 celle «libero» su un ostacolo; per condizione $[0;\,1{,}3\,\%]$ & \prova{S} & \S\ref{sec:map}\\
«Libero» nell'ombra vera & da scheda, variante che conserva il «libero» per tre fotogrammi; 135 sequenze, 131 delle quali con i filtri della camera spenti & 0 su $2\,907\,032$ celle\,$\cdot$\,fotogramma; per sequenza, sulle 131 con i filtri spenti, $[0;\,2{,}9\,\%]$. In una seconda serie di 381 sequenze, 5 su $3\,992\,186$ (una al livello di rumore nominale): non è risolto se dipendano dalla regola o dal bordo della maschera di riferimento & \prova{S} & \S\ref{sec:val}\\
Contenimento in $\Pres(t)$ e riduzione solo per osservazione & profondità valida corretta; posa esatta, o libero eroso del margine di posa; suolo piano; lettura per punto esatta (margine di lettura nullo); velocità al più $\vmax$ & proposizioni; con il margine di lettura e con la mappa a celle il contenimento poggia sulla simulazione & \prova{D} & \S\ref{sec:occ}\\
Fascia di guardia di $0{,}40\unita{m}$ & tempo di rilevamento assunto $t_a=0{,}133\unita{s}$, latenza stimata $0{,}155\unita{s}$, $\Delta=0{,}067\unita{s}$; ostacolo fermo & copre, per il centro del soggetto, $\vmax=1{,}13\unita{m/s}$; $0{,}56\unita{m/s}$ se il bordo dell'ombra sbaglia di $0{,}20\unita{m}$ & \prova{D}\prova{C} & \S\ref{sec:occ}\\
Continuità & ideale; 12 combinazioni (3 disposizioni, $P_D=0{,}9$ e $0{,}7$, corpo esteso e puntiforme) per 5 comportamenti; premessa: ingresso nell'ombra riconosciuto & 0 passaggi con un ciclo senza regione dopo l'ingresso, su $23\,402$ passaggi nascosti (6000 camminatori distinti: meno di 7 su $10\,000$). La premessa manca in 928 passaggi ($3{,}97\,\%$; 852 con $P_D=0{,}7$), senza regione: la soglia 0 della Tabella~\ref{tab:intro-requisiti} è rispettata solo dove l'ingresso è riconosciuto. Con l'avviso acceso anche per le tracce chiuse entro $0{,}75\unita{m}$ da un'ombra (regola attiva solo dopo la prova d'ingresso sul suolo vuoto) restano senza avviso 24 passaggi ($0{,}10\,\%$). Con corpo puntiforme: 0 su $24\,835$ $[0;\,0{,}016\,\%]$. Da scheda la regola della regione persistente non è stata eseguita & \prova{D}\prova{S}\prova{N} & \S\ref{sec:occ}, \S\ref{sec:val}\\
Copertura di $\Reg_\alpha$ al livello $0{,}90$ & ideale; ostacolo largo $2\unita{m}$; corpo esteso; soggetto sul proprio percorso; $P_D=0{,}9$; bordo dell'ombra noto & $\ell^->0{,}90$ fino a $2\unita{s}$ e $>0{,}85$ a $3\unita{s}$ in 30 combinazioni su 30; peggiori $2125/2279$ ($\ell^-=0{,}921$) e $1754/1922$ ($\ell^-=0{,}899$); da 608 a 2359 casi per combinazione & \prova{S} & \S\ref{sec:val}\\
\quad pannello largo $1\unita{m}$ & come sopra & a $1\unita{s}$ $942/1045$ ($\ell^-=0{,}882$; $0{,}87$ su cinque combinazioni insieme); oltre $1\unita{s}$ nessuna cifra. Con corpo puntiforme: $826/879$ a $2\unita{s}$ ($\ell^-=0{,}922$) & \prova{S} & \S\ref{sec:val}\\
\quad fuori dal dominio & ideale; chi svolta o esce dal lato lontano; bordo dell'ombra stimato più stretto di 10--20\,cm & a 2 e $3\unita{s}$: $0{,}13$--$0{,}70$ con corpo puntiforme ($N$ da 538 a 1791), $0{,}57$--$0{,}79$ con corpo esteso e ostacolo largo $2\unita{m}$; bordo stimato: $0{,}30$--$0{,}46$ a $1\unita{s}$ (estremo basso: chi prosegue). Resta la sola area d'avviso & \prova{S} & \S\ref{sec:val}\\
Avviso di possibile presenza & ideale; un soggetto per ombra & soggetto nascosto dentro l'area d'avviso per costruzione; avviso non acceso in 29 passaggi su $95\,399$ (ombra esatta) e in 15 su 2000 (ombra stimata più stretta di $0{,}20\unita{m}$). Da scheda, con i gruppi a celle forti e l'avviso che resta dopo il riaggancio: soggetto nascosto entro $3\unita{m}$ senza avviso in 0 fotogrammi su 4324 a 424 pixel di larghezza e in 0 su 4325 a 848 (con la configurazione di riferimento, sugli stessi 150 passaggi, in 180 e in 271), e in 0 su 3608 per risoluzione in una replica. Dietro un ostacolo fermo l'avviso resta acceso; che l'avviso termini quando il suolo è rivisto dopo la rimozione dell'ostacolo non è sostenuto: nella realizzazione provata resta acceso dopo $2\unita{s}$ in 10 prove su 10 (72--100\,\%), per una causa non determinata & \prova{D}\prova{S} & \S\ref{sec:occ}\\
\bottomrule
\end{tabularx}
\end{table}

\begin{table}[p]
\centering\footnotesize
\setlength{\tabcolsep}{4pt}
\caption{Risultati principali (seguito).}\label{tab:intro-sintesi-due}
\begin{tabularx}{\textwidth}{@{}>{\RaggedRight\arraybackslash}p{2.8cm}>{\RaggedRight\arraybackslash}p{4.1cm}L>{\RaggedRight\arraybackslash}p{1.2cm}>{\RaggedRight\arraybackslash}p{1.05cm}@{}}
\toprule
Risultato & Condizioni & Valore & Prova & Dove\\
\midrule
Identità alla ricomparsa & ideale, sul percorso, tre disposizioni & $10\,711/10\,757$; $9228/9240$; $6772/6967$ ($\ell^-=0{,}994$; $0{,}998$; $0{,}968$) & \prova{S} & \S\ref{sec:sti}\\
\quad su profondità simulata & da scheda; configurazione di riferimento / gruppi a celle forti e condizioni aggiuntive del riaggancio, soggetto a $0{,}5\unita{m}$ dietro il pannello & 0 su 90 $[0;\,4{,}1\,\%]$ / 176 su 180 $[94{,}4;\,99{,}1\,\%]$ a 424 pixel di larghezza e 174 su 180 $[92{,}9;\,98{,}5\,\%]$ a 848: un tasso in un dominio ristretto, non il superamento di una soglia per ciascun comportamento. Un secondo soggetto nella stessa ombra riceve l'identità del primo (19 casi su 20 da scheda; 2400 su 2400 con modello ideale). L'avviso resta per costruzione; da scheda, in una scena con due persone, 0 fotogrammi su 336 con il primo soggetto fuori dall'avviso; con modello ideale la verifica non è stata eseguita & \prova{S}\prova{N} & \S\ref{sec:val}\\
Modello di moto & traiettorie pubbliche di pedoni \cite{thor}; 219 tratti, 9 persone; $q=0{,}2\unita{m^2/s^3}$ & copertura della regione al 90\,\% del filtro: $0{,}977$; $1{,}000$; $1{,}000$ a 1, 2, $3\unita{s}$ ($[0{,}948;\,0{,}990]$ a $1\unita{s}$, tratti contati come indipendenti) & \prova{P} & \S\ref{sec:val}\\
Regola «mobile» & ideale, $0{,}5\unita{m/s}$ da $2\unita{m}$, $N=2000$; da scheda, pannello fermo, 90 passaggi & «in movimento» entro $1\unita{s}$: 100\,\% ($\ell^-=99{,}81\,\%$); pannello fermo dichiarato mobile in 79 su 90 $[79{,}4;\,93{,}0\,\%]$; in 5 su 90 con la soglia minima di punti portata da 3 a 40, non adottata perché una scatola alta $15\unita{cm}$ a $2\unita{m}$ è allora vista in 3 prove su 10; con i gruppi a celle forti in 0 su 140 nel dominio dell'identità al ritorno e in 7 su 10 fuori da esso & \prova{S} & \S\ref{sec:sti}\\
Ostacoli fermi & da scheda; una sequenza per caso & entro $15\unita{cm}$ dalla posizione vera: cassa alta $1\unita{m}$ a $1{,}68\unita{m}$, 84 fotogrammi su 84; scatola di $15\unita{cm}$ a $2{,}0\unita{m}$, 10 su 10; la stessa a $0{,}8\unita{m}$ e una lastra di vetro, 0 su 10. Descrivono un caso, non una frequenza: la soglia di riferimento (5 prove su 5 fra $1{,}1$ e $3\unita{m}$) non è stata provata in questa forma & \prova{S} & \S\ref{sec:val}\\
Osservato distinto dal previsto & da scheda; configurazione di riferimento / gruppi a celle forti & uscita «osservato» con il soggetto nascosto in 8 passaggi su 90 in una prima serie; in una seconda, riagganci a soggetto ancora nascosto in 11 e 6 passaggi su 140 (424 e 848 pixel di larghezza) / in 1 e 1 su 140: la soglia 0 non è rispettata & \prova{S} & \S\ref{sec:sti}\\
Massa e accelerazione & punto materiale con forza ignota; misure di posizione, $\sigma_p=3\unita{cm}$ assunto, $15\unita{Hz}$ & massa non identificabile; scarto tipo dell'accelerazione $16{,}5\unita{m/s^2}$ contro $0{,}67\unita{m/s^2}$ di valore efficace & \prova{D}\prova{C} & \S\ref{sec:sti}\\
Costo per ciclo & da scheda, 9152 cicli, un filo di CPU da portatile; una registrazione pubblica affollata, 320 cicli & media $6{,}35\unita{ms}$, 95º percentile $18{,}6\unita{ms}$, 4 cicli oltre $\Delta$; media $81\unita{ms}$ sulla registrazione pubblica & \prova{S}\prova{P} & \S\ref{sec:cos}\\
Latenza & somma delle medie delle voci note, nessuna misurata sulla catena & $155\unita{ms}$; il 95º percentile non è disponibile; il tempo dal primo fotogramma all'uscita sonora non è misurato & \prova{C}\prova{N} & \S\ref{sec:lim}\\
Affidabilità persa & simulazione, rumore dell'unità inerziale ipotizzato & segnale in $0{,}13\unita{s}$; 0 falsi in $60\unita{s}$ & \prova{S} & \S\ref{sec:lim}\\
Classe e avviso & per costruzione; sequenze casuali di etichette sulla sola fascia di guardia & 0 riduzioni dell'avviso su $20\,000$, e 0 su altre $20\,000$; mai eseguito sulla catena completa & \prova{D}\prova{S}\prova{N} & \S\ref{sec:occ}\\
Avviso tempestivo per classe & portata $3\unita{m}$, avvicinamento frontale, reazione di $1{,}0\unita{s}$ ipotizzata; velocità delle classi in parte ipotizzate (\S\ref{sec:sti}) & fino a $1{,}44\unita{m/s}$: 3 classi su 15 alla velocità tipica, 0 su 15 alla massima & \prova{C} & \S\ref{sec:occ}, \S\ref{sec:lim}\\
\bottomrule
\end{tabularx}
\end{table}

\subsection{Che cosa l'allegato non dimostra}\label{sec:intro-limiti}

\begin{riquadro}{Limiti di validità}
\textbf{Nessuna misura su un sensore fisico.} Ogni risultato è una dimostrazione in un modello dichiarato, un calcolo, una simulazione (con modello ideale del sensore o su profondità simulata dalle schede tecniche) o un'elaborazione di dati pubblici. Copertura della regione, errore del bordo dell'ombra, identità alla ricomparsa e latenza non sono verificati su un sensore \prova{N}. Il simulatore da scheda produce meno pixel senza dato di due registrazioni pubbliche di camere a stereo attivo con cui è stato confrontato ($7{,}8\,\%$ al livello di rumore più alto, contro $24{,}1\,\%$ e $34{,}4\,\%$; camere e scene diverse da quelle simulate, \S\ref{sec:val}): le cifre al livello nominale sono ottimistiche.

\textbf{Dominio.} Camera ferma, a posa nota entro $0{,}5^\circ$; al chiuso; fino a $3\unita{m}$; suolo piano; ostacolo fermo, opaco, che restituisce profondità e nasconde il soggetto per intero; un soggetto per ombra, visto per almeno 4 fotogrammi prima di scomparire. Il livello della regione è misurato solo in simulazione con modello ideale del sensore, e vale con un ostacolo largo $2\unita{m}$, il bordo dell'ombra noto e per chi resta sul proprio percorso; con un ostacolo largo $1\unita{m}$ è «almeno 87 su 100» fino a $1\unita{s}$; non è dichiarata alcuna cifra oltre $3\unita{s}$.

\textbf{Classificazione.} Il classificatore è un componente richiesto dal sistema; in questo lavoro non è realizzato né valutato: è simulato con etichette assegnate a mano, eventualmente con errori iniettati per misurare la sensibilità. Nessuna classe diversa da «persona» è stata simulata nella catena completa. Un'etichetta non produce un «libero» né riduce un avviso acceso; può cambiare l'esito di un riaggancio, e quindi gli ingressi successivi (\S\ref{sec:occ-classe}).

Il sistema non è un dispositivo di protezione e non sostituisce il bastone, il cane guida o le procedure di sicurezza di un impianto.
\end{riquadro}

\subsection{Ambiti d'uso}\label{sec:intro-ambiti}

La sfida ammette due applicazioni. La prima è il monitoraggio di attività e oggetti in ambienti aerospaziali e industriali, con il sensore su una stazione fissa o su un veicolo senza equipaggio fermo, e un'uscita a schermo e su registro degli eventi. La seconda è l'avviso di presenza per una persona non vedente ferma, con il sensore su un supporto accanto alla persona e un'uscita sonora. Hanno in comune la catena, la mappa, il trattamento dell'occlusione e il principio di sicurezza; cambiano il profilo di chi porta il sensore, il canale d'uscita e il comportamento alla perdita di affidabilità (\S\ref{sec:lim}). In entrambe il sistema svolge solo osservazione, monitoraggio e avviso: nessuna prova è stata fatta con il sensore in moto o indossato.

\subsection{Struttura del documento}\label{sec:intro-struttura}

\S\ref{sec:geo} descrive il sensore e la geometria d'osservazione; \S\ref{sec:map} la rappresentazione dello spazio e la sicurezza del «libero»; \S\ref{sec:sti} il rilevamento, la classificazione, l'associazione e la stima dello stato; \S\ref{sec:occ} l'occlusione; \S\ref{sec:val} il metodo di valutazione e i risultati; \S\ref{sec:cos} il costo di calcolo, le alternative considerate e i lavori correlati; \S\ref{sec:lim} l'uscita, il profilo di chi porta il sensore, la perdita di affidabilità e i limiti; \S\ref{sec:con} riassume che cosa è dimostrato, che cosa è verificato in simulazione e che cosa resta non verificato sperimentalmente. L'Appendice~\ref{sec:par} raccoglie i parametri e l'Appendice~\ref{sec:rip} le indicazioni per rifare i conti.

\section{Sensore e geometria d'osservazione}\label{sec:geo}

Questo capitolo fissa che cosa la camera può osservare dal punto di montaggio e con quale incertezza. Tutti i risultati sono dimostrazioni in un modello dichiarato, calcoli, letture di schede dei costruttori, misure pubblicate da terzi o simulazioni: nessuna grandezza è stata misurata su un sensore fisico, e dove un risultato dipende dal sensore reale lo si scrive come limite di validità. Ai simboli comuni si aggiungono: $\Phi_h$, $\Phi_v$ ampiezze orizzontale e verticale del campo, $T_h=\tan(\Phi_h/2)$, $T_v=\tan(\Phi_v/2)$; $\beta_-=\theta+\Phi_v/2$ e $\beta_+=\Phi_v/2-\theta$, angoli del bordo basso e del bordo alto del campo rispetto all'orizzontale; $a$ altezza di un ostacolo, $s$ statura di una persona, $z_o$ e $z_p$ coordinate in avanti dell'ostacolo e della persona.

\subsection{Modello della camera}\label{sec:geo-camera}

\begin{definizione}[Camera]\label{def:geo-camera}
La camera è un foro stenopeico senza distorsione, con centro $\mathbf c=(0,h,0)$, asse ottico $\mathbf e_d=(0,-\sin\theta,\cos\theta)$, asse d'immagine verso il basso $\mathbf e_\eta=(0,-\cos\theta,-\sin\theta)$ e verso destra $\mathbf e_\xi=(1,0,0)$. Un punto $P=(x,y,z)$ ha profondità e coordinate d'immagine normalizzate
\begin{equation}\label{eq:geo-proiezione}
d=(P-\mathbf c)\cdot\mathbf e_d=(h-y)\sin\theta+z\cos\theta,\qquad \xi=\frac{x}{d},\qquad \eta=\frac{(h-y)\cos\theta-z\sin\theta}{d},
\end{equation}
ed è nel campo se $d>0$, $|\xi|\le T_h$ e $|\eta|\le T_v$; il suo pixel è $(c_x+f\,\xi,\;c_y+f\,\eta)$. Posa di riferimento: $h=1{,}2\unita{m}$, $\theta=15^\circ$.
\end{definizione}

\begin{table}[htb]\centering\small
\caption{Profili di sensore considerati: tre camere di profondità a stereo attivo. Valori dichiarati dai costruttori per il flusso a 848 × 480 pixel e 30 fotogrammi al secondo; nessun valore è misurato su un esemplare. I valori senza richiamo bibliografico sono letti dalle schede tecniche.}\label{tab:geo-profili}
\begin{tabularx}{\textwidth}{@{}p{4.1cm}p{2.9cm}p{3.3cm}L@{}}
\toprule
 & RealSense D455 & RealSense D435i & Orbbec Gemini 335\\
\midrule
Campo $\Phi_h\times\Phi_v$ & 87° × 58° \cite{realsense-d455} & 87° × 58° \cite{realsense-confronto-camere} & 90° × 60° (90° × 65° a 1280 × 800 \cite{orbbec-gemini-335})\\
Base stereo $b$ & 95 mm \cite{realsense-d455} & 50 mm \cite{realsense-taratura-profondita} & 50 mm \cite{orbbec-gemini-335}\\
Focale $f$ (conti di rumore) & 424 px & 424 px \cite{realsense-taratura-profondita} & 410,7 px\\
Distanza minima & 0,35 m & 0,195 m & 0,18 m\\
Intervallo ideale & 0,6--6 m \cite{realsense-d455} & 0,3--3 m \cite{realsense-confronto-camere} & 0,3--3 m (scheda); 0,26--3 m (pagina del prodotto \cite{orbbec-gemini-335})\\
$\sigma_d$ nominale / limite & 0,08 / 0,20 px & 0,08 / 0,21 px & 0,082 / 0,154 px\\
Tolleranza sul campo & ±3° & ±3° & ±3°\\
\bottomrule
\end{tabularx}
\end{table}

\begin{osservazione}[Focale e campo]\label{oss:geo-focale}
Nel foro stenopeico la focale in pixel è la semilarghezza dell'immagine divisa per $T_h$: con 848 px, 87° danno 446,8 px, e 58° su 480 px danno 433,0 px; $f=424\unita{px}$ corrisponde invece a 90° × 59,0°. I 424 px sono il valore con cui il costruttore traccia la curva di rumore (campo dell'imager di 90° \cite{realsense-taratura-profondita}); 87° × 58° è il campo utile della profondità. In tutto il capitolo la geometria del campo usa gli angoli dichiarati (87° × 58°) e il rumore usa il prodotto $f\,b$ con $f=424\unita{px}$; lo scarto fra le due scelte è il 5 \% sulle grandezze laterali e 2 cm sulla distanza minima a cui si vede il suolo (1,22 m con 59,0° contro 1,24 m). I parametri intrinseci veri si leggono dall'esemplare. \prova{C}
\end{osservazione}

\subsection{Visibilità del suolo}\label{sec:geo-suolo}

\begin{proposizione}[Campo in coordinate del suolo]\label{prop:geo-campo}
Nel modello della Definizione~\ref{def:geo-camera}, con $\beta_\pm\in(0,90^\circ)$, un punto $P=(x,y,z)$ con $z>0$ è nel campo se e solo se
\begin{equation}\label{eq:geo-campo}
h-z\tan\beta_-\;\le\;y\;\le\;h+z\tan\beta_+\qquad\text{e}\qquad |x|\le T_h\,\bigl[(h-y)\sin\theta+z\cos\theta\bigr].\ \text{\prova{D}}
\end{equation}
\end{proposizione}
\begin{proof}
Sia $\beta=\operatorname{atan}\bigl((h-y)/z\bigr)$ la depressione del punto vista dal centro. Dividendo numeratore e denominatore di $\eta$ per $z\cos\theta$ si ha $\eta=(\tan\beta-\tan\theta)/(1+\tan\beta\tan\theta)=\tan(\beta-\theta)$. Quindi $|\eta|\le T_v$ equivale a $\theta-\Phi_v/2\le\beta\le\theta+\Phi_v/2$, cioè $-\tan\beta_+\le(h-y)/z\le\tan\beta_-$, che è la prima condizione; in questo intervallo $d>0$. La seconda è $|\xi|\le T_h$ con $d$ della~\eqref{eq:geo-proiezione}.
\end{proof}

I due bordi verticali del campo sono dunque due piani che dipendono solo da $z$ e non dall'azimut. Con la posa di riferimento e $\Phi_v=58^\circ$ si ha $\beta_-=44^\circ$ e $\beta_+=14^\circ$ (figura~\ref{fig:geo-sezione}).

\begin{corollario}[Suolo]\label{cor:geo-suolo}
Il suolo ($y=0$) è nel campo per
\begin{equation}\label{eq:geo-zmin}
z\;\ge\;z_{\min}=\frac{h}{\tan\beta_-}=\frac{1{,}2}{\tan 44^\circ}=1{,}243\unita{m}\qquad\text{e}\qquad |x|\le T_h\,(z\cos\theta+h\sin\theta)=0{,}917\,z+0{,}295\unita{m}.
\end{equation}
Il confine vicino è la retta $z=z_{\min}$: la zona cieca vicina è una striscia, non un anello. In coordinate polari $\rho_{\min}(\varphi)=z_{\min}/\cos\varphi$: 1,24 m sull'asse, 1,43 m a 30°, 1,71 m a 43,5°; lo spigolo del campo tocca il suolo in $x=\pm1{,}43\unita{m}$, cioè a $\varphi=\pm49{,}1^\circ$ e $\rho=1{,}90\unita{m}$. Con $f=424\unita{px}$ ($\Phi_v=59{,}0^\circ$) si ha $z_{\min}=1{,}221\unita{m}$, e 1,70 m a 44°.
\end{corollario}

\begin{corollario}[Ostacolo basso]\label{cor:geo-basso}
La cima di un ostacolo alto $a<h$ è nel campo solo per $z\ge(h-a)/\tan\beta_-$: 1,087 m per $a=0{,}15\unita{m}$, 1,139 m per $a=0{,}10\unita{m}$. Sotto questa distanza l'ostacolo è interamente fuori campo.
\end{corollario}

\begin{corollario}[Tetto del campo]\label{cor:geo-tetto}
La quota massima visibile è $y_{\max}(z)=h+z\tan\beta_+=1{,}2+0{,}249\,z$ (metri); una quota $y>h$ è vista solo per $z\ge(y-h)/\tan\beta_+$.
\end{corollario}

\begin{figure}[htb]\centering
\begin{tikzpicture}[x=1.95cm,y=1.95cm,font=\scriptsize,>=Stealth]
  \fill[accent!10] (0,1.2)--(1.243,0)--(3.3,0)--(3.3,2.023)--cycle;
  \draw[->] (0,0)--(3.5,0) node[right]{$z$};
  \draw[->] (0,0)--(0,2.2) node[above]{$y$};
  \foreach \t in {0.5,1,...,3}{\draw (\t,0)--(\t,-0.03);}
  \foreach \t in {1,2}{\draw (0,\t)--(-0.03,\t) node[left]{\t};}
  \draw[accent,thick] (0,1.2)--(1.243,0) node[pos=0.45,sloped,below,text=black]{$\beta_-=44^\circ$};
  \draw[accent,thick] (0,1.2)--(3.3,2.023) node[pos=0.2,sloped,above,text=black]{$\beta_+=14^\circ$};
  \draw[muted,dashed] (0,1.2)--(1.3,0.852);
  \node[muted] at (1.42,0.90) {asse};
  \draw[red!70!black,line width=1.6pt] (0,0)--(1.243,0);
  \fill (0,1.2) circle (1.6pt);
  \node[left] at (-0.03,1.2) {$h$};
  \fill[black!65] (1.087,0) rectangle (1.117,0.15);
  \draw[muted] (1.085,0.13)--(0.50,0.13) node[left,text=black]{15 cm};
  \fill[black!65] (1.5,0) rectangle (1.54,0.5);
  \draw[densely dotted] (0,1.2)--(2.571,0);
  \draw[orange!85!black,line width=1.6pt] (1.54,0)--(2.571,0);
  \node[orange!60!black] at (1.86,0.09) {ombra};
  \draw[dashed,muted] (3,0)--(3,2.12);
  \draw[densely dotted] (0,1.70)--(2.006,1.70)--(2.006,0);
  \foreach \z/\y in {1.7/1.624,2.3/1.773,3.0/1.948}{\fill[accent] (\z,\y) circle (1.1pt);}
  \node[above] at (1.62,1.63) {1,62};
  \node[above] at (2.26,1.78) {1,77};
  \node[above left] at (3.0,1.95) {1,95};
  \node[left] at (0,1.72) {1,70};
  \node[below] at (0.62,-0.03) {\color{red!70!black}suolo non visto};
  \node[below] at (1.243,-0.03) {1,24};
  \node[below] at (2.006,-0.03) {2,01};
  \node[below] at (2.571,-0.03) {2,57};
  \node[below] at (3,-0.03) {3};
  \node at (1.65,-0.36) {(a)};
\end{tikzpicture}\hfill
\begin{tikzpicture}[x=1.95cm,y=1.95cm,font=\scriptsize,>=Stealth]
  \fill[accent!10] (0,1.2)--(1.243,0)--(1.8,0)--(1.8,1.649)--cycle;
  \fill[black!12] (1.83,0)--(3.3,0)--(3.3,2.023)--(1.83,1.656)--cycle;
  \draw[->] (0,0)--(3.5,0) node[right]{$z$};
  \draw[->] (0,0)--(0,2.2) node[above]{$y$};
  \foreach \t in {0.5,1,...,3}{\draw (\t,0)--(\t,-0.03);}
  \foreach \t in {1,2}{\draw (0,\t)--(-0.03,\t) node[left]{\t};}
  \draw[accent,thick] (0,1.2)--(1.243,0);
  \draw[accent,thick] (0,1.2)--(3.3,2.023);
  \fill (0,1.2) circle (1.6pt);
  \fill[black!75] (1.8,0) rectangle (1.83,1.7);
  \draw[white,line width=0.8pt] (1.795,1.5)--(1.835,1.5);
  \draw[densely dashed] (0,1.2)--(2.3,1.583);
  \draw[rounded corners=2pt,fill=white] (2.2,0) rectangle (2.4,1.75);
  \draw[red!75!black,line width=1.8pt] (2.3,1.583)--(2.3,1.75);
  \draw[muted] (2.32,1.67)--(2.72,1.42) node[below,text=black,align=center]{in vista sopra\\un pannello di 1,5 m};
  \node[below] at (1.8,-0.03) {$z_o=1{,}8$};
  \node[below] at (2.42,-0.03) {$z_p=2{,}3$};
  \node[above] at (1.72,1.72) {$a=1{,}7$};
  \node at (2.3,0.8) {$s$};
  \node at (2.88,0.5) {ombra};
  \node at (1.65,-0.36) {(b)};
\end{tikzpicture}
\caption{Sezione verticale nel piano $x=0$ (lunghezze in metri). (a)~Campo della camera ($h=1{,}2$ m, $\theta=15^\circ$, $\Phi_v=58^\circ$): il suolo è visto da 1,24 m; un ostacolo di 15 cm è fuori campo sotto 1,09 m; il tetto del campo vale 1,62, 1,77 e 1,95 m a 1,7, 2,3 e 3 m, e una quota di 1,70 m è vista solo da 2,01 m; un ostacolo alto 0,5 m a 1,5 m nasconde il suolo fino a 2,57 m; la linea tratteggiata verticale è la portata di 3 m. (b)~Un pannello alto 1,7 m a 1,8 m nasconde per intero una persona a 2,3 m; sopra un pannello di 1,5 m resterebbe in vista la parte alta del corpo (tratto rosso).}\label{fig:geo-sezione}
\end{figure}

\begin{table}[htb]\centering\small
\caption{Limiti verticali del campo ($h=1{,}2$ m, $\theta=15^\circ$). In alto: quota massima visibile e distanza minima da cui una quota è vista, per $\Phi_v=58^\circ$. In basso: gli stessi limiti per profilo. \prova{C}}\label{tab:geo-verticale}
\begin{tabular}{@{}lcccccc@{}}
\toprule
$z$ (m) & 1,24 & 1,7 & 1,9 & 2,0 & 2,3 & 3,0\\
$y_{\max}$ (m) & 1,51 & 1,62 & 1,67 & 1,70 & 1,77 & 1,95\\
\bottomrule
\end{tabular}\qquad
\begin{tabular}{@{}lccccc@{}}
\toprule
quota (m) & 1,60 & 1,70 & 1,80 & 1,90 & 2,00\\
vista da $z\ge$ (m) & 1,60 & 2,01 & 2,41 & 2,81 & 3,21\\
\bottomrule
\end{tabular}

\smallskip
\begin{tabular}{@{}lcccc@{}}
\toprule
Profilo ($\Phi_v$) & $z_{\min}$ del suolo & $z$ minima per $a=0{,}15$ m & $\tan\beta_+$ & quota 1,80 m vista da\\
\midrule
D455, D435i (58°) & 1,243 m & 1,087 m & 0,249 & 2,41 m\\
Gemini 335 (60°) & 1,200 m & 1,050 m & 0,268 & 2,24 m\\
Gemini 335 (65°) & 1,100 m & 0,962 m & 0,315 & 1,90 m\\
\bottomrule
\end{tabular}
\end{table}

Ne segue che entro 3 m la colonna da 0 a 2 m di quota non è mai osservata per intero: una testa a 1,70 m non è vista sotto 2,01 m, e un ostacolo sospeso a 1,80 m non è visto sotto 2,41 m (tabella~\ref{tab:geo-verticale}). Lo stato «libero» della mappa (\S\ref{sec:map}) è quindi un'affermazione sul suolo e sul volume fino a $y_{\max}(z)$, non sulla colonna intera. La distanza $z_{\min}$ varia da 1,18 a 1,31 m per $\Phi_v=58^\circ\pm3^\circ$ e da 1,20 a 1,29 m per $\theta=15^\circ\pm1^\circ$. I valori sono verificati per calcolo e con un lancio di raggi su 848 × 480 pixel (minimo di $z$ al suolo 1,245 m ai centri dei pixel, uguale su tutte le colonne) \prova{C}. Valgono con suolo piano, posa nota e assenza di distorsione: su una rampa, o con la camera più inclinata, i limiti si spostano secondo il \S\ref{sec:geo-posa}.

\begin{proposizione}[Frazione del disco coperta]\label{prop:geo-disco}
Sia $D$ il disco di raggio 3 m centrato sulla proiezione della camera sul suolo, di area 28,27 m². (i)~Approssimando la regione vista con il settore d'anello di apertura $\Phi_h$ fra $\rho=z_{\min}$ e 3 m, l'area è $\tfrac12\Phi_h(3^2-z_{\min}^2)=5{,}66\unita{m^2}$, il 20,0 \% di $D$. (ii)~Con la regione esatta del Corollario~\ref{cor:geo-suolo} intersecata con $D$, l'area è 5,79 m², il 20,5 \% (integrazione numerica su una griglia di 2 mm). In entrambi i conti quattro quinti del disco non sono osservati dalla camera. \prova{C}
\end{proposizione}

\subsection{Ombra di un ostacolo}\label{sec:geo-ombra}

\begin{proposizione}[Ombra sul suolo]\label{prop:geo-ombra}
Un ostacolo opaco alto $a$, a distanza $\rho_o$ lungo una direzione $\varphi$, nasconde il suolo in quella direzione per
\begin{equation}\label{eq:geo-ombra}
\rho_o<\rho<\rho_o\,\frac{h}{h-a}\quad\text{se } a<h,\qquad\text{e per ogni }\rho>\rho_o\ \text{se } a\ge h.\ \text{\prova{D}}
\end{equation}
\end{proposizione}
\begin{proof}
Nel piano verticale per il centro e la direzione $\varphi$, il raggio che sfiora la cima $(\rho_o,a)$ ha equazione $y=h-(h-a)\rho/\rho_o$ e tocca il suolo in $\rho=\rho_o h/(h-a)$; i punti del suolo più vicini hanno il raggio intercettato dall'ostacolo. Se $a\ge h$ il raggio non scende.
\end{proof}

Esempi: $a=0{,}3$ m a 1,7 m dà ombra fino a 2,27 m; $a=0{,}5$ m a 1,5 m fino a 2,57 m; $a=1{,}0$ m a 1,5 m fino a 9,0 m. In pianta l'ombra è il cuneo compreso fra gli azimut dei due bordi dell'ostacolo (figura~\ref{fig:geo-pianta}); per un ostacolo più basso della camera il cuneo è troncato alla distanza della~\eqref{eq:geo-ombra} e non arriva al bordo della portata. Il risultato vale in ottica geometrica, per un ostacolo opaco su suolo piano.

\begin{proposizione}[Altezza che nasconde una persona]\label{prop:geo-amin}
Un pannello opaco a $z_o$, alto $a$, nasconde per intero dall'alto una persona di statura $s$ a $z_p>z_o$ se e solo se
\begin{equation}\label{eq:geo-amin}
a\;\ge\;a_{\min}=h+z_o\,\min\!\Bigl(\frac{s-h}{z_p},\;\tan\beta_+\Bigr).\ \text{\prova{D}}
\end{equation}
Viceversa, sopra un pannello alto $a$ resta visibile a $z_p$ la fascia di quota fra $h+(a-h)\,z_p/z_o$ e $\min\bigl(s,\,y_{\max}(z_p)\bigr)$.
\end{proposizione}
\begin{proof}
Il raggio dal centro alla sommità $(z_p,s)$ passa a $z_o$ alla quota $h+(s-h)z_o/z_p$; se il pannello la raggiunge, ogni punto più basso della persona è nascosto, perché il suo raggio passa più in basso. Se $s>y_{\max}(z_p)$, la parte sopra il tetto del campo non è comunque vista e basta raggiungere il tetto a $z_o$, cioè $h+z_o\tan\beta_+$.
\end{proof}

\begin{table}[htb]\centering\small
\caption{Altezza minima $a_{\min}$ (m) del pannello che nasconde per intero una persona a $z_p=2{,}3$ m, secondo la~\eqref{eq:geo-amin}, per statura $s$ e distanza $z_o$ del pannello. \prova{C}}\label{tab:geo-amin}
\begin{tabular}{@{}lcccccc@{}}
\toprule
$s$ \textbackslash{} $z_o$ & 1,15 m & 1,5 m & 1,7 m & 1,8 m & 1,9 m & 2,0 m\\
\midrule
1,60 m & 1,40 & 1,46 & 1,50 & 1,51 & 1,53 & 1,55\\
1,70 m & 1,45 & 1,53 & 1,57 & 1,59 & 1,61 & 1,63\\
1,75 m & 1,48 & 1,56 & 1,61 & 1,63 & 1,65 & 1,68\\
$\ge$ 1,77 m (tetto del campo) & 1,49 & 1,57 & 1,62 & 1,65 & 1,67 & 1,70\\
\bottomrule
\end{tabular}
\end{table}

Un pannello alto 1,5 m a 1,7--1,9 m dalla camera non nasconde una persona a 2,3 m: resta in vista ciò che supera 1,61 m (pannello a 1,7 m) o 1,56 m (a 1,9 m), fino al tetto di 1,77 m. A quelle distanze servono fra 1,61 m (statura 1,75 m) e 1,67 m (qualunque statura). Un pannello di 1,7 m nasconde ogni statura fino a $z_o=2{,}0$ m, a margine nullo a 2,0 m e con 3--8 cm a 1,9--1,7 m; un pannello di 1,8 m dà da 10 a 18 cm di margine contro l'errore di posa e la tolleranza del campo, che da sola sposta il tetto di circa ±5 cm a 1,9 m. Il pannello deve inoltre coprire in larghezza, come precisa la proposizione seguente.

\begin{proposizione}[Durata dell'occlusione totale]\label{prop:geo-durata}
Un pannello largo $L$ a $z_o$, parallelo al piano dell'immagine, proietta alla distanza $z_p$ un'ombra larga $L\,z_p/z_o$. Un corpo largo $w_c$ che la attraversa di traverso a velocità $v$ è nascosto per intero per
\begin{equation}\label{eq:geo-durata}
T_{\mathrm{tot}}=\frac{L\,z_p/z_o-w_c}{v}.\ \text{\prova{D}}
\end{equation}
\end{proposizione}
\begin{proof}
In pianta i raggi per i due bordi del pannello sono rette per l'origine: a $z_p$ distano $L z_p/z_o$ (triangoli simili). Il corpo è tutto nascosto finché entrambi i suoi estremi stanno nell'ombra, cioè lungo un tratto $L z_p/z_o-w_c$.
\end{proof}

A passo normale (1,0--1,4 m/s) un pannello da 1 m dà 0,54--0,90 s di occlusione totale, uno da 2 m dà 1,4--2,3 s (tabella~\ref{tab:geo-durata}).

\begin{figure}[htb]\centering
\begin{tikzpicture}[x=1.3cm,y=1.3cm,font=\scriptsize,>=Stealth]
  \draw[muted] (-3,0) arc (180:0:3);
  \draw[muted] (-3,0)--(3,0);
  \fill[accent!13] (-1.435,1.243)--(1.435,1.243)--(2.182,2.058) arc (43.3:136.7:3)--cycle;
  \begin{scope}
    \clip (-1.435,1.243)--(1.435,1.243)--(2.182,2.058) arc (43.3:136.7:3)--cycle;
    \foreach \r in {1.5,2,2.5}{\draw[white,line width=0.4pt] (0,0) circle (\r);}
    \foreach \g in {50,60,...,130}{\draw[white,line width=0.4pt] (0,0)--(\g:3);}
  \end{scope}
  \draw[accent] (1.435,1.243)--(2.182,2.058);
  \draw[accent] (-1.435,1.243)--(-2.182,2.058);
  \draw[dashed,muted] (46.5:1.243) arc (46.5:133.5:1.243);
  \draw[red!70!black,thick] (-1.435,1.243)--(1.435,1.243);
  \begin{scope}
    \clip (-3,0) arc (180:0:3)--cycle;
    \draw[orange!40,line width=0.78cm,line join=round] (-0.5,1.8)--(0.5,1.8)--(0.9,3.24)--(-0.9,3.24)--cycle;
    \fill[black!35] (-0.5,1.8)--(0.5,1.8)--(0.9,3.24)--(-0.9,3.24)--cycle;
    \fill[red!65] (-0.5,1.8)--(-0.9,3.24)--(-1.03,3.24)--cycle;
  \end{scope}
  \draw[densely dotted] (0,0)--(-0.5,1.8);
  \draw[densely dotted] (0,0)--(0.5,1.8);
  \draw[line width=1.8pt] (-0.5,1.8)--(0.5,1.8);
  \draw[white,line width=0.5pt] (-0.639,2.3)--(0.639,2.3);
  \draw[fill=white] (0.12,2.3) ellipse [x radius=0.225,y radius=0.11];
  \fill (0,0) circle (1.6pt) node[below]{camera};
  \draw[muted] (0.62,2.3)--(1.75,2.62) node[right,text=black]{ombra larga $L\,z_p/z_o$; corpo largo $w_c$};
  \draw[muted] (0.5,1.8)--(1.75,1.62) node[right,text=black]{pannello largo $L$ a $z_o$};
  \draw[muted] (0.93,2.05)--(1.75,2.12) node[right,text=black]{fascia di guardia, larga $w$};
  \draw[muted] (-0.74,2.6)--(-1.75,2.62) node[left,text=black]{banda senza dato $\Delta x$};
  \draw[muted] (-1.0,1.243)--(-1.75,1.5) node[left,text=black]{$z=z_{\min}$};
  \draw[muted] (-0.7,1.03)--(-1.75,0.9) node[left,text=black]{$\rho=z_{\min}$ (anello)};
  \node[left] at (-2.6,2.2) {portata 3 m};
  \node at (0,2.66) {ombra};
\end{tikzpicture}
\caption{Pianta. In azzurro la parte del suolo nel campo entro 3 m, con la griglia polare accennata: il confine vicino è la retta $z=z_{\min}$ (in rosso), non l'arco $\rho=z_{\min}$ (tratteggiato). Un pannello largo $L$ a $z_o$ proietta il cuneo d'ombra (grigio), largo $L\,z_p/z_o$ alla distanza $z_p$; attorno al cuneo, la fascia di guardia (arancione), la cui larghezza $g$ è fissata nel \S\ref{sec:occ}; sul lato sinistro del bordo, la banda senza disparità della Proposizione~\ref{prop:geo-banda} (disegno non in scala per $g$ e $\Delta x$).}\label{fig:geo-pianta}
\end{figure}

\begin{proposizione}[Banda senza dato accanto a un bordo]\label{prop:geo-banda}
In una coppia stereo con riferimento sull'imager sinistro e base $b$, accanto a un bordo verticale a profondità $d_{\mathrm{bordo}}$, dal lato sinistro, lo sfondo a profondità $d_{\mathrm{sfondo}}$ non ha disparità su una banda larga
\begin{equation}\label{eq:geo-banda}
\Delta x=b\,\frac{d_{\mathrm{sfondo}}-d_{\mathrm{bordo}}}{d_{\mathrm{bordo}}}\ \ \text{(sullo sfondo)},\qquad \Delta u=f\,b\Bigl(\frac1{d_{\mathrm{bordo}}}-\frac1{d_{\mathrm{sfondo}}}\Bigr)\ \ \text{(in pixel)}.\ \text{\prova{D}}
\end{equation}
\end{proposizione}
\begin{proof}
In pianta, con l'imager sinistro in $(0,0)$ e il destro in $(b,0)$, l'ostacolo occupi $x\ge x_e$ alla profondità $d_{\mathrm{bordo}}$. Un punto $(x,d_{\mathrm{sfondo}})$ è visto dal sinistro se $x<x_e\,d_{\mathrm{sfondo}}/d_{\mathrm{bordo}}$; dal destro se il segmento da $(b,0)$ passa a sinistra del bordo, $b+(x-b)\,d_{\mathrm{bordo}}/d_{\mathrm{sfondo}}<x_e$, cioè $x<x_e\,d_{\mathrm{sfondo}}/d_{\mathrm{bordo}}-b\,(d_{\mathrm{sfondo}}-d_{\mathrm{bordo}})/d_{\mathrm{bordo}}$. La differenza dei due limiti è $\Delta x$; in pixel $\Delta u=f\,\Delta x/d_{\mathrm{sfondo}}$.
\end{proof}

\begin{table}[htb]\small
\begin{minipage}[t]{0.52\textwidth}\centering
\caption{Durata dell'occlusione totale $T_{\mathrm{tot}}$ (s) secondo la~\eqref{eq:geo-durata}, con $z_p=2{,}3$ m e $w_c=0{,}45$ m. \prova{C}}\label{tab:geo-durata}
\begin{tabular}{@{}lccccc@{}}
\toprule
 & ombra & tratto & \multicolumn{3}{c}{velocità $v$ (m/s)}\\
$L$; $z_o$ & (m) & nascosto (m) & 1,4 & 1,0 & 0,5\\
\midrule
1 m; 1,7 m & 1,35 & 0,90 & 0,64 & 0,90 & 1,81\\
1 m; 1,9 m & 1,21 & 0,76 & 0,54 & 0,76 & 1,52\\
2 m; 1,7 m & 2,71 & 2,26 & 1,61 & 2,26 & 4,51\\
2 m; 1,9 m & 2,42 & 1,97 & 1,41 & 1,97 & 3,94\\
\bottomrule
\end{tabular}
\end{minipage}\hfill
\begin{minipage}[t]{0.45\textwidth}\centering
\caption{Banda senza disparità secondo la~\eqref{eq:geo-banda}, sullo sfondo (cm) e in pixel. \prova{C}}\label{tab:geo-banda}
\begin{tabular}{@{}lccc@{}}
\toprule
bordo; sfondo & D455 & D435i & Gemini 335\\
\midrule
1 m; 3 m & 19,0; 27 & 10,0; 14 & 10,0; 14\\
1,5 m; 3 m & 9,5; 13 & 5,0; 7 & 5,0; 7\\
2 m; 3 m & 4,8; 7 & 2,5; 4 & 2,5; 3\\
\bottomrule
\end{tabular}
\end{minipage}
\end{table}

La banda è sistematica e sta da un lato solo; i bordi orizzontali non ne hanno. Per la stessa ragione manca il dato su una striscia di $f\,b/d$ pixel al bordo sinistro dell'immagine (20 px a 2 m con $b=95$ mm, 11 px con 50 mm). Chi ricompare dal lato della banda è misurato $\Delta x/v$ più tardi: 0,16 s a 1,2 m/s con una banda di 19 cm. Il risultato è dimostrato nella geometria epipolare ideale e verificato numericamente (0,1900 m). Una misura di terzi su una camera della stessa famiglia descrive un'ombra che peggiora sul bordo sinistro del bersaglio quando questo ruota, cioè dal lato previsto, e riporta un rumore laterale al 90° percentile di 4,3 cm sul bordo sinistro e di 4,1 cm sul destro \cite{ahn2019}: le due cifre sono vicine e non misurano l'ampiezza della banda; l'ampiezza della banda su un esemplare reale non è verificata sperimentalmente \prova{N}, e i filtri di riempimento del costruttore, se attivi, la mascherano copiando lo sfondo \cite{realsense-filtri-post}.

\subsection{Rumore di profondità dello stereo}\label{sec:geo-rumore}

\begin{proposizione}[Rumore di profondità]\label{prop:geo-rumore}
Se la disparità $\delta=f\,b/d$ è misurata con un errore di scarto tipo $\sigma_d$ (pixel), al primo ordine
\begin{equation}\label{eq:geo-sigmaz}
\sigma_z(d)=\frac{d^2}{f\,b}\,\sigma_d=k\,d^2,\qquad k=\frac{\sigma_d}{f\,b}.\ \text{\prova{D}}
\end{equation}
\end{proposizione}
\begin{proof}
Da $d=f\,b/\delta$ segue $\partial d/\partial\delta=-f\,b/\delta^2=-d^2/(f\,b)$; la tesi è la propagazione lineare dell'errore.
\end{proof}

Il costruttore dichiara che l'errore di profondità cresce con il quadrato della distanza e traccia la curva teorica con $\sigma_d=0{,}08$ px \cite{realsense-taratura-profondita}, indicando meno di 0,1 px per una camera ben tarata \cite{realsense-bkm-taratura}: è il valore qui detto «nominale». Il valore «limite» si ricava dal massimo dichiarato: scarto quadratico medio non superiore al 2 \% fino a 4 m per la D455 e fino a 2 m per la D435i, sull'80 \% della regione utile e a risoluzione HD \cite{realsense-d400-accuratezza} ($k=0{,}02\cdot4/16=0{,}005\unita{m^{-1}}$ e $k=0{,}010\unita{m^{-1}}$; per la D455 la pagina del prodotto indica un errore di profondità inferiore al 2 \% a 4 m \cite{realsense-d455}), e all'1,5 \% a 2 m per la Gemini 335, a 1280 × 800 e sul 90 \% × 90 \% della regione \cite{orbbec-gemini-335} ($k=0{,}0075\unita{m^{-1}}$); queste condizioni non sono quelle d'uso del sistema (848 × 480 o 424 × 240); per quest'ultima il nominale ($k=0{,}004\unita{m^{-1}}$) è letto da un grafico della scheda.

\begin{table}[htb]\centering\small
\caption{Rumore di profondità $\sigma_z$ (cm) secondo la~\eqref{eq:geo-sigmaz}, e passo di quantizzazione (mm) per una disparità quantizzata a 1/32 px. I valori nominali e limite sono calcolati da dati di scheda. \prova{C}}\label{tab:geo-rumore}
\begin{tabular}{@{}llccccc@{}}
\toprule
Profilo ($f\,b$ in px·m) & & 1 m & 2 m & 3 m & 4 m & $\sigma_z=5$ cm a\\
\midrule
D455 (40,28) & nominale & 0,20 & 0,79 & 1,79 & 3,18 & 5,0 m\\
 & limite & 0,50 & 2,00 & 4,50 & 8,00 & 3,2 m\\
D435i (21,20) & nominale & 0,38 & 1,51 & 3,40 & 6,04 & 3,6 m\\
 & limite & 1,00 & 4,00 & 9,00 & 16,0 & 2,2 m\\
Gemini 335 (20,54) & nominale & 0,40 & 1,60 & 3,60 & 6,40 & 3,5 m\\
 & limite & 0,75 & 3,00 & 6,75 & 12,0 & 2,6 m\\
\midrule
Passo di quantizzazione (mm) & D455 & 0,8 & 3,1 & 7,0 & 12,4 & \\
 & D435i & 1,5 & 5,9 & 13,3 & 23,6 & \\
\bottomrule
\end{tabular}
\end{table}

\emph{Quantizzazione.} Un passo di disparità $\epsilon_\delta$ dà un passo di profondità $d^2\epsilon_\delta/(f\,b)$ e uno scarto tipo pari a $1/\sqrt{12}$ di quel passo (2,0 mm a 3 m per la D455, 3,8 mm per la D435i): è inferiore al rumore nominale. Il valore $\epsilon_\delta=1/32$ px è un'ipotesi, non confermata da una scheda.

\emph{Confronto con una misura indipendente.} Su una D435 a 848 × 480, con un bersaglio bianco frontale fra 0,4 e 2,1 m, è stato misurato $\sigma_z=0{,}00106+0{,}00073\,d+0{,}00395\,d^2$ (metri) \cite{ahn2019}: 1,83 cm a 2 m, cioè 1,21 volte il nominale (0,097 px equivalenti). A 3 m la stessa formula dà 3,9 cm, ma è un'estrapolazione fuori dall'intervallo misurato.

\emph{Limiti di validità.} I valori dichiarati valgono su un bersaglio chiaro e frontale, nella parte centrale del campo. Il rumore è correlato nello spazio, quindi la media su $N$ pixel non lo riduce come $\sqrt N$. Dal punto di montaggio il suolo è visto con un'incidenza $\operatorname{atan}(\rho/h)$ dalla normale: 46° a 1,24 m, 60° a $\rho=h\tan 60^\circ=2{,}08$ m, 68° a 3 m; oltre 60° non esiste una misura del rumore sul suolo. Il comportamento sul suolo radente, su tessuti scuri e su superfici lucide non è verificato sperimentalmente \prova{N}.

\emph{Risoluzione e rumore laterali.} Un pixel copre $d/f$: con il campo dichiarato di 87° ($f=446{,}8$ px) sono 4,5 mm a 2 m e 6,7 mm a 3 m. Un oggetto che debba coprire 5 pixel di lato è quindi osservabile a 3 m solo se è largo almeno 3,4 cm (6,7 cm a 6 m), e un oggetto di 1 cm solo entro 0,89 m; con $f=424$ px i valori crescono del 5 \% (3,5 cm). Il criterio dei 5 pixel è una convenzione. La sfocatura da moto vale la velocità per il tempo d'esposizione: 12 mm a 1,5 m/s con 8 ms, 45 mm con 30 ms. \prova{C}

\subsection{Effetto di un errore di posa}\label{sec:geo-posa}

\begin{proposizione}[Errore di quota]\label{prop:geo-posa}
Se l'inclinazione creduta differisce dalla vera di $\delta\theta$, un punto del suolo a distanza $\rho$ sull'asse è ricostruito alla quota
\begin{equation}\label{eq:geo-posa}
\Delta y=\rho\sin\delta\theta+h\,(1-\cos\delta\theta)\;\approx\;\rho\,\delta\theta.\ \text{\prova{D}}
\end{equation}
\end{proposizione}
\begin{proof}
Rispetto al centro il punto è $(\rho,-h)$ nel piano verticale; una rotazione di $\delta\theta$ attorno all'asse trasversale porta la sua quota a $-h\cos\delta\theta+\rho\sin\delta\theta$, e la differenza da $-h$ è $\Delta y$.
\end{proof}

Con $\delta\theta=1^\circ$ l'errore è 3,5 cm a 2 m, 4,9 cm a 2,8 m e 5,3 cm a 3 m (5,2 cm con la sola approssimazione lineare). Sia $y_s=0{,}07$ m la soglia di quota sopra la quale un ritorno è «occupato». Con un errore che alza il suolo, il suolo supera la soglia oltre $\rho\approx y_s/\delta\theta$: con 2° il suolo risulta occupato da 1,99 m in poi. Con un errore che lo abbassa, un ostacolo alto $a$ scende sotto la soglia, ed è preso per suolo, quando
\begin{equation}\label{eq:geo-suolo-falso}
\rho\,\delta\theta>a-y_s :
\end{equation}
per $a=0{,}10$ m bastano 0,86° a 2 m, 0,67° a 2,55 m e 0,57° a 3 m; per $a=0{,}15$ m servono 2,3° a 2 m.

\emph{Riscontro in simulazione} (catena di elaborazione completa su profondità simulata dalle schede tecniche; il primo caso su una scena senza rumore, il secondo con il profilo nominale della D435i). Un ostacolo di 10 cm a 2,55 m risulta «occupato» in 7 celle su 7 con errore nullo e «libero» in 7 su 7 con $+1^\circ$ (soglia prevista: 0,67°); uno di 12 cm cambia stato fra 1,0° e 1,5° (previsione 1,12°), uno di 15 cm fra 1,5° e 2,0° (previsione 1,78°). Per una scatola di 15 cm a 2,0--2,15 m, i fotogrammi con almeno una cella dell'impronta «libero» sono 0 su 10 a 0°, 1 su 10 a 1°, 2 su 10 a 2°, 9 su 10 a 2,5° (intervallo di Wilson al 95 \%: 0,60--0,98) e 10 su 10 a 3° (0,72--1); la soglia prevista è 2,1--2,3°. Le soglie previste dalla~\eqref{eq:geo-suolo-falso} sono compatibili con i cambi di stato osservati; le numerosità ($N=7$ celle, $N=10$ fotogrammi) sono piccole e non bastano a misurare una soglia \prova{S}. La proprietà «mai libero su un ostacolo» (\S\ref{sec:map}) vale quindi solo a posa nota entro questi margini. I requisiti sulla posa per un sensore in moto, che usano la~\eqref{eq:geo-posa}, sono analizzati nel \S\ref{sec:lim}.

\subsection{Sensori a tempo di volo multizona e LiDAR a scansione su un piano}\label{sec:geo-altri}

Sono stati considerati, come complemento laterale, i sensori a tempo di volo multizona. Il VL53L7CX ha un campo di 60° × 60° su una matrice di 8 × 8 zone, con portata di 2,0 m al buio e 0,5 m con 5 klx (chilolux) di luce ambiente a 15 Hz \cite{st-vl53l7cx}; il VL53L8CX ha un campo di 45° × 45°, con 1,55 m a 5 klx \cite{st-vl53l8cx}. Una zona sottende 7,5° (5,625° nel secondo), e a distanza $D$ il suo lato è $2D\tan(\text{passo}/2)$ (tabella~\ref{tab:geo-tof}).

\begin{table}[htb]\small
\begin{minipage}[c]{0.57\textwidth}\centering
\begin{tabular}{@{}lcc@{}}
\toprule
 & VL53L7CX & VL53L8CX\\
\midrule
Lato di una zona a 2 m; diagonale & 26 cm; 37 cm & 20 cm; 28 cm\\
Campo intero a 2 m & 2,31 m & 1,66 m\\
Tre sensori affiancati & 180° × 60° & 135° × 45°\\
Sensori necessari per 360° & 6 & 8\\
Portata: buio / 5 klx & 2,0 m / 0,5 m & 4,0 m / 1,55 m\\
Persona larga 0,5 m a 4 m & 0,95 zone & 1,27 zone\\
\bottomrule
\end{tabular}
\caption{Sensori a tempo di volo multizona: geometria delle zone e copertura. Campi e portate (matrice 8 × 8 a 15 Hz, bersaglio bianco) dalle schede \cite{st-vl53l7cx,st-vl53l8cx}; il resto per calcolo. \prova{C}}\label{tab:geo-tof}
\end{minipage}\hfill
\begin{minipage}[c]{0.40\textwidth}\centering
\begin{tikzpicture}[x=0.56cm,y=0.56cm,font=\scriptsize]
  \draw[muted] (0,0) circle (3);
  \fill[accent!18] (46.5:1.243) arc (46.5:133.5:1.243)--(133.5:3) arc (133.5:46.5:3)--cycle;
  \foreach \g in {0,60,120}{\fill[orange!45,opacity=0.55] (0,0)--(\g:2) arc (\g:\g+60:2)--cycle; \draw[orange!70!black] (0,0)--(\g:2) arc (\g:\g+60:2)--cycle;}
  \foreach \g in {0,60,120}{\fill[orange!80!black] (0,0)--(\g:0.5) arc (\g:\g+60:0.5)--cycle;}
  \fill (0,0) circle (1.2pt);
  \node at (0,2.35) {campo};
  \node at (0,-1.5) {180° scoperti};
  \node[orange!50!black] at (-1.25,0.75) {2 m};
  \begin{scope}[shift={(3.6,-2.0)},x=1.0cm,y=1.0cm]
    \fill[black!22] (0.905,0) rectangle (1.405,1.655);
    \draw[step=0.28875,muted] (0,0) grid (2.31,2.31);
    \draw (0,0) rectangle (2.31,2.31);
    \node[above] at (1.155,2.31) {8 × 8 a 2 m: 2,31 m};
  \end{scope}
\end{tikzpicture}
\captionof{figure}{Copertura in pianta entro il disco di 3 m: settore della camera (azzurro, schematico), tre sensori a tempo di volo da 60° con portata 2 m al buio e 0,5 m a 5 klx (arancione). A destra, la matrice a 2 m (zone di 26 cm) con la sagoma di una persona larga 0,5 m.}\label{fig:geo-tof}
\end{minipage}
\end{table}

Tre sensori affiancati coprono dunque 180°, non 360° (figura~\ref{fig:geo-tof}). Per regola di progetto il canale a tempo di volo è solo un indicatore laterale di prossimità e non genera mai lo stato «libero», per tre ragioni geometriche. (a)~Una zona restituisce una sola distanza per un cono largo 26 cm a 2 m, contro celle della mappa di 1° × 0,1 m (3,5 cm × 10 cm a 2 m): non distingue il suolo visto alla quota attesa da un ostacolo basso. (b)~Da 1,2 m di altezza la distanza obliqua dal suolo è almeno 1,20 m e vale 2,33 m a $\rho=2$ m: è oltre la portata con luce ambiente (0,5 m a 5 klx su un bersaglio bianco). (c)~Una zona senza bersaglio valido significa «nessun dato», non «vuoto»: la guida del costruttore assegna al primo campione dopo un'assenza di bersaglio uno stato con fiducia inferiore al 50 \% \cite{st-um3038}. Come copertura, tre sensori da 60° con portata 2 m aggiungono $3\cdot\tfrac12\cdot\tfrac{\pi}{3}\cdot2^2=6{,}28\unita{m^2}$, il 22,2 \% del disco di 3 m, al buio; con portata 0,5 m l'1,4 \%. L'interferenza ottica fra il proiettore della camera e i sensori a tempo di volo non è verificata sperimentalmente \prova{N}.

Un LiDAR (\emph{light detection and ranging}) a scansione su un piano, con passo angolare di 0,72° \cite{slamtec-rplidar-c1} o di 0,8° (valore derivato da 4500 misure al secondo a 10 giri al secondo \cite{ldrobot-ld19}), vede solo ciò che interseca il piano alla quota di montaggio: non gradini, ostacoli bassi, tavoli od oggetti sospesi. Il passo fra due punti a 3 m è 3,8--4,2 cm: 12--13 punti su una persona larga 0,5 m a 3 m, meno di 2 su una gamba di 12 cm a 5 m. Ogni punto ha inoltre il proprio istante: a 30°/s il sensore ruota di 3° dentro un giro. \prova{C}

\begin{riquadro}{Condizioni di validità del capitolo}
Foro stenopeico senza distorsione; suolo piano; posa $(h,\theta)$ nota e camera ferma (per il sensore in moto valgono i requisiti del \S\ref{sec:lim}); ostacoli opachi che restituiscono profondità (vetro, specchi, nero assorbente e pavimenti lucidi sono fuori dal modello); campi e basi come da scheda, con ±3° di tolleranza; rumore gaussiano in disparità. La fascia di guardia attorno all'ombra è trattata nel \S\ref{sec:occ}, il bilancio di latenza nel \S\ref{sec:lim}. Tutte le cifre sono conti o simulazioni: la loro validità su un sensore reale non è verificata sperimentalmente.
\end{riquadro}

\section{Rappresentazione dello spazio e sicurezza del «libero»}\label{sec:map}

Il capitolo definisce la mappa del suolo e la regola con cui una cella è dichiarata «libero», enuncia che cosa quella dichiarazione garantisce e in quale modello, elenca i casi in cui la garanzia cade e misura quanto «libero» vero la regola rinuncia a dichiarare. Tutti i risultati numerici sono dimostrazioni, conti o simulazioni: nessuno è stato ottenuto su un sensore fisico.

\subsection{Mappa polare del suolo a quattro stati}\label{sec:map-stati}

Il suolo davanti alla camera è diviso in celle polari $c=(\varphi,\rho)$ di $1^\circ\times0{,}1\unita{m}$ fino alla portata di $3\unita{m}$: con un campo orizzontale di $87^\circ$ sono $N_c=87\cdot30=2\,610$ celle \prova{C}. Sia $X$ la regione coperta. La camera è ferma ($h=1{,}2\unita{m}$, $\theta=15^\circ$, \S\ref{sec:geo}), quindi la corrispondenza fra pixel e celle si calcola una sola volta: per ogni pixel $p$ il cui raggio tocca il suolo entro $3\unita{m}$ si tabulano la cella toccata $c(p)$ e la profondità attesa del suolo $z_0(p)$. L'\emph{impronta} di una cella è $I(c)=\{p: c(p)=c\}$, con $n(c)=|I(c)|$ pixel; nella camera usata in simulazione ($320\times240$ pixel) una cella ha circa 12 pixel a $2\unita{m}$ e 4 a $3\unita{m}$. A ogni fotogramma ogni pixel valido è assegnato alla sua cella con queste tabelle; nessun raggio viene percorso.

\begin{definizione}[Stati della cella]\label{def:map-stati}
Sia $D_t$ l'immagine di profondità del fotogramma $t$. Una cella $c\in X$ è:
\begin{enumerate}[label=(\roman*)]
\item \emph{occupato osservato}, se almeno 3 pixel validi di $D_t$ hanno il punto ricostruito dentro $c$ a quota compresa fra $0{,}07$ e $2{,}0\unita{m}$;
\item \emph{in ombra}, cioè $c\in\Omb(t)$, se non è occupata e sta oltre la prima cella occupata del suo azimut;
\item \emph{libero}, cioè $c\in\Oss(t)$, se soddisfa la regola della Definizione~\ref{def:map-libero};
\item \emph{non osservato}, in ogni altro caso.
\end{enumerate}
I quattro stati formano una partizione di $X$ e sono ricalcolati a ogni fotogramma.
\end{definizione}

I due stati senza osservazione diretta hanno significati diversi. «In ombra» dice che il fotogramma contiene un ostacolo e che la cella sta dietro di esso lungo la linea di vista: il contenuto non è noto, la causa sì. «Non osservato» dice che il fotogramma non porta alcuna informazione sufficiente sulla cella: fuori campo, oltre la portata, pixel non validi, suolo letto da troppo pochi pixel, oppure cella adiacente a una cella non letta. La distinzione serve in \S\ref{sec:occ}: la «possibile presenza» nasce sull'ombra in cui è entrato un oggetto seguito, non su tutto il non osservato. In nessun caso «non osservato» o «in ombra» diventano «libero».

\begin{figure}[H]\centering
\begin{tikzpicture}[scale=1.35,font=\footnotesize]
  % non osservato: tutto il settore
  \fill[muted!22] (0,0) -- (46.5:3) arc (46.5:133.5:3) -- cycle;
  % libero: centro da 1,3 m, lati da 1,7 m
  \fill[accent!22] (61.5:1.3) arc (61.5:118.5:1.3) -- (118.5:3) arc (118.5:61.5:3) -- cycle;
  \fill[accent!22] (46.5:1.7) arc (46.5:61.5:1.7) -- (61.5:3) arc (61.5:46.5:3) -- cycle;
  \fill[accent!22] (118.5:1.7) arc (118.5:133.5:1.7) -- (133.5:3) arc (133.5:118.5:3) -- cycle;
  % cella di guardia attorno a occupato e ombra (torna «non osservato»)
  \fill[muted!22] (91:1.88) arc (91:112:1.88) -- (112:3) arc (112:91:3) -- cycle;
  \draw[muted,densely dotted] (91:3) -- (91:1.88) arc (91:112:1.88) -- (112:3);
  % ombra
  \fill[white] (94:2.12) arc (94:109:2.12) -- (109:3) arc (109:94:3) -- cycle;
  \fill[pattern=north east lines,pattern color=muted] (94:2.12) arc (94:109:2.12) -- (109:3) arc (109:94:3) -- cycle;
  % occupato
  \fill[accent] (94:2.0) arc (94:109:2.0) -- (109:2.12) arc (109:94:2.12) -- cycle;
  % reticolo
  \foreach \r in {1,2,3} {\draw[muted!70,very thin] (46.5:\r) arc (46.5:133.5:\r); \node[muted,anchor=west] at (45:\r) {$\r\unita{m}$};}
  \draw[muted] (0,0) -- (46.5:3) arc (46.5:133.5:3) -- cycle;
  \draw[muted,very thin,dashed] (94:0) -- (94:2.0) (109:0) -- (109:2.0);
  \fill (0,0) circle (1.4pt) node[below] {camera};
  \draw[-{Stealth}] (0,0) -- (90:0.55) node[right] {$z$};
  \draw[-{Stealth}] (0,0) -- (0:0.55) node[below] {$x$};
  % legenda
  \begin{scope}[shift={(2.9,2.55)}]
    \fill[accent!22] (0,0) rectangle (0.32,0.2); \node[anchor=west] at (0.4,0.1) {libero, $\Oss(t)$};
    \fill[accent] (0,-0.38) rectangle (0.32,-0.18); \node[anchor=west] at (0.4,-0.28) {occupato osservato};
    \draw[muted] (0,-0.76) rectangle (0.32,-0.56); \fill[pattern=north east lines,pattern color=muted] (0,-0.76) rectangle (0.32,-0.56); \node[anchor=west] at (0.4,-0.66) {in ombra, $\Omb(t)$};
    \fill[muted!22] (0,-1.14) rectangle (0.32,-0.94); \node[anchor=west] at (0.4,-1.04) {non osservato};
    \draw[muted,densely dotted] (0,-1.52) rectangle (0.32,-1.32); \node[anchor=west,text width=4.4cm] at (0.4,-1.42) {cella di guardia (non osservato)};
  \end{scope}
\end{tikzpicture}
\caption{Mappa polare del suolo attorno a un ostacolo, in pianta (schema non in scala: celle e cella di guardia ingrandite). Il suolo vicino alla camera, i lati del campo dove il suolo entra più lontano, l'esterno del campo e ciò che sta oltre $3\unita{m}$ sono «non osservato».}\label{fig:map-polare}
\end{figure}

\subsection{La regola del «libero»: impronta del suolo con cella di guardia}\label{sec:map-regola}

Un pixel $p$ \emph{legge il suolo} nel fotogramma $t$ se $D_t(p)$ è valido e
\begin{equation}\label{eq:map-margine}
|D_t(p)-z_0(p)|\le m\big(z_0(p)\big),\qquad m(z)=0{,}03+3\,k\,z^2\ \ [\mathrm{m}],
\end{equation}
dove il termine $k\,z^2$ segue la crescita quadratica del rumore di profondità con la distanza (\S\ref{sec:geo}) e $k$ dipende dal sensore.

\begin{definizione}[Regola del «libero» della cella]\label{def:map-libero}
A ogni fotogramma, senza alcuna memoria dei fotogrammi precedenti:
\begin{enumerate}
\item $\mathrm{buoni}(c)$ è il numero di pixel di $I(c)$ che leggono il suolo secondo la \eqref{eq:map-margine};
\item $c$ è \emph{letta} se $\mathrm{buoni}(c)\ge n(c)/2$ (e almeno 1), e $c$ non è occupata né in ombra;
\item $c$ è \emph{libero} se è letta e sono lette tutte le sue 8 celle vicine con $n>0$ (\emph{cella di guardia}).
\end{enumerate}
Occupato e ombra del fotogramma corrente prevalgono; nessun «libero» è ricordato, dedotto o prodotto da un contatore di tempo.
\end{definizione}

\textbf{Parametri.} Cella $1^\circ\times0{,}1\unita{m}$; quota d'ostacolo $0{,}07$–$2{,}0\unita{m}$; 3 pixel minimi per l'occupato; $k=0{,}004$ per il modello nominale del sensore e $k=0{,}0100$ per quello al limite; soglia di metà impronta; una cella di guardia; inclinazione nota entro $0{,}5^\circ$; memoria di 0 fotogrammi (Appendice~\ref{sec:par}). La soglia di metà impronta e la cella di guardia sono state scelte dopo un primo confronto fra varianti sullo stesso insieme di scene (Tabella~\ref{tab:map-regole}); la verifica è stata poi ripetuta su semi diversi.

\subsection{Proposizione di sicurezza}\label{sec:map-sicurezza}

\begin{ipotesi}[Profondità valida = giusta]\label{ip:map-valida} Un pixel valido con profondità $d$ garantisce che il suo raggio è privo di superfici fino a $d$, meno la tolleranza di misura.\end{ipotesi}
\begin{ipotesi}[Camera ferma, posa nota entro una tolleranza]\label{ip:map-posa} L'errore di posa sposta un punto a distanza $\rho$ di al più $\varepsilon_p(\rho)=\varepsilon_t+\rho\,\delta_\theta$, con $\varepsilon_t$ errore di traslazione e $\delta_\theta$ errore angolare in radianti.\end{ipotesi}
\begin{ipotesi}[Suolo piano]\label{ip:map-piano} Il suolo è un piano.\end{ipotesi}

\begin{proposizione}[Sicurezza del «libero»]\label{prop:map-sicurezza}
Si assuma la regola: lo stato «libero» è scritto solo dalla funzione che legge il fotogramma corrente, e solo per le celle di $\Oss(t)$. Allora:
\begin{enumerate}[label=(\alph*)]
\item nessuna cella in ombra, occupata, fuori campo o oltre portata è dichiarata libera al fotogramma $t$, e nemmeno una cella in cui meno di metà dei pixel dell'impronta legge il suolo (in particolare, con più di metà dei pixel non validi) o che ha una cella vicina in questa condizione;
\item sotto le Ipotesi~\ref{ip:map-valida}, \ref{ip:map-posa} (con $\varepsilon_p=0$) e \ref{ip:map-piano}, se una cella è dichiarata libera al fotogramma $t$, all'istante d'acquisizione nessuna superficie interseca i raggi dei pixel che ne leggono il suolo a profondità inferiore a $z_0(p)-m(z_0(p))$, meno la tolleranza di misura, dove $z_0(p)$ è la profondità attesa del suolo lungo il raggio del pixel $p$: non è esclusa una superficie più vicina al suolo del margine \eqref{eq:map-margine}.
\end{enumerate}
\end{proposizione}
\begin{proof}
(a) Per la Definizione~\ref{def:map-libero}, $\Oss(t)$ contiene solo celle il cui suolo è letto da pixel validi del fotogramma $t$ e che non sono occupate né in ombra; gli stati sono una partizione, quindi una cella in ombra o non osservata non sta in $\Oss(t)$, e nessun'altra funzione scrive «libero». (b) Un pixel che legge il suolo ha profondità non inferiore a $z_0(p)-m(z_0(p))$, e $z_0(p)$ è la distanza del suolo lungo il raggio (Ipotesi~\ref{ip:map-piano} e posa esatta); per l'Ipotesi~\ref{ip:map-valida} il raggio è vuoto fino a quella profondità.
\end{proof}
\emph{Livello di prova.} (a) \prova{D} per costruzione; (b) \prova{D} nel modello delle tre ipotesi; \prova{N} su un sensore fisico.

\begin{osservazione}[Che cosa la proposizione non dice]\label{oss:map-nondice}
(i) La garanzia riguarda il suolo e i raggi che lo raggiungono, non la colonna sopra la cella: un oggetto sospeso sopra quei raggi lascia la cella «libero». L'uscita corretta è «suolo visto libero», non «spazio libero». (ii) La parte (b) copre i raggi dei pixel che leggono il suolo, non l'intera cella: con la soglia di metà impronta gli altri pixel non sono coperti dalla dimostrazione, e la sicurezza della cella intera è verificata in simulazione (\S\ref{sec:map-banco}), non dimostrata. (iii) La garanzia vale all'istante d'acquisizione, non a quello d'uso (\S\ref{sec:occ}). In simulazione, con modello ideale del sensore (150 fotogrammi, finestra d'uso di $0{,}20\unita{s}$), 537 celle·fotogramma dichiarate libere su 88\,418 risultano occupate entro la finestra; un margine di $0{,}33\unita{m}$ (per $1{,}4\unita{m/s}$) attorno a ombra e non osservato porta gli errori a 0 ma conserva lo $0{,}4\,\%$ del libero \prova{S}.
\end{osservazione}

\begin{lemma}[Margine di posa]\label{lem:map-posa}
Sotto l'Ipotesi~\ref{ip:map-posa}, intesa come spostamento nel piano del suolo (traslazione e rotazione attorno alla verticale), e supposto corretto nel riferimento stimato lo stato calcolato, sia $\hat{\Oss}(t)$ il libero calcolato con la posa stimata e $x$ un punto a distanza $\rho$. Se il disco di centro $x$ e raggio $\varepsilon_p(\rho)$ è contenuto in $\hat{\Oss}(t)$, allora $x$ è libero nel riferimento vero.
\end{lemma}
\begin{proof} La posa vera porta $x$ in un punto $x'$ con $|x-x'|\le\varepsilon_p(\rho)$; quindi $x'$ appartiene al disco, che è contenuto in $\hat{\Oss}(t)$.\end{proof}
A $\rho=3\unita{m}$: $\delta_\theta=0{,}5^\circ$ dà $3\cdot0{,}00873=2{,}6\unita{cm}$; $\varepsilon_t=3\unita{cm}$ e $\delta_\theta=1^\circ$ danno $3+5{,}2=8{,}2\unita{cm}$; $\varepsilon_t=5\unita{cm}$ e $\delta_\theta=2^\circ$ danno $5+10{,}5=15{,}5\unita{cm}$ \prova{C}. In simulazione su una sezione piana (200 scene per caso), senza margine le scene con almeno un falso «libero» sono 114 su 200 e 162 su 200; con un margine non inferiore a quello teorico 0 su 200; con $10\unita{cm}$ dove ne servono $15{,}5$, 1 su 200 \prova{S}. Il margine non copre l'errore d'inclinazione, che agisce sulla quota ricostruita e quindi sullo stato: per esso vale un conto diretto (piano e raggio ideali): un oggetto alto $10\unita{cm}$ resta sopra la soglia di quota di $0{,}07\unita{m}$ se l'errore è sotto $0{,}69^\circ$ a $2{,}5\unita{m}$ e sotto $0{,}57^\circ$ a $3\unita{m}$; con $1^\circ$ la sua quota letta scende a $5{,}6$ e $4{,}8\unita{cm}$ \prova{C}. Da qui la tolleranza d'inclinazione di $0{,}5^\circ$ del dominio.

\subsection{Verifica in simulazione della regola}\label{sec:map-banco}

\textbf{Insieme di prova.} In simulazione, con modello ideale del sensore: camera ferma $320\times240$ pixel, campo $87^\circ\times58^\circ$, $h=1{,}2\unita{m}$, $\theta=15^\circ$; griglia $87\times30$. Dieci scene; $k\in\{0{,}004;\,0{,}0100\}$; cinque regimi di pixel non validi (nessuno; sparsi al $5\,\%$ e al $30\,\%$; a chiazze di $20\times20$ pixel al $5\,\%$ e al $30\,\%$); cinque errori d'inclinazione ($0$, $\pm0{,}5^\circ$, $\pm1^\circ$: la camera è inclinata di $15^\circ+e$, la regola assume $15^\circ$; in queste prove l'errore che abbassa la quota ricostruita, e fa scendere un ostacolo basso sotto la soglia, è $e<0$, mentre le prove del \S\ref{sec:map-schede} riportano lo stesso effetto con il segno opposto): 500 condizioni di 30 fotogrammi, cioè 15\,000 fotogrammi, generati due volte con semi diversi. Si conta una \emph{violazione} quando una cella «libero» ha il centro dentro la pianta di un corpo appoggiato, o quando il suolo del suo centro è nascosto da un corpo.

\begin{table}[t]\centering\small
\caption{Regole candidate per il «libero» della cella: fotogrammi con almeno una violazione su 7\,200 (8 scene, 2 valori di $k$, 5 regimi di pixel non validi, errore d'inclinazione $0$ e $\pm0{,}5^\circ$); fra parentesi la ripetizione con semi diversi. Ultima colonna: mediana per settore della distanza di suolo libero in una stanza vuota con il $5\,\%$ di pixel non validi sparsi (scala da $1{,}3$ a $3{,}0\unita{m}$; uguale per i due valori di $k$). In simulazione, con modello ideale del sensore \prova{S}.}\label{tab:map-regole}
\begin{tabularx}{\linewidth}{Lrr}
\toprule
Regola & Fotogrammi con violazione & Distanza mediana \\
\midrule
Lettura di un punto sopra il centro della cella (a $0{,}10\unita{m}$) & 932 (899) & $1{,}3\unita{m}$\\
Lettura di un punto a $0{,}20\unita{m}$ dal suolo & 1\,732 (1\,742) & $1{,}3\unita{m}$\\
Lettura a $0{,}20\unita{m}$ con veto delle celle occupate e quota verificata & 163 (179) & $1{,}3\unita{m}$\\
Impronta: basta un pixel che legge il suolo & 98 (126) & $3{,}0\unita{m}$\\
Impronta: metà dei pixel, senza cella di guardia & 6 (9) & $3{,}0\unita{m}$\\
\textbf{Impronta: metà dei pixel e cella di guardia} (Def.~\ref{def:map-libero}) & \textbf{0 (0)} & $3{,}0\unita{m}$\\
\bottomrule
\end{tabularx}
\end{table}

\textbf{Sicurezza.} La Tabella~\ref{tab:map-regole} confronta la regola adottata con le varianti più semplici. La lettura di un solo punto per cella cade perché la cella è estesa (a $0{,}20\unita{m}$: 15\,648 celle libere dentro un corpo); la variante con veto cade con i pixel non validi a chiazze, che tolgono i pixel che testimoniano l'ostacolo; con un solo pixel dell'impronta basta un raggio che sfiora il bordo del corpo. La regola adottata dà 0 violazioni su 9\,000 fotogrammi (10 scene, 300 condizioni con errore d'inclinazione entro $\pm0{,}5^\circ$), in entrambe le generazioni; le 300 condizioni comprendono un corpo che non rimanda segnale (0 su 900 fotogrammi, dove la regola a un pixel cade in 900 su 900) e una scatola alta $10\unita{cm}$ a $2{,}6$–$2{,}9\unita{m}$ (0 su 900). Una seconda esecuzione, con semi scelti in modo indipendente, dà ancora 0 su 9\,000 \prova{S}. Lo zero va letto per condizione, non per fotogramma, perché i fotogrammi di una condizione condividono la scena: 0 su 300 condizioni ha limite superiore al $95\,\%$ pari a $1{,}3\,\%$ (Wilson~\cite{wilson1927}); per scena, 0 su 10, il limite è $28\,\%$. Le scene sono state scelte, non campionate: il risultato è «verificato sull'insieme dichiarato», e la garanzia resta la Proposizione~\ref{prop:map-sicurezza} nelle sue ipotesi.

\textbf{Il prezzo in resa.} La Tabella~\ref{tab:map-resa} riporta la quota del libero vero che la regola dichiara. La cella di guardia costa $2{,}5$ punti senza pixel non validi, 23 punti con il $30\,\%$ sparso e 19 con il $30\,\%$ a chiazze. La perdita è sempre dal lato prudente: libero vero dichiarato «non osservato». In una stanza vuota, con $k=0{,}004$, lo stato di una cella cambia $5{,}7$ volte al secondo con il $30\,\%$ di pixel non validi sparsi, $2{,}7$ con il $5\,\%$ a chiazze e $6{,}7$ con il $30\,\%$ a chiazze, contro $0{,}01$ con il $5\,\%$ sparso \prova{S}. Quale forma abbiano i pixel non validi di un sensore reale non è verificato sperimentalmente \prova{N}.

\begin{table}[t]\centering\small
\caption{Resa del libero vero (media su 8 scene e 2 valori di $k$, errore d'inclinazione nullo) e distanza mediana di suolo libero per settore (stanza vuota, $k=0{,}004$). In simulazione, con modello ideale del sensore; rapporti di conteggi, senza intervallo \prova{S}.}\label{tab:map-resa}
\begin{tabular}{lrrr}
\toprule
Pixel non validi & Metà impronta, senza guardia & Regola adottata & Distanza per settore\\
\midrule
nessuno & $89{,}3\,\%$ & $86{,}8\,\%$ & —\\
$5\,\%$ sparsi & $89{,}3\,\%$ & $86{,}8\,\%$ & $3{,}0\unita{m}$\\
$30\,\%$ sparsi & $86{,}0\,\%$ & $62{,}6\,\%$ & $1{,}5\unita{m}$\\
$5\,\%$ a chiazze & $85{,}0\,\%$ & $78{,}3\,\%$ & $3{,}0\unita{m}$\\
$30\,\%$ a chiazze & $64{,}1\,\%$ & $45{,}4\,\%$ & $1{,}4\unita{m}$\\
\bottomrule
\end{tabular}
\end{table}

Con il $30\,\%$ di pixel non validi la distanza mediana scende a $1{,}4$–$1{,}5\unita{m}$ e il $43$–$47\,\%$ dei settori risponde «sconosciuto». Anche l'errore d'inclinazione dentro il dominio costa resa: in una stanza vuota, con $k=0{,}0100$ e $+0{,}5^\circ$ la resa è $80{,}6\,\%$ e la distanza mediana $2{,}4\unita{m}$; con $+1^\circ$, $33{,}7\,\%$ e $1{,}6\unita{m}$; con $k=0{,}004$ e $-1^\circ$, $82{,}4\,\%$ e $2{,}3\unita{m}$ \prova{S}.

\subsection{Controesempi e limiti del dominio}\label{sec:map-controesempi}

\begin{controesempio}[Errore d'inclinazione di $1^\circ$]\label{cex:map-inclinazione}
Errore d'inclinazione oltre la tolleranza di $0{,}5^\circ$ del dominio (il Lemma~\ref{lem:map-posa} non lo copre). Con errore di $-1^\circ$ la scatola alta $10\unita{cm}$ a $2{,}6$–$2{,}9\unita{m}$ è dichiarata «libero» in 150 fotogrammi su 300 (503 celle); in altre due esecuzioni con altri semi, in 145 e in 146 su 300. Sulle 8 scene di base la regola resta a 0 violazioni anche a $\pm1^\circ$ (4\,800 fotogrammi). In simulazione, con modello ideale del sensore \prova{S}.
\end{controesempio}

\begin{controesempio}[Vetro: profondità valida ma sbagliata]\label{cex:map-vetro}
Cade l'Ipotesi~\ref{ip:map-valida}: il pixel è valido, ma misura ciò che sta dietro la lastra. In simulazione, con una lastra trasparente di $1\times2\unita{m}$ a $2{,}0\unita{m}$ dal sensore, la lettura del fotogramma corrente dà per liberi 654 punti della lastra su 800 (600 su 800 con la lastra a $1{,}5\unita{m}$) \prova{S}. Nessuna regola sulla mappa risolve questo limite.
\end{controesempio}

\begin{controesempio}[Ostacolo che non rimanda segnale]\label{cex:map-muto}
La parte (a) della Proposizione~\ref{prop:map-sicurezza} protegge solo dove i pixel dell'ostacolo risultano non validi su più di metà dell'impronta della cella o di una cella vicina. Un corpo esteso che non rimanda segnale è retto dalla regola adottata (0 su 900 fotogrammi, \S\ref{sec:map-banco}). Un palo di $3\unita{cm}$ a $2{,}05\unita{m}$ con i pixel posti a zero, in una scena di prova letta con la variante che conserva il «libero» per tre fotogrammi, ha il piede «libero» in 6 fotogrammi su 6 \prova{S}; con la regola adottata il caso non è stato provato \prova{N}.
\end{controesempio}

\textbf{Altri limiti} (in simulazione, con modello ideale del sensore \prova{S}). Sotto un ostacolo sospeso il suolo visto resta «libero» (374 celle nell'insieme di prova): è il caso (i) dell'Osservazione~\ref{oss:map-nondice}. Una barra orizzontale di $1\unita{cm}$ a quota $0{,}8\unita{m}$ e a $2{,}0\unita{m}$ di distanza lascia la distanza per settore oltre la barra nel $4{,}7$–$7{,}3\,\%$ dei settori·fotogramma (7–11 su 150); con $2\unita{cm}$ o più, 0 su 150. Un oggetto alto $0{,}35\unita{m}$ a $0{,}5$–$0{,}8\unita{m}$, o alto $0{,}60\unita{m}$ a $0{,}3$–$0{,}6\unita{m}$, non produce alcun pixel, e il settore riporta $3{,}0\unita{m}$: l'uscita dichiara perciò in modo esplicito che il suolo è visto solo da $1{,}3\unita{m}$ in poi (\S\ref{sec:map-uscita}). Restano fuori dal dominio il suolo non piano e il sensore in moto (\S\ref{sec:lim}).

\subsection{Misure su profondità simulata dalle schede tecniche}\label{sec:map-schede}

Su profondità simulata dalle schede tecniche di tre camere di profondità (modelli di rumore nominale e al limite), a posa esatta e con ostacoli che restituiscono profondità, le celle dichiarate «libero» dentro l'ombra vera di un ostacolo sono 0 su 2\,907\,032 celle·fotogramma, in 135 sequenze (131 con i filtri della camera spenti e 4 con i filtri accesi); per sequenza, 0 su 131 con i filtri spenti (limite superiore di Wilson $2{,}9\,\%$, arrotondato per eccesso da $2{,}85\,\%$) \prova{S}. La misura è stata fatta con la variante della regola di lettura che conserva il «libero» per tre fotogrammi: sostiene la parte (a) su una profondità con rumore e pixel non validi realistici, non la Definizione~\ref{def:map-libero}. Una seconda serie di 381 sequenze, con la stessa variante, dà 5 celle·fotogramma «libero» nell'ombra vera su 3\,992\,186 (4 con il modello al limite e una scatola a $2\unita{m}$, 1 con il modello nominale e la scatola a $2{,}8\unita{m}$); non è risolto se dipendano dalla regola o dal bordo della maschera di riferimento. Fuori dominio, con la stessa variante: in una scena costruita a mano, con $+1^\circ$ di errore di posa, un oggetto di $10\unita{cm}$ a $2{,}55\unita{m}$ è «libero» in 7 celle su 7; su profondità simulata, con $+2{,}5^\circ$, una scatola di $15\unita{cm}$ a $2\unita{m}$ lo è in 9 fotogrammi su 10; con i filtri temporali della camera accesi e il sensore coperto restano 989–995 celle «libero» per altri 3 fotogrammi, e per questo i filtri vanno spenti \prova{S}.

La regola adottata è stata applicata a tre gruppi di 6 fotogrammi consecutivi della stessa profondità simulata (senza riferimento per il «libero»): con il modello nominale, una cassa e una persona, dichiara libero il $50{,}8\,\%$ delle celle con impronta fino a $3{,}0\unita{m}$, e nessuna cella dentro le piante note; con il modello al limite e il suolo vuoto il $15{,}9\,\%$ e fino a $2{,}2\unita{m}$, perché il $21\,\%$ dei pixel validi dell'impronta legge fuori dal margine \eqref{eq:map-margine}; con il suolo lucido lo $0{,}1\,\%$ \prova{S}. Su 6 fotogrammi di una sequenza pubblica ripresa con un'altra camera di profondità \cite{tum-rgbd}, la regola non dichiara quasi niente ($0{,}0$–$0{,}1\,\%$): il $37\,\%$ dei pixel validi dell'impronta legge più vicino del suolo atteso, e non è risolto se per gli oggetti in scena o per la posa \prova{P}. Il margine $m(z)$ va quindi tarato sul suolo reale prima di dichiarare una distanza \prova{N}.

\subsection{Perché non una mappa persistente o probabilistica a sensore fermo}\label{sec:map-persistente}

Le rappresentazioni in uso nella letteratura (griglie di occupazione, octree, voxel sparsi con tabella hash) sono richiamate fra i lavori correlati (\S\ref{sec:cos-correlati}). Sono state considerate, oltre alla mappa polare, la lettura diretta nell'immagine di profondità, la griglia densa di voxel, i voxel sparsi a blocchi, l'octree, la mappa a fasce di quota e la griglia bayesiana in logaritmo delle quote; costi e memoria sono in \S\ref{sec:cos}. Qui interessa il solo argomento di sicurezza.

\begin{osservazione}[Sensore fermo]\label{oss:map-fermo}
A sensore fermo ogni fotogramma osserva gli stessi raggi: una mappa persistente contiene o ciò che il fotogramma corrente dice già, o un valore ricordato che il fotogramma corrente non conferma. Per lo stato «libero», in un mondo che cambia, un valore ricordato non è una garanzia. È una deduzione, non un teorema; la sostengono i controesempi seguenti.
\end{osservazione}

\begin{controesempio}[Regola con memoria]\label{cex:map-memoria}
La regola «libero se il suolo è stato visto in almeno uno degli ultimi 3 fotogrammi» non soddisfa la Proposizione~\ref{prop:map-sicurezza}. Caso minimo, una cella e due fotogrammi: al primo il suolo è visto; al secondo un corpo che non rimanda segnale entra nella cella e i pixel diventano non validi. La regola con memoria dichiara «libero» una cella occupata; la regola del solo fotogramma corrente dichiara «non osservato». L'età massima del dato ricordato è $2\Delta=0{,}133\unita{s}$ a $15\unita{Hz}$, cioè $0{,}33\unita{m}$ a $2{,}5\unita{m/s}$ \prova{C}. In simulazione, con due programmi indipendenti: 1\,186 punti del suolo sotto il corpo dati per liberi, con almeno un errore in 87 fotogrammi su 87; e 58 fotogrammi su 90 con almeno un errore ($0{,}64$; Wilson $0{,}54$–$0{,}74$) \prova{S}.
\end{controesempio}

\begin{proposizione}[Latenza della griglia bayesiana]\label{prop:map-logodds}
Con i parametri predefiniti di una libreria diffusa (probabilità $0{,}7$ al colpo e $0{,}4$ al mancato colpo, saturazione a $0{,}1192$ e $0{,}971$ \cite{octomap}) l'incremento al colpo è $\ln(0{,}7/0{,}3)=0{,}847$ e il limite inferiore $\ln(0{,}1192/0{,}8808)=-2{,}000$: un voxel saturo «libero» torna positivo dopo $\lceil 2{,}000/0{,}847\rceil=3$ colpi, cioè $0{,}2\unita{s}$ a $15\unita{Hz}$. Un voxel interno a un corpo appena comparso non riceve mai un colpo, perché i ritorni stanno sulla superficie: il suo valore resta negativo senza limite di tempo. Il primo enunciato è un calcolo \prova{C}; il secondo è verificato in simulazione \prova{S} (Tabella~\ref{tab:map-varianti}, quarta riga).
\end{proposizione}

\begin{table}[t]\centering\small
\caption{Varianti che ricordano o deducono un «libero», con il controesempio che le esclude. In simulazione, con modello ideale del sensore (camera ferma, $h=1{,}2\unita{m}$, $\theta=15^\circ$, $15\unita{Hz}$); fra parentesi quadre l'intervallo di Wilson al $95\,\%$ dove il conteggio è per fotogrammi \prova{S}.}\label{tab:map-varianti}
\begin{tabularx}{\linewidth}{>{\hsize=0.78\hsize}L>{\hsize=0.72\hsize}L>{\hsize=1.5\hsize}L}
\toprule
Variante & Scena di prova & Esito \\
\midrule
Due stati: «non occupato $\Rightarrow$ libero» su blocchi sparsi & scena ferma, blocchi non allocati & 80\,505 falsi su 146\,901 punti dati per liberi ($54{,}8\,\%$); con tre stati: 0\\\addlinespace[2pt]
Decadimento del ricordo verso il libero & pannello che scorre davanti a un soggetto fermo già visto & corpo dato per libero in 830 fotogrammi su 897 ($92{,}5\,\%$ $[90{,}6;\,94{,}1]$)\\\addlinespace[2pt]
Blocchi a memoria limitata, espulsione del meno recente & soggetto dietro un ostacolo & soggetto fuori da $\Pres(t)$ in 132 fotogrammi su 134 ($98{,}5\,\%$ $[94{,}7;\,99{,}6]$)\\\addlinespace[2pt]
Logaritmo delle quote letto come stato corrente & 40 scatole compaiono; lettura 14 fotogrammi dopo & 5\,553 voxel interni «liberi» su 16\,851 ($33{,}0\,\%$); lettura del fotogramma corrente: 0\\\addlinespace[2pt]
«Attraversato = libero» senza età & sensore coperto; scena con ostacolo & 71\,006 punti restano «libero»; falsi dentro i corpi in 147 fotogrammi su 147\\\addlinespace[2pt]
Raggi decimati (un pixel ogni $4\times4$) & scena tipica, voxel di $5\unita{cm}$ & 2\,096 voxel con superficie dati per liberi su 6\,648 ($31{,}5\,\%$)\\\addlinespace[2pt]
Finestra di tre fotogrammi senza condizione sul pixel centrale & soggetto che non rimanda segnale & falsi dentro i corpi in 44 fotogrammi su 87 ($50{,}6\,\%$ $[40{,}3;\,60{,}8]$)\\\\
\bottomrule
\end{tabularx}
\end{table}

La Tabella~\ref{tab:map-varianti} raccoglie le varianti provate. Sistemi industriali di sorveglianza tridimensionale dichiarano un principio analogo (\S\ref{sec:cos-correlati}).

\textbf{Lettura di un punto nell'immagine di profondità.} Per interrogare un punto 3D $p$, di profondità $d_p$ e pixel di proiezione $(u,v)$ con intorno $3\times3$ $\mathcal N$, si usa la lettura diretta: «fuori campo» se $(u,v)$ è fuori dall'immagine; «non valido» se un pixel di $\mathcal N$ non è valido; «libero» se $d_p<\min_{\mathcal N}D_t-m(d_p)$; «in ombra» se $d_p>\max_{\mathcal N}D_t+m(d_p)$; «occupato» altrimenti. Se lo stato è «libero» allora $D_t(u,v)\ge d_p+m(d_p)$, e se la misura valida differisce meno di $m(d_p)$ dalla profondità della prima superficie lungo il raggio, nessuna superficie sta fra il sensore e $p$ \prova{D} (nel modello: assetto noto, misura valida entro il margine, una profondità per pixel). Una variante \emph{con conferma} accetta il «libero» di uno dei due fotogrammi precedenti, purché nel fotogramma corrente il pixel centrale sia valido e nessun pixel valido di $\mathcal N$ sia più vicino di $d_p+m(d_p)$. Per entrambe, i punti dati per liberi dentro un corpo a più di $5\unita{cm}$ dalla superficie sono 0 in 321 fotogrammi ($0\,\%$ $[0;\,1{,}18]$) con un primo programma e 0 in 900 fotogrammi ($0\,\%$ $[0;\,0{,}43]$; 10 condizioni di 90 fotogrammi, reticolo di 208\,000 punti a $5\unita{cm}$) con un secondo programma indipendente; i fotogrammi di una sequenza non sono indipendenti \prova{S}. Le rese delle due letture sono distinte: senza pixel non validi $97{,}1\,\%$ per entrambe ($88{,}4$–$91{,}8\,\%$ con il secondo programma); con il $5\,\%$ sparso $61{,}1\,\%$ a un fotogramma ($55{,}7$–$57{,}9\,\%$) e $87{,}7\,\%$ con conferma ($79{,}4$–$83{,}1\,\%$); con il $30\,\%$ sparso $3{,}6\,\%$ e $8{,}0\,\%$ (secondo programma) \prova{S}. La lettura per punto non basta però a dichiarare libera una \emph{cella} (prime righe della Tabella~\ref{tab:map-regole}): per le celle vale la Definizione~\ref{def:map-libero}, che non usa alcuna conferma.

\subsection{Uscita per settore: «suolo visto libero da 1,3 m a $d$ m»}\label{sec:map-uscita}

I settori dell'uscita raggruppano 5 azimut ($5^\circ$). Per ogni settore, $d$ è il minimo, sui 5 azimut, del raggio della prima cella non «libero» incontrata da $1{,}3\unita{m}$ in fuori; $d$ si ferma quindi una cella ($0{,}1\unita{m}$) prima di un ostacolo, ed è limitata a $2{,}6\unita{m}$, cioè alla portata meno la fascia di guardia di $0{,}40\unita{m}$ (\S\ref{sec:occ}). Se $d=1{,}3\unita{m}$, o se il settore non vede il suolo da $1{,}3\unita{m}$, l'uscita è «sconosciuto»: con la camera dell'insieme di prova il suolo si vede da $1{,}3\unita{m}$ in 9 settori su 17, e gli 8 laterali rispondono «sconosciuto» per geometria \prova{C}. La frase non porta mai una quota e non equivale a «nessun ostacolo»: per l'Osservazione~\ref{oss:map-nondice} e per i limiti del \S\ref{sec:map-controesempi} parla del suolo, non del volume. Le parole dell'uscita e la loro priorità sono in \S\ref{sec:lim}.

\begin{riquadro}{Condizioni di validità del «libero»}
Camera ferma; inclinazione nota entro $0{,}5^\circ$; suolo piano; ostacoli che restituiscono profondità valida; filtri temporali della camera spenti. Entro queste condizioni: parte (a) dimostrata per costruzione; parte (b) dimostrata nel modello a posa esatta; regola della cella, e tenuta con errore d'inclinazione entro $0{,}5^\circ$, verificate in simulazione, con modello ideale del sensore, su 300 condizioni dichiarate, non dimostrate. Fuori: vetro, ostacoli sottili che non rimandano segnale, errore d'inclinazione di $1^\circ$, ostacoli sospesi o sotto il campo, suolo lucido. Nessuna verifica su un sensore fisico.
\end{riquadro}

\section{Rilevamento, classificazione, associazione e stima dello stato}\label{sec:sti}

Il capitolo descrive come i punti di profondità diventano gruppi, quando un gruppo è dichiarato «mobile», quale classe e quale profilo cinematico porta una traccia, e come un filtro di Kalman a velocità quasi costante \cite{kalman1960} ne stima posizione e velocità. Tutti i risultati vengono da dimostrazioni nel modello dichiarato \prova{D}, da calcoli \prova{C}, da simulazioni \prova{S} o da traiettorie pubbliche di pedoni \prova{P}; nessuno viene da un sensore fisico, e ciò che richiederebbe un sensore è marcato \prova{N}. Le frequenze portano la numerosità $N$ e l'intervallo di Wilson al 95\,\%.

\begin{proposizione}[Ordine della catena]\label{prop:sti-ordine}
Siano $Z_k$ i rilevamenti non etichettati del ciclo $k$ e, per ogni traccia $i$, un filtro con previsione $(\hat{\mathbf x}_i^-,P_i^-)=(F\hat{\mathbf x}_i,\,FP_iF^{\top}+Q)$ e aggiornamento funzione di $(\hat{\mathbf x}_i^-,P_i^-,\mathbf z)$ per una misura $\mathbf z$ attribuita a $i$. Ogni esecuzione corretta rispetta l'ordine parziale: rilevamento e previsione $\to$ associazione $\to$ aggiornamento $\to$ decisione. \prova{D}
\end{proposizione}
\begin{proof}
L'aggiornamento richiede i rilevamenti e la corrispondenza fra tracce e misure; l'associazione usa la porta $\nu^{\top}S^{-1}\nu\le\gamma$, con $\nu=\mathbf z-H\hat{\mathbf x}_i^-$ e $S=HP_i^-H^{\top}+R$, quindi richiede la previsione; la decisione legge stato e covarianza aggiornati.
\end{proof}
La previsione precede dunque l'associazione, e l'aggiornamento consuma misure già associate: è lo schema degli inseguitori per rilevamento, con filtro di Kalman e assegnamento ungherese a valle del rilevatore \cite{bewley2016}.


\subsection{Dai punti ai gruppi}\label{sec:sti-gruppi}

\begin{definizione}[Cella occupata, gruppo, gruppo tagliato]\label{def:sti-gruppo}
Una cella $(\varphi,\rho)$ della mappa polare ($1^\circ\times0{,}1\unita{m}$, $\rho\le3\unita{m}$) è «occupata osservata» se contiene almeno $n_{\min}=3$ punti di profondità con quota $y\in[0{,}07;\,2]\unita{m}$ ed è il primo ritorno lungo il suo azimut. Un \emph{gruppo} è una componente connessa di celle occupate, con vicinanza a otto celle negli indici $(\varphi,\rho)$; il suo centro $\mathbf c$ è la media delle coordinate $(x,z)$ dei suoi punti. Un gruppo è \emph{tagliato} quando confina con celle in ombra di un gruppo più vicino alla camera o tocca il bordo del campo.
\end{definizione}
Il centro della parte visibile non è il centro dell'oggetto: un gruppo tagliato non aggiorna il filtro, non entra nella regola «mobile», non fa nascere tracce e non riceve etichette nuove.

\paragraph{Rumore del centro.} Per una camera stereo $d=f\,b/\text{disparità}$, quindi un errore di disparità $\sigma_d$ dà
\begin{equation}\label{eq:sti-sigmaz}
\sigma_z(d)=\frac{d^2}{f\,b}\,\sigma_d,\qquad
\sigma_t(d)=\frac{d\,\pi/180}{\sqrt{12}},\qquad
R(d)=\sigma_z^2(d)\,\mathbf u_r\mathbf u_r^{\top}+\sigma_t^2(d)\,\mathbf u_t\mathbf u_t^{\top},
\end{equation}
con $\mathbf u_r$ versore dalla camera al centro, $\mathbf u_t$ il suo ortogonale sul suolo e $\sigma_t$ lo scarto tipo di un errore uniforme sulla cella di $1^\circ$. Con $\sigma_d=0{,}08\unita{px}$ e crescita quadratica con la distanza \cite{realsense-taratura-profondita}, e con i valori assunti $f=424\unita{px}$ e $b=0{,}05\unita{m}$, si ha $\sigma_z=1{,}51/3{,}40/6{,}04\unita{cm}$ a $2/3/4\unita{m}$ (curva \emph{nominale}) e $\sigma_t=1{,}01/1{,}51/2{,}02\unita{cm}$ \prova{C}. Come analisi di sensibilità si usano due curve più severe: $4\,(d/2)^2\unita{cm}$ (\emph{pessimista}: $4/9/16\unita{cm}$) e $5\,(d/2)^2\unita{cm}$ (\emph{severa}: $5/11{,}25/20\unita{cm}$). Il rumore del centro di un gruppo di $n$ punti sta fra $\sigma/\sqrt n$ (errori indipendenti) e $\sigma$ (errori del tutto correlati), se gli errori hanno uguale varianza e correlazioni non negative \prova{D}; il progetto adotta l'estremo prudente $\sigma_c=\sigma_z(d)$. Il valore reale di $\sigma_c$ è non verificato sperimentalmente \prova{N}: su profondità simulata dalle schede tecniche, per una persona ferma (20 semi), il rumore radiale mediano è $0{,}36/0{,}99/1{,}25\unita{cm}$ a $2/3/4\unita{m}$; l'unico dato da una sequenza pubblica reale \cite{depthtrack} (altro sensore della stessa famiglia, 288 fotogrammi, senza verità di riferimento) è $5{,}8$--$6{,}8\unita{cm}$ sotto $3\unita{m}$. I due valori stanno ai lati opposti della curva nominale.

\subsection{La regola «mobile»}\label{sec:sti-mobile}

\begin{definizione}\label{def:sti-mobile}
Sia $\bar{\mathbf c}_k=(\mathbf c_k+\mathbf c_{k-1}+\mathbf c_{k-2})/3$ la media di $m=3$ centri consecutivi non tagliati. Al fotogramma $k$ si confronta $\bar{\mathbf c}_k$ con la media più vecchia disponibile $\bar{\mathbf c}_j$, con $m\le k-j\le1/\Delta$ (medie disgiunte, finestra fino a $1\unita{s}$). Il gruppo è «mobile» se, lungo l'asse radiale o lungo quello trasversale,
\begin{equation}\label{eq:sti-soglia}
|\bar c_k-\bar c_j|>k_s\,\sigma\,\sqrt{2/m},\qquad k_s=5,
\end{equation}
con $\sigma=\sigma_z(d)$ o $\sigma_t(d)$ alla maggiore delle distanze ai due capi.
\end{definizione}

\begin{proposizione}[Falso allarme e rilevamento]\label{prop:sti-mobile}
Siano gli errori del centro gaussiani, indipendenti fra i fotogrammi, di scarto tipo $\sigma$ lungo un asse. (a) Per un oggetto fermo la differenza di due medie disgiunte di $m$ fotogrammi è $\mathcal N(0,\sigma_\delta^2)$ con $\sigma_\delta=\sigma\sqrt{2/m}$; la soglia \eqref{eq:sti-soglia} vale $k_s\sigma_\delta$ e la probabilità di falso allarme per prova e per asse è $P_{fa}=2\,[1-\Phi(k_s)]$. (b) Per un oggetto con componente di velocità $v$ lungo l'asse, nella finestra $T_w=(k-j)\Delta$ lo spostamento è $\delta=v\,T_w$ e
\[P_{\mathrm{mob}}=\Phi\!\Big(\frac{\delta}{\sigma_\delta}-k_s\Big)+\Phi\!\Big(-\frac{\delta}{\sigma_\delta}-k_s\Big),\]
quindi nella prova singola la velocità rilevata con probabilità $0{,}5$ è $v_{0,5}=k_s\sigma_\delta/T_w$ e quella rilevata con probabilità $0{,}9$ è $v_{0,9}=(k_s+1{,}2816)\,\sigma_\delta/T_w$. \prova{D}
\end{proposizione}
\begin{proof}
Ogni media ha varianza $\sigma^2/m$; le due medie sono indipendenti perché disgiunte, quindi la differenza ha varianza $2\sigma^2/m$ e media $0$ (oggetto fermo) o $\delta$ (oggetto in moto). $P_{\mathrm{mob}}$ è la massa oltre $\pm k_s\sigma_\delta$; il secondo termine è trascurabile e $P_{\mathrm{mob}}=0{,}9$ quando $\delta/\sigma_\delta-k_s=1{,}2816$.
\end{proof}
Con $m=3$: $\sigma_\delta=0{,}8165\,\sigma$, soglia $4{,}08\,\sigma$ ($6{,}2\unita{cm}$ in radiale e $4{,}1\unita{cm}$ in trasversale a $2\unita{m}$, curva nominale), $P_{fa}=5{,}7\cdot10^{-7}$. Nel primo secondo di una traccia a $15\unita{Hz}$ le prove sono al più 13 per asse: per la disuguaglianza dell'unione la probabilità di almeno un'etichetta falsa è $\le2\cdot13\cdot5{,}7\cdot10^{-7}=1{,}5\cdot10^{-5}$. Con la finestra massima $T_w=13\Delta=0{,}867\unita{s}$, in radiale $v_{0,5}/v_{0,9}$ valgono $0{,}071/0{,}089\unita{m/s}$ a $2\unita{m}$ e $0{,}160/0{,}201$ a $3\unita{m}$ con la curva nominale, $0{,}424/0{,}533$ a $3\unita{m}$ con la pessimista e $0{,}530/0{,}666$ con la severa \prova{C}: un pedone a $0{,}5\unita{m/s}$ in avvicinamento da $3\unita{m}$ sta sopra la soglia con la prima, vicino con la seconda, sotto con la terza (Tabella~\ref{tab:sti-mobile}).

\begin{table}[H]\centering\footnotesize
\caption{Regola «mobile», oggetto a $0{,}5\unita{m/s}$: quota di tracce etichettate entro $1\unita{s}$. In simulazione, con modello ideale del sensore (centro puntiforme, rumore gaussiano indipendente e noto, $15\unita{Hz}$) \prova{S}. Da $2\unita{m}$: di traverso o in avvicinamento, tre curve di rumore; da $3\unita{m}$: in avvicinamento. Prova singola: valori del calcolo d'origine con la Proposizione~\ref{prop:sti-mobile}; con la curva severa sovrastima, perché la regola usa la maggiore delle distanze \emph{misurate}, che il rumore gonfia. Rifacendo il conto con $T_w=0{,}867\unita{s}$ e lo scarto tipo a $3\unita{m}$ si ottengono $0{,}815$ e $0{,}389$ per le curve pessimista e severa: lo scarto dai valori in tabella non è spiegato. Le percentuali della prima colonna sono quelle stampate dalla simulazione; i conteggi sono dedotti.}\label{tab:sti-mobile}
\begin{tabular}{lccc}
\toprule
Caso & Simulazione, $N=2000$ & Ricalcolo, $N=20\,000$ & Prova singola\\
\midrule
da 2 m, sei casi & 100\,\% in ciascuno ($\ge99{,}81\,\%$) & 100\,\% ($\ge99{,}98\,\%$) & 1,000\\
da 3 m, curva nominale & 100\,\% & 100\,\% & 1,000\\
da 3 m, curva pessimista & 85,85\,\% (84,3--87,3) & 86,2\,\% (85,7--86,7) & 0,853\\
da 3 m, curva severa & 26,6\,\% (24,7--28,6) & 26,9\,\% (26,3--27,5) & 0,371\\
fermo: tracce con etichetta falsa & 0 su 4000 in ciascuno ($\le0{,}10\,\%$) & 1 su 120\,000 & $\le1{,}5\cdot10^{-5}$\\
\bottomrule
\end{tabular}
\end{table}
\paragraph{Condizioni di validità.} (i) Rumore noto: se quello creduto è la metà del vero la soglia scende a $2{,}5\,\sigma_\delta$ e, in simulazione con modello ideale del sensore, il $10{,}9\,\%$ delle tracce ferme riceve un'etichetta falsa nel primo secondo ($N=20\,000$, curva pessimista, $2\unita{m}$) \prova{S}. (ii) Errori indipendenti fra i fotogrammi, senza valori aberranti. (iii) Il centro cambia solo per il rumore di misura, non perché cambia l'insieme delle celle del gruppo.

\begin{controesempio}[Pixel misti e fusione dei gruppi]\label{ctr:sti-misti}
Su profondità simulata dalle schede tecniche, con i \emph{pixel misti} ai bordi delle sagome (profondità intermedie fra l'oggetto vicino e lo sfondo), un pannello fermo dietro cui passa una persona è etichettato «mobile» in 79 passaggi su 90 (79,4--93,0\,\%) \prova{S}. Cade la condizione (iii): quando le sagome si toccano i pixel misti riempiono celle «ponte» e i due gruppi diventano uno (90 passaggi su 90); il centro del gruppo fuso si sposta di decine di centimetri contro una soglia di 4--6\,cm, e la regola legge come moto un cambio di segmentazione. Senza pixel misti la fusione passa da 10 casi su 10 a 1 su 10. L'entità dei pixel misti è un'ipotesi del modello del sensore: l'effetto è non verificato sperimentalmente \prova{N}.
\end{controesempio}

\paragraph{Una soglia di punti assoluta non rimedia.} Con $n_{\min}=40$ a $424\times240$ pixel le etichette false scendono a 5 passaggi su 90 (2,4--12,4\,\%), ma una scatola alta 15\,cm a $2\unita{m}$ è vista in 3 prove su 10 anziché 10 su 10; a $848\times480$ la stessa soglia dà la continuità dell'identità in 0 passaggi su 10, e 160 in 10 su 10 \prova{S}: i punti per cella scalano con $f_xf_y/d^2$ e con l'area della cella. Le celle ponte hanno 33--125 punti e quelle della scatola 13--43: nessuna soglia per cella le separa. Una soglia proporzionale ai pixel attesi applicata alle celle della mappa produce inoltre «libero» nell'ombra vera della scatola (98 e 91 celle$\cdot$fotogramma a 424 e 848 pixel di larghezza, contro 0 e 4 con la mappa invariata) \prova{S}: le soglie vanno poste sui gruppi, non sulle celle che decidono ombra e «libero» (\S\ref{sec:map}).

\paragraph{Regola dei gruppi a celle forti.} Questa regola, le condizioni (c)--(e) del riaggancio (\S\ref{sec:sti-riaggancio}) e l'avviso per le tracce chiuse vicino a un'ombra (\S\ref{sec:occ-risultati}) non fanno parte della \emph{configurazione di riferimento} (gruppi per componenti connesse della Definizione~\ref{def:sti-gruppo}, riaggancio senza le condizioni (c)--(e)): si attivano all'avvio solo se le prove d'ingresso del \S\ref{sec:lim-salute} sono superate, e l'uscita dichiara quali sono attive. La mappa non cambia. Sia $n(c)$ il numero di punti della cella $c$ e $N_{\mathrm{att}}(\rho)=\max\{n_{\min},\lceil f_xf_y\,\Delta\varphi\,h_{\mathrm{rif}}/\rho\rceil\}$, con $h_{\mathrm{rif}}=0{,}05\unita{m}$ e $\Delta\varphi=1^\circ$ in radianti (20 punti a $2\unita{m}$ a 424 pixel di larghezza, 79 a 848). In ogni componente connessa $K$ di celle occupate sono \emph{forti} le celle con $n(c)\ge\eta\max_K n$ e $n(c)\ge N_{\mathrm{att}}(\rho_c)$, $\eta=0{,}30$; i gruppi sono le componenti connesse delle celle forti, e un gruppo a distanza $d_G$ è scartato se ha meno di $f_xf_y\,A_{\mathrm{rif}}/d_G^2$ punti, $A_{\mathrm{rif}}=0{,}004\unita{m^2}$. Le celle deboli restano «occupato» e fanno ombra. Il valore di $\eta$ è tarato su sequenze simulate; le cifre seguenti sono su altri semi. Su profondità simulata dalle schede tecniche, 424 e 848 pixel di larghezza \prova{S}: pannello «mobile» in 0 passaggi su 140 nel dominio dell'identità al ritorno, con rumore nominale (130 e 94 con i gruppi della Definizione~\ref{def:sti-gruppo}); fuori da quel dominio, in una replica con il soggetto sulla linea a $2{,}1\unita{m}$ dalla camera, in 7 passaggi su 10 a entrambe le risoluzioni (40--89\,\%): lo zero vale solo nel dominio; 0 tracce confermate sul suolo vuoto (9 sequenze; 137 e 206 con i gruppi per componenti connesse); scatola di 15\,cm a $2\unita{m}$ vista in 190 fotogrammi su 190 (19 sequenze, 3 delle quali al livello degradato del modello di rumore). \emph{Limiti di validità}: un cubo di 10\,cm a $2\unita{m}$ resta «occupato» ma non è mai un gruppo (0 fotogrammi su 60), e la scatola di 15\,cm a $2{,}8\unita{m}$ lo è in 45 su 60 (63--84\,\%): la taglia minima inseguibile è 15\,cm entro $2\unita{m}$. Al livello degradato del modello di rumore la regola non regge (105--106 tracce confermate su 3 pavimenti vuoti): l'avvio prevede perciò una prova d'ingresso (suolo vuoto, zero gruppi).

\subsection{Classificazione e profili cinematici di classe}\label{sec:sti-classi}

\begin{riquadro}{Stato dichiarato del componente}
Il sistema \emph{richiede} un classificatore su più classi di oggetti in movimento. In questo lavoro il classificatore \emph{non è realizzato né valutato}: nessun modello è stato eseguito o addestrato. È simulato con etichette assegnate a mano nelle prove (da un copione della scena) e, per misurare la sensibilità del sistema, con errori iniettati da una matrice di confusione e da un ritardo dichiarati, a seme fisso. Nessuna classe diversa da «persona» è stata simulata nella catena completa. La sicurezza delle indicazioni non dipende dall'etichetta (\S\ref{sec:occ}).
\end{riquadro}

\paragraph{Interfaccia.} Per ogni traccia confermata il componente riceve i ritagli di profondità degli ultimi fotogrammi con la maschera del gruppo e la marca «tagliato», le dimensioni metriche, la storia della traccia e lo stato di visibilità. Restituisce, come campi della traccia: \emph{classe}; \emph{livello} (tipo, categoria, non classificato); \emph{fiducia}; \emph{origine} dell'etichetta (a mano; a mano con errori iniettati; modello); \emph{età}; \emph{smentita}. Con origine «a mano» la fiducia è «non applicabile»; l'origine «modello» qui non compare mai. Il componente è asincrono (1--2\,Hz o a evento) e il ciclo a $15\unita{Hz}$ non lo attende. L'etichetta è una proprietà della traccia: non entra nel costo d'associazione né nella porta; in occlusione resta con la sua età; al riaggancio è «presunta», come l'identità (\S\ref{sec:sti-riaggancio}); con identità ambigua decade a «non classificato». Il componente scrive solo i campi dell'etichetta: non scrive celle, non rilascia avvisi, non abbassa livelli.

\paragraph{Tassonomia con ripiego.} Sei categorie: persona; veicolo terrestre; veicolo ferroviario; aeromobile; oggetto mobile generico (carrello a mano, porta o barriera, carico sospeso, animale); non classificato. I tipi sono le righe della Tabella~\ref{tab:sti-profili}. Vale il tipo se la sua probabilità è almeno $0{,}7$ (ipotesi, da tarare); altrimenti la categoria, se la somma dei suoi tipi è almeno $0{,}7$; altrimenti «non classificato». Il profilo di una categoria è il caso peggiore dei suoi tipi; quello di «non classificato» è l'inviluppo dello \emph{scenario}, cioè delle classi ammesse nel sito, dichiarato all'avvio. Salire nella gerarchia può quindi solo allargare. È «non classificato» anche chi non ha etichetta, il gruppo tagliato e l'etichetta smentita.

\begin{table}[!htb]\centering\footnotesize
\caption{Profili cinematici a priori delle classi: valori iniziali di configurazione, non un riferimento. Origine dei valori: $^{a}$ letto in una fonte pubblica di terza mano, non riscontrato su una fonte primaria e per questo senza riferimento bibliografico; $^{b}$ valore non verificato su una fonte primaria; senza marca: ipotesi. Solo 9 velocità su 30 hanno una fonte consultata.}\label{tab:sti-profili}
\begin{tabularx}{\linewidth}{@{}p{4.6cm}p{2.0cm}p{2.0cm}L@{}}
\toprule
Classe & $v$ tipica / massima (m/s) & acc. / frenata (m/s$^2$) & Gradi di libertà e modello di moto\\
\midrule
persona a piedi & 1,4$^{a}$ / 2,5$^{a}$ & 3,0 / 4,0 & 2, quasi olonomo; si ferma e inverte in un passo\\
persona di corsa & 3,0 / 8,0 & 4,0 / 4,0 & come sopra\\
bicicletta & 4,3$^{a}$ / 8,3 & 2,0 / 5,0 & non olonomo, raggio $\ge2$ m\\
monopattino & 4,2 / 5,6$^{a}$ & 2,0 / 4,0 & non olonomo, raggio $\ge1{,}5$ m\\
carrello elevatore & 2,2$^{b}$ / 4,5$^{b}$ & 1,5 / 2,5 & non olonomo, sterzo posteriore\\
veicolo a guida automatica & 1,0 / 2,0$^{a}$ & 1,0$^{b}$ / 2,0 & uniciclo o omnidirezionale\\
auto o furgone, area di lavoro & 2,8 / 8,3 & 3,0 / 6,9$^{a}$ & non olonomo, raggio $\approx5{,}5$ m$^{b}$\\
auto o furgone, strada urbana & 8,3 / 13,9$^{a}$ & 3,0 / 6,9$^{a}$ & come sopra\\
trattore di piazzale con rimorchi & 1,7$^{a}$ / 8,3$^{a}$ & 1,5 / 4,0 & articolato: i rimorchi tagliano la curva\\
aereo in rullaggio & 9,0$^{a}$ / 15,4$^{b}$ & 1,0 / 3,0 & non olonomo, raggio $\ge20$ m; l'ala spazza oltre le ruote\\
aereo al traino & 1,4$^{b}$ / 4,2$^{b}$ & 0,5 / 1,5 & articolato, grande raggio\\
treno in arrivo in banchina & 11,0 / 16,7 & 1,0 / 1,5$^{a}$ & 1, sul binario\\
tram a una fermata & 5,6 / 13,9 & 1,3 / 2,5$^{b}$ & 1, sul binario\\
treno o carri in manovra & 4,2 / 8,3$^{b}$ & 0,5 / 1,5$^{a}$ & 1, sul binario; spesso spinto\\
multirotore & 3,0$^{b}$ / 19,0$^{b}$ & 6,9$^{b}$ / 6,9$^{b}$ & 3 più l'imbardata; non rappresentato dalla mappa del suolo\\
non classificato & inviluppo dello scenario & inviluppo & olonomo: il caso peggiore\\
\bottomrule
\end{tabularx}
\end{table}
Nella Tabella~\ref{tab:sti-profili} il massimo della persona a piedi è la media del passo massimo misurato, non un estremo. Treni e aerei non sono etichettabili da questo sensore: lunghezze di $200\unita{m}$ e $37{,}6\unita{m}$ contro una portata di $3\unita{m}$ danno sempre un gruppo tagliato; le loro righe servono alla configurazione dello scenario e al dominio per classe (\S\ref{sec:occ}).

\paragraph{Controllo di coerenza cinematica.} Dal quarto ciclo con misura, un'etichetta è \emph{smentita} se in 4 degli ultimi 6 cicli vale una delle condizioni: (i) $|\hat{\mathbf v}|-2\sigma_v>\vmax(\text{classe})$; (ii) per le classi a un grado di libertà, distanza dalla linea configurata maggiore di $0{,}30\unita{m}+2\sigma_p$; (iii) per le classi non olonome, curvatura maggiore di $1{,}5$ volte la massima della classe, o inversione del verso in meno di $1\unita{s}$. Le soglie sono scelte, non tarate. Un'etichetta smentita decade a «non classificato», per quella traccia non torna, e la regione di previsione perde il livello dichiarato; l'insieme di possibile presenza e la fascia di guardia non si restringono. In simulazione, con modello ideale del sensore (4000 tracce per caso) \prova{S}: con l'etichetta giusta gli scatti falsi sono 0, 1 e 0 su 4000 per pedone, carrello e treno (un limite inferiore: le tracce rispettano per costruzione i limiti della classe); una «persona» a $6\unita{m/s}$ o a $3{,}5\unita{m/s}$ è smentita sempre, dopo $0{,}47\unita{s}$ (a $6\unita{m/s}$ sono $2{,}8\unita{m}$, quasi tutta la portata), e a $2{,}8\unita{m/s}$ nel $10{,}6\,\%$ dei casi; un carrello lento (1--$2{,}2\unita{m/s}$) etichettato «pedone» nel $4{,}1\,\%$, un pedone etichettato «carrello» nel $5{,}2\,\%$; un veicolo fermo etichettato «persona» in 0 casi su 2000. Il controllo toglie le etichette palesemente incompatibili con il moto osservato: non convalida un'etichetta e non è una barriera di sicurezza.

\paragraph{Requisiti posti al classificatore.} Le soglie sono ipotesi di progetto \prova{N}. (1)~Classi: la tassonomia almeno a livello di categoria, con «non classificato» come uscita effettiva. (2)~Fiducia tarata: errore di taratura atteso $\le0{,}05$ per classe su dati dello stesso sito; senza taratura la fiducia non si mostra. (3)~Prima etichetta entro $0{,}5\unita{s}$ dalla conferma della traccia: la portata di $3\unita{m}$ si attraversa in $0{,}75\unita{s}$ a $4\unita{m/s}$ e in $0{,}21\unita{s}$ a $14\unita{m/s}$, e la conferma chiede $0{,}27\unita{s}$: per le classi veloci protegge lo scenario dichiarato, non l'etichetta. (4)~Errori asimmetrici: tollerati gli scambi fra profili uguali e da lento a veloce ($\le10\,\%$); la probabilità che la classe vera sia più veloce di un'etichetta lenta è $\le1\,\%$; mai da libero a vincolato. (5)~Un tipo mai visto esce «non classificato»; la quota di oggetti ignoti con un'etichetta sicura è $\le5\,\%$. (6)~Al più un cambio d'etichetta ogni $2\unita{s}$ per traccia; nessuna etichetta nuova su un gruppo tagliato. (7)~Valutazione con partizioni per sito e per sequenza, mai per fotogramma, su un sito mai visto, per classe e con intervalli.

\subsection{Modello del moto e filtro}\label{sec:sti-filtro}

Lo stato è $\mathbf x=[p_x,p_z,v_x,v_z]^{\top}$. Per asse, $\dot p=v$ e $\dot v=a(t)$, con $a$ rumore bianco di densità $q$ (m$^2$/s$^3$), indipendente fra gli assi. Campionando a passo $\Delta$:
\begin{equation}\label{eq:sti-modello}
F=\begin{bmatrix}I_2&\Delta I_2\\0&I_2\end{bmatrix},\quad
Q=q\begin{bmatrix}\tfrac{\Delta^3}{3}I_2&\tfrac{\Delta^2}{2}I_2\\[2pt]\tfrac{\Delta^2}{2}I_2&\Delta I_2\end{bmatrix},\quad
H=\begin{bmatrix}I_2&0\end{bmatrix},\quad R=R(d).
\end{equation}
Predizione: $\hat{\mathbf x}^-=F\hat{\mathbf x}$, $P^-=FPF^{\top}+Q$. Innovazione: $\nu=\mathbf z-H\hat{\mathbf x}^-$, $S=HP^-H^{\top}+R$. Correzione: $K=P^-H^{\top}S^{-1}$, $\hat{\mathbf x}=\hat{\mathbf x}^-+K\nu$, $P=(I-KH)P^-(I-KH)^{\top}+KRK^{\top}$ (forma di Joseph, che conserva simmetria e positività). La correzione usa solo gruppi non tagliati che cadono nella porta (\S\ref{sec:sti-porta}).

\begin{proposizione}[Forma chiusa di $Q$ e crescita senza misure]\label{prop:sti-crescita}
Per asse, con $A=\begin{bsmallmatrix}0&1\\0&0\end{bsmallmatrix}$ e $G=[0,1]^{\top}$: (a) $Q(\Delta)=\int_0^{\Delta}e^{As}G\,q\,G^{\top}e^{A^{\top}s}\,ds=q\begin{bsmallmatrix}\Delta^3/3&\Delta^2/2\\\Delta^2/2&\Delta\end{bsmallmatrix}$, e $Q(\Delta_1+\Delta_2)=F(\Delta_2)Q(\Delta_1)F(\Delta_2)^{\top}+Q(\Delta_2)$: la previsione non dipende dal numero di passi in cui è spezzata. (b) Se all'ultima correzione la covarianza per asse ha elementi $P_{pp},P_{pv},P_{vv}$, dopo un tempo $\tocc$ senza misure
\begin{equation}\label{eq:sti-crescita}
\sigma_p^2(\tocc)=P_{pp}+2\,\tocc\,P_{pv}+\tocc^2\,P_{vv}+\frac{q\,\tocc^3}{3}.
\end{equation}
\prova{D}
\end{proposizione}
\begin{proof}
(a) $A^2=0$, quindi $e^{As}=I+As$ e $e^{As}G=[s,1]^{\top}$: l'integrando è $q\begin{bsmallmatrix}s^2&s\\s&1\end{bsmallmatrix}$ e si integra termine a termine. La seconda identità si verifica sostituendo. (b) Per (a), $P(\tocc)=F(\tocc)PF(\tocc)^{\top}+Q(\tocc)$, il cui elemento di posizione è la tesi.
\end{proof}
Con covarianza isotropa la regione al livello $\alpha$ è un cerchio centrato in $\hat{\mathbf p}$ di raggio $r_\alpha=\sigma_p\sqrt{-2\ln(1-\alpha)}$, perché $\chi^2_2$ ha funzione di ripartizione $1-e^{-x/2}$: $r_{0,50}=1{,}177\,\sigma_p$, $r_{0,90}=2{,}146\,\sigma_p$, $r_{0,95}=2{,}448\,\sigma_p$.

\begin{table}[H]\centering\footnotesize
\caption{$\sigma_p(\tocc)$ per asse, in metri, dalla covarianza stazionaria a $15\unita{Hz}$ con rumore di misura di 5\,cm (per $q=0{,}2$: $\sqrt{P_{pp}}=3{,}3\unita{cm}$, $\sqrt{P_{vv}}=0{,}205\unita{m/s}$) \prova{C}.}\label{tab:sti-crescita}
\begin{tabular}{llcccc}
\toprule
$q$ (m$^2$/s$^3$) & & 0,5 s & 1 s & 2 s & 3 s\\
\midrule
0,2 & $\sigma_p$ & 0,156 & 0,344 & 0,848 & 1,485\\
 & $r_{0,95}$ & 0,38 & 0,84 & 2,08 & 3,64\\
0,1 & $\sigma_p$ & 0,121 & 0,257 & 0,618 & 1,072\\
0,05 & $\sigma_p$ & 0,096 & 0,194 & 0,453 & 0,777\\
\bottomrule
\end{tabular}
\end{table}
Con $q=0{,}2$ il raggio al 95\,\% supera $1\unita{m}$ dopo $1{,}15\unita{s}$ e $2\unita{m}$ dopo $1{,}95\unita{s}$: oltre un paio di secondi la regione copre gran parte di un'ombra larga uno o due metri e smette di localizzare. L'accumulo dell'errore senza misure è il punto di partenza anche di metodi di inseguimento recenti \cite{cao2023}.

\subsection{Massa e accelerazione non si stimano}\label{sec:sti-massa}

\begin{proposizione}[Non identificabilità della massa]\label{prop:sti-massa}
Nel modello del punto materiale $M\ddot p(t)=u(t)$, con forza $u$ ignota e misure di sola posizione $y_k=p(t_k)+n_k$, la massa $M$ non è identificabile: per ogni $\lambda>0$ le coppie $(M,u)$ e $(\lambda M,\lambda u)$ producono la stessa legge delle misure. \prova{D}
\end{proposizione}
\begin{proof}
$\ddot p=u/M=(\lambda u)/(\lambda M)$: a parità di condizioni iniziali le due coppie danno la stessa traiettoria e quindi la stessa distribuzione di $(y_k)$. Con forza costante in una dimensione, $p(t)=p_0+v_0t+\tfrac{u}{2M}t^2$; il gradiente rispetto a $(u,M)$ è $\tfrac{t^2}{2}\,g$ con $g=(1/M,-u/M^2)^{\top}$, e l'informazione di Fisher $\sigma^{-2}\sum_k\tfrac{t_k^4}{4}\,gg^{\top}$ ha rango 1, con direzione nulla $(u,M)^{\top}$.
\end{proof}
Il calcolo in aritmetica razionale ($u=120\unita{N}$, $M=75\unita{kg}$, 30 misure a $15\unita{Hz}$) dà determinante esattamente 0 \prova{C}. Un sensore di profondità non misura forze: ciò che si osserva è l'ingombro.

\paragraph{Accelerazione.} Con errori di posizione indipendenti di scarto tipo $\sigma_p$, la differenza seconda $\hat a_k=(p_{k+1}-2p_k+p_{k-1})/\Delta^2$ ha scarto tipo $\sigma_a=\sqrt6\,\sigma_p/\Delta^2$, poiché $\mathrm{Var}(\hat a_k)=(1+4+1)\,\sigma_p^2/\Delta^4$ \prova{D}. Con $\sigma_p=3\unita{cm}$ (valore assunto: è la curva nominale a circa $2{,}8\unita{m}$) si ha $\sigma_a=66{,}1\unita{m/s^2}$ a $30\unita{Hz}$, $16{,}5$ a $15\unita{Hz}$ e $7{,}3$ a $10\unita{Hz}$, contro un'accelerazione efficace di $0{,}67\unita{m/s^2}$ nel modello di pedone usato in simulazione \prova{C}. Il limite non è algebrico (lo stato $[p,v,a]$ con misura della sola posizione ha matrice di osservabilità di rango 3) ma di precisione: ai minimi quadrati, per scendere sotto $1\unita{m/s^2}$ serve circa mezzo secondo di accelerazione \emph{costante} ($\sigma_a=0{,}71\unita{m/s^2}$ su $0{,}5\unita{s}$ a $30\unita{Hz}$) \prova{C}. In simulazione, con modello ideale del sensore (un solo modello di pedone che manovra, $15\unita{Hz}$, rumore di 5\,cm, 1500 prove per caso) \prova{S}, un filtro ad accelerazione costante sbaglia l'accelerazione di $0{,}82\unita{m/s^2}$ contro $0{,}67$ di valore efficace vero e, dopo $2\unita{s}$ senza misure, sbaglia la posizione di $2{,}29\unita{m}$ contro $1{,}18\unita{m}$ del modello a velocità quasi costante: lo stato si ferma a posizione e velocità. \emph{Con le classi}, massa e accelerazione entrano come \emph{valori a priori della classe} (intervalli della Tabella~\ref{tab:sti-profili}), marcati come tali, mai come stime dalle misure. Per la massa una fonte è stata consultata per tre sole classi; per le altre il campo resta vuoto.

\subsection{Associazione, porta di validazione e conferma}\label{sec:sti-porta}

\paragraph{Porta e assegnamento.} Una misura $\mathbf z$ è candidata per una traccia se $\nu^{\top}S^{-1}\nu\le\gamma$. Se l'innovazione è gaussiana di covarianza $S$ la forma quadratica è $\chi^2_2$, e la misura vera cade in porta con probabilità $P_G$ per $\gamma=-2\ln(1-P_G)$: $\gamma=9{,}210$ al 99\,\%, valore adottato \prova{D}. L'area della porta, $\pi\gamma\sqrt{\det S}$, cresce con $\sigma_p^2$: senza misure si allarga, e rende possibili gli scambi d'identità (\S\ref{sec:sti-riaggancio}). Fra tracce e gruppi non tagliati si costruisce la matrice dei costi $c_{ij}=\nu_{ij}^{\top}S_i^{-1}\nu_{ij}$ se in porta, vietato altrimenti, e si sceglie l'assegnamento di costo totale minimo con il metodo ungherese \cite{kuhn1955}, prima per le tracce confermate e poi, sui gruppi rimasti, per quelle in attesa di conferma.

\begin{proposizione}[Conferma]\label{prop:sti-conferma}
Una traccia è confermata con 3 rilevamenti in 5 cicli consecutivi. Con rilevamenti indipendenti di probabilità $P_D$ per fotogramma, un oggetto visibile per $n\le5$ fotogrammi è confermato con probabilità $\Pr[\mathrm{Bin}(n,P_D)\ge3]$ e, dal primo rilevamento, la conferma arriva al ciclo $+j$ ($j=2,3,4$) con probabilità $(j-1)\,P_D^2\,(1-P_D)^{j-2}$. \prova{D}
\end{proposizione}
\begin{proof}
Al ciclo $+j$ cade il terzo rilevamento, e fra i $j-1$ cicli precedenti ne cade esattamente uno.
\end{proof}
Con $P_D=0{,}9$ resta non confermato il $27{,}1\,\%$ degli oggetti visti per 3 fotogrammi, il $5{,}2\,\%$ con 4 e lo $0{,}9\,\%$ con 5: le proprietà enunciate per le tracce confermate valgono per chi è stato visto almeno 4 fotogrammi ($0{,}27\unita{s}$ a $15\unita{Hz}$). Il ritardo dal primo rilevamento è di $+2$ cicli ($0{,}133\unita{s}$) con probabilità $0{,}81$; media $0{,}147\unita{s}$ \prova{C}. La conferma protegge dai falsi indipendenti da un ciclo all'altro, non da un ritorno spurio che persiste nella stessa cella. Su profondità simulata dalle schede tecniche, con i gruppi della Definizione~\ref{def:sti-gruppo}, i ritorni spuri sul suolo a grande incidenza hanno prodotto da 7 a 66 tracce false per passaggio \prova{S}; la frequenza di falsi di un sensore reale è non verificata \prova{N}.

\subsection{Il valore di \texorpdfstring{$q$}{q}}\label{sec:sti-q}

\paragraph{Scelta.} In simulazione, con modello ideale del sensore (un solo modello di pedone: salti di velocità a istanti casuali con rilassamento, $15\unita{Hz}$, rumore assunto 5\,cm) \prova{S}, $q=0{,}2\unita{m^2/s^3}$ è il valore per cui la media dell'innovazione normalizzata è $1{,}95$ contro il valore atteso 2. Dopo un tratto senza misure il cerchio al 95\,\% contiene la posizione vera nello $0{,}936/0{,}917/0{,}903/0{,}920$ dei casi a $0{,}5/1/2/3\unita{s}$ (1500 prove per caso); con $q=0{,}01$ la copertura a $2\unita{s}$ scende a $0{,}42$.

\paragraph{Confronto con traiettorie pubbliche di pedoni.} Si usano le traiettorie dell'insieme pubblico THÖR \cite{thor} (persone al chiuso, con posizioni di riferimento accurate): ricampionate a $15\unita{Hz}$, con rumore gaussiano aggiunto di 2\,cm, $1\unita{s}$ di filtro con misure e poi $3\unita{s}$ di sola previsione. Una registrazione dà 227 tratti di 9 persone (taratura), un'altra 219 tratti di 9 persone (prova). L'occlusione è artificiale: questi dati misurano solo il modello del moto, non il rilevamento, l'ombra o l'identità.

\begin{table}[H]\centering\footnotesize
\caption{Traiettorie pubbliche di pedoni, $N=219$ tratti di 9 persone \prova{P}: copertura della regione al 90\,\% (tratti contenuti; intervallo di Wilson, che tratta i tratti come indipendenti). $^{\ast}$ Riga calcolata dopo aver visto i risultati della registrazione di prova: non vale come verifica.}\label{tab:sti-thor}
\begin{tabular}{lccccc}
\toprule
$q$ & a 1 s & a 2 s & a 3 s & $r_{0,95}$ a 1/2/3 s (m) & Errore mediano (m)\\
\midrule
0,2 & 0,977 (214; 0,948--0,990) & 1,000 (219; 0,983--1) & 1,000 (219; 0,983--1) & 0,75 / 1,96 / 3,49 & 0,21 / 0,49 / 0,81\\
0,1$^{\ast}$ & 0,945 (207; 0,907--0,968) & 0,945 & 0,945 & 0,55 / 1,41 / 2,50 & ---\\
0,05 & 0,795 (174; 0,736--0,843) & 0,804 (176; 0,746--0,851) & 0,831 (182; 0,776--0,875) & 0,41 / 1,02 / 1,80 & 0,19 / 0,44 / 0,72\\
\bottomrule
\end{tabular}
\end{table}
Una seconda elaborazione degli stessi dati (altro codice, altra estrazione del rumore) dà $0{,}982/0{,}995/1{,}000$ per $q=0{,}2$, $0{,}922/0{,}950/0{,}941$ per $0{,}1$ e $0{,}836/0{,}822/0{,}840$ per $0{,}05$, con errore mediano $0{,}22/0{,}44/0{,}73\unita{m}$ per $q=0{,}2$. I tratti di una stessa persona non sono indipendenti (per persona la copertura con $q=0{,}05$ va da $0{,}43$ a $1{,}00$; ricampionando le persone l'intervallo è $1{,}27$ volte quello sui tratti), e la stessa regola di scelta dà $q=0{,}05$ su una registrazione e $0{,}1$ sull'altra.

\paragraph{Copertura attesa se il modello fosse esatto.} Se il moto vero seguisse il modello \eqref{eq:sti-modello} con lo stesso $q$ del filtro, con rumori e condizione iniziale gaussiani ed esatti, $P$ sarebbe la covarianza dell'errore, $\mathbf e^{\top}P^{-1}\mathbf e\sim\chi^2_2$, e la regione al livello $\alpha$ conterrebbe la posizione vera con probabilità esattamente $\alpha$ a ogni $\tocc$; su $N$ tratti indipendenti i contenuti sarebbero $\mathrm{Bin}(N,\alpha)$ \prova{D}. Con $N=219$ e $\alpha=0{,}90$ si attendono $197{,}1\pm4{,}4$ tratti (intervallo centrale al 95\,\%: da 188 a 205). Con $q=0{,}2$ se ne osservano 214: averne almeno tanti ha probabilità $8\cdot10^{-6}$ sotto il modello esatto, quindi $q=0{,}2$ \emph{copre troppo}, nel verso prudente. Con $q=0{,}05$ se ne osservano 174 (probabilità $2{,}5\cdot10^{-6}$): copre troppo poco. I 207 con $q=0{,}1$ stanno appena sopra l'intervallo (probabilità $0{,}012$); questa riga è stata calcolata dopo aver visto i risultati della registrazione di prova e non vale come verifica. Il valore calibrato sta fra $0{,}05$ e $0{,}1\unita{m^2/s^3}$; la densità che riproduce gli errori mediani osservati è $q_{\mathrm{eff}}=0{,}058/0{,}050/0{,}044$ a $1/2/3\unita{s}$ \prova{C}. Il sistema mantiene $q=0{,}2$ come scelta prudente, non come taratura: a $3\unita{s}$ il cerchio al 95\,\% ha un raggio di $3{,}5\unita{m}$ e informa poco.

\subsection{Il cerchio di covarianza non è una dichiarazione di precisione in occlusione}\label{sec:sti-cerchio}

In simulazione, con modello ideale del sensore (soggetto puntiforme, $15\unita{Hz}$, rumore di 5\,cm, comportamenti generati da più simulatori scritti indipendentemente, $q=0{,}2$) \prova{S}: al livello nominale 90\,\% il cerchio contiene il soggetto nello $0{,}88$--$0{,}96$ dei casi fino a $5\unita{s}$ ($N=6000$), ma al livello nominale 50\,\% contiene lo $0{,}63$--$0{,}85$. Per comportamento, al 90\,\%: chi prosegue $1{,}00$; chi si ferma $0{,}86$--$1{,}00$; chi svolta $0{,}72$--$0{,}90$; chi torna indietro, a $2\unita{s}$, da $0{,}29$ a $0{,}77$ secondo il simulatore. Il rapporto fra i quantili 90\,\% e 50\,\% dell'errore va da $1{,}8$ a $4{,}5$; per una gaussiana è $\sqrt{\ln10/\ln2}=1{,}82$. Se la regione viene mostrata solo finché copre meno di metà dell'ombra, a $2\unita{s}$ senza osservazione essa contiene chi prosegue nello $0{,}00$--$0{,}06$ dei passaggi (tre disposizioni, 400 passaggi ciascuna), perché a quel punto non viene più emessa.

La spiegazione è nel modello: il filtro assume un'accelerazione bianca, con errore gaussiano a una sola scala, $\sqrt{q\tocc^3/3}$, mentre il moto vero in occlusione è una miscela di comportamenti. Chi prosegue ha un errore molto più piccolo di $r_\alpha$; chi si ferma o torna indietro ha un errore grande, lungo la traiettoria e dietro il punto previsto. Il valore quasi giusto al 90\,\% in media nasce da due errori di segno opposto, e nessun valore di $q$ lo corregge, perché il rapporto fra i quantili di una gaussiana è fisso. Di conseguenza il cerchio del filtro serve come porta di associazione e come centro della previsione; la precisione dichiarata durante l'occlusione è affidata alla regione $\Reg_\alpha(t)$ del \S\ref{sec:occ}, e va misurata per fascia di $\tocc$ e per comportamento (\S\ref{sec:val}).

\emph{Consistenza del filtro.} Nella simulazione del pedone che manovra (rumore 5\,cm, $q=0{,}2$, da 10 a $30\unita{Hz}$) l'innovazione normalizzata $\nu^{\top}S^{-1}\nu$ ha media $1{,}93$--$1{,}97$ (attesa 2), l'errore normalizzato dello stato $3{,}29$--$3{,}41$ (attesa 4) e l'autocorrelazione dell'innovazione a ritardo 1 è $-0{,}02$ \prova{S}. Poiché $R(d)$ non è misurata (\S\ref{sec:sti-gruppi}), la consistenza con un sensore reale è non verificata sperimentalmente \prova{N}: la deciderebbe la stessa statistica su un bersaglio fermo e su uno in moto noto.

\subsection{Riaggancio dell'identità alla ricomparsa}\label{sec:sti-riaggancio}

\paragraph{Regola.} Una traccia confermata che entra in un'ombra $\Omb(t)$ diventa una \emph{traccia in previsione} in quell'ombra. Un nuovo gruppo ne riprende l'identità solo se: (a) cade entro una fascia dall'ombra d'ingresso; (b) dista più della fascia dal bordo della portata e non tocca il bordo del campo; (c) non è tagliato; (d) non è il gruppo dell'ostacolo (non dà il primo ritorno su un azimut dell'ombra); (e) entro la fascia non ha celle non osservate o in ombra estranee all'ombra d'ingresso e alla propria; (f) la traccia in previsione è una sola. Le condizioni (c)--(e) si attivano solo dopo le prove d'ingresso (\S\ref{sec:lim-salute}). Altrimenti l'uscita è «identità ambigua» e nasce una traccia nuova. La fascia vale $0{,}40\unita{m}$ nelle prove su 27 configurazioni e su profondità simulata, $0{,}75\unita{m}$ nelle prove d'insieme con modello ideale del sensore. Al riaggancio il filtro \emph{riparte dalla misura}, con covarianza di posizione $R+(\ell/2)^2I$ ($\ell$ larghezza del gruppo) e scarto tipo di velocità di $2{,}5\unita{m/s}$: su una scena costruita a mano con una sosta di $5\unita{s}$ lo scarto massimo dopo il riaggancio scende da $5{,}00\unita{m}$ a $0{,}17\unita{m}$ ($N=1$) \prova{S}. L'identità ripresa è dichiarata \emph{presunta}, e l'avviso di possibile presenza non dipende da essa: resta finché le celle non sono viste libere (\S\ref{sec:occ}).

\paragraph{Risultati in simulazione, con modello ideale del sensore} (soggetto rigido, rumore gaussiano, bordo d'ombra noto entro $\pm20\unita{cm}$, sensore e ostacolo fermi) \prova{S}. Ricomparse legittime riagganciate, con $P_D=1$: 400 su 400 in ciascuna di 27 configurazioni (limite inferiore di Wilson $99{,}05\,\%$; il conteggio è dedotto dalla quota stampata); con $P_D=0{,}9$, dal $99{,}7$ al 100\,\%. Nelle prove d'insieme, per chi ricompare lungo il percorso: 10\,711 su 10\,757, 9228 su 9240 e 6772 su 6967 in tre disposizioni (limiti inferiori $0{,}994/0{,}998/0{,}968$); per chi svolta o esce lontano dal punto previsto, $0{,}72$--$0{,}80$. Secondo soggetto mai visto che entra dal bordo del campo: 0 riagganci sbagliati su 2400 ($\le0{,}16\,\%$).

\paragraph{Risultati su profondità simulata dalle schede tecniche} \prova{S}. Con i gruppi della Definizione~\ref{def:sti-gruppo} e la regola senza le condizioni (c)--(e), la stessa identità senza riagganci spuri si ottiene in 0 passaggi su 90 (limite superiore $4{,}1\,\%$): la persona si fonde con il pannello (Controesempio~\ref{ctr:sti-misti}), un ritorno spurio allo spigolo dell'ostacolo viene riagganciato a soggetto ancora nascosto, e le tracce false sul suolo rendono falsa la condizione (f). Con $n_{\min}=40$ a $424\times240$ pixel: 41 passaggi su 45 (79,3--96,5\,\%), e 0 su 6 al rumore degradato. Con i gruppi a celle forti (\S\ref{sec:sti-mobile}), le condizioni (a)--(f) e la ripartenza dalla misura: stessa identità presunta in 176 passaggi su 180 (94--99\,\%) a 424 pixel di larghezza e in 174 su 180 (93--98\,\%) a 848, su semi non usati per la taratura, compresa una replica su altri semi (chi passa a $1{,}2$ e a $0{,}5\unita{m/s}$ e chi si ferma $3\unita{s}$ dietro il pannello); i riagganci a persona ancora nascosta scendono da 11 e 6 a 1 e 1 su 140 passaggi.

\begin{riquadro}{Dominio e limiti di validità del riaggancio}
\emph{Dominio in cui il tasso vale} (solo simulazione): rumore nominale; un adulto solo ($\ge1{,}2\unita{m}$) a $0{,}5$--$1{,}6\unita{m/s}$; pannello fermo fra $1{,}5$ e $1{,}8\unita{m}$; persona $0{,}5\unita{m}$ dietro il pannello ed entro $2{,}6\unita{m}$ dalla camera; profili di tre camere; 424--848 pixel di larghezza; posa nota entro meno di $0{,}5^\circ$.
\emph{Fuori}: persona a $2{,}8\unita{m}$, 0 passaggi su 10; a $0{,}3\unita{m}$ dal pannello, 0 e 1 su 10; bambino di $1{,}0\unita{m}$, 2 su 10 (6--51\,\%); errore di posa di $+0{,}5^\circ$, 6 su 10 (31--83\,\%); rumore degradato, 0 su 10. Chi ricompare sul bordo del campo riceve un'identità nuova (10 su 10).
\emph{Secondo soggetto nella stessa ombra.} Un soggetto mai visto che esce dall'ombra in cui ne è entrato un altro riceve l'identità del primo: 2400 casi su 2400 (limite inferiore $99{,}84\,\%$) e 193 su 193 con modello ideale del sensore; 19 su 20 (76--99\,\%) su profondità simulata. Il limite è strutturale: due oggetti che escono dalla stessa fascia non sono distinguibili per posizione. L'identità presunta significa «dalla stessa ombra è uscito un oggetto».
\emph{Fuori dal dominio della regola}: più tracce in previsione, ostacolo mobile, bordo d'ombra errato oltre $-20\unita{cm}$. Su un sensore reale nessuna cifra è verificata \prova{N}.
\end{riquadro}

\section{Occlusione: possibile presenza e previsione con livello dichiarato}\label{sec:occ}

Quando un oggetto seguito scompare dietro un ostacolo il sistema continua a produrre un'uscita per tutta la durata dell'occlusione. L'uscita ha due strati con compiti distinti. Il primo, l'\emph{area d'avviso di possibile presenza}, è un insieme prudente: sotto le ipotesi dichiarate contiene le posizioni che l'oggetto nascosto può occupare e non si riduce se non per osservazione diretta del suolo. Il secondo, la \emph{regione di previsione}, è un sottoinsieme dell'area d'avviso che contiene l'oggetto con una frequenza dichiarata, misurata in simulazione; serve a localizzare, non a proteggere. Trattare l'occlusione come una previsione a breve termine è la linea di \cite{khurana2021}; un filtro lasciato senza misure accumula errore per tutta l'occlusione \cite{cao2023}, e per questo la precisione dichiarata non è letta dalla covarianza del filtro (\S\ref{sec:occ-indicatori}). Nessun risultato di questo capitolo è stato ottenuto su un sensore fisico: sono dimostrazioni in un modello dichiarato \prova{D}, calcoli \prova{C}, simulazioni \prova{S}; ciò che non è stato esercitato è marcato \prova{N}.

\subsection{Notazione e ipotesi}\label{sec:occ-ipotesi}

Sia $X$ la regione del suolo coperta dalla mappa, $B(r)$ il disco chiuso di raggio $r$ e $\dilata(A,r)=A\oplus B(r)=\{x:\operatorname{dist}(x,A)\le r\}$ la dilatazione di $A$. Si indicano con $\hat\Omb$ l'ombra stimata dell'ostacolo e con $\varepsilon_\Omb$ l'errore del suo bordo, con $w$ la larghezza della fascia di guardia, con $W(t)$ l'area d'avviso, con $t_a$ il tempo entro cui un soggetto fuori dall'ombra, in campo e in portata, è riportato come osservato, con $\varepsilon_p(\rho)=\varepsilon_t+\rho\,\delta_\theta$ il margine di posa a distanza $\rho$ (errore di traslazione $\varepsilon_t$, errore angolare $\delta_\theta$ in radianti) e con $c(t)$ il centro del soggetto. Gli stati della cella (libero, occupato osservato, in ombra, non osservato) sono quelli di \S\ref{sec:map} e formano una partizione di $X$.

\begin{ipotesi}[Modello]\label{ip:occ-modello}
Ogni enunciato dichiara quali di queste ipotesi usa.
\begin{enumerate}[label=(\roman*),itemsep=0pt]
\item \emph{Profondità valida uguale a profondità giusta:} un pixel valido con profondità $d$ garantisce che il raggio è privo di superfici fino a $d$, a meno della tolleranza di misura.
\item \emph{Superfici che rimandano segnale:} ogni superficie nel campo e in portata produce un pixel valido; un soggetto che è fuori dall'ombra, in campo e in portata da un tempo $t_a$ è riportato come osservato.
\item \emph{Camera ferma, posa nota entro una tolleranza:} l'errore di posa sposta un punto a distanza $\rho$ di al più $\varepsilon_p(\rho)$.
\item \emph{Velocità limitata:} ogni punto di ogni oggetto si muove con continuità a velocità non superiore a $\vmax$.
\item \emph{Suolo planare.}
\item \emph{Un solo soggetto per ombra.}
\end{enumerate}
\end{ipotesi}

\subsection{Insieme di possibile presenza}\label{sec:occ-presenza}

\begin{definizione}[Insieme di possibile presenza]\label{def:occ-presenza}
$\Pres(0)=X$ e, a ogni fotogramma,
\begin{equation}\label{eq:occ-presenza}
\Pres(t)=\dilata\bigl(\Pres(t-\Delta),\ \vmax\Delta+\varepsilon_p\bigr)\setminus\Oss(t).
\end{equation}
Sulla griglia la dilatazione si calcola per eccesso; ciò che è fuori dalla griglia o non è mai stato valutato appartiene a $\Pres$.
\end{definizione}

La ricorsione \eqref{eq:occ-presenza} è il passo di previsione di un filtro a insiemi con dinamica a velocità limitata: è l'insieme raggiungibile dagli oggetti non osservati, nella linea dei metodi che sovrastimano lo stato degli ostacoli nascosti \cite{orzechowski2018} e della previsione insiemistica dei partecipanti rilevati e potenzialmente nascosti \cite{koschi2021}.

\begin{proposizione}[Contenimento]\label{prop:occ-contenimento}
Sia $\Oss(t)$ l'insieme dei punti del suolo il cui raggio, nel fotogramma $t$, è letto da un pixel valido alla profondità attesa del suolo (lettura per punto \emph{esatta}: margine di lettura nullo nella \eqref{eq:map-margine} e tolleranza di misura nulla nella (i)). Sotto le ipotesi (i), (iv) e (v), a posa esatta ($\varepsilon_p=0$), ogni punto interno all'impronta al suolo di un oggetto al fotogramma $t$ (punto \emph{occupato}) appartiene a $\Pres(t)$. In particolare la posizione di un soggetto entrato in un'ombra appartiene a $\Pres(t)$ per ogni $t$ fino alla riosservazione, senza alcuna ipotesi sul numero dei soggetti né sull'esistenza di una traccia. Con posa nota entro $\varepsilon_p$ (ipotesi (iii), intesa come spostamento nel piano del suolo, come nel Lemma~\ref{lem:map-posa}) l'enunciato vale nel riferimento stimato; nel riferimento vero vale se nella \eqref{eq:occ-presenza} a $\Oss(t)$ si sostituisce l'insieme dei punti $x$ il cui disco di raggio $\varepsilon_p$ è contenuto nel libero stimato $\hat{\Oss}(t)$. \prova{D}
\end{proposizione}
\begin{proof}
Per induzione su $t$. Base: $\Pres(0)=X$. Passo: sia $x$ occupato al tempo $t$. Per (iv) il punto materiale che sta in $x$ si trovava al tempo $t-\Delta$ in un punto $y$ con $|x-y|\le\vmax\Delta$; $y$ era occupato, quindi $y\in\Pres(t-\Delta)$ per ipotesi induttiva, e $x\in\dilata(\Pres(t-\Delta),\vmax\Delta)$. Per (i) e (v) un raggio letto alla profondità attesa del suolo è privo di superfici fino al suolo; il raggio diretto a un punto interno all'impronta incontra invece l'oggetto prima del suolo, quindi un punto occupato non è letto come suolo: $x\notin\Oss(t)$. Dunque $x\in\Pres(t)$. Con posa stimata l'oggetto compare in $x'$ con $|x-x'|\le\varepsilon_p$ e $x'\notin\hat{\Oss}(t)$; se il disco di centro $x$ e raggio $\varepsilon_p$ fosse contenuto in $\hat{\Oss}(t)$ si avrebbe $x'\in\hat{\Oss}(t)$: dunque $x$ non appartiene all'insieme sottratto.
\end{proof}

Con il margine di lettura $m>0$ della \eqref{eq:map-margine} il passo «un punto occupato non è letto come suolo» cade su una striscia dell'impronta: il raggio diretto a un punto che sta a distanza $\delta$, in pianta, dietro la faccia vista dell'oggetto incontra la faccia alla profondità $z_0\,(1-\delta/\rho)$ e, se $\delta\le\rho\,m(z_0)/z_0$, è letto come suolo, qualunque sia l'altezza dell'oggetto (sull'asse del campo, con $k=0{,}004$: $11\unita{cm}$ a $2{,}5\unita{m}$ e $14\unita{cm}$ a $3\unita{m}$ \prova{C}). Nemmeno per la regola della cella (Definizione~\ref{def:map-libero}) il passo è dimostrato: una cella può essere «libero» con meno di metà dell'impronta coperta da un oggetto. La Proposizione è quindi dimostrata per la lettura per punto esatta \prova{D}; con il margine di lettura e con la mappa a celle il contenimento dell'intera impronta non è dimostrato e poggia sulla verifica in simulazione del \S\ref{sec:map-banco} \prova{S}.

\begin{corollario}[«Libero adesso»]\label{cor:occ-libero-adesso}
Sotto le ipotesi della Proposizione~\ref{prop:occ-contenimento}, se l'uscita del fotogramma $t$ è in uso in $[t+\tau,\,t+\tau+\Delta]$, una cella $V$ con $\operatorname{dist}(V,\Pres(t))>\vmax(\tau+\Delta)+\varepsilon_p$ è vuota in tutto $[t,\,t+\tau+\Delta]$. \prova{D}
\end{corollario}
\begin{proof}
Ogni punto occupato in un istante $s$ dell'intervallo dista al più $\vmax(s-t)\le\vmax(\tau+\Delta)$ da un punto occupato al tempo $t$, che per la Proposizione~\ref{prop:occ-contenimento} sta in $\Pres(t)$.
\end{proof}

\begin{proposizione}[Monotonia: nessun rilascio a tempo]\label{prop:occ-monotonia}
$\Pres(t)\supseteq\Pres(t-\Delta)\setminus\Oss(t)$: un punto esce da $\Pres$ solo in un fotogramma in cui è osservato libero. Né un contatore di tempo, né la covarianza di un filtro, né la chiusura di una traccia riducono $\Pres$. Inoltre, per oggetti puntiformi, $\varepsilon_p=0$ e dilatazione esatta (senza l'eccesso di griglia), $\Pres(t)$ è minimo: per ogni $x\in\Pres(t)$ esiste un mondo ammesso dalle ipotesi (i) e (iv), coerente con tutti i fotogrammi, con un oggetto in $x$ al tempo $t$. \prova{D}
\end{proposizione}
\begin{proof}
Poiché $0\in B(r)$, $A\subseteq\dilata(A,r)$; togliendo $\Oss(t)$ ai due membri si ha l'inclusione, e nella ricorsione non compare altro termine sottrattivo. Minimalità, per induzione. Base: $\Pres(0)=X$ e nessun fotogramma vincola l'istante 0. Passo: se $x\in\Pres(t)$ allora $x\notin\Oss(t)$ ed esiste $y\in\Pres(t-\Delta)$ con $|x-y|\le\vmax\Delta$; per ipotesi induttiva una traiettoria ammessa e coerente con i fotogrammi fino a $t-\Delta$ arriva in $y$; la si prolunga con il segmento da $y$ a $x$ percorso in $\Delta$ a velocità non superiore a $\vmax$, coerente con il fotogramma $t$ perché $x\notin\Oss(t)$.
\end{proof}

Dalla minimalità segue che ogni regola che dichiari libera una cella esclusa dal Corollario~\ref{cor:occ-libero-adesso} sbaglia in almeno un mondo ammesso: una probabilità d'occupazione accumulata nel tempo non può estendere il «libero» senza sbagliare in almeno un mondo ammesso.

\emph{Verifiche.} Per calcolo esaustivo su griglie in una dimensione (passo di al più una cella per fotogramma, tutte le successioni di insiemi osservati): 4 celle e 4 fotogrammi, 65\,536 successioni; 5 celle e 3 fotogrammi, 32\,768 successioni; in tutte l'insieme delle posizioni finali delle traiettorie coerenti coincide con $\Pres(t)$, senza violazioni della monotonia \prova{C}. In simulazione su una sezione bidimensionale (2--4 oggetti, 300 fotogrammi per corsa), con $\vmax$ dichiarata pari a quella vera ($2{,}5$ e $1{,}4\unita{m/s}$; misure mancanti 0 e 5\,\%), le corse con almeno una cella «libera adesso» occupata durante l'uso sono 0 su 96 (intervallo di Wilson al 95\,\%: 0--0,038); con la regola «suolo visto, senza margine» sono 24 su 24 (45\,829 celle-fotogramma) \prova{S}. Il controllo esaustivo è in una dimensione: copre la ricorsione, non la geometria del sensore.

\begin{controesempio}[Velocità oltre quella dichiarata]\label{ce:occ-velocita}
Cade (iv). Tre celle e due fotogrammi: al primo sono viste libere le celle 1 e 2, quindi $\Pres=\{3\}$; al secondo non si vede nulla, $\Pres=\{2,3\}$; un oggetto che percorre due celle per fotogramma può essere nella cella 1 (8 successioni su 64; 17\,696 su 65\,536 con 4 celle, contando la violazione all'ultimo fotogramma) \prova{C}. In simulazione, con velocità vera doppia della dichiarata ($0{,}9$ contro $1{,}8\unita{m/s}$ e $0{,}4$ contro $0{,}8\unita{m/s}$), le corse con almeno un errore sono 23 su 24 (0,80--0,99) e 21 su 24 (0,69--0,96) \prova{S}. Se cade (i), cioè con profondità valida e sbagliata, la sottrazione di $\Oss(t)$ toglie punti occupati e il contenimento non vale.
\end{controesempio}

\emph{Costo.} Il «libero adesso» garantito è quasi vuoto a passo d'uomo: nella stessa simulazione resta dichiarabile lo 0,1\,\% del suolo visto libero con $\vmax=2{,}5\unita{m/s}$ e il 3,2\,\% con $1{,}4\unita{m/s}$ \prova{S}. Per questo l'uscita non afferma mai «via libera» ma «suolo visto libero fino a $d$» all'istante d'acquisizione (\S\ref{sec:lim}). L'insieme $\Pres$ è il riferimento teorico; il sistema realizza l'area d'avviso $W$ del \S\ref{sec:occ-avviso}, che nasce solo quando l'ingresso di una traccia in un'ombra è riconosciuto.

\subsection{Ciclo di vita dell'oggetto seguito e ingresso nell'ombra}\label{sec:occ-ciclo}

Ogni traccia confermata (\S\ref{sec:sti}) è in uno e un solo stato, ciascuno con la propria uscita: \emph{osservato} (una misura non tagliata dentro la porta nel ciclo corrente; disco pieno); \emph{visto in parte} (il gruppo è tagliato da un bordo di profondità più vicino o dal bordo del campo; contorno senza velocità); \emph{non confermato} (senza misura in campo aperto: anello vuoto con l'età per $0{,}2\unita{s}$, poi la traccia è chiusa come «non più rilevato»); \emph{previsto} (entrata in un'ombra, con un livello dichiarabile: anello vuoto, regione tratteggiata, etichetta con livello, larghezza ed età); \emph{previsto ampio} (regione senza punto e senza livello; non torna a «previsto»). Non esiste uno stato «perso»: una previsione non esce mai come osservazione, e un disco pieno compare solo con una misura nel ciclo.

\emph{Regola d'ingresso.} Per una traccia con almeno 4 fotogrammi di misura, rimasta senza misura non tagliata nel ciclo, si considera l'ombra di un ostacolo diverso dal gruppo della traccia. L'ingresso scatta (a) subito, se il gruppo è tagliato in quel ciclo; (b) dopo 2 cicli senza misura, se l'ultima misura dista meno di $w$ dall'ombra \emph{oppure} il punto previsto dal filtro sta nell'ombra o entro $w$ da essa. Se l'ultima misura dista meno di $w$ dal bordo del campo o dalla zona cieca vicina alla camera, la traccia è chiusa come «fuori dall'area osservata» e l'avviso resta sul non osservato entro $w$. Dall'ingresso partono l'area d'avviso (\S\ref{sec:occ-avviso}) e la regione di previsione (\S\ref{sec:occ-regione}); l'età «non visto da» si conta dall'ultima misura non tagliata.

\emph{Chiusure in ombra}, elenco chiuso: riaggancio dell'identità; ombra vista libera oltre $w$; tutti i membri della previsione fuori dall'area osservata. Nessun contatore di tempo chiude una traccia. «Identità ambigua» e «affidabilità persa» sono eventi, non chiusure: la traccia resta «previsto ampio».

In simulazione, con modello ideale del sensore, la soglia di 4 fotogrammi è scelta sul conto binomiale del \S\ref{sec:sti-porta} (con 3 fotogrammi la traccia manca nel 27,1\,\% dei casi, con 4 nel 5,2\,\%) \prova{C}; con la sola condizione sull'ultima misura il soggetto resta nascosto senza traccia in previsione in 87 passaggi su 2\,800, con la disgiunzione completa in 9 su 2\,800 (0,3\,\%; Wilson 0,2--0,6\,\%); l'attesa di 2 cicli porta le false entrate da 92 a 7 passaggi su 400 \prova{S}.

\subsection{Area d'avviso e fascia di guardia}\label{sec:occ-avviso}

\begin{definizione}[Area d'avviso]\label{def:occ-avviso}
Quando all'istante $t_0$ una traccia entra nell'ombra di un ostacolo nasce l'area $W(t_0)=\dilata(\hat\Omb,w)$, tagliata alla griglia, con $w=0{,}40\unita{m}$ e senza limite di durata. A ogni fotogramma $W(t)=W(t-\Delta)\setminus L(t)$, dove $L(t)$ è l'insieme delle celle viste libere per 3 fotogrammi consecutivi e distanti più di $w$ dall'ombra corrente dell'ostacolo d'ingresso.
\end{definizione}

L'avviso finisce solo per osservazione, cella per cella: non lo tolgono il riaggancio, l'identità, la chiusura di una traccia, un contatore di tempo o l'operatore. Dopo un riaggancio l'area resta come «possibile presenza residua», senza traccia, distinta a schermo da quella con una traccia in previsione. Più tracce nella stessa ombra danno un solo avviso con contatore ed età. Un'ombra in cui non è entrata alcuna traccia resta «in ombra»: non è mai «libero» e non porta l'avviso. Con l'ostacolo fermo e l'ombra stimata costante nessuna cella di $W(t_0)$ dista più di $w$ dall'ombra corrente: $L(t)\cap W(t_0)=\emptyset$ e $W(t)=W(t_0)$.

\begin{proposizione}[Fascia di guardia]\label{prop:occ-fascia}
Sia l'ostacolo fermo, $\Omb$ l'ombra vera e $\hat\Omb$ quella stimata, con $\Omb\subseteq\dilata(\hat\Omb,\varepsilon_\Omb)$, e sia $W=\dilata(\hat\Omb,w)$. Sotto le ipotesi (ii) e (iv), con la (ii) riferita al centro del soggetto ($c\notin\Omb$), se
\begin{equation}\label{eq:occ-fascia}
w\;\ge\;\varepsilon_\Omb+\vmax\,(t_a+\tau+\Delta),
\end{equation}
allora in ogni istante $s$ in cui è in uso un'uscita, prodotta da un fotogramma $t$, che non riporta il soggetto come osservato, se il soggetto è rimasto nel campo e in portata per tutto $[t-t_a,\,s]$, si ha $c(s)\in W$. Per contenere l'intero corpo di raggio $r_c$ si aggiunge $r_c$ al secondo membro; se la (ii) vale solo quando tutto il corpo è fuori dall'ombra, se ne aggiunge un altro. \prova{D}
\end{proposizione}
\begin{proof}
L'uscita in uso in $s$ viene da un fotogramma $t$ con $s\le t+\tau+\Delta$, che non riporta il soggetto come osservato. Poiché il soggetto è nel campo e in portata in tutto $[t-t_a,t]$ e non è riportato come osservato in $t$, per (ii) esiste $t'\in[t-t_a,t]$ con $c(t')\in\Omb$: altrimenti il centro sarebbe fuori dall'ombra da un tempo $t_a$ e il soggetto sarebbe riportato come osservato. Per (iv), $|c(s)-c(t')|\le\vmax(t_a+\tau+\Delta)$; poiché $c(t')\in\dilata(\hat\Omb,\varepsilon_\Omb)$, $\operatorname{dist}(c(s),\hat\Omb)\le\varepsilon_\Omb+\vmax(t_a+\tau+\Delta)\le w$.
\end{proof}

La velocità coperta dalla fascia è $\vmax=(w-\varepsilon_\Omb)/(t_a+\tau+\Delta)$. La \emph{tenuta} è il termine $\Delta$ della \eqref{eq:occ-fascia}: il ciclo durante il quale un'uscita resta in uso. La Tabella~\ref{tab:occ-fascia} riporta i conti con i termini che ciascuno include; la latenza $\tau$ è stimata, non misurata, e $t_a=0{,}133\unita{s}$ (due cicli) è un valore assunto \prova{N}: con $P_D<1$ nessun tempo finito è un limite certo, e la (ii) è un'idealizzazione (\S\ref{sec:sti-porta}). La Proposizione riguarda l'insieme $W$, non il fatto che l'avviso sia acceso: la premessa è che l'ingresso nell'ombra sia stato riconosciuto.

\begin{table}[H]\centering\small
\caption{Velocità coperta dalla fascia $w=0{,}40\unita{m}$ con $\Delta=1/15\unita{s}$ \prova{C}. L'ultima colonna dice quali termini della \eqref{eq:occ-fascia} il conto include.}\label{tab:occ-fascia}
\begin{tabularx}{\linewidth}{ccccc L}
\toprule
$t_a$ (s) & $\tau$ (s) & $\varepsilon_\Omb$ (m) & Tempo totale (s) & $\vmax$ coperta (m/s) & Che cosa include\\
\midrule
0 & 0,150 & 0 & 0,217 & 1,85 & latenza (qui $0{,}150\unita{s}$) e tenuta, con $t_a=0$\\
0,133 & 0,155 & 0 & 0,288 & meno di 1,4 & $t_a$ e latenza, senza tenuta ($0{,}40/0{,}288$): $2{,}5\cdot0{,}288=0{,}72\unita{m}$, oltre la fascia\\
0,133 & 0,155 & 0 & 0,355 & 1,13 & conto completo: $t_a$ (assunto: due cicli), latenza, tenuta\\
0,133 & 0,155 & 0,20 & 0,355 & 0,56 & conto completo con errore del bordo d'ombra\\
\bottomrule
\end{tabularx}
\end{table}

La condizione \eqref{eq:occ-fascia} è sufficiente, non necessaria: il caso peggiore è un soggetto che si allontana dall'ombra alla velocità massima senza essere riportato come osservato per tutto il tempo $t_a$. Il valore da ritenere per il sistema è quello del conto completo: $1{,}13\unita{m/s}$, e $0{,}56\unita{m/s}$ se il bordo dell'ombra sbaglia di $20\unita{cm}$; il passo d'uomo ($1{,}4\unita{m/s}$) non è coperto nemmeno trascurando la tenuta. La tenuta su un sensore fisico, dove $\varepsilon_\Omb$ e $t_a$ vanno misurati, non è verificata \prova{N}. Fuori dalle ipotesi: ostacolo mobile ($\Omb$ non costante); velocità doppia di quella dichiarata (raddoppia il termine cinematico); soggetto che non rimanda segnale ($t_a$ non è limitato).

\emph{Differenza fra $\Pres$ e $W$.} $\Pres$ non dipende dalle tracce e vale per qualunque numero di soggetti (Proposizione~\ref{prop:occ-contenimento}); $W$ esiste solo se l'ingresso di una traccia è stato riconosciuto ed è dimensionata dalla fascia $w$ anziché dalla ricorsione \eqref{eq:occ-presenza}. La variante in cui anche il riaggancio dell'identità chiude $W$ non è adottata: quella chiusura non è un'osservazione del suolo, ed è corretta solo sotto l'ipotesi (vi). Il caso che la smentisce è riportato nel \S\ref{sec:occ-risultati}; per questo nella regola adottata l'avviso non dipende dall'identità.

\subsection{Regione di previsione con livello dichiarato}\label{sec:occ-regione}

\emph{Costruzione.} Il predittore riceve lo stato della traccia all'ultima misura non tagliata e, a ogni ciclo, la maschera delle celle non viste libere; restituisce a ogni ciclo, fino a una chiusura, un insieme di celle, un punto, l'età, il livello e lo stato. Il predittore principale usa una libreria di 400 percorsi simulati, prodotti dallo stesso generatore di camminate usato per i passaggi di prova, trapiantati sul punto e sulla velocità d'ingresso: 200 proseguono, 200 si fermano a un istante estratto in modo uniforme. Un membro che si trova su una cella vista libera a più di $w_p=0{,}20\unita{m}$ dall'ombra ha il peso moltiplicato per $0{,}1$ a ogni ciclo in cui vi resta: è l'uso dell'informazione negativa nel seguimento \cite{koch2007,wyffels2015}, che qui restringe solo la regione e mai l'area d'avviso. L'insieme $C_\alpha(t)$ è formato dalle celle più dense fino alla quota $\alpha=0{,}9$ del peso, dopo 3 lisciature. Il predittore di ripiego è un disco attorno alla previsione a velocità costante, con raggio da tabella, ritagliato sulla maschera.

\begin{definizione}[Regione mostrata]\label{def:occ-regione}
$\Reg_\alpha(t)=C_\alpha(t)\cap W(t)$ per $\alpha<1$; se questo insieme è vuoto, e in ogni altro caso di «previsto ampio», $\Reg_\alpha(t)=\bigl(\hat\Omb(t)\setminus\Oss(t)\bigr)\cap W(t)$, cioè tutto il non visto libero dell'ombra (\emph{regione di ripiego}). Si pone $\Reg_1(t)=W(t)$.
\end{definizione}

\emph{Punto.} È quello del filtro finché cade sulla maschera; altrimenti resta fermo per 8 cicli e poi si muove verso il baricentro dei membri sulla maschera, di al più $0{,}167\unita{m}$ per ciclo. In «previsto ampio» non è mostrato.

\emph{Livello a schermo.} Il livello non viene dalla covarianza del filtro: è la frequenza con cui la regione contiene il soggetto, misurata in simulazione per comportamento e per durata dell'occlusione (\S\ref{sec:val}), ed è dichiarato con la parola «almeno». Con il predittore principale e un'ombra larga almeno $2\unita{m}$: «almeno 9 su 10» fino a $\tocc=2\unita{s}$, «almeno 85 su 100» fra 2 e $3\unita{s}$; con un'ombra più stretta (pannello da $1\unita{m}$): «almeno 87 su 100» fino a $1\unita{s}$, poi «previsto ampio». Con il predittore di ripiego non è dichiarata alcuna cifra. Una garanzia di copertura ottenuta per taratura è marginale, cioè una media sulla popolazione di taratura \cite{angelopoulos2021}, e la copertura condizionata esatta non è ottenibile senza ipotesi sulla distribuzione \cite{barber2019}: per questo il livello è verificato per comportamento e dichiarato con il suo dominio (\S\ref{sec:occ-risultati}).

\emph{Passaggio a «previsto ampio».} Avviene se la regione supera l'80\,\% delle celle dell'ombra, se scatta la regione di ripiego della Definizione~\ref{def:occ-regione}, oltre l'orizzonte del livello, se la velocità d'ingresso è inferiore a $0{,}3\unita{m/s}$, se il filtro è stato sotto $0{,}3\unita{m/s}$ in un ciclo degli ultimi $2\unita{s}$ prima dell'ultima misura, o se profondità o posa sono degradate. In «previsto ampio» la regione è tutta la maschera dell'ombra, tratteggiata, con l'età; non sparisce e non torna a «previsto». La soglia dell'80\,\% e l'orizzonte di $3\unita{s}$ sono stati scelti dopo aver visto i dati di simulazione e non sono stati riprovati su dati indipendenti dalla scelta.

\begin{figure}[H]\centering
\begin{tikzpicture}[scale=0.92,font=\footnotesize]
\foreach \k/\tit in {0/{(a) appena nascosto: «previsto»},1/{(b) $\tocc=1\unita{s}$: «previsto», livello dichiarato},2/{(c) oltre l'orizzonte: «previsto ampio»}}{
\begin{scope}[xshift=\k*5.45cm]
  \begin{scope}
  \clip (-2.45,-0.15) rectangle (2.45,3.12);
  % campo e portata
  \draw[muted,very thin] (0,0)--(46.5:3) arc[start angle=46.5,end angle=133.5,radius=3]--cycle;
  % ombra stimata
  \fill[black!18] (-1,1.5)--(1,1.5)--(56.31:3) arc[start angle=56.31,end angle=123.69,radius=3]--cycle;
  % area d'avviso
  \draw[accent,dashed,thick] (-1,1.1)--(1,1.1) arc[start angle=-90,end angle=-33.69,radius=0.4]--(1.982,2.252) arc[start angle=48.65,end angle=131.35,radius=3]--(-1.333,1.278) arc[start angle=213.69,end angle=270,radius=0.4]--cycle;
  % ostacolo
  \draw[line width=2.2pt] (-1,1.5)--(1,1.5);
  % camera
  \fill (0,0) circle[radius=0.06];
  % ultima misura
  \fill[muted] (-1.75,2.2) circle[radius=0.07];
  \ifnum\k=0
    \draw[pattern=north east lines,pattern color=accent] (-1.28,2.2) ellipse[x radius=0.27,y radius=0.19];
    \draw[thick,fill=white] (-1.38,2.2) circle[radius=0.07];
  \fi
  \ifnum\k=1
    \draw[pattern=north east lines,pattern color=accent] (-0.65,2.2) ellipse[x radius=0.95,y radius=0.38];
    \draw[thick,fill=white] (-0.35,2.2) circle[radius=0.07];
  \fi
  \ifnum\k=2
    \draw[pattern=north east lines,pattern color=accent] (-1,1.5)--(1,1.5)--(56.31:3) arc[start angle=56.31,end angle=123.69,radius=3]--cycle;
  \fi
  \end{scope}
  \node[below] at (0,-0.12) {\tit};
\end{scope}}
\end{tikzpicture}
\caption{Pianta schematica. In basso la camera; il segmento spesso è l'ostacolo; in grigio l'ombra stimata $\hat\Omb$ fino alla portata; la linea tratteggiata è il bordo dell'area d'avviso $W=\dilata(\hat\Omb,w)$; il disco grigio è l'ultima misura; la zona a righe è la regione $\Reg_{0,9}$; l'anello è il punto mostrato. In (c) la regione è tutto il non visto libero dell'ombra, senza punto; $W$ è la stessa nei tre istanti.}\label{fig:occ-pianta}
\end{figure}

\begin{proposizione}[Annidamento, interferenza monotòna, continuità]\label{prop:occ-annidamento}
(a) $\Reg_\alpha(t)\subseteq\Reg_1(t)=W(t)\subseteq\dilata(\hat\Omb,w)$ per ogni $t$ e ogni $\alpha$. (b) Il predittore non scrive alcuno stato di cella e non può far emettere «osservato» in un ciclo senza misura: a parità di fotogrammi, due predittori diversi danno gli stessi stati delle celle. Gli eventi «osservato» coincidono fino al primo riaggancio; dopo possono differire per l'identità attribuita e per l'istante di conferma, perché la porta del riaggancio usa $\Reg_{0,9}$. (c) Per ogni ingresso in un'ombra, a parità di stato delle tracce in quell'istante, l'area d'avviso con il predittore contiene quella senza: la previsione può accendere o mantenere l'avviso, mai toglierlo. (d) Con l'ostacolo fermo e l'ombra stimata costante, dall'ingresso nell'ombra alla chiusura della traccia, $\Reg_\alpha(t)\neq\emptyset$ a ogni ciclo. \prova{D}
\end{proposizione}
\begin{proof}
(a) Entrambe le espressioni della Definizione~\ref{def:occ-regione} sono sottoinsiemi di $W(t)$, e $W(t)\subseteq W(t_0)$ perché la ricorsione solo sottrae. (b) Lo stato della cella è funzione dei soli fotogrammi di profondità; «osservato» richiede nel ciclo una misura dentro la porta di una traccia confermata; il predittore entra nell'attribuzione delle misure alle tracce solo attraverso la porta del riaggancio. (c) La condizione d'ingresso è una disgiunzione, e il termine che usa il punto previsto è un disgiunto in più: non toglie ingressi. Il rilascio $L(t)$ non dipende dal predittore, e il riaggancio, che usa la regione nella sua porta, non chiude l'avviso. (d) $\hat\Omb\cap\Oss(t)=\emptyset$ per definizione d'ombra e, con l'ostacolo fermo e l'ombra stimata costante, $\hat\Omb\subseteq W(t_0)=W(t)$; quindi la regione di ripiego è $\hat\Omb\neq\emptyset$, e negli altri casi $C_\alpha(t)\cap W(t)\neq\emptyset$.
\end{proof}

L'enunciato «la previsione non interferisce con l'avviso» va quindi letto nella forma (c): l'interferenza esiste, perché la regola d'ingresso usa il punto previsto, ed è monotòna dal lato prudente. La continuità (d) è condizionata all'ingresso: se la traccia è chiusa prima che l'ingresso sia riconosciuto, non c'è né avviso né regione. I punti (b) e (c) sono dimostrati per costruzione della regola descritta; la conformità di un programma alla regola non è oggetto della dimostrazione, e il confronto lungo una successione di occlusioni non è dimostrato, perché lo stato delle tracce dipende dal riaggancio precedente. La proprietà «$\Reg_{0,9}$ non copre celle viste libere oltre $w_p$» è ottenuta con il peso descritto sopra e non è garantita dalla sola regola: è verificata in simulazione (0 cicli; entro la fascia $w_p$ la regione tocca celle viste libere nel 92--100\,\% dei cicli) \prova{S}.

\begin{proposizione}[Livello di una regione ritagliata]\label{prop:occ-ritaglio}
Sia $G$ una misura di probabilità sul suolo, $E_p$ una regione con $G(E_p)\ge p$ e $S$ l'insieme delle celle ammesse, con $G(S)>0$. Il livello della regione ritagliata $E_p\cap S$ rispetto alla distribuzione condizionata a $S$ è
\begin{equation}\label{eq:occ-ritaglio}
\alpha_S=\frac{G(E_p\cap S)}{G(S)}\;\ge\;1-\frac{1-p}{G(S)},
\end{equation}
con uguaglianza se $S\supseteq X\setminus E_p$ e $G(E_p)=p$. \prova{D}
\end{proposizione}
\begin{proof}
$G(E_p\cap S)=G(S)-G(S\setminus E_p)$ e $G(S\setminus E_p)\le G(X\setminus E_p)\le1-p$; si divide per $G(S)$. Se $S\supseteq X\setminus E_p$ la prima disuguaglianza è un'uguaglianza.
\end{proof}

Il risultato non richiede che $G$ sia gaussiana ed è informativo solo se $G(S)>1-p$. Con $p=0{,}95$ e $G(S)=0{,}23$ si ha $\alpha_S\ge1-0{,}05/0{,}23=0{,}78$; con $G(S)=0{,}08$, $\alpha_S\ge0{,}375$ (0 violazioni su 2\,000 casi casuali con gaussiana discreta e $S$ semipiano \prova{C}). Ritagliare un cerchio «al 95\,\%» sulle celle ammesse senza rinormalizzare non conserva dunque il livello dichiarato: per questo la regione è estratta dalla distribuzione dei membri già pesata sulla maschera, e il suo livello è misurato, non dedotto.

\emph{Riaggancio.} Alla ricomparsa, candidata è la prima misura non tagliata che non tocca il bordo del campo; l'identità è riagganciata se la misura è entro $w$ da $\Reg_{0,9}$, entro $0{,}75\unita{m}$ dall'ombra d'ingresso e la traccia è la più recente in previsione per quell'ombra. L'identità riagganciata è marcata «presunta»; negli altri casi si ha l'evento «identità ambigua», con una traccia nuova e la vecchia in «previsto ampio». In entrambi i casi l'avviso resta. È la forma usata nelle prove d'insieme con modello ideale del sensore; le condizioni complete del riaggancio e le cifre sull'identità sono nel \S\ref{sec:sti-riaggancio}.

\subsection{Risultati che spiegano il meccanismo}\label{sec:occ-risultati}

Le tavole complete sono nel \S\ref{sec:val}; qui si riportano i risultati che mostrano perché la regione è costruita così. Le disposizioni di prova sono tre: $G_1$ con un pannello largo $1\unita{m}$, $G_2$ e $G_3$ con un ostacolo largo $2\unita{m}$. I comportamenti «sul percorso» sono cinque: prosegue, rallenta, si ferma, torna indietro, sosta lunga. Tutte le prove sono in simulazione, con modello ideale del sensore (posizioni al suolo con rumore gaussiano di $5\unita{cm}$, rilevamento con probabilità fissa $P_D$), camera ferma e un soggetto per volta salvo dove detto.

\emph{Cerchio del filtro e regione persistente.} La Tabella~\ref{tab:occ-prima-dopo} confronta due costruzioni. Nella prima la regione mostrata è il cerchio d'incertezza del filtro e viene spenta, con la traccia dichiarata persa, quando supera una soglia d'ampiezza; nella seconda è la regione persistente della Definizione~\ref{def:occ-regione}. Con la prima, in $G_1$ lo stato «previsto» dura $0{,}67\unita{s}$ ed è spento prima della ricomparsa in 332 passaggi su 355.

\begin{table}[H]\centering\small
\caption{Continuità e precisione con il cerchio del filtro e con la regione persistente, in simulazione con modello ideale del sensore \prova{S}. Terne: $G_1$ / $G_2$ / $G_3$. Gli intervalli della seconda e della terza riga sono in \S\ref{sec:val}.}\label{tab:occ-prima-dopo}
\begin{tabularx}{\linewidth}{L L L}
\toprule
Misura & Cerchio del filtro, spento oltre soglia & Regione persistente con ripiego\\
\midrule
Cicli senza osservazione con una regione mostrata & 19,0 / 38,1 / 33,3\,\% ($N=105\,301$ / $118\,875$ / $139\,132$ cicli) & 100\,\% ($N=785\,645$ / $342\,552$ / $395\,479$ cicli)\\
Passaggi con almeno un ciclo senza regione & 1\,958 su 1\,958; 1\,991 su 2\,000; 1\,999 su 2\,000 & 0 su 13\,638; 0 su 5\,600; 0 su 5\,597; su semi nuovi 0 su 23\,402 (dopo l'ingresso)\\
Soggetto dentro la regione a $\tocc=2\unita{s}$ & 0,00--0,06 & limite inferiore di Wilson sopra 0,90 nel dominio della Tabella~\ref{tab:occ-dominio}\\
Errore mediano del punto in $G_2$ a 1 / 2 / $3\unita{s}$ ($N=2\,400$ / 2\,077--2\,115 / 1\,674--1\,734) & punto fermo all'ultima misura: 0,97--1,00 / 1,48--1,54 / 1,59--$1{,}66\unita{m}$ & 0,18--0,19 / 0,37--0,39 / 0,54--$0{,}60\unita{m}$\\
\bottomrule
\end{tabularx}
\end{table}

Lo zero della seconda riga vale per costruzione (Proposizione~\ref{prop:occ-annidamento}, d) ed è confermato dal conteggio. I 23\,402 passaggi nascosti della prova su semi nuovi (12 combinazioni: 3 disposizioni, $P_D=0{,}9$ e $0{,}7$, corpo esteso e puntiforme) vengono da 6\,000 camminatori distinti: il limite superiore al 95\,\% è «meno di 7 passaggi su 10\,000» (6,4 se ogni camminatore ha almeno un passaggio con l'ingresso riconosciuto, 6,8 nel caso peggiore), non quello che si otterrebbe trattando i passaggi come indipendenti. La premessa è l'ingresso: nella stessa prova la traccia è chiusa prima dell'ombra in 928 passaggi su 23\,402 (3,97\,\%: 76 su 11\,701 con $P_D=0{,}9$ e 852 su 11\,701 con $P_D=0{,}7$), e lo zero vale per i 22\,474 restanti; accendendo l'avviso anche per le tracce non confermate o chiuse la cui ultima misura dista meno di $0{,}75\unita{m}$ da un'ombra, i passaggi scoperti scendono a 24 su 23\,402 e non a zero (con $P_D=0{,}9$ e corpo esteso 1 su 5\,704; con $P_D=0{,}7$ fino a 5 su 1\,000); la causa del residuo non è determinata. Il punto mostrato cade su celle viste libere dopo il quinto ciclo in 13 / 0 / 23 cicli: la soglia zero non è raggiunta.

\emph{Dominio dichiarato e fuori dominio.} La Tabella~\ref{tab:occ-dominio} riporta la frequenza con cui $\Reg_{0,9}$ contiene il soggetto, finché resta nascosto.

\begin{table}[!htb]\centering\footnotesize
\caption{Quota di istanti con il soggetto dentro $\Reg_{0,9}$, in simulazione con modello ideale del sensore \prova{S}. In alto il dominio dichiarato (corpo esteso, $P_D=0{,}9$, bordo dell'ombra noto, comportamenti sul percorso, semi nuovi): limite inferiore di Wilson al 95\,\% della combinazione peggiore (una combinazione è una terna di disposizione, comportamento e durata dell'occlusione, \S\ref{sec:val-misure}). In basso, fuori dominio: stime puntuali. Terne: $G_1$ / $G_2$ / $G_3$.}\label{tab:occ-dominio}
\begin{tabularx}{\linewidth}{>{\RaggedRight\arraybackslash}p{4.3cm} L}
\toprule
Condizione & Quota dentro $\Reg_{0,9}$\\
\midrule
$G_2$, ostacolo da $2\unita{m}$ & fino a $2\unita{s}$: 0,955 (1\,222 su 1\,265); a $3\unita{s}$: 0,936 (666 su 698)\\
$G_3$, ostacolo da $2\unita{m}$ & fino a $2\unita{s}$: 0,921 (2\,125 su 2\,279); a $3\unita{s}$: 0,899 (1\,754 su 1\,922)\\
$G_2$ e $G_3$ insieme & sopra 0,90 fino a $2\unita{s}$ e sopra 0,85 a $3\unita{s}$ in 30 combinazioni su 30, anche con la correzione di Bonferroni; da 608 a 2\,359 casi per combinazione\\
$G_1$, pannello da $1\unita{m}$ & a $1\unita{s}$: 0,882 (942 su 1\,045), 0,87 sulle cinque combinazioni giudicate insieme; oltre $1\unita{s}$ nessun livello («previsto ampio» nel 92\,\% degli istanti)\\
\midrule
Chi svolta, a 2 e $3\unita{s}$ & corpo esteso: 0,86 / 0,67 / 0,79 e 0,84 / 0,58 / 0,67; soggetto puntiforme, altri semi: 0,54--0,70 e 0,38--0,52\\
Chi esce dal lato lontano, a 2 e $3\unita{s}$ & corpo esteso: 0,82 / 0,61 / 0,79 e 0,76 / 0,57 / 0,70; soggetto puntiforme, altri semi: 0,32--0,67 e 0,13--0,47\\
Ombra stimata più stretta di 10 e $20\unita{cm}$ ($G_2$, prosegue, $1\unita{s}$) & 0,414 e 0,297; con il bordo di profondità riconosciuto (marca «tagliato») a ogni ciclo 0,998 e 1,000; riconosciuto in 8 cicli su 10: 0,742 e 0,757\\
Chi parte da fermo, a $1\unita{s}$ & regione senza ritocchi: 0,146 / 0,049 / 0,043; con la sola soglia di $0{,}3\unita{m/s}$: 0,787 / 0,745 / 0,717\\
Camminatori meno regolari della libreria ($G_3$), a 1 / 2 / $3\unita{s}$ & 323 su 342 / 283 su 339 / 226 su 282 = 0,944 / 0,835 / 0,801 (Wilson a 2 e $3\unita{s}$: 0,792--0,871 e 0,751--0,844); soggetto puntiforme, contati i soli istanti con il soggetto nascosto dentro l'area osservata\\
Oltre l'orizzonte: $5\unita{s}$, $G_1$, prosegue & 174 su 221 = 0,787 (0,729--0,836)\\
\bottomrule
\end{tabularx}
\end{table}

Lettura. (1) Il livello «almeno 9 su 10 fino a $2\unita{s}$» regge con l'ostacolo largo; con il pannello da $1\unita{m}$ cade, e la caduta dipende dal corpo esteso: sugli stessi 400 semi per comportamento, in $G_1$, il soggetto puntiforme dà 0,975 / 0,967 / 0,951 a 1 / 2 / $3\unita{s}$ e quello esteso 0,938 / 0,924 / 0,913. (2) Per chi svolta o esce dal lato lontano le due serie non concordano e la differenza non è spiegata; in entrambe la regione non raggiunge il livello: questi casi sono coperti solo dall'area d'avviso. (3) Il bordo dell'ombra è la grandezza critica: su un sensore il livello non va dichiarato finché non è misurata la quota dei cicli in cui il bordo si riconosce. (4) Chi parte da fermo non è nel dominio: la regola sugli ultimi $2\unita{s}$ lo porta a «previsto ampio», dove il soggetto è dentro la regione in 256 su 256, 342 su 342 e 249 su 249 istanti con il corpo tutto nascosto ($N<400$ per combinazione), ma in 0,832 e 0,818 di tutti gli istanti con il centro nascosto; nessun livello è dichiarato. (5) In nessuna prova la libreria viene da un generatore di camminate e i passaggi da un altro; una libreria di 30 percorsi non basta (0,88 e 0,82 a 2 e $3\unita{s}$). (6) Il centro del soggetto nascosto è dentro $W$ per costruzione finché l'errore dell'ombra non supera $w$ più mezzo corpo: lo «0 su 11,9 milioni di istanti» della Tabella~\ref{tab:val-sicurezza} non è una misura del sistema. Con l'avviso effettivamente acceso come requisito, i passaggi scoperti sono 29 su 95\,399 con la maschera esatta, e 2 / 15 / 1\,704 / 1\,979 su 2\,000 con l'ombra stimata più stretta di 0,10 / 0,20 / 0,45 / $0{,}60\unita{m}$. Su profondità simulata dalle schede tecniche, con il predittore a velocità costante, il soggetto resta nella regione in 36 passaggi su 46 e, con la soglia di punti adattata alla risoluzione, in 83 su 90.

\emph{Secondo soggetto nella stessa ombra (limite di validità).} Il soggetto seguito entra e sosta dietro l'ostacolo; un secondo soggetto, mai visto, parte da dentro l'ombra 0,5--$3\unita{s}$ dopo ed esce da un lato a caso. Alla sua prima misura valida, l'identità del primo gli è attribuita in 3\,027 casi su 3\,041, 3\,493 su 3\,616 e 3\,239 su 3\,546 (0,995 / 0,966 / 0,913; Wilson 0,992--0,997 / 0,960--0,971 / 0,904--0,922), in 2\,124, 2\,713 e 2\,984 dei quali il primo è ancora tutto nascosto; con una porta di riaggancio più stretta ($0{,}40\unita{m}$) le quote sono 0,481 / 0,711 / 0,676 \prova{S}. La porta sulla regione non lo impedisce perché all'uscita del secondo soggetto la regione è già «previsto ampio», cioè tutta l'ombra, nel 98 / 75 / 90\,\% dei casi. Nessuna delle regole provate ripara l'identità: con due soggetti nella stessa ombra non è determinabile, ed è per questo marcata «presunta». Con la variante in cui il riaggancio chiude l'avviso, ogni scambio rilascia l'avviso senza che il primo soggetto sia stato visto (in altre due prove: 193 casi su 193 e 347 su 357, 0,97; Wilson 0,95--0,98). Con la regola adottata i rilasci senza osservazione sono zero per costruzione; la verifica sul caso del secondo soggetto con modello ideale del sensore non è stata eseguita \prova{N}. Su profondità simulata dalle schede tecniche, in una scena costruita con due persone, i fotogrammi con la persona nascosta fuori dall'avviso passano da 282 su 336 a 0 su 336 (0 su 169 con il rumore) \prova{S}. Il costo è dichiarato: dietro un ostacolo fermo l'avviso non si spegne (140 passaggi su 140; 0 celle rilasciabili su 808) e finisce solo quando l'ombra si sposta o l'ostacolo è tolto; che l'avviso termini quando il suolo è rivisto dopo la rimozione dell'ostacolo non è sostenuto: nella realizzazione provata resta acceso dopo $2\unita{s}$ in 10 prove su 10 (72--100\,\%), per una causa non determinata \prova{S}.

L'alternativa con due ipotesi di moto pesate (transito e sosta) è esclusa perché attribuisce una sosta inesistente nel 38--100\,\% dei passaggi; il confronto è nel \S\ref{sec:cos}.

\subsection{Tre indicatori che non misurano ciò che dichiarano}\label{sec:occ-indicatori}

\begin{osservazione}[Il «$\pm\sigma$» in due dimensioni]\label{oss:occ-sigma}
Per $x\sim\mathcal N(\mu,\Sigma)$ in $\Reali^2$, $P\bigl((x-\mu)^{\top}\Sigma^{-1}(x-\mu)\le k^2\bigr)=1-e^{-k^2/2}$: con il cambio di variabile che rende $\Sigma$ l'identità e in coordinate polari, $\int_0^k r\,e^{-r^2/2}\,dr=1-e^{-k^2/2}$. Per $k=1$ si ha $0{,}393$: un'etichetta «posizione $\pm\sigma$» contiene il soggetto nel 39\,\% dei casi anche con un modello esatto (in una dimensione sarebbe il 68,3\,\%); il 95\,\% richiede $k=\sqrt{-2\ln 0{,}05}=2{,}45$ \prova{D}; Monte Carlo con $N=2\cdot10^6$: 0,3932 \prova{C}.
\end{osservazione}

\begin{osservazione}[Probabilità d'intrusione su una distribuzione ritagliata]\label{oss:occ-intrusione}
Sia $\pi_t$ la distribuzione di posizione di un soggetto nascosto, con supporto nelle celle non viste libere $S(t)=X\setminus\Oss(t)$, e $V$ un volume protetto. Se $V\subseteq\Oss(t)$ allora $\pi_t(V)=0$ a ogni ciclo, qualunque sia la distanza del soggetto da $V$, perché $V\cap S(t)=\emptyset$ \prova{D}. Con il soggetto fermo nell'ombra a $9\unita{cm}$ da $V$, il tempo per entrarvi a $1{,}4\unita{m/s}$ è $0{,}09/1{,}4=64\unita{ms}$, meno della latenza. L'indicatore dice dov'è il soggetto secondo il modello, non se può entrare: può solo abbassare un allarme, e non è usato come criterio.
\end{osservazione}

\begin{osservazione}[Criterio d'accettazione non stratificato]\label{oss:occ-criterio}
Si consideri il criterio a due stadi su 30 passaggi (almeno 28 dentro: superato; al più 22: respinto; altrimenti altri 30 e almeno 52 su 60), con 21 passaggi di chi prosegue, 5 di chi sosta e 4 di chi torna. Se i primi due comportamenti sono sempre coperti e il terzo mai, i passaggi dentro sono 26 con certezza: si va al secondo stadio, dove $26+26=52$ supera il criterio con probabilità 1 \prova{C}. Con 30 passaggi per comportamento (almeno 28: superato; al più 23: respinto) una copertura di 0,5 è respinta al primo stadio con probabilità 0,999 e una di 0,90 è respinta per errore al primo stadio con probabilità 0,026. Un criterio che conta l'identità al ritorno, inoltre, misura l'identità e non la presenza del soggetto nella regione: la precisione dichiarata si verifica sulla regione e per comportamento (\S\ref{sec:val}).
\end{osservazione}

\subsection{Dipendenza dalla classe e dominio per classe}\label{sec:occ-classe}

Il classificatore su più classi è un componente richiesto dal sistema (\S\ref{sec:sti}). In questo lavoro non è realizzato né valutato: è simulato, con le etichette assegnate a mano; la sensibilità a etichette errate è misurata solo in prove brevi, su sequenze di etichette generate a caso e fuori dalla catena completa; nessuna classe diversa da «persona a piedi» è stata simulata nella catena completa, e in quella catena nessuna funzione legge ancora un'etichetta. Questo paragrafo descrive come la classe entra nelle tre costruzioni del capitolo e con quale livello di prova.

\emph{Possibile presenza e fascia.} La velocità $\vmax$ dell'ipotesi (iv) deve limitare ogni oggetto che può stare nell'ombra, non quello etichettato (la Proposizione~\ref{prop:occ-contenimento} non suppone una traccia). La fissa quindi lo \emph{scenario} dichiarato all'avvio, cioè l'insieme delle classi ammesse nel sito: $v_{\mathrm{scen}}$ è la massima $\vmax$ delle classi dello scenario e, dalla \eqref{eq:occ-fascia}, $w_{\mathrm{scen}}=\max\{w,\ \varepsilon_\Omb+(v_{\mathrm{scen}}+v_e)(t_a+\tau+\Delta)\}$, con $t_a+\tau+\Delta=0{,}355\unita{s}$ e $v_e$ velocità di chi porta il sensore (zero nel dominio provato). Per ogni ombra con un'area d'avviso la fascia è $w_\Omb(t)=\max\{w_\Omb(t-\Delta),\ w_{\mathrm{scen}},\ 0{,}355\unita{s}\cdot\vmax(c)\}$ per ogni classe $c$ non smentita entrata in quell'ombra; $W=\dilata(\hat\Omb,w_\Omb)$ e il rilascio richiede la distanza $w_\Omb$ dall'ombra corrente. Nel prototipo descritto la fascia è fissa a $w=0{,}40\unita{m}$ e l'uscita dichiara la velocità coperta ($1{,}13\unita{m/s}$) e le classi dello scenario che la superano; la fascia di scenario è una regola di progetto non eseguita nella catena \prova{N}.

\emph{Regione di previsione.} La forma della regione dipende dal modello di moto della classe. Solo per «persona a piedi» esiste una libreria con cifre, e il livello è dichiarato «se l'etichetta è giusta»; per ogni altra classe, per una categoria o per «non classificato» la regione è il disco a velocità costante, senza livello, e così resta fino alla chiusura se l'etichetta cambia o è smentita durante l'occlusione. La forma non tocca mai $W$.

\begin{proposizione}[Un'etichetta non riduce l'avviso né produce un «libero»]\label{prop:occ-etichetta}
Si supponga che (1) lo stato delle celle sia funzione dei soli fotogrammi di profondità, (2) associazione, porta, riaggancio e ciclo di vita non leggano l'etichetta né grandezze che ne dipendono (in particolare la porta del riaggancio usi una regione che non dipende dalla classe), (3) la fascia segua la regola del capoverso «Possibile presenza e fascia». Siano date due esecuzioni sugli stessi fotogrammi e con lo stesso scenario, che differiscono solo nelle etichette (assenti, tardive, errate, variabili durante l'occlusione). Allora le celle «libero» sono identiche; per ogni ingresso in un'ombra l'area d'avviso con etichette contiene, in ogni istante, quella senza etichette; l'insieme di rilascio $L(t)$ con etichette è contenuto in quello senza. \prova{D}
\end{proposizione}
\begin{proof}
Per (1) e (2) stati delle celle, tracce e ingressi coincidono. Senza etichette $w_\Omb=w_{\mathrm{scen}}$; con etichette $w_\Omb(t)\ge w_{\mathrm{scen}}$ in ogni istante e non decresce. La dilatazione è monotòna nel raggio, quindi $W(t_0)$ con etichette contiene quella senza; la condizione di rilascio «oltre $w_\Omb$ dall'ombra corrente» è più restrittiva con $w_\Omb$ maggiore, quindi $L(t)$ con etichette è contenuto in $L(t)$ senza; per induzione su $W(t)=W(t-\Delta)\setminus L(t)$ l'inclusione si conserva.
\end{proof}

Nella costruzione descritta la forma di $\Reg_{0,9}$ dipende dalla classe e la porta del riaggancio usa $\Reg_{0,9}$: la premessa (2) non è soddisfatta alla lettera. L'identità delle celle «libero» vale comunque, per la sola (1); l'inclusione delle aree d'avviso vale per ogni ingresso riconosciuto in entrambe le esecuzioni, non per la successione degli ingressi.

\emph{Livello di prova e limiti.} L'enunciato vale per costruzione della regola. In simulazione è stato esercitato sulla sola fascia, non su mappa, ombra e rilascio: 0 sequenze su 20\,000 con la fascia più stretta che senza etichette (limite superiore al 95\,\% a una coda: 1,5 su 10\,000), e 0 su 20\,000 con un secondo programma \prova{S}; sulla catena completa non è stato eseguito \prova{N}. La regola alternativa «fascia dalla $\vmax$ dell'etichetta, mai sotto $w$» lo viola: con scenario \{persona, carrello elevatore\} la fascia senza etichetta è $4{,}5\cdot0{,}355=1{,}60\unita{m}$; con l'etichetta «persona» scende a $0{,}89\unita{m}$, e un carrello a $4{,}0\unita{m/s}$, che ne chiede 1,42, resta scoperto di $0{,}53\unita{m}$; su 20\,000 sequenze di etichette la violazione compare in una quota di 0,353 e di 0,791 con due generatori diversi \prova{S}. L'enunciato vale rispetto allo scenario dichiarato, non contro uno scenario dichiarato male: un'auto fuori scenario resta scoperta nel 4,8\,\% di 20\,000 uscite dall'ombra anche con la fascia di scenario.

\emph{Effetto di un'etichetta sbagliata sulla regione} (prove brevi: soggetto puntiforme, suolo senza ostacoli, senza ritaglio sulla maschera, un solo generatore; 6\,000 prove) \prova{S}: una persona è dentro la regione costruita per «carrello» in 0,70--0,74 dei casi a 1--$3\unita{s}$ e dentro quella per «treno» in 0,11--0,24; un carrello a $3\unita{m/s}$ è dentro la regione «persona» in 0,02--0,10. Se la fascia fosse presa dall'etichetta, un carrello etichettato «persona» resterebbe scoperto nel 29,1\,\% di 20\,000 uscite. Il controllo di coerenza cinematica (\S\ref{sec:sti}) toglie le etichette incompatibili con il moto ma non le convalida: «persona» a $6\unita{m/s}$ è smentita in 4\,000 tracce su 4\,000, dopo $0{,}47\unita{s}$, cioè $2{,}8\unita{m}$ di percorso; un carrello lento etichettato «persona» è smentito nel 4,1\,\% delle tracce. Per questo nessun livello numerico è dichiarato per una classe senza libreria provata.

\emph{Dominio per classe.} Per un destinatario fermo l'avviso è tempestivo se l'oggetto è rilevato a una distanza non inferiore a $d_{\mathrm{serve}}(v)=v\,(t_{\mathrm{conf}}+\tau+T_r)+C$ (\S\ref{sec:lim}), con tempo di conferma $t_{\mathrm{conf}}=0{,}267\unita{s}$, latenza stimata $\tau=0{,}155\unita{s}$, tempo di reazione $T_r=1{,}0\unita{s}$ (ipotesi) e margine $C=0{,}95\unita{m}$. Con la portata di $3\unita{m}$, avvicinamento frontale e oggetto che non frena, la velocità coperta è $1{,}44\unita{m/s}$ \prova{C}; oltre $7{,}1\unita{m/s}$ i $3\unita{m}$ sono percorsi in meno di $t_{\mathrm{conf}}+\tau=0{,}42\unita{s}$ e l'avviso uscirebbe dopo il contatto. La Tabella~\ref{tab:occ-classi} applica il conto ai profili di classe del \S\ref{sec:sti}.

\begin{table}[H]\centering\footnotesize
\caption{Dominio per classe con camera ferma e portata di $3\unita{m}$ \prova{C}. Fascia a $\vmax$: $0{,}355\unita{s}\cdot\vmax$. Nessuna classe diversa da «persona a piedi» è stata simulata nella catena \prova{N}; le velocità dei profili sono in parte ipotesi (\S\ref{sec:sti}).}\label{tab:occ-classi}
\begin{tabularx}{\linewidth}{>{\RaggedRight\arraybackslash}p{3.4cm} c c c L c}
\toprule
Classe & $v$ tipica / $\vmax$ (m/s) & Fascia (m) & $d_{\mathrm{serve}}$ (m) & Avviso tempestivo entro $3\unita{m}$ & Portata necessaria (m)\\
\midrule
persona a piedi & 1,4 / 2,5 & 0,89 & 2,94 / 4,51 & sì a passo normale ($6\unita{cm}$ di margine); no a passo svelto & 4,5\\
veicolo a guida automatica & 1,0 / 2,0 & 0,71 & 2,37 / 3,79 & sì a velocità tipica; no alla massima & 3,8\\
carrello elevatore & 2,2 / 4,5 & 1,60 & 4,08 / 7,35 & no & 4--7,5\\
persona di corsa & 3,0 / 8,0 & 2,84 & 5,22 / 12,33 & no; alla massima dopo il contatto & 13\\
trattore di piazzale con rimorchi & 1,7 / 8,3 & 2,95 & 3,37 / 12,75 & no; alla massima dopo il contatto & 3,4--13\\
auto, strada urbana & 8,3 / 13,9 & 4,93 & 12,75 / 20,72 & no: dopo il contatto già a velocità tipica & 13--21\\
aereo in rullaggio & 9,0 / 15,4 & 5,47 & 13,75 / 22,85 & no: dopo il contatto & 14--23\\
treno in arrivo in banchina & 11,0 / 16,7 & 5,93 & 16,59 / 24,70 & no: dopo il contatto & 17--25\\
\bottomrule
\end{tabularx}
\end{table}

Sulle 15 classi profilate, l'avviso è tempestivo entro $3\unita{m}$ per 3 a velocità tipica (la terza, l'aereo al traino, non è etichettabile da questo sensore) e per nessuna a velocità massima; l'avviso uscirebbe dopo il contatto per 3 classi a velocità tipica e per 10 alla massima. Con lo scenario «sole persone a piedi» la fascia di scenario è $0{,}89\unita{m}$ e la distanza massima del «suolo visto libero» scende da 2,6 a $2{,}11\unita{m}$; con una classe di $\vmax\ge4{,}8\unita{m/s}$ (11 su 15) la fascia è almeno $1{,}7\unita{m}$ e nessun settore può più dichiarare «suolo visto libero». Per traffico urbano, banchine e rullaggio servirebbero da 13 a $25\unita{m}$ di portata; per treni, tram e aerei il caso rilevante non è l'avvicinamento frontale ma la distanza dal binario o dall'ala, una zona fissa di configurazione esterna a questo lavoro.

\begin{riquadro}{Condizioni di validità del capitolo}
\textbf{Verificato in simulazione, con modello ideale del sensore:} continuità, dato l'ingresso nell'ombra; livello «almeno 9 su 10 fino a $2\unita{s}$, almeno 85 su 100 fino a $3\unita{s}$» con camera e ostacolo fermi, ostacolo largo $2\unita{m}$, un soggetto per ombra visto almeno 4 fotogrammi, ingresso fra 0,3 e $1{,}8\unita{m/s}$, comportamento sul percorso, bordo dell'ombra noto esattamente, libreria dello stesso tipo di camminatore.
\textbf{Fuori dal livello dichiarato} (resta la sola area d'avviso): pannello da $1\unita{m}$ oltre $1\unita{s}$; chi svolta, esce dal lato lontano o parte da fermo; oltre $3\unita{s}$; bordo dell'ombra stimato; secondo soggetto nella stessa ombra; ostacolo mobile; camera in moto; classi diverse da «persona a piedi».
\textbf{Non verificato sperimentalmente:} ogni enunciato su un sensore fisico, in particolare $\varepsilon_\Omb$, $t_a$, $\tau$ e la quota dei cicli in cui il bordo dell'ombra è riconosciuto.
\end{riquadro}

\section{Metodo di valutazione e risultati}\label{sec:val}

Ogni risultato porta condizioni, valore, numerosità $N$, intervallo e \emph{tipo di dato}: (I)~modello ideale del sensore (posizioni al suolo, o immagini di profondità sintetiche, con rumore gaussiano e misure mancanti assegnati); (II)~profondità simulata dai parametri delle schede tecniche di tre camere a stereo attivo \cite{realsense-confronto-camere,realsense-d455,orbbec-gemini-335}; (III)~dati pubblici. Livello di prova: \prova{D} dimostrato nel modello dichiarato, \prova{C} calcolo, \prova{S} simulazione, \prova{P} dati pubblici, \prova{N} non verificato sperimentalmente. \textbf{Nessun risultato è stato ottenuto su un sensore fisico.}

\subsection{Misure e denominatori}\label{sec:val-misure}

L'unità statistica è il \emph{passaggio}: una persona attraversa il campo, scompare dietro un ostacolo fermo e di norma ricompare. Un passaggio è \emph{valido} se la traccia ha ricevuto almeno 4 misure prima di scomparire. I fotogrammi di uno stesso passaggio non sono indipendenti: gli intervalli si calcolano sui passaggi, salvo dove è detto. Si usano tre disposizioni di prova $G_1$, $G_2$, $G_3$, con ombra stimata di $1{,}36$, $7{,}82$ e $14{,}70\unita{m^2}$ ($G_1$: pannello largo $1\unita{m}$; $G_2$ e $G_3$: $2\unita{m}$). Simboli: $k$ eventi su $N$ prove, $\hat p=k/N$; $P_D$ probabilità di rilevamento per fotogramma; $e_\Omega$ errore con segno del bordo dell'ombra stimata (negativo: ombra stimata più stretta della vera).

\begin{definizione}[Misure]\label{def:val-misure}
\begin{enumerate}[label=(\alph*),itemsep=0pt]
\item \emph{Continuità per passaggio}: quota dei passaggi validi con almeno un \emph{vuoto}, cioè un ciclo senza osservazione del soggetto e senza regione $\Reg_\alpha(t)\neq\emptyset$. \emph{Per ciclo}: quota dei cicli senza osservazione in cui una regione è emessa.
\item \emph{Copertura della regione}: per un tempo dall'ultima osservazione $\tocc$ e un comportamento $b$,
\begin{equation}\label{eq:val-copertura}
c(\tocc,b)=\frac{\#\{\text{passaggi di tipo } b \text{ ancora nascosti a } \tocc \text{ con il soggetto in } \Reg_\alpha\}}{\#\{\text{passaggi di tipo } b \text{ ancora nascosti a } \tocc\}} .
\end{equation}
\item \emph{Area relativa}: $|\Reg_\alpha(\tocc)|/|\Omb|$, media sui passaggi ancora nascosti a $\tocc$.
\item \emph{Errore del punto}: distanza al suolo fra il punto mostrato e la posizione vera a $\tocc$ (mediana sui passaggi ancora nascosti).
\item \emph{Sicurezza}: (i) quota degli istanti con il soggetto nascosto il cui centro sta fuori dall'area d'avviso $W(t)$; (ii) numero di celle dichiarate «libero» e contenute per intero nell'ombra vera, sul totale delle celle in ombra vera per fotogramma (celle·fotogramma). La soglia per entrambe è zero.
\item \emph{Identità al ritorno}: quota dei passaggi con ricomparsa in cui la traccia ha lo stesso identificativo di prima e nessuna associazione intermedia a un altro oggetto.
\item \emph{Falsi «in movimento»}: quota dei passaggi in cui un oggetto fermo riceve almeno un'etichetta «in movimento».
\end{enumerate}
\end{definizione}

\emph{Scelta del denominatore.} Il denominatore della~\eqref{eq:val-copertura} comprende tutti i passaggi ancora nascosti, anche quelli in cui il soggetto è uscito dalla portata di $3\unita{m}$ senza essere rivisto: è la scelta prudente. Il denominatore alternativo («ancora nascosti e dentro l'area osservata») è il solo misurabile sul posto con un'occlusione virtuale; dà valori uguali o maggiori e differisce in modo sensibile solo in $G_1$ (stessa cella: 461 su 513 contro 461 su 480). Il livello dichiarato è uno solo, $\alpha=0{,}90$. Una \emph{combinazione}, detta anche \emph{cella} nelle tavole di questo capitolo, è una terna (disposizione, comportamento, $\tocc$) con $\tocc\in\{0{,}5;1;2;3\}\unita{s}$. I comportamenti \emph{sul percorso} sono cinque (prosegue, rallenta, si ferma, torna indietro, sosta lunga); \emph{fuori dal percorso}: svolta, esce dal lato lontano. Una copertura alta con area relativa vicina a 1 dice «è dietro l'ostacolo», non «dove».

\emph{Regola d'esito.} Detti $\ell^-$ e $\ell^+$ gli estremi dell'intervallo, una soglia $s$ su una quota è \emph{rispettata} se $\ell^->s$, \emph{violata} se $\ell^+<s$, altrimenti \emph{non decisa}. Soglie e numerosità delle prove principali sono state fissate prima dei lanci; dove un parametro è stato scelto dopo aver visto i dati lo si dice nella riga.

\subsection{Intervalli}\label{sec:val-intervalli}

Per $k$ eventi su $N$ prove indipendenti si usa l'intervallo di Wilson \cite{wilson1927}, raccomandato per numerosità piccole \cite{brown2001}; con $z=1{,}960$ (95\,\% a due code):
\begin{equation}\label{eq:val-wilson}
\ell^\pm=\frac{\hat p+\dfrac{z^2}{2N}\pm z\sqrt{\dfrac{\hat p(1-\hat p)}{N}+\dfrac{z^2}{4N^2}}}{1+\dfrac{z^2}{N}} .
\end{equation}
\emph{Esempio} ($k=826$, $N=879$): $\hat p=0{,}93970$; centro $(0{,}93970+0{,}002185)/1{,}004370=0{,}93779$; semiampiezza $1{,}960\cdot0{,}008106/1{,}004370=0{,}01582$; intervallo $[0{,}9220;\,0{,}9536]$. \prova{C}

\begin{proposizione}[Zero eventi su $N$]\label{prop:val-zero}
Se in $N$ prove indipendenti con probabilità d'evento $p$ si osservano 0 eventi: (i) il limite superiore esatto a una coda al 95\,\% è $p^+=1-0{,}05^{1/N}$; (ii) $p^+<3/N$; (iii) l'estremo superiore di Wilson a due code vale $z^2/(N+z^2)=3{,}84/(N+3{,}84)$. \prova{D}
\end{proposizione}
\begin{proof}
(i) La probabilità di 0 eventi è $(1-p)^N$, decrescente in $p$; i valori non rifiutati al livello $0{,}05$ sono quelli con $(1-p)^N\ge0{,}05$. (ii) $1-0{,}05^{1/N}=1-e^{-x}$ con $x=2{,}9957/N$, e $1-e^{-x}<x$ per $x>0$. (iii) Con $\hat p=0$ nella~\eqref{eq:val-wilson} centro e semiampiezza coincidono, quindi $\ell^-=0$ e $\ell^+=\frac{z^2/N}{1+z^2/N}$.
\end{proof}
Valori \prova{C}: $N=90$: $3{,}27\,\%$ esatto, $4{,}09\,\%$ Wilson; $N=131$: $2{,}26\,\%$ e $2{,}85\,\%$; $N=400$: $0{,}746\,\%$ e $0{,}951\,\%$; $N=24\,835$: $0{,}0121\,\%$ e $0{,}0155\,\%$. Per simmetria, $N$ successi su $N$ danno il limite inferiore $N/(N+z^2)$: con 400 su 400, $99{,}05\,\%$ ($99{,}25\,\%$ con il limite esatto a una coda \cite{clopper1934}). \emph{Convenzione:} nelle tavole l'intervallo per «0 su $N$» è quello di Wilson a due code; il limite a una coda, più stretto, è citato solo dove è dichiarato.

\begin{proposizione}[Dichiarazioni «in tutte le celle»]\label{prop:val-celle}
Se si giudicano insieme $m$ celle e ogni intervallo ha copertura non inferiore a $1-0{,}05/m$, ed è costruito con $z_m=\Phi^{-1}\bigl(1-\tfrac{0{,}05}{2m}\bigr)$, la probabilità che almeno un intervallo non contenga il proprio valore vero è al più $0{,}05$, qualunque sia la dipendenza fra le celle. \prova{D}
\end{proposizione}
\begin{proof}
Se $E_j$ è l'evento «l'intervallo $j$ non copre», $P(E_j)\le0{,}05/m$ e $P(\bigcup_jE_j)\le\sum_jP(E_j)\le0{,}05$.
\end{proof}
Con l'intervallo di Wilson, la cui copertura è approssimata, il limite $0{,}05$ è a sua volta approssimato.
Valori \prova{C}: $z_{12}=2{,}865$, $z_{30}=3{,}144$, $z_{45}=3{,}261$, $z_{60}=3{,}341$; con 90 su 100 l'intervallo passa da $[0{,}826;0{,}945]$ ($m=1$) a $[0{,}782;0{,}958]$ ($m=12$).

\begin{osservazione}[Unità d'analisi]\label{oss:val-unita}
Un conteggio «0 su $M$» fatto su fotogrammi, istanti o celle·fotogramma non autorizza il limite $3/M$: le unità dello stesso passaggio o della stessa sequenza sono dipendenti. Il limite si calcola sul numero di passaggi, di sequenze o di soggetti simulati distinti; il conteggio fine resta come descrizione dell'estensione del controllo.
\end{osservazione}

\subsection{Numerosità e lettura di una prova sul posto «dentro $k$ su $n$»}\label{sec:val-numerosita}

Per distinguere una copertura $p_0$ da una $p_1$ si usa la prova binomiale esatta a una coda: si rifiuta $H_0\colon p=p_0$ quando $k$ supera (o non raggiunge) una soglia $c$ scelta in modo che il livello della prova sia al più $0{,}05$; la potenza è la probabilità di rifiuto sotto $p_1$. Poiché $c$ è intero, la potenza in funzione di $n$ è a dente di sega: la tabella~\ref{tab:val-numerosita} riporta il primo $n$ che raggiunge la potenza voluta e l'$n$ \emph{stabile}, dal quale la raggiunge ogni $n$ maggiore. Per dimensionare una prova vale l'$n$ stabile.

\begin{table}[H]\centering\footnotesize
\caption{Numerosità per distinguere due coperture (livello $0{,}05$, potenza $0{,}80$) e lettura di un conteggio «dentro $k$ su $n$». \prova{C}}\label{tab:val-numerosita}
\begin{tabular}{@{}lll@{\hspace{2.2em}}lll@{}}
\toprule
$H_0$ contro alternativa & primo $n$ (soglia) & $n$ stabile & osservato & Wilson 95\,\% & esatto 95\,\% (due code)\\
\midrule
0,90 contro 0,80 & 78 ($k\le65$) & \textbf{89} & 9 su 10 & 0,596–0,982 & 0,555–0,997\\
0,80 contro 0,90 & 82 ($k\ge72$) & 94 & 27 su 30 & 0,744–0,965 & 0,735–0,979\\
0,90 contro 0,85 & 270 ($k\le234$) & \textbf{299} & 30 su 30 & 0,886–1 & 0,884–1\\
0,85 contro 0,90 & 281 ($k\ge249$) & 304 & 46 su 50 & 0,812–0,968 & 0,808–0,978\\
0,90 contro 0,95 & 179 ($k\ge168$) & \textbf{203} & 90 su 100 & 0,826–0,945 & 0,824–0,951\\
\bottomrule
\end{tabular}
\end{table}

Servono dunque circa 90 passaggi per cella per distinguere $0{,}90$ da $0{,}80$, circa 300 per distinguerlo da $0{,}85$ e circa 200 da $0{,}95$. \textbf{«27 su 30»} dice «fra $0{,}74$ e $0{,}97$»: non esclude $0{,}80$. Con $n=30$ la prova di $H_0\colon p=0{,}90$ contro $0{,}80$ rifiuta per $k\le23$ e ha potenza $0{,}39$; «30 su 30» dice «almeno $0{,}89$», non «sempre», e con copertura vera $0{,}90$ capita nel $4{,}2\,\%$ dei casi: un conteggio senza errori su trenta prove è più probabile con una regione troppo larga che con una tarata. Una regione che copre $0{,}80$ mostra «almeno 9 su 10» nel $37{,}6\,\%$ delle prove da dieci passaggi; una che copre $0{,}90$ nel $73{,}6\,\%$. In forma operativa \prova{C}: «dentro almeno 82 su 94» esclude $0{,}80$ con potenza $0{,}86$ se il valore vero è $0{,}90$; «almeno 262 su 296» esclude $0{,}85$ con potenza $0{,}83$; per scrivere «almeno 9 su 10» (limite inferiore di Wilson $\ge0{,}90$) servono almeno 96 su 100, con potenza $0{,}44$ se il vero è $0{,}95$ e $0{,}82$ se è $0{,}97$.

\emph{Regole di lettura.} (1)~Si mostra «dentro $k$ su $n$» con l'intervallo, mai una percentuale sola; sotto $n=10$ solo i conteggi. (2)~L'evento contato è «la persona è dentro $\Reg_\alpha$ a $\tocc=1\unita{s}$ e a $2\unita{s}$», non «la prima misura dopo la ricomparsa cade vicino al punto previsto»: le due grandezze divergono (in simulazione, con $e_\Omega=-0{,}2\unita{m}$ e senza margine, la seconda segna 1216 su 1244 $=0{,}977$ mentre la cella peggiore della prima vale $0{,}68$). (3)~I cicli in cui la regione coincide con quasi tutta l'ombra si contano a parte. (4)~Comportamento e $n$ si fissano prima; nessun azzeramento; i passaggi usati per la taratura sono esclusi. (5)~Passaggi della stessa persona si somigliano: con $m$ passaggi a testa e correlazione $\rho_c$ la numerosità efficace è $n/(1+(m-1)\rho_c)$; con $m=10$ e $\rho_c=0{,}2$, 40 passaggi ne valgono 14.

\subsection{Taratura della regione per quantili}\label{sec:val-taratura}

Questo paragrafo analizza la costruzione per quantili, che non è quella adottata (Definizione~\ref{def:occ-regione}): dice che cosa una taratura può garantire e che cosa no.

Per ogni fascia di $\tocc$ si calcola, su $n_c$ passaggi di taratura, un punteggio di non conformità $s_i$ (per esempio la distanza fra la posizione vera e il punto previsto) e si prende come soglia la statistica d'ordine $s_{(k)}$ con $k=\lceil\alpha(n_c+1)\rceil$; $\Reg_\alpha$ è l'insieme delle posizioni con punteggio $\le s_{(k)}$. Se $k>n_c$ la soglia è infinita.

\begin{proposizione}[Garanzia marginale sotto scambiabilità]\label{prop:val-marginale}
Se i punteggi $s_1,\dots,s_{n_c},s_{n_c+1}$ (taratura più il passaggio nuovo) sono scambiabili, allora $P\bigl(s_{n_c+1}\le s_{(k)}\bigr)\ge\alpha$; se inoltre sono quasi certamente distinti, la probabilità vale $k/(n_c+1)\le\alpha+1/(n_c+1)$. \prova{D}
\end{proposizione}
\begin{proof}
Per scambiabilità e assenza di parità il rango di $s_{n_c+1}$ fra gli $n_c+1$ punteggi è uniforme su $\{1,\dots,n_c+1\}$; l'evento equivale a «il rango è $\le k$», di probabilità $k/(n_c+1)\in[\alpha,\,\alpha+1/(n_c+1))$. Con parità l'uguaglianza diventa $\ge$.
\end{proof}
È la copertura \emph{marginale} della previsione conforme a insieme separato \cite{angelopoulos2021}: una media sulla popolazione da cui vengono taratura e prova.

\begin{proposizione}[Copertura realizzata con una taratura data]\label{prop:val-beta}
Se i punteggi sono indipendenti con la stessa distribuzione continua $F$, la copertura condizionata all'insieme di taratura, $C=F(s_{(k)})$, ha distribuzione $\mathrm{Beta}(k,\,n_c+1-k)$. \prova{D}
\end{proposizione}
\begin{proof}
$U_i=F(s_i)$ sono uniformi indipendenti e $F$ è monotona, quindi $F(s_{(k)})=U_{(k)}$, la $k$-esima statistica d'ordine di $n_c$ uniformi.
\end{proof}
Valori al livello $0{,}90$ \prova{C}: con $n_c=30$ ($k=28$) il 5º percentile della copertura realizzata è $0{,}805$; con $n_c=100$ ($k=91$) è $0{,}848$; il primo $n_c$ per cui la copertura realizzata supera $0{,}85$ con probabilità $0{,}95$ è 68 (stabilmente da 110); sotto $n_c=9$ la regione è tutto lo spazio.

\begin{osservazione}[Dove la garanzia non vale]\label{oss:val-nonvale}
La Proposizione~\ref{prop:val-marginale} non dice nulla (a)~sulla copertura \emph{per comportamento}, perché una media $\sum_bw_bc_b\ge\alpha$ non implica $c_b\ge\alpha$ per ogni $b$; (b)~su una popolazione diversa da quella della taratura, perché cade la scambiabilità. In simulazione, con modello ideale del sensore \prova{S}: una soglia unica per fascia di tempo copre $0{,}897$–$0{,}907$ nel complesso, ma solo $0{,}455$ a $2\unita{s}$ e $0{,}307$ a $3\unita{s}$ per chi torna indietro; una taratura su una miscela 90/10/0/0 (prosegue, si ferma, torna, svolta) provata su 50/20/20/10 copre $0{,}72$, $0{,}66$ e $0{,}64$ a 1, 2 e $3\unita{s}$; la regione adottata, con camminatori meno regolari di quelli della taratura (variabilità di velocità 15\,\% anziché 6\,\%, di rotta $8^\circ$ anziché $3^\circ$) in $G_3$, copre 323/342 $=0{,}944$, 283/339 $=0{,}835$ $[0{,}792$–$0{,}871]$ e 226/282 $=0{,}801$ $[0{,}751$–$0{,}844]$ fino a 1, 2 e $3\unita{s}$.
\end{osservazione}

\begin{proposizione}[Soglia del comportamento peggiore]\label{prop:val-suolo}
Se per ogni comportamento $b$ di un insieme finito $B$ si tara una soglia $s^{(b)}$ con la costruzione sopra e si adotta $s^\ast=\max_{b\in B}s^{(b)}$, allora per ogni $b\in B$ la copertura sul comportamento $b$ è almeno $\alpha$, sotto scambiabilità dentro ciascun comportamento. \prova{D}
\end{proposizione}
\begin{proof}
La regione è monotona nella soglia: $\{s\le s^{(b)}\}\subseteq\{s\le s^\ast\}$; si applica la Proposizione~\ref{prop:val-marginale} a ogni $b$.
\end{proof}
Il prezzo è l'area \prova{S}: il disco con il raggio massimo sui comportamenti sul percorso, ritagliato sull'ombra, occupa a 2 e $3\unita{s}$ una quota dell'ombra di $0{,}99$ e $1{,}21$ in $G_2$ e di $0{,}79$ e $1{,}12$ in $G_3$ (sopra 1: la regione comprende il margine attorno all'ombra); la regione adottata, costruita da una libreria di percorsi simulati (\S\ref{sec:occ}), ne occupa $0{,}24$ e $0{,}31$ in $G_2$ e $0{,}16$ e $0{,}27$ in $G_3$.

\subsection{Indipendenza dei campioni: generatore e semi}\label{sec:val-semi}

Le simulazioni usano un generatore a incremento costante: stato $s\in\mathbb Z_{2^{64}}$, a ogni estrazione $s\leftarrow s+G\pmod{2^{64}}$ ($G$ dispari) e uscita $\mu(s)$, con $\mu$ una funzione di mescolamento biiettiva.

\begin{lemma}[Semi contigui]\label{lem:val-semi}
Se lo stato iniziale del seme $\sigma$ è $s_0(\sigma)=\sigma G+c\pmod{2^{64}}$, con $c$ indipendente da $\sigma$, la successione prodotta dal seme $\sigma+1$ è quella del seme $\sigma$ avanzata di un'estrazione: $x_i(\sigma+1)=x_{i+1}(\sigma)$ per ogni $i\ge1$. \prova{D}
\end{lemma}
\begin{proof}
$s_0(\sigma+1)=s_0(\sigma)+G$; lo stato dopo $i$ estrazioni è $s_0+iG$, quindi $x_i(\sigma+1)=\mu(s_0(\sigma)+(i+1)G)=x_{i+1}(\sigma)$.
\end{proof}
\emph{Avvertenza di metodo.} Con lo schema del Lemma~\ref{lem:val-semi} (verificato per calcolo: 1000 valori uguali su 1000 fra due semi contigui; 0 su 1000 con lo stato iniziale $\mu(\sigma G+c)$) i passaggi di un lancio a semi contigui non sono prove indipendenti. Su dieci blocchi da 400 semi contigui in $G_1$, nella cella «prosegue» a $3\unita{s}$ la quota va da $0{,}726$ a $0{,}977$ fra i blocchi e la varianza fra blocchi è $6{,}0$ volte quella binomiale; sulle 20 celle sul percorso fino a $3\unita{s}$ il rapporto vale in media $2{,}67$ ($1{,}62$ con il denominatore «dentro l'area osservata»). Gli intervalli binomiali su quei lanci sono quindi troppo stretti di un fattore fra $1{,}3$ e $2{,}5$; con i semi riparati il rapporto torna a $0{,}94$. Le tavole dicono quali cifre vengono da semi indipendenti; non sono stati rifatti con semi riparati 71\,671 dei 175\,802 istanti del controllo sull'area d'avviso né le prove fuori dal dominio con il soggetto puntiforme, eseguite con un altro programma e un altro schema di semi. Le prove su profondità simulata e su dati pubblici non usano quel generatore. Le prove d'insieme usano semi distanziati e uno stato iniziale mescolato (0 flussi spostati su 1000 coppie di semi vicini), ma riusano gli stessi semi nelle diverse configurazioni: i totali contano perciò i soggetti simulati distinti (Osservazione~\ref{oss:val-unita}). In ogni prova taratura e verifica usano insiemi di semi disgiunti; in nessuna prova la libreria di taratura viene da un generatore di camminate e i passaggi di verifica da un altro.

\subsection{Tavole dei risultati}\label{sec:val-tavole}

Condizioni comuni del tipo (I), salvo diversa indicazione: sensore e ostacolo fermi; un soggetto per ombra; velocità d'ingresso $0{,}3$–$1{,}8\unita{m/s}$, rotta entro $\pm15^\circ$; rumore di misura gaussiano di $5\unita{cm}$ per asse, noto al filtro; $P_D=0{,}9$; $\Delta=1/15\unita{s}$; $q=0{,}2\unita{m^2/s^3}$; maschera dell'ombra esatta. Il soggetto è \emph{puntiforme} nelle prove di base e un \emph{corpo esteso} (disco di $0{,}45\unita{m}$, misura nel baricentro della parte visibile) nelle prove d'insieme. Tipo (II): tre profili di camera (RealSense D435i e D455, Orbbec Gemini 335) a $424\times240$ (salvo indicazione), a due livelli di rumore («nominale» e «limite»), posa nota esattamente. Intervalli: Wilson al 95\,\% a due code.

\subsubsection*{Sicurezza}

I risultati sono nella tabella~\ref{tab:val-sicurezza}. Sulle 240 condizioni comuni a tutte le regole (7200 fotogrammi, delle 300 condizioni della regola adottata) le regole del «libero» non adottate danno violazioni in: 932 fotogrammi la lettura al centro della cella; 98 un solo pixel sul suolo; 6 metà dell'impronta senza cella di guardia. La regola adottata vede libero fino a $3{,}0\unita{m}$ con il 5\,\% di misure mancanti sparse e fino a $1{,}5\unita{m}$ con il 30\,\%; con rumore al livello limite e inclinazione errata di $+0{,}5^\circ$ fino a $2{,}4\unita{m}$; con inclinazione errata di $-1^\circ$ dà violazioni in 150 fotogrammi su 300 (scatola di 10 cm a 2,6–2,9 m). \emph{Che cosa il risultato non dice.} Lo zero della regola adottata è «verificato sull'insieme dichiarato» (scene scelte, non campionate; regola scelta fra quattro dopo un primo lancio), non una frequenza; vale a posa nota entro $0{,}5^\circ$ e per ostacoli che restituiscono profondità. Lo zero su 2\,907\,032 celle·fotogramma vale per la prima serie di sequenze e non si ripete in altre due serie: 4 celle·fotogramma «libero» in ombra vera a 848 px (livello limite, scatola di 15 cm a $2\unita{m}$) e 5 in una serie indipendente (4 al livello limite con misure mancanti 3,5 volte più frequenti, 1 al livello nominale con la scatola a $2{,}8\unita{m}$); non è stabilito se dipendano dalla catena o dal bordo della maschera di riferimento. Lo zero della prima riga è dimostrato solo nelle ipotesi della Proposizione~\ref{prop:occ-fascia} (bordo dell'ombra esatto, ingresso riconosciuto, un solo soggetto per ombra, ostacolo fermo); altrove è un conteggio. L'avviso che sopravvive al riaggancio non si spegne dietro un pannello fermo (140 passaggi su 140) e resta acceso $2\unita{s}$ dopo la rimozione del pannello in 10 prove su 10.

\begin{table}[H]\centering\footnotesize
\caption{Sicurezza: area d'avviso e «libero».}\label{tab:val-sicurezza}
\begin{tabularx}{\textwidth}{@{}>{\RaggedRight}p{3.3cm}L>{\RaggedRight}p{1.25cm}>{\RaggedRight}p{2.75cm}>{\RaggedRight}p{2.45cm}l@{}}
\toprule
Misura & Condizioni & Valore & $N$ & Intervallo 95\,\% & Dato\\
\midrule
Soggetto nascosto fuori da $W$ & $G_1,G_2,G_3$; 7 comportamenti; $\tocc=0{,}5$–$5\unita{s}$ & 0 & 175\,802 istanti, di cui 104\,131 in 24\,835 passaggi a semi indipendenti & per passaggio: 0–0,016\,\% & I \prova{D}\prova{S}\\
idem, prove d'insieme & corpo esteso; anche $e_\Omega=-0{,}2\unita{m}$ & 0 & 11\,946\,289 istanti (dipendenti) & — & I \prova{S}\\
Traccia chiusa prima dell'ombra (avviso non acceso) & senza / con l'avviso mantenuto per le tracce chiuse entro $0{,}75\unita{m}$ dall'ombra & 928 / 24 & 23\,402 passaggi & quote: 3,97\,\% / 0,10\,\% & I \prova{S}\\
idem, con la regola & corpo esteso; $P_D=0{,}9$ / $0{,}7$ & 1 / 17 & 5704 & $<0{,}99$ per mille / 1,9–4,8 per mille & I \prova{S}\\
Avviso mancante & $P_D=1$, $e_\Omega\in\{-0{,}2;0;+0{,}2\}\unita{m}$, margine $0{,}40\unita{m}$ / $P_D=0{,}9$, $e_\Omega=-0{,}2\unita{m}$ & 0 / 13 & 14\,400 / 4800 & 0–0,027\,\% / 0,16–0,46\,\% & I \prova{S}\\
Celle «libero» nell'ombra vera, variante che conserva il «libero» per tre fotogrammi & tre camere; 135 sequenze, di cui 4 con i filtri del costruttore accesi & 0 & 2\,907\,032 celle·fotogramma & per sequenza, sulle 131 senza filtri: 0–2,9\,\% & II \prova{S}\\
idem, altre due serie & nominale e limite, 424 e 848 px; con buchi 3,5 volte più frequenti & 0 e 4; 5 & 3\,588\,836; 3\,992\,186 celle·fot. & — & II \prova{S}\\
«Libero» con ostacolo sopra, regola adottata (\S\ref{sec:map}) & immagine sintetica $320\times240$; 10 scene, 300 condizioni; inclinazione entro $\pm0{,}5^\circ$ & 0 & 9000 fotogrammi, due lanci & per condizione: 0–1,3\,\% & I \prova{S}\\
Soggetto nascosto fuori dall'avviso, prima serie & pannello; configurazione di riferimento & 99 & 3038 fotogrammi & 15 passaggi su 90: 10,4–25,7\,\% & II \prova{S}\\
Soggetto nascosto entro $3\unita{m}$ fuori dall'avviso, seconda serie (424 e 848 px) & configurazione di riferimento / con l'avviso che resta dopo il riaggancio & 180 e 271 / 0 e 0 & 4324 e 4325 fotogrammi (3921 e 3922 senza la disposizione in cui il soggetto si allontana nell'ombra); replica: 0 su 3608 per risoluzione & — & II \prova{S}\\
\bottomrule
\end{tabularx}
\end{table}

\subsubsection*{Continuità}

\begin{table}[H]\centering\footnotesize
\caption{Continuità (tipo I, in simulazione). Terne: $G_1$ / $G_2$ / $G_3$.}\label{tab:val-continuita}
\begin{tabularx}{\textwidth}{@{}>{\RaggedRight}p{3.9cm}L>{\RaggedRight}p{2.5cm}>{\RaggedRight}p{2.9cm}>{\RaggedRight}p{3.3cm}@{}}
\toprule
Misura & Condizioni & Valore & $N$ & Intervallo 95\,\%\\
\midrule
Passaggi con almeno un vuoto, senza regione persistente & 5 comportamenti sul percorso; semi contigui & 1958 / 1991 / 1999 & 1958 / 2000 / 2000 & 99,8–100 / 99,1–99,8 / 99,7–100\,\% (ottimistici)\\
Cicli senza osservazione con una regione, senza regione persistente & idem & 19,0 / 38,1 / 33,3\,\% & 105\,301 / 118\,875 / 139\,132 cicli & non calcolato (cicli dipendenti)\\
Passaggi con almeno un vuoto, con regione persistente & 7 comportamenti; semi indipendenti; soggetto puntiforme & 0 / 0 / 0 & 13\,638 / 5600 / 5597 & 0–0,028 / 0–0,069 / 0–0,069\,\%; insieme 0–0,016\,\%\\
idem, prove d'insieme & corpo esteso; 12 configurazioni $\times$ 5 comportamenti & 0 & 23\,402 passaggi (6000 soggetti distinti) & 0–0,064\,\% (su 6000)\\
\bottomrule
\end{tabularx}
\end{table}

Sull'insieme di tutti i lanci d'insieme i vuoti sono 0 su 217\,861 passaggi (al più 70\,400 soggetti distinti). \prova{D} per costruzione, data la premessa, e \prova{S}. \emph{Che cosa il risultato non dice.} La continuità è garantita solo dopo che la traccia è entrata nell'ombra: se la traccia si chiude prima (tabella~\ref{tab:val-sicurezza}) non c'è regione. Lo zero non riguarda l'identità al ritorno né una camera in moto.

\subsubsection*{Precisione della regione nel dominio dichiarato}

Dominio: sensore e ostacolo fermi; un soggetto per ombra, visto per almeno 4 fotogrammi; ingresso fra $0{,}3$ e $1{,}8\unita{m/s}$; comportamenti sul percorso; rumore noto; libreria di taratura dello stesso tipo di camminatore (400 percorsi simulati); maschera dell'ombra esatta. Tabella~\ref{tab:val-copertura}: copertura~\eqref{eq:val-copertura}, soggetto puntiforme, semi indipendenti; semiampiezza degli intervalli $\le0{,}03$ in ogni cella.

\begin{table}[H]\centering\footnotesize
\caption{Copertura della regione al 90\,\%, soggetto puntiforme, $G_1$ / $G_2$ / $G_3$ (tipo I). \prova{S}}\label{tab:val-copertura}
\begin{tabular}{@{}llll@{}}
\toprule
Comportamento & $\tocc=1\unita{s}$ & $\tocc=2\unita{s}$ & $\tocc=3\unita{s}$\\
\midrule
prosegue & 0,980 / 0,995 / 0,990 & 0,940 / 0,964 / 0,960 & 0,899 / 0,968 / 0,924\\
rallenta & 0,988 / 0,991 / 0,996 & 0,959 / 0,959 / 0,959 & 0,935 / 0,926 / 0,919\\
si ferma & 0,985 / 0,991 / 0,989 & 0,965 / 0,977 / 0,959 & 0,956 / 0,968 / 0,939\\
torna indietro & 0,985 / 0,993 / 0,987 & 0,972 / 0,969 / 0,956 & 0,977 / 0,976 / 0,935\\
sosta lunga & 0,991 / 0,991 / 0,993 & 0,970 / 0,973 / 0,971 & 0,961 / 0,966 / 0,948\\
\midrule
$N$ per cella & 1872–1935 / 800 / 799–800 & 879–1844 / 529–783 / 774–800 & 513–1818 / 284–770 / 669–791\\
area relativa & 0,72 / 0,12 / 0,06 & 0,95 / 0,24 / 0,16 & 1,00 / 0,31 / 0,27\\
\bottomrule
\end{tabular}
\end{table}

\begin{table}[H]\centering\footnotesize
\caption{Riepilogo per soglia (tipo I). $\ell^-$: limite inferiore di Wilson; fra parentesi con la correzione della Proposizione~\ref{prop:val-celle}: $m=45$ nella prima parte ($z_{45}=3{,}261$), $m=30$ nella seconda ($z_{30}=3{,}144$, le 30 celle di $G_2$ e $G_3$ insieme), $m=5$ dove è scritto. \prova{S}}\label{tab:val-soglie}
\begin{tabularx}{\textwidth}{@{}>{\RaggedRight}p{5.4cm}>{\RaggedRight}p{2.6cm}L@{}}
\toprule
Enunciato & Celle sopra soglia & Cella peggiore\\
\midrule
\multicolumn{3}{@{}l}{\emph{Soggetto puntiforme, semi indipendenti}}\\
$\ell^->0{,}90$ fino a $2\unita{s}$ ($\tocc=0{,}5;1;2$) & 45 su 45 (45 con $z_{45}$) & $G_1$, prosegue, 2 s: 826/879 $=0{,}940$ $[0{,}922$–$0{,}954]$ ($0{,}908$)\\
$\ell^->0{,}85$ fino a $3\unita{s}$ & 60 su 60 & $G_1$, prosegue, 3 s: 461/513 $=0{,}899$ $[0{,}869$–$0{,}922]$\\
$\ell^->0{,}90$ a $3\unita{s}$ & 12 su 15; 3 non decise & 461/513; 437/472 $[0{,}899$–$0{,}946]$; 711/774 $[0{,}897$–$0{,}936]$\\
\midrule
\multicolumn{3}{@{}l}{\emph{Corpo esteso, semi nuovi, prove d'insieme (2400 passaggi per comportamento in $G_2$ e $G_3$, 4000 in $G_1$)}}\\
$G_2$: $\ell^->0{,}90$ a 1 e 2 s, $>0{,}85$ a 3 s & 15 su 15 (15) & fino a 2 s: 1222/1265, $\ell^-=0{,}955$ ($0{,}946$); 3 s: 666/698, $0{,}936$ ($0{,}922$)\\
$G_3$: idem & 15 su 15 (15) & fino a 2 s: 2125/2279, $0{,}921$ ($0{,}914$); 3 s: 1754/1922, $0{,}899$ ($0{,}890$)\\
$G_1$, solo istanti in cui l'uscita dichiara il livello & 1 su 5 a 1 s (0 su 5) & 1 s: 942/1045, $0{,}882$ ($0{,}875$ su 5 celle); a 2 e 3 s meno di 400 casi per cella\\
$G_1$, tutti gli istanti & 5 su 5 a 1 s; 4 su 5 a 2 s e a 3 s & 2 s: 1257/1388, $0{,}889$; 3 s: 845/984, $0{,}836$\\
\bottomrule
\end{tabularx}
\end{table}

Errore mediano del punto mostrato in $G_2$ (miscela di comportamenti, $N=2400$ / 2077–2115 / 1674–1734): $0{,}18$–$0{,}19$, $0{,}37$–$0{,}39$ e $0{,}54$–$0{,}60\unita{m}$ a 1, 2 e $3\unita{s}$, contro $0{,}97$–$1{,}00$, $1{,}48$–$1{,}54$ e $1{,}59$–$1{,}66\unita{m}$ di un punto fermo all'ultima misura. \prova{S}

\emph{Forma dichiarabile.} In simulazione, con modello ideale del sensore, ostacolo largo $2\unita{m}$, bordo dell'ombra noto, una persona alla volta che resta sul percorso: la persona nascosta sta nella regione in almeno 9 casi su 10 fino a $2\unita{s}$ dall'ultima misura e in almeno 85 su 100 fino a $3\unita{s}$, finché resta nascosta (30 celle su 30; limiti peggiori $0{,}921$ e $0{,}899$). Con il pannello da $1\unita{m}$ ($G_1$) il risultato del soggetto puntiforme non si ritrova con il corpo esteso: sugli stessi 400 semi per comportamento il primo dà $0{,}975$ / $0{,}967$ / $0{,}951$ a 1, 2 e $3\unita{s}$, il secondo $0{,}938$ / $0{,}924$ / $0{,}913$; lì la soglia di $0{,}90$ non è raggiunta: si dichiara solo «almeno 87 su 100 a $1\unita{s}$» (limite $0{,}875$ sulle cinque celle giudicate insieme), poi l'uscita dice «è dietro l'ostacolo» senza livello (area relativa $0{,}95$–$1{,}00$).

\subsubsection*{Precisione della regione fuori dal dominio dichiarato}

\begin{table}[H]\centering\footnotesize
\caption{Controesempi alla copertura dichiarata. \prova{S}}\label{tab:val-fuori}
\begin{tabularx}{\textwidth}{@{}>{\RaggedRight}p{5.6cm}L>{\RaggedRight}p{2.6cm}l@{}}
\toprule
Che cosa cambia & Copertura a 1 / 2 / 3 s & $N$ & Dato\\
\midrule
chi svolta, soggetto puntiforme ($G_1$ / $G_2$ / $G_3$) & 0,864 / 0,871 / 0,941; 0,567 / 0,541 / 0,698; 0,451 / 0,383 / 0,517 & 538–1911 per cella & I\\
chi esce dal lato lontano, idem & 0,771 / 0,839 / 0,920; 0,320 / 0,448 / 0,665; 0,131 / 0,262 / 0,466 & 725–1933 per cella & I\\
chi svolta; chi esce lontano, corpo esteso & a 2 s: 0,86 / 0,67 / 0,79; 0,82 / 0,61 / 0,79. A 3 s: 0,84 / 0,58 / 0,67; 0,76 / 0,57 / 0,70 & — & I\\
oltre l'orizzonte: $\tocc=5\unita{s}$, $G_1$, prosegue & 174/221 $=0{,}787$ $[0{,}729$–$0{,}836]$ & 221 & I\\
bordo dell'ombra stimato, $e_\Omega=-0{,}2\unita{m}$, corpo esteso, sul percorso, senza la marca «tagliato» ($G_2$; $G_3$) & 0,459 / 0,649 / 0,765; 0,350 / 0,375 / 0,482; per il solo «prosegue», in $G_2$ a $1\unita{s}$: 0,414 con $e_\Omega=-0{,}1\unita{m}$ e 0,297 con $-0{,}2\unita{m}$ & — & I\\
idem, con la marca decisa sul bordo vero a ogni ciclo & 0,996 / 0,974 / 0,968; 0,986 / 0,960 / 0,928 & — & I\\
idem, con la marca che manca 2 cicli su 10 & 0,801 / 0,835 / 0,884; 0,750 / 0,708 / 0,733 & — & I\\
chi parte da fermo accanto al pannello, $G_2$ & 15/141 $=0{,}11$; 21/134 $=0{,}16$; 14/54 $=0{,}26$ & 141; 134; 54 & I\\
profondità simulata, configurazione di riferimento: il soggetto resta sempre nella regione & 36 passaggi su 46 con una previsione (su 90) $=0{,}78$ $[0{,}64$–$0{,}88]$ & 46 & II\\
idem, variante con soglia di punti per cella portata da 3 a 40 & 83 su 90 $=0{,}92$ $[0{,}85$–$0{,}96]$ & 90 & II\\
\bottomrule
\end{tabularx}
\end{table}

\emph{Che cosa il risultato non dice.} Fuori dal dominio (svolte, uscite laterali, bordo dell'ombra stimato male di 10–20 cm, partenza da fermi, popolazione diversa) la copertura scende fino a $0{,}1$–$0{,}8$ e resta solo l'area d'avviso: in nessuna di queste prove un soggetto nascosto esce da $W$. La riparazione del bordo stimato è ottimistica per costruzione (la marca è decisa sul bordo vero) e fragile: sul sensore il livello non si dichiara finché non è misurata la quota dei cicli in cui il bordo si riconosce. Chi è fermo negli ultimi $2\unita{s}$ riceve «previsto ampio» (256/256, 342/342, 249/249 a $1\unita{s}$, meno di 400 casi per cella). Le due righe del tipo (II) si riferiscono a due varianti diverse e non sono confrontabili fra loro; la seconda variante non è adottata perché introduce «libero» nell'ombra vera (2 celle·fotogramma, più 45 al bordo). Nessuna cifra è dichiarata oltre $3\unita{s}$, né per una camera in moto.

\subsubsection*{Identità al ritorno}

\begin{table}[H]\centering\footnotesize
\caption{Identità al ritorno. Terne: $G_1$ / $G_2$ / $G_3$.}\label{tab:val-identita}
\begin{tabularx}{\textwidth}{@{}>{\RaggedRight}p{4.6cm}L>{\RaggedRight}p{3.2cm}>{\RaggedRight}p{2.6cm}l@{}}
\toprule
Misura & Condizioni & Valore su $N$ & Intervallo 95\,\% & Dato\\
\midrule
Stessa identità & una persona, $P_D=1$; 27 configurazioni & 400 su 400 in ciascuna & 99,05–100\,\% & I \prova{S}\\
Stessa identità, riaggancio esteso & corpo esteso, sul percorso & 10\,711/10\,757; 9228/9240; 6772/6967 & $\ell^-$: 0,994 / 0,998 / 0,968 & I \prova{S}\\
Stessa identità, configurazione di riferimento & D435i nominale; pannello, 3 disposizioni $\times$ 30 semi & 0 su 90 & 0–4,1\,\% & II \prova{S}\\
Stessa identità (dichiarata «presunta»), gruppi a celle forti e condizioni (c)–(e) del riaggancio & adulto solo, $0{,}5\unita{m}$ dietro il pannello, entro $2{,}6\unita{m}$; nominale; 424 px / 848 px & 176 su 180 / 174 su 180 & 94–99\,\% / 93–98\,\% & II \prova{S}\\
idem, fuori dal dominio & persona a $0{,}3\unita{m}$ dal pannello; livello limite; oltre $2{,}6\unita{m}$; bambino di $1{,}0\unita{m}$; posa errata di $+0{,}5^\circ$ & 0 e 1 su 10; 0 su 10; 0 su 10; 2 su 10; 6 su 10 & 0–28\,\% (0 su 10); 2–40\,\% (1 su 10); 6–51\,\%; 31–83\,\% & II \prova{S}\\
Identità data a un secondo soggetto, mai visto, uscito dalla stessa ombra & prova di base, 2 disposizioni $\times$ 3 valori di $e_\Omega$ & 2400 su 2400 & 99,84–100\,\% & I \prova{S}\\
idem, prove d'insieme & riaggancio entro $0{,}40\unita{m}$ / riaggancio esteso & 1457/3029; 2567/3609; 2393/3540 / 3027/3041; 3493/3616; 3239/3546 & quote 0,481; 0,711; 0,676 / 0,992–0,997; 0,960–0,971; 0,904–0,922 & I \prova{S}\\
idem & scene costruite a mano / con rumore & 12 su 12 / 19 su 20 & 76–100\,\% / 76–99\,\% & II \prova{S}\\
\bottomrule
\end{tabularx}
\end{table}

\emph{Che cosa il risultato non dice.} Il risultato del modello ideale non si trasferisce alla profondità simulata: con la configurazione di riferimento la sagoma della persona si fonde con quella del pannello in 90 passaggi su 90. La correzione che porta a 176 su 180 è stata scelta dopo aver visto i dati (rapporto col picco $0{,}30$) e i semi nuovi cambiano rumore e percorso, non la scena: è una prova di ripetibilità, non di generalità; il valore è un tasso (circa 97\,\%, su tre comportamenti e due serie di semi), non il superamento di una soglia fissata per ciascun comportamento (in una serie uno dei tre dà 27 su 30). Con due persone nella stessa ombra l'identità passa alla persona sbagliata e, con il modello ideale, nella variante non adottata in cui il riaggancio chiude l'avviso, l'avviso viene rilasciato insieme all'identità (quote della tabella); con la regola adottata i rilasci sono zero per costruzione, non contati. Su profondità simulata, con l'avviso che resta dopo il riaggancio, il primo soggetto resta nell'avviso (0 fotogrammi fuori su 264, 0–1,4\,\%); le due regole non sono state eseguite insieme nello stesso programma \prova{N}. Chi ricompare sul bordo del campo riceve un'identità nuova (10 su 10).

\subsubsection*{Elementi mobili e ostacoli}

\begin{table}[H]\centering\footnotesize
\caption{Etichetta «in movimento», ostacoli fermi, tracce spurie.}\label{tab:val-mobili}
\begin{tabularx}{\textwidth}{@{}>{\RaggedRight}p{4.2cm}L>{\RaggedRight}p{2.1cm}>{\RaggedRight}p{1.9cm}>{\RaggedRight}p{2.3cm}l@{}}
\toprule
Misura & Condizioni & Valore & $N$ & Intervallo 95\,\% & Dato\\
\midrule
Soggetto a $0{,}5\unita{m/s}$ «in movimento» entro $1\unita{s}$ & da $2\unita{m}$; tre curve di rumore / da $3\unita{m}$, in avvicinamento, due curve peggiori & 100\,\% / 85,85\,\%; 26,60\,\% & 2000 per cella & $\ge99{,}81\,\%$ / $\pm1{,}5$; $\pm1{,}9$ punti & I \prova{S}\\
Parete ferma «in movimento» mentre passa una persona & con / senza lo scarto delle misure al bordo di un'ombra & 0 / 400 & 400 & 0–0,95\,\% / 99,05–100\,\% & I \prova{S}\\
Pannello fermo «in movimento» & configurazione di riferimento, 3 disposizioni $\times$ 30 semi & 79 & 90 passaggi & 79,4–93,0\,\% & II \prova{S}\\
idem, gruppi a celle forti & semi nuovi; adulto solo, $0{,}5\unita{m}$ dietro il pannello, entro $2{,}6\unita{m}$, nominale; 424 e 848 px / fuori da questo dominio, in una replica & 0 e 0 / 7 e 7 & 140 passaggi / 10 passaggi & 0–2,7\,\% / 40–89\,\% & II \prova{S}\\
Tracce confermate su suolo vuoto & 9 sequenze: gruppi per componenti connesse / a celle forti (424 e 848 px) & 137 e 206 / 0 e 0 & 270 fotogrammi & — & II \prova{S}\\
Scatola di 15 cm a $2\unita{m}$ vista come ostacolo & 19 sequenze (3 al livello limite) & 190 & 190 fotogrammi & — & II \prova{S}\\
Oggetto piccolo che forma un gruppo inseguibile & cubo di 10 cm a $2\unita{m}$ / scatola di 15 cm a $2{,}8\unita{m}$ & 0 / 45 & 60 fotogrammi & 0–6\,\% / 63–84\,\% & II \prova{S}\\
Ostacolo entro 15 cm dalla posizione vera, casi limite & scatola di 15 cm a $0{,}8\unita{m}$; lastra di vetro; ostacolo sospeso a $2{,}0\unita{m}$ & 0; 0; 7 & 10 fotogrammi ciascuno & — & II \prova{S}\\
\bottomrule
\end{tabularx}
\end{table}

\emph{Che cosa il risultato non dice.} Le righe sugli ostacoli contano fotogrammi di poche sequenze: descrivono casi, non frequenze. Fuori dal dominio dell'identità al ritorno il pannello fermo torna «in movimento» (7 passaggi su 10 alle due risoluzioni). Con i gruppi a celle forti il cubo di 10 cm resta «occupato» nella mappa (60 fotogrammi su 60) ma non è un oggetto seguito: la taglia minima inseguibile è 15 cm entro $2\unita{m}$. Al livello limite del rumore le tracce spurie tornano (105–106 tracce confermate su 3 pavimenti vuoti) e l'identità al ritorno cade (tabella~\ref{tab:val-identita}).

\subsubsection*{Profondità simulata dalle schede tecniche di tre camere}

Errore del bordo dell'ombra stimata, al livello nominale, con un pannello: da $+2{,}0$ a $+14{,}6\unita{cm}$, mediana $+6{,}2\unita{cm}$ (verso il largo), con tutto l'errore entro $\pm20\unita{cm}$ in 115 sequenze su 115 (96,8–100\,\%); al livello limite in 0 su 8 (0–32,4\,\%; errori da $-65$ a $+73\unita{cm}$). \prova{S} \emph{Che cosa il risultato non dice.} Il simulatore è verificato contro il proprio modello, non contro una camera: confrontato con registrazioni pubbliche, sottostima i pixel senza dato (7,8\,\% al livello limite contro 24–34\,\% osservati). Le cifre del tipo (II) al livello nominale sono quindi ottimistiche e non sostituiscono una misura sul sensore.

\subsubsection*{Traiettorie pubbliche di pedoni}

Tracciati di persone in cammino al chiuso, con verità di riferimento per la posizione e senza immagini di profondità \cite{thor}: una registrazione per scegliere $q$ (227 tratti), un'altra per la prova (219 tratti, 9 persone). Si simula un'occlusione e si verifica se la posizione vera sta nella regione al 90\,\% del filtro a velocità quasi costante.

\begin{table}[H]\centering\footnotesize
\caption{Su traiettorie pubbliche di pedoni (tipo III). \prova{P}}\label{tab:val-pubblici}
\begin{tabularx}{\textwidth}{@{}>{\RaggedRight}p{4.4cm}>{\RaggedRight}p{2.6cm}L>{\RaggedRight}p{4.3cm}@{}}
\toprule
Misura & Condizioni & Valore a 1 / 2 / 3 s & Intervallo 95\,\% (tratti come indipendenti)\\
\midrule
Copertura della regione al 90\,\% & $q=0{,}2$; $N=219$ & 0,977 / 1,000 / 1,000 (214, 219, 219 su 219); con un secondo programma 0,982 / 0,995 / 1,000 & 0,948–0,990 / 0,983–1 / 0,983–1\\
idem & $q=0{,}05$; $q=0{,}01$ & 0,795 / 0,804 / 0,831; 0,534 / 0,521 / 0,516 & $\pm0{,}05$; $\pm0{,}07$\\
Errore mediano della previsione & $q=0{,}2$ & 0,21 / 0,49 / 0,81 m (secondo programma: 0,22 / 0,44 / 0,73 m) & —\\
\bottomrule
\end{tabularx}
\end{table}

\emph{Che cosa il risultato non dice.} Il filtro parte da posizioni misurate con precisione millimetrica fino all'ultimo istante: la prova riguarda il modello di moto, non la catena su profondità, e non conferma il livello $0{,}90$ della regione adottata. L'errore mediano di $0{,}2\unita{m}$ a $1\unita{s}$ non è una prestazione del prototipo: con misure peggiori è $0{,}32$--$0{,}52\unita{m}$, e per chi si ferma è circa $1\unita{m}$ a $2\unita{s}$ ($N=12$). I tratti di una stessa persona sono correlati (da 7 a 42 a testa): la numerosità efficace sta fra 9 e 219. Con $q=0{,}2$ la regione copre più del livello dichiarato, cioè è più larga del necessario.

\subsubsection*{Classificazione e profili di classe}

Il classificatore è un componente richiesto dal sistema (\S\ref{sec:sti}); in questo lavoro \textbf{non è realizzato né valutato} \prova{N}: nelle prove le etichette sono assegnate a mano, e nessuna classe diversa da «persona» è stata simulata nella catena completa. Le prove fatte riguardano le regole che usano l'etichetta, su modelli ridotti: (i)~«un'etichetta non riduce l'avviso né produce un libero» vale per costruzione (\S\ref{sec:occ}) e, in una simulazione della sola fascia di guardia, dà 0 violazioni su 20\,000 sequenze di etichette casuali (limite superiore a una coda $1{,}5$ su 10\,000) e 0 su 20\,000 in una seconda simulazione scritta in modo indipendente \prova{S}; una regola alternativa, scartata, viola in $0{,}353$ e $0{,}791$ delle sequenze; (ii)~se l'etichetta potesse restringere la fascia, un carrello etichettato «pedone» resterebbe scoperto nel 29,1\,\% di 20\,000 uscite \prova{S}; (iii)~il controllo di coerenza cinematica non protegge da un'etichetta lenta finché l'oggetto è fermo: 0 scatti su 2000 da fermo; con partenza a $3{,}0\unita{m/s^2}$ lo scatto precede il contatto in 2 casi su 2000 \prova{S}; (iv)~con l'etichetta sbagliata la persona sta nella regione di un'altra classe in $0{,}70$–$0{,}74$ (carrello) e $0{,}11$–$0{,}24$ (treno) dei casi, e un carrello a $3\unita{m/s}$ nella regione «pedone» in $0{,}02$–$0{,}10$ (6000 prove) \prova{S}: per questo nessun livello è dichiarato per classi diverse da «persona». La regola (i) non è mai stata eseguita sulla catena completa \prova{N}.

\subsection{Limiti di validità comuni}\label{sec:val-limiti}

Il costo per ciclo (tempi e memoria) è riportato in \S\ref{sec:cos}.
\begin{riquadro}{Limiti di validità}
Tutte le frequenze di questo capitolo sono condizionate al generatore che le ha prodotte: modello ideale del sensore (I), profondità simulata da scheda tecnica a posa nota (II), traiettorie pubbliche prive di profondità (III). La continuità e la copertura della regione sono misurate solo con il tipo (I): su profondità simulata la regola della regione persistente non è stata eseguita (\S\ref{sec:lim-limiti}). La copertura della regione, l'errore del bordo dell'ombra, il rumore del centro, l'identità al ritorno e la latenza \textbf{non sono verificati sperimentalmente su un sensore} \prova{N}. Le riconte dei risultati condividono con le simulazioni i programmi di simulazione e i loro modelli: rilevano errori di somma, di denominatore e d'intervallo, non un difetto del simulatore. Sul posto vale il conteggio «dentro $k$ su $n$» letto con le regole di \S\ref{sec:val-numerosita}: con 10–50 passaggi per riga esso può smentire una cifra, non confermarla al centesimo.
\end{riquadro}

\section{Costo di calcolo, alternative considerate e lavori correlati}\label{sec:cos}

Questo capitolo dà il costo di calcolo del sistema e delle rappresentazioni dello spazio con cui è stato confrontato, elenca le alternative considerate ed escluse con il motivo tecnico e con il controesempio o la misura che le esclude, e colloca il lavoro rispetto alla letteratura. Tutte le cifre di tempo e di memoria sono misure fatte su un calcolatore portatile, su profondità sintetica o su una registrazione pubblica: nessuna viene da un sensore fisico né dalla scheda di calcolo di bordo. La scelta della rappresentazione è decisa dalla sicurezza dello stato «libero» (\S\ref{sec:map}) prima che dal costo.

\subsection{Rappresentazioni dello spazio e loro complessità}\label{sec:cos-rappresentazioni}

Simboli aggiunti alla notazione comune: $W\times H$ pixel dell'immagine di profondità; $k$ passo di sottocampionamento dei pixel; $N_r\le WH/k^2$ raggi (pixel validi letti) per fotogramma; $L_x\times L_z\times L_y$ regione rappresentata (pianta e altezza); $s$ lato del voxel; $n_x,n_z,n_y$ numero di voxel per lato; $N_v=n_xn_yn_z$; $c$ byte per voxel; $\beta$ lato di un blocco, in voxel; $B$ blocchi allocati; $K$ passi di raggio per fotogramma (voxel attraversati, sommati sui raggi); $N_c$ celle della griglia polare; $n_T$ oggetti seguiti in un ciclo; $T=\Delta=1/15\unita{s}=66{,}7\unita{ms}$ periodo del ciclo.

Sono state considerate sette rappresentazioni. (i)~La \emph{griglia polare del suolo a quattro stati} (\S\ref{sec:map}): celle di $1^\circ\times0{,}1\unita{m}$ fino a $3\unita{m}$, $N_c=87\cdot30=2\,610$ con campo orizzontale di $87^\circ$; ogni pixel valido è assegnato alla sua cella con tabelle precalcolate e nessun raggio viene percorso. (ii)~La \emph{lettura diretta nell'immagine di profondità}: nessuna mappa; lo stato di un punto si legge per proiezione sugli ultimi fotogrammi (\S\ref{sec:map}). (iii)~La \emph{griglia densa di voxel}: ogni raggio è percorso voxel per voxel fino al ritorno. (iv)~I \emph{voxel sparsi con tabella hash a blocchi} di $\beta^3=8^3$ voxel, allocati alla prima scrittura, schema proposto per la fusione volumetrica in tempo reale su hardware grafico comune \cite{niessner2013}. (v)~L'\emph{octree}, di profondità 6 a $s=10\unita{cm}$ e 7 a $5\unita{cm}$, 8 byte per nodo, con o senza fusione dei figli uguali. (vi)~La \emph{mappa a fasce di quota}: griglia in pianta con tre fasce per colonna (0,15--0,60 m, 0,60--1,60 m, 1,60--2,0 m), e la variante a colonne di bit (30 livelli da 10 cm, 12 byte per colonna). (vii)~La \emph{griglia di occupazione bayesiana}, che è il contenuto di un voxel (logaritmo delle quote di occupazione) e si appoggia a una delle tre strutture precedenti. La tabella~\ref{tab:cos-complessita} ne riassume la complessità.

\begin{table}[H]\centering\small
\caption{Complessità delle rappresentazioni considerate.}\label{tab:cos-complessita}
\begin{tabular}{@{}llll@{}}\toprule
Rappresentazione & Memoria & Aggiornamento per fotogramma & Lettura di un punto\\\midrule
Griglia polare del suolo & $O(N_c)$ & $O(N_r+N_c)$ & $O(1)$\\
Lettura nell'immagine & $O(WH)$ per fotogramma tenuto & $O(WH)$ & $O(1)$: nove pixel\\
Griglia densa di voxel & $c\,N_v=c\,L_xL_yL_z/s^3$ & $O(K)$, $K=O(N_rL/s)$ & $O(1)$\\
Voxel sparsi con hash & $c\,\beta^3B$ più la tabella & $O(K)$ atteso, costante maggiore & $O(1)$ atteso\\
Octree & 8 byte per nodo & $O(K\log_2 n)$ & $O(\log_2 n)$\\
Fasce di quota & $O(n_xn_z)$ per fascia & $O(K)$ & $O(1)$\\
Occupazione bayesiana & quella del supporto & quella del supporto & $O(1)$\\\bottomrule
\end{tabular}
\end{table}

\begin{lemma}[Numero di voxel e memoria]\label{lem:cos-voxel}
Per la regione $6\times4\times3\unita{m}$ si ha $N_v=60\cdot40\cdot30=72\,000$ a $s=10\unita{cm}$ e $N_v=120\cdot80\cdot60=576\,000$ a $s=5\unita{cm}$; la memoria densa $c\,N_v$ vale $288\,000$ e $2\,304\,000$ byte con $c=4$ byte, il doppio con $c=8$ byte. I blocchi da $8^3$ voxel sono al più $160$ e $1\,200$. \prova{C}
\end{lemma}

\begin{lemma}[Tetto dei passi di raggio]\label{lem:cos-passi}
In una griglia di $n_x\times n_y\times n_z$ voxel ogni raggio attraversa al più $n_x+n_y+n_z-2$ voxel; quindi $K\le N_r\,(n_x+n_y+n_z-2)$, cioè $128\,N_r$ a $10\unita{cm}$ e $258\,N_r$ a $5\unita{cm}$ nella regione del lemma~\ref{lem:cos-voxel}, qualunque sia il numero di oggetti nella scena. \prova{D}
\end{lemma}
\begin{proof}
Passando da un voxel al successivo il raggio attraversa esattamente una superficie di separazione della griglia. Un segmento di retta attraversa ciascuna di queste superfici al più una volta, e quelle interne alla regione sono $n_x-1$, $n_y-1$ e $n_z-1$ nelle tre direzioni: i voxel toccati sono al più $1+(n_x-1)+(n_y-1)+(n_z-1)$. Un oggetto in più può solo accorciare un raggio.
\end{proof}

Nelle misure del \S\ref{sec:cos-misure} i passi per raggio sono in media 46,8 a 10 cm e 94,1 a 5 cm: dimezzando $s$ il costo raddoppia, mentre $N_v$ cresce di otto volte. Il caso peggiore per il tempo è la scena vuota, non quella affollata.

\begin{lemma}[Limite del guadagno di una struttura sparsa a tre stati]\label{lem:cos-sparsa}
Se «non allocato» deve valere «non osservato» (lo impone il controesempio della lettura a due stati, \S\ref{sec:map}), ogni voxel dichiarabile libero va allocato, e la memoria cresce con il volume visto libero, non con le superfici. Entro una portata $R$ il sensore vede al più la piramide
\begin{equation}\label{eq:cos-piramide}
V_p=\tfrac43\,\tan\!\Big(\tfrac{87^\circ}{2}\Big)\tan\!\Big(\tfrac{58^\circ}{2}\Big)R^3=\tfrac43\cdot0{,}9490\cdot0{,}5543\cdot27\unita{m^3}=18{,}9\unita{m^3}\qquad(R=3\unita{m}),
\end{equation}
cioè il 26,3\,\% dei $72\unita{m^3}$ della regione: nel caso peggiore, con tutta la piramide vista libera, il rapporto fra memoria densa e memoria sparsa è $72/18{,}9=3{,}8$, prima del costo dell'indice; è il guadagno su cui si può contare, perché con ostacoli in scena il volume visto libero è minore e il rapporto maggiore. \prova{C}
\end{lemma}

\subsection{Misure di tempo e memoria}\label{sec:cos-misure}

\textbf{Condizioni di misura.} Calcolatore portatile con processore AMD Ryzen AI 9 HX 370 (quattro nuclei fino a 5,16 GHz, otto fino a 3,29 GHz), Linux, programmi in C compilati con gcc 16.2.1 e \texttt{-O2}, un solo filo d'esecuzione fissato su un nucleo. Profondità \emph{sintetica} $640\times480$, campo $87^\circ\times58^\circ$, $h=1{,}2\unita{m}$, $\theta=15^\circ$, valida fino a $6\unita{m}$, rumore $\sigma_z=0{,}0025\,d^2$ (metri), 1,2\,\% di pixel non validi; regione $6\times4\times3\unita{m}$, $c=4$ byte; scena con parete di fondo, tre scaffali e venti oggetti; $N=300$ fotogrammi per configurazione. Il tempo comprende deproiezione, percorso dei raggi e scritture, non la generazione dell'immagine. Con $k=4$ i raggi validi sono $16\,557$, con $k=1$ sono $264\,709$.

\begin{table}[H]\centering\footnotesize
\caption{Memoria e tempo per fotogramma (media / 95° percentile, in ms) sul nucleo veloce; profondità sintetica, $N=300$ fotogrammi per casella. Per le strutture sparse la memoria va dalla scena con il solo suolo a quella ingombra. La griglia polare e la lettura nell'immagine non dipendono da $s$: le due colonne sono due esecuzioni. \prova{S}}\label{tab:cos-tempi}
\begin{tabularx}{\textwidth}{@{}Lrrcccc@{}}\toprule
 & \multicolumn{2}{c}{Memoria (byte)} & \multicolumn{2}{c}{$k=4$} & \multicolumn{2}{c}{$k=1$ (tutti i raggi)}\\\cmidrule(lr){2-3}\cmidrule(lr){4-5}\cmidrule(l){6-7}
 & 10 cm & 5 cm & 10 cm & 5 cm & 10 cm & 5 cm\\\midrule
Griglia polare, 3 m & \multicolumn{2}{c}{$10\,788$ (più 1,5 MB di tabelle)} & 0,03 / 0,04 & uguale & 0,42 / 0,55 & 0,37 / 0,40\\
Lettura nell'immagine & \multicolumn{2}{c}{$2\,457\,600$ (due fotogrammi)} & 0,51 / 0,52 & 0,51 / 0,52 & 0,74 / 0,88 & 0,51 / 0,53\\
Colonne di bit & $28\,800$ & $115\,200$ & 1,30 / 1,34 & 2,47 / 2,55 & 19,5 / 20,0 & 37,4 / 38,4\\
Tre fasce di quota & $48\,000$ & $192\,000$ & 3,29 / 3,37 & 6,23 / 6,41 & 44,5 / 45,7 & 85,5 / 97,6\\
Griglia densa & $288\,000$ & $2\,304\,000$ & 2,03 / 2,08 & 4,21 / 4,35 & 22,4 / 22,8 & 44,5 / 49,1\\
Hash a blocchi $8^3$ & $159\,744$--$241\,664$ & $763\,904$--$1\,515\,520$ & 3,02 / 3,12 & 6,71 / 7,45 & 32,6 / 47,9 & 72,4 / 76,4\\
Octree & $163\,264$--$377\,856$ & $1\,004\,672$--$2\,691\,008$ & 5,02 / 5,39 & 11,2 / 12,1 & 72,2 / 75,7 & 161 / 171\\\bottomrule
\end{tabularx}
\end{table}

\begin{figure}[H]\centering
\begin{tikzpicture}[x=2.7cm,y=0.40cm,font=\footnotesize]
  \foreach \e/\lab in {-1/{0{,}1},0/{1},1/{10},2/{100}}{
    \draw[muted!50] ({\e+1},-0.3)--({\e+1},-7.7);
    \node[below] at ({\e+1},-7.7) {$\lab$};}
  \node[below] at (3.75,-7.7) {ms};
  \foreach \nome/\v/\lab [count=\i] in {griglia polare/0.40/{0{,}40},lettura nell'immagine/0.53/{0{,}53},colonne di bit/38.4/{38{,}4},griglia densa/49.1/{49{,}1},hash a blocchi/76.4/{76{,}4},tre fasce di quota/97.6/{97{,}6},octree/171/{171}}{
    \pgfmathsetmacro{\x}{log10(\v)+1}
    \fill[accent!65] (0,-\i+0.3) rectangle (\x,-\i-0.3);
    \node[left] at (0,-\i) {\nome};
    \node[left,text=white,inner sep=2pt] at (\x,-\i) {$\lab$};}
  \pgfmathsetmacro{\xt}{log10(66.7)+1}
  \draw[dashed,thick] (\xt,0) -- (\xt,-7.7);
  \node[above] at (\xt,0) {$T=66{,}7$ ms};
\end{tikzpicture}
\caption{95° percentile del tempo per fotogramma con tutti i raggi ($k=1$) e $s=5\unita{cm}$, in scala logaritmica, contro il periodo $T$ (dati della tabella~\ref{tab:cos-tempi}).}\label{fig:cos-tempi}
\end{figure}

Dalle misure seguono quattro constatazioni.
\begin{enumerate}
\item \emph{Il costo segue il prodotto raggi $\times$ passi, non la struttura.} La griglia polare costa 0,03--0,6 ms; la lettura nell'immagine 0,5 ms a fotogramma più 5,9 ns a punto interrogato (636 µs per $108\,000$ punti); la griglia densa 2--4 ms con un raggio su 16 e 45--67 ms con tutti i raggi a 5 cm (44,5 / 49,1 ms sul nucleo veloce; 65,2 / 67,4 ms su quello lento, che in mediana costa 1,49 volte).
\item \emph{La sparsità costa tempo e risparmia poco.} Sulle medie di 20 configurazioni la tabella hash costa 1,39--1,76 volte la griglia densa e l'octree 2,2--3,8 volte (2,3--3,9 con la fusione dei figli uguali, che non riduce il picco dei nodi: almeno il 99\,\% del non fuso). La tabella hash occupa dal 33 all'84\,\% della densa (rapporto 1,2--3,0, con profondità valida fino a 6 m; il limite del lemma~\ref{lem:cos-sparsa} vale a 3 m): il risparmio massimo è $2\,304\,000-763\,904=1\,540\,096$ byte, cioè 1,54 MB. L'octree nella scena ingombra occupa più della densa (117--131\,\%).
\item \emph{Il linguaggio pesa più della struttura.} Lo stesso aggiornamento denso (10 cm, $19\,200$ raggi) scritto in Python senza librerie numeriche costa 102 ms in mediana, contro 2,0 ms in C; la griglia polare 1,6 ms.
\item \emph{Confronto con il periodo.} Con $k=4$ tutte le strutture stanno sotto 12,8 ms al 95° percentile (20,4 ms sul nucleo lento). Con tutti i raggi a 5 cm (figura~\ref{fig:cos-tempi}) la griglia densa occupa il 74\,\% del periodo sul nucleo veloce e lo supera su quello lento (67,4 ms); tabella hash e octree lo superano sempre.
\end{enumerate}
Una ripetizione delle misure sulla stessa macchina, con altri processi in esecuzione, ha dato medie entro il $\pm10\,\%$ in 29 configurazioni su 32. Un secondo programma indipendente (voxel da 8 byte, stanza chiusa, 33,3 milioni di passi) dà 88,9 ms per la griglia densa a 5 cm con tutti i raggi, cioè 2,67 ns per passo, contro 2,6 ns del primo sullo stesso tipo di nucleo.

\subsection{Costo della catena completa}\label{sec:cos-catena}

\textbf{Previsione durante l'occlusione.} Il calcolo di $\Reg_\alpha(t)$ (\S\ref{sec:occ}) per gli oggetti senza osservazione costa, in C, 0,6--1,4 ms al 95° percentile con dieci oggetti nascosti insieme (due programmi, due geometrie dell'ostacolo) e 5,4 ms con 120, cioè l'8,2\,\% di $T$: cresce in modo lineare, circa 0,05 ms a oggetto ($N=1\,500$ cicli, scena sintetica). In Python senza librerie numeriche esistono due misure, entrambe su $N=150$ cicli con dieci oggetti: 5,3 ms di mediana e 6,2 ms al 95° percentile in una prima misura; 6,4 ms di mediana e 15,9 ms al 95° percentile in una ripetizione con altri processi in esecuzione sulla stessa macchina. Il tetto assegnato a questo calcolo è 6,7 ms, un decimo di $T$: la prima misura lo rispetta, la seconda no; il valore in Python non è stabile sotto carico. \prova{S}

\textbf{Catena di riferimento} (mappa polare, gruppi, associazione, filtro di Kalman, regole, uscita; \S\ref{sec:sti}), in Python con una libreria numerica, un solo filo, stesso calcolatore: tabella~\ref{tab:cos-catena}.

\begin{table}[H]\centering\small
\caption{Tempo per ciclo della catena di riferimento (ms). La terza riga è riportata senza ricalcolo dai dati per ciclo.}\label{tab:cos-catena}
\begin{tabularx}{\textwidth}{@{}Lrrrrc@{}}\toprule
Ingresso & $N$ cicli & Media & 95° perc. & Massimo & Livello\\\midrule
Profondità simulata dalle schede tecniche, $424\times240$, 162 sequenze & $9\,152$ & 6,35 & 18,6 & 105,0 & \prova{S}\\
La stessa, con la soglia minima di punti per ostacolo portata da 3 a 40 & $9\,152$ & 0,95 & 1,28 & 7,7 & \prova{S}\\
Scene senza rumore, $848\times480$ & 48 & 2,5 & 2,9 & 4,7 & \prova{S}\\
Registrazione pubblica affollata, $640\times360$ \cite{depthtrack} & 320 & 81 & 152 & 383 & \prova{P}\\\bottomrule
\end{tabularx}
\end{table}

Sulle sequenze simulate 4 cicli su $9\,152$ superano 66 ms (0,04\,\%; intervallo di Wilson al 95\,\%: $[0{,}02;\,0{,}11]$\,\%). Sulla registrazione pubblica la media supera il periodo: in 21,3 s la catena conferma 254 oggetti, e ne tiene insieme 152 in mediana (le due cifre misurano cose diverse: gli oggetti confermati nell'intera sequenza e quelli presenti nello stesso ciclo).

\textbf{Dipendenza dal numero di oggetti seguiti.} A ogni ciclo la catena confronta ogni oggetto seguito con ogni gruppo (porta $\gamma$ sull'innovazione $\nu$ con covarianza $S$), risolve un'assegnazione ottima \cite{kuhn1955} e aggiorna un filtro per oggetto: lavoro almeno lineare in $n_T$ e polinomiale nell'assegnazione. Sui $9\,152$ cicli simulati il tempo medio cresce con $n_T$: 0,9 ms con 0--2 oggetti, 2,6 ms con 6--10, 5,3 ms con 11--20, 8,1 ms con 21--40, 17,6 ms con 41--80, 34,0 ms oltre 80; ai minimi quadrati
\begin{equation}\label{eq:cos-retta}
t\approx 0{,}33+0{,}30\,n_T\unita{ms}\qquad(\text{termine quadratico } 1{,}7\cdot10^{-4}\,n_T^2).
\end{equation}
Con $n_T=152$ la retta dà circa 46 ms, 50 ms con il termine quadratico, lo stesso ordine degli 81 ms misurati sulla registrazione pubblica; l'estrapolazione oltre i valori di $n_T$ simulati non è stata verificata. Gli stessi fotogrammi, con la soglia a 40 punti (mediana 2 oggetti seguiti invece di 17), costano 0,95 ms. Il costo si governa quindi limitando gli oggetti spuri, non cambiando la rappresentazione dello spazio. \prova{S}

\begin{riquadro}{Limiti di validità delle misure di costo}
(1)~Tutti i tempi sono di un calcolatore portatile, con un solo filo, su profondità sintetica o registrata: la scheda di bordo non è misurata; il solo cambio di tipo di nucleo vale 1,49 volte. (2)~La catena completa nel linguaggio finale non è misurata: le strutture sono misurate in C, la catena in Python, e il divario fra i due linguaggi sull'aggiornamento denso è di circa 50 volte. (3)~La latenza $\tau$ da un estremo all'altro non è misurata (\S\ref{sec:lim}). (4)~Octree e strutture sparse non sono stati misurati con una libreria esistente né su regioni grandi. \prova{N}
\end{riquadro}

\subsection{Alternative considerate ed escluse}\label{sec:cos-alternative}

La tabella~\ref{tab:cos-alternative} elenca le alternative e ciò che le esclude. Dove non è detto altro le prove sono in simulazione, con modello ideale del sensore, camera ferma a $h=1{,}2\unita{m}$ e $\theta=15^\circ$, 15 Hz.

\begin{table}[H]\centering\footnotesize
\caption{Alternative considerate ed escluse. Le frequenze sono date come conteggi; dove la fonte della misura non riporta la numerosità è scritto «$N$ non riportata».}\label{tab:cos-alternative}
\begin{tabularx}{\textwidth}{@{}>{\RaggedRight\arraybackslash}p{3.3cm}>{\RaggedRight\arraybackslash}p{4.4cm}Lc@{}}\toprule
Alternativa & Motivo tecnico & Controesempio o misura & Liv.\\\midrule
Tre sensori a tempo di volo multizona per coprire $360^\circ$ (\S\ref{sec:geo}) & Il campo di un sensore è $60^\circ\times60^\circ$ \cite{st-vl53l7cx}; una zona non distingue il suolo da un ostacolo basso & Tre sensori affiancati coprono $3\cdot60^\circ=180^\circ$: per $360^\circ$ ne servono sei. Una zona è larga 26 cm a 2 m. Portata dichiarata a 5 klx nel modo $8\times8$: 0,5 m \cite{st-vl53l7cx} & \prova{C}\\
Stima della massa (\S\ref{sec:sti}) & Non identificabile da misure di sola posizione con forza ignota & $(M,u)$ e $(\lambda M,\lambda u)$ danno la stessa legge delle misure; informazione di Fisher di rango 1 & \prova{D}\\
Stima dell'accelerazione (\S\ref{sec:sti}) & L'accelerazione fotogramma per fotogramma sta sotto il rumore & $\sigma_a=\sqrt6\,\sigma_p/\Delta^2=16{,}5\unita{m/s^2}$ a 15 Hz con $\sigma_p=3$ cm (valore assunto), contro $0{,}67\unita{m/s^2}$ efficaci nel modello di pedone. Filtro ad accelerazione costante dopo 2 s senza misure: errore 2,29 m contro 1,18 m ($1\,500$ prove per caso, rumore assunto 5 cm) & \prova{C} \prova{S}\\
Catena filtro $\to$ rete neurale $\to$ decisione & Gli enunciati sul «libero» e sull'avviso sono dimostrati, nel modello dichiarato e sotto le sue ipotesi, su regole scritte (\S\ref{sec:map}, \S\ref{sec:occ}) & Scelta di progetto: nessuna misura la esclude. Un componente appreso assegna solo un'etichetta (vedi testo) & \prova{N}\\
Mappa 3D a voxel: densa, sparsa a hash, octree & A sensore fermo non contiene più del fotogramma corrente; le strutture sparse costano tempo & Hash: 1,39--1,76 volte il tempo della densa per al più 1,54 MB; octree: 2,2--3,8 volte. Densa a 5 cm con tutti i raggi: 44,5 / 49,1 ms (tabella~\ref{tab:cos-tempi}) & \prova{S}\\
Varianti che ricordano o deducono un «libero» (\S\ref{sec:map}) & Un «libero» ricordato non è garantito in un mondo che cambia & Sette varianti, ciascuna con il suo controesempio, nella Tabella~\ref{tab:map-varianti}; per esempio la lettura a due stati dà $80\,505$ falsi su $146\,901$ punti dati per liberi (54,8\,\%). Lettura del fotogramma corrente, stesse scene: 0 & \prova{S}\\
Mappa a tre fasce di quota & Parte da 0,15 m; non costa meno della griglia densa & Dichiara il 30,6\,\% del libero vero (0,0\,\% con il 5\,\% di pixel non validi sparsi); 32 celle «libere» sopra una scatola alta 10 cm (griglia polare: 0); con la fascia alta, 0 celle libere su $2\,610$ entro 3 m; 95° percentile 3,4--6,4 ms contro 2,1--4,4 & \prova{S} \prova{C}\\
Previsione a due ipotesi di moto con pesi aggiornati a ogni ciclo, regione a ellisse & I pesi contano la stessa evidenza a ogni ciclo; l'ellisse non è una regione tarata & Sosta dichiarata nel 38--100\,\% dei passaggi normali (cinque simulatori). Regione assente per almeno un ciclo in 374 e $1\,312$ passaggi su $2\,000$, e in $1\,976$ e $1\,987$ su $2\,000$ (due geometrie con ombra larga, due letture della regola). Ellisse al 95\,\%: vedi testo & \prova{S}\\
Scelta automatica della configurazione sotto carico & Ridurre la mappa in esercizio crea «libero» non osservato; il tempo dipende dalla configurazione scelta & «Libero» sul 13,8\,\% e sul 57,0\,\% della zona (due simulatori); 200 cambi in 10 minuti con l'isteresi; 17,5--28 cambi al minuto (quattro simulatori); $N$ non riportata & \prova{S}\\
Probabilità d'intrusione come criterio d'allarme & Scende quando il dato peggiora; non dice se il soggetto \emph{può} entrare & Su rotta di collisione esatta, a 2 s, vale 0,16--0,18; vale 0 con il soggetto nascosto a 9 cm dal volume protetto (\S\ref{sec:occ}) & \prova{D} \prova{S}\\
Sensori laterali come preallarme (\S\ref{sec:geo}) & Il margine temporale guadagnato è di centesimi di secondo & A 1,4 m/s restano 0,07 s; chi passa a 0,5 m di lato entra nel campo a 0,73 m, dentro la zona vicina che la camera non vede. Resta possibile un canale di sola prossimità laterale, che non dichiara mai «libero» & \prova{C}\\
Predittori appresi delle traiettorie & Sull'orizzonte breve il guadagno sul punto è piccolo; la taratura dell'incertezza non è garantita & Nessuna prova in questo lavoro; dati di letteratura \cite{scholler2020,cheraghipouria2026} (vedi testo) & \prova{N}\\\bottomrule
\end{tabularx}
\end{table}

\textbf{Mappe volumetriche.} L'argomento di sicurezza contro una mappa persistente a sensore fermo (Osservazione~\ref{oss:map-fermo}), i controesempi che lo sostengono (Tabella~\ref{tab:map-varianti}) e il conto sulla latenza della griglia bayesiana (Proposizione~\ref{prop:map-logodds}) sono nel \S\ref{sec:map-persistente}; qui restano i conti che riguardano il costo. La decimazione dei raggi non salta voxel solo se $k\le s\,f/d_{\max}$: con $f=337$ pixel e $d_{\max}=6\unita{m}$, $k\le2$ a 5 cm e $k\le5$ a 10 cm \prova{C}. La mappa a tre fasce di quota, con la fascia alta, non dichiara mai libero entro 3 m perché la quota di 2,0 m entra tutta nel campo solo da $z=(2{,}0-h)/\tan(58^\circ/2-\theta)=0{,}8/\tan14^\circ=3{,}21\unita{m}$ \prova{C}. Una mappa persistente tornerebbe necessaria con il sensore in moto, caso non verificato né sperimentalmente né in simulazione (\S\ref{sec:lim}): la forma compatibile con i risultati del \S\ref{sec:map} è una griglia densa locale a tre stati con l'età del dato, che aggiunge avvisi «occupato ricordato» e non produce mai «libero»; eseguita su un altro nucleo, in simulazione porta il 95° percentile della griglia polare da 0,84 a 0,88 ms. Le strutture sparse diventano convenienti quando la regione cresce: $20\times20\times5\unita{m}$ a 5 cm sono $16\cdot10^6$ voxel, 128 MB a 8 byte.

\textbf{Previsione a due ipotesi.} L'alternativa mantiene due ipotesi di moto (transito e sosta), ne moltiplica i pesi a ogni ciclo, dichiara la regione come ellisse al 95\,\% della miscela e smette di mostrarla oltre una soglia fissa di incertezza (0,8 m). In simulazione i pesi moltiplicati a ogni ciclo contano più volte la stessa evidenza e dichiarano una sosta in una quota dei passaggi normali che va dal 38 al 100\,\% secondo il simulatore. La soglia a 0,8 m non scatta mai con un pannello largo 1 m (0 passaggi su $3\,600$) e, in due geometrie con ombra larga, fa mancare la regione nei passaggi della tabella~\ref{tab:cos-alternative}. La prima misura alla ricomparsa cade dentro l'ellisse dichiarata al 95\,\%, per il soggetto che prosegue, in 0,28 / 0,30 / 0,00 dei passaggi con una lettura della regola e in 0,997 / 0,71 / 0,00 con l'altra, in tre geometrie dell'ostacolo (la prima con pannello largo 1 m, la terza con pannello largo 2 m); la dicitura «$\pm\sigma$» corrisponde a una regione che contiene il 39,3\,\%. Alla prova di taratura la regola, applicata alla lettera, copre il 76,4\,\% $[70{,}7;\,81{,}3]$ contro il 90--95\,\% richiesto (nella lettura più favorevole: 100\,\% per chi prosegue, 46,6\,\% per chi torna indietro); anche con i pesi corretti chi esce dal lato lontano dell'ombra è contenuto nel 28\,\% e nel 7\,\% dei casi a 2 e 3 s. Dell'alternativa restano il ritaglio della regione sul suolo visto libero, con la fascia di guardia, e le regole di disegno, fra cui quella per cui la regione non sparisce (\S\ref{sec:occ}). \prova{S}

\textbf{Scelta automatica della configurazione.} Un selettore che in esercizio cambia la configurazione sotto carico e riduce la mappa fa nascere «libero» su suolo non osservato e, se la griglia viene riallocata, perde l'intero insieme di possibile presenza $\Pres(t)$; l'isteresi nel tempo non ferma l'oscillazione quando il tempo di ciclo dipende dalla configurazione scelta: ne fissa solo il ritmo. Nel sistema la configurazione si sceglie e si controlla all'avvio, è dichiarata in ogni uscita, e in esercizio la mappa non cambia (\S\ref{sec:lim}). \prova{S}

\textbf{Rete neurale e classificazione.} Nel sistema decide una catena di rilevamento geometrico, associazione, filtro di Kalman e regole scritte; nessun modello appreso produce un «libero», toglie un avviso o abbassa un livello. Il componente di classificazione su più classi che il sistema richiede (\S\ref{sec:sti}) sta fuori dal percorso che decide: assegna un'etichetta all'oggetto seguito, in modo asincrono, e l'etichetta può solo allargare l'avviso (\S\ref{sec:occ}). In questo lavoro il classificatore non è realizzato né valutato: è simulato con etichette assegnate a mano (\S\ref{sec:sti}). Un conto mostra che per le classi veloci l'etichetta arriverebbe comunque tardi: i 3 m di portata si attraversano in 0,75 s a 4 m/s, e la sola conferma di un oggetto seguito chiede 4 cicli, cioè 0,27 s; per queste classi protegge la configurazione iniziale, non l'etichetta \prova{C}. Nessun rilevatore di oggetti è stato eseguito in questo lavoro. Dalla documentazione pubblica, i rilevatori pronti addestrati sull'insieme pubblico di immagini COCO coprono, secondo un conto fatto sull'elenco delle sue classi, 9 dei circa 25 tipi di una tassonomia più fine di quella della Tabella~\ref{tab:sti-profili}, che distingue per esempio la persona su carrozzina (mancano carrello elevatore, tram, multirotore, carrozzina), e chiedono l'immagine a colori; la letteratura riporta che le reti moderne danno probabilità mal tarate e propone una correzione a posteriori. Queste affermazioni di letteratura non sono state riscontrate su fonti presenti in bibliografia e sono date senza riferimento. Se un modello di questo tipo basti per queste classi, a 3 m e sulla sola profondità, non è verificato. \prova{N}

\textbf{Predittori appresi delle traiettorie.} Sull'orizzonte di 4,8 s, sulla media dei cinque scenari degli insiemi pubblici ETH e UCY \cite{pellegrini2009,lerner2007}, il modello a velocità costante ha errore medio e finale di 0,39 m e 0,83 m, contro 0,44 m e 0,91 m del miglior predittore appreso deterministico e 0,32 m e 0,60 m di un modello generativo valutato sul migliore di 20 campioni \cite{scholler2020}; i predittori addestrati sulla verosimiglianza tendono a distribuzioni troppo o troppo poco sicure \cite{cheraghipouria2026}. Il confronto riguarda il punto previsto, non una regione, non comprende occlusioni, e quei dati non vengono da un sensore di profondità frontale. Per questo il sistema usa il modello a velocità quasi costante e dichiara la regione con un livello misurato a parte (\S\ref{sec:val}); nessun predittore appreso è stato provato in questo lavoro.

\subsection{Lavori correlati}\label{sec:cos-correlati}

\textbf{Sensori di profondità.} Per le camere a stereo attivo la documentazione del costruttore dichiara che l'errore di profondità cresce con il quadrato della distanza \cite{realsense-taratura-profondita}; una caratterizzazione indipendente su un bersaglio planare bianco fra 0,4 e 2,1 m dà per il rumore assiale un polinomio di secondo grado nella distanza \cite{ahn2019}. I sensori a tempo di volo multizona offrono al più $8\times8$ zone a 15 Hz, oppure $4\times4$ zone a 60 Hz \cite{st-um2884}, con una portata che la luce ambiente riduce \cite{st-vl53l7cx}; l'uso all'aperto delle camere a tempo di volo resta limitato dall'interferenza solare \cite{qiu2023}; un LiDAR 2D rotante economico misura un solo strato orizzontale, a circa 10 Hz e con risoluzione angolare di $0{,}72^\circ$ \cite{slamtec-rplidar-c1}. Sono dichiarazioni dei costruttori o misure di terzi, non misure fatte sul sistema descritto; il loro uso nei calcoli è nel \S\ref{sec:geo}.

\textbf{Mappe di occupazione e stato «sconosciuto».} Le griglie di occupazione sono la rappresentazione classica per la percezione dei robot mobili \cite{elfes1989}. OctoMap rappresenta in modo esplicito, oltre allo spazio occupato, anche quello libero e quello sconosciuto \cite{hornung2013}; nella sua realizzazione lo sconosciuto è codificato in modo implicito, come assenza di informazione \cite{octomap}. UFOMap dichiara una rappresentazione esplicita di tutti e tre gli stati \cite{duberg2020}. Le griglie evidenziali trattano l'ignoranza come una sorgente d'incertezza distinta dall'errore di misura \cite{moras2011}, e le griglie dinamiche di occupazione sono state realizzate in tempo reale con un filtro a particelle in parallelo \cite{nuss2018}. Le mappe volumetriche a distanza con segno funzionano in tempo reale su un solo nucleo di CPU \cite{oleynikova2017} o su GPU \cite{millane2024,niessner2013}, e per i veicoli in volo è stata proposta una griglia locale, in un buffer circolare, che si muove con il veicolo \cite{usenko2017}. Rispetto a questi lavori il sistema adotta una mappa polare del suolo, centrata su una camera ferma, con quattro stati: separa lo sconosciuto in «in ombra» $\Omb(t)$ (dietro un ostacolo osservato) e «non osservato», e non aggiorna una probabilità per cella. Non è stata confrontata sperimentalmente con le librerie citate: i confronti del \S\ref{sec:cos-misure} usano realizzazioni proprie delle stesse strutture.

\textbf{Seguimento di più oggetti e occlusioni.} Lo schema «rilevamento, filtro di Kalman \cite{kalman1960}, assegnazione ottima \cite{kuhn1955}» è quello di SORT, che dato l'insieme dei rilevamenti aggiorna gli oggetti seguiti a 260 Hz \cite{bewley2016}, e della sua estensione in tre dimensioni \cite{weng2020}; la teoria della stima per il seguimento è materia di manuale \cite{barshalom2001}. I perfezionamenti noti agiscono sull'associazione (uso di quasi tutti i rilevamenti, anche a basso punteggio \cite{zhang2022}; compensazione del moto della camera \cite{aharon2022}) e sulla gestione del ciclo di vita, isolata come modulo a sé \cite{pang2021}. Durante un'occlusione il filtro, privo di misure, accumula errore \cite{cao2023}: è il punto in cui si colloca questo lavoro. La letteratura segue tre strade. La prima tratta l'occlusione come un problema di previsione a breve termine \cite{khurana2021} o la impara da dati che portano le annotazioni degli oggetti invisibili \cite{tokmakov2021}; i lavori di riferimento riconoscono però che i metodi faticano sotto occlusione pesante \cite{hsieh2023} e che resta un divario notevole rispetto a una vera permanenza dell'oggetto \cite{vanhoorick2023}. La seconda usa l'informazione «negativa», cioè la mancata osservazione là dove il sensore vede \cite{koch2007,wyffels2015,krejci2026}. La terza, nata nella guida automatica, sovrastima con insiemi lo stato di ciò che è nascosto \cite{orzechowski2018} e predice in forma insiemistica i comportamenti ammissibili dei partecipanti rilevati e di quelli potenzialmente nascosti \cite{koschi2021}; il problema affine dell'inseguimento con visibilità considera un evasore imprevedibile, di posizione iniziale ignota e senza limite di velocità \cite{guibas1997}. Il sistema combina la prima e la terza strada: per l'oggetto già seguito mantiene una regione di previsione $\Reg_\alpha(t)$ dichiarata a un livello $\alpha$; per l'ombra in cui è entrato un oggetto seguito mantiene un'area d'avviso $W(t)$, che realizza l'insieme di possibile presenza $\Pres(t)$ del \S\ref{sec:occ} con una velocità massima assunta $\vmax$; un'ombra in cui non è entrato alcun oggetto seguito resta «in ombra», senza avviso. La percezione umana offre un criterio per la resa: l'identità di un oggetto nascosto è mantenuta se scompare e ricompare lungo il bordo fisso dell'ostacolo \cite{scholl1999}.

\textbf{Incertezza tarata.} Le previsioni probabilistiche vanno valutate per affidabilità e nitidezza \cite{kress2021} con punteggi propri \cite{gneiting2007}; un'analisi recente conferma che la velocità costante approssima bene la dinamica dei pedoni nella maggior parte dei casi \cite{uhlemann2024}. La previsione conforme dà regioni con copertura marginale garantita senza ipotesi sulla distribuzione, purché i dati di taratura e di prova siano scambiabili \cite{angelopoulos2021}; è stata applicata alla pianificazione in ambienti dinamici, con regioni su 20 passi e 498 traiettorie di prova su 500 dentro la regione al 95\,\% \cite{lindemann2023}, resa meno prudente \cite{cleaveland2024} e adattata in linea quando la distribuzione cambia \cite{gibbs2021,dixit2022}. Due limiti noti valgono anche qui: la garanzia è marginale, e una copertura condizionata esatta è impossibile senza ipotesi sulla distribuzione \cite{barber2019}; una probabilità di collisione calcolata con dati peggiori può apparire più bassa \cite{balch2019}, il che motiva l'uso di regioni e insiemi anziché di una sola probabilità.

\textbf{Sicurezza dei sensori di presenza.} Nella sorveglianza di velocità e separazione dei robot collaborativi la distanza di protezione ha forma lineare, somma dei percorsi dell'operatore e del robot nei tempi di reazione e di arresto e di termini d'incertezza di posizione, con una velocità dell'operatore di 1,6 m/s \cite{marvel2017}. I sistemi 3D industriali di questo tipo classificano i volumi in vuoti, occupati e sconosciuti e trattano come occupato ciò che nessun sensore vede \cite{therobotreport-veo,us-brevetto-11846916}; il segnalamento ferroviario applica lo stesso principio a guasto sicuro, imponendo l'aspetto più restrittivo quando manca l'energia \cite{us-49cfr-236-205}. Il sistema ne riprende due principi: non dichiarare mai libero ciò che non è osservato, e segnalare in modo esplicito la perdita di affidabilità. Non è un dispositivo di protezione certificato, e nessuna conformità a norme di sicurezza delle macchine è rivendicata; le norme sulla distanza di separazione non sono state consultate nel testo originale.

\textbf{Ausili per persone con disabilità visiva.} Le rassegne mostrano che efficacia e sicurezza sono la prima preoccupazione degli utenti \cite{xu2023} e che solo il 22,9\,\% di 70 lavori valuta il dispositivo con persone con disabilità visiva \cite{romlay2021}. Un sistema di avviso con previsione per estrapolazione lineare e soglie a 5 s e 2,5 s è stato provato con sei utenti non vedenti \cite{kayukawa2019}; esistono dispositivi indossabili con sensori di profondità e riscontro aptico \cite{dossantos2025,kozlov2026}. Sul canale d'uscita gli studi sono discordi: per chi usa il bastone il riscontro sonoro ha dato un carico cognitivo 22 volte maggiore di quello aptico, effetto assente nei partecipanti bendati \cite{martinez2014}; in un altro studio, con otto partecipanti, gli stimoli aptici sono risultati meno intuitivi di quelli sonori \cite{lupu2020}; la guida con suono spazializzato regge un compito secondario meglio del linguaggio spaziale \cite{giudice2008}. Il sistema non risolve questa scelta: espone gli stati e l'incertezza in forma indipendente dal canale (\S\ref{sec:lim}).

\section{Uscita, profilo di chi porta il sensore, perdita di affidabilità, condizioni di validità e limiti}\label{sec:lim}

Questo capitolo descrive ciò che il sistema dichiara verso l'esterno e ciò che non dichiara mai, come usa la configurazione di chi porta il sensore, come segnala che non è più affidabile, e dentro quali condizioni valgono gli enunciati dei capitoli precedenti. Nessuna uscita è un comando: il sistema osserva, sorveglia lo stato della scena e avvisa; non ha funzioni d'ingaggio né di puntamento. Tutte le cifre sono dimostrazioni in un modello dichiarato, calcoli o simulazioni a sensore fermo: nessuna è una misura su un sensore fisico.

\subsection{Contratto d'uscita}\label{sec:lim-uscita}

\paragraph{Che cosa si dichiara.} Di ogni cella della mappa polare il sistema dichiara uno fra quattro stati: \emph{libero}, \emph{occupato osservato}, \emph{in ombra}, \emph{non osservato} (\S\ref{sec:map}); «non osservato» non diventa mai «libero». Di ogni oggetto seguito dichiara se lo vede, se non lo vede ma indica dove può essere (con un livello, quando è dichiarabile), o se non lo sa. Gli stati di una traccia e i loro segni sono nella tabella~\ref{tab:lim-stati} e nella figura~\ref{fig:lim-segni}. Valgono due regole di forma.
\begin{enumerate}[label=(\alph*)]
\item \emph{Contorno pieno solo con una misura.} Il disco pieno compare soltanto nei cicli in cui la traccia ha una misura non tagliata dentro la porta; una previsione non esce mai con il segno di un'osservazione, e il passaggio avviene in un solo verso (da previsto a osservato serve una misura).
\item \emph{Regione tratteggiata con l'età del dato.} Ogni regione di previsione porta l'etichetta «previsto · livello · larga $X{,}X\unita{m}$ · non visto da $Y{,}Y\unita{s}$»; l'età è misurata dall'ultima misura non tagliata e non è mai scritta come nulla.
\end{enumerate}
Ogni distinzione sta nella forma, nella trama e nella parola; il colore la ripete soltanto e non è mai verde. Ogni riga d'uscita porta contatore del fotogramma, età, livello, configurazione e l'indicazione «osservato» o «dedotto».

\begin{table}[H]\centering\small
\caption{Stati di una traccia e uscita corrispondente. Le soglie sono in secondi; a $15\unita{Hz}$ tre fotogrammi valgono $0{,}2\unita{s}$.}\label{tab:lim-stati}
\begin{tabularx}{\linewidth}{@{}p{2.5cm}Lp{4.6cm}L@{}}
\toprule
Stato & Quando & Segno e parole & Finisce quando\\
\midrule
osservato & misura non tagliata in porta nel ciclo & disco pieno; «visto · $d\unita{m}$» & manca la misura\\
visto in parte & il gruppo è tagliato da un bordo d'ombra o dal bordo del campo & contorno pieno, senza velocità, insieme a regione e avviso & torna intero o scompare\\
non confermato & traccia confermata, senza misura, in campo aperto & anello vuoto con l'età, per $0{,}2\unita{s}$ & evento «non più rilevato»\\
previsto & entrata in un'ombra, livello dichiarabile & anello vuoto, regione tratteggiata $\Reg_\alpha(t)$, età & diventa «previsto ampio»; non torna indietro\\
previsto ampio & regione oltre l'$80\,\%$ dell'ombra, oltre l'orizzonte del livello, ingresso da fermo, sensore non in salute & tutta la parte dell'ombra non vista libera, a tratteggio rado, senza punto & riaggancio; ombra vista libera; uscita dall'area osservata\\
possibile presenza & una traccia è entrata nell'ombra & ombra più fascia $w$, quadrettata, con l'età; «residua» dopo un riaggancio & solo dove il suolo è rivisto libero oltre $w$ dall'ombra corrente\\
\bottomrule
\end{tabularx}
\end{table}

\begin{figure}[H]\centering
\begin{tikzpicture}[x=1cm,y=1cm,font=\footnotesize]
\fill (0,0) circle (0.16); \node[anchor=west] at (0.3,0) {osservato};
\draw[line width=0.9pt] (2.55,-0.18) rectangle (2.95,0.18); \node[anchor=west] at (3.1,0) {visto in parte};
\draw[dashed,pattern=north east lines,pattern color=muted] (6.0,-0.3) ellipse (0.55 and 0.32);
\draw[line width=0.8pt,fill=white] (6.0,-0.3) circle (0.13); \node[anchor=west] at (6.7,-0.05) {previsto: anello vuoto,};
\node[anchor=west] at (6.7,-0.42) {regione tratteggiata, età};
\draw[dotted,line width=0.8pt,pattern=north west lines,pattern color=muted!60] (10.6,-0.6) -- (11.9,-0.6) -- (11.6,0.15) -- (10.9,0.15) -- cycle;
\node[anchor=west] at (12.0,-0.05) {previsto ampio:};
\node[anchor=west] at (12.0,-0.42) {nessun punto};
\draw[line width=1pt,pattern=crosshatch,pattern color=muted] (0,-1.55) -- (1.6,-1.55) -- (1.25,-0.75) -- (0.35,-0.75) -- cycle;
\node[anchor=west] at (1.8,-0.97) {possibile presenza: ombra dell'ostacolo più la fascia $w$, con l'età;};
\node[anchor=west] at (1.8,-1.34) {non dipende dall'identità e finisce solo dove il suolo è rivisto libero};
\end{tikzpicture}
\caption{Segni dell'uscita visiva. Il segno pieno è riservato a ciò che ha una misura nel ciclo corrente.}\label{fig:lim-segni}
\end{figure}

\paragraph{Parole per settore.} La zona valutata è divisa in 17 settori d'azimut; ogni settore riceve un solo livello, il più alto fra: «attenzione» (un elemento mobile in avvicinamento, detto per primo, poi un ostacolo occupato a $d$ metri); «zona non verificata» (possibile presenza, regione di previsione, ombra, o un buco di dati prima della portata); «suolo visto libero da $1{,}3\unita{m}$ a $d\unita{m}$», dove $d$ è il raggio della prima cella non libera del settore e non supera la portata meno la fascia, $3-0{,}40=2{,}6\unita{m}$; «sconosciuto» se il settore non vede il suolo da $1{,}3\unita{m}$. Sopra tutti, fuori scala, sta «affidabilità persa» (\S\ref{sec:lim-salute}). Con la geometria del capitolo~\ref{sec:geo}, nove settori su 17 vedono il suolo da $1{,}3\unita{m}$; gli 8 laterali dicono «sconosciuto» per geometria \prova{S}. La frase per settori è una regola di progetto: l'ordine fra i primi due livelli non è stato provato \prova{N}.

\paragraph{Che cosa non si dichiara mai.} (i) «Via libera» o «nessun ostacolo osservato fino a $d$»: la frase parla del suolo letto, non del volume, perché un oggetto basso vicino, una barra sottile o un ostacolo sospeso non interrompono la lettura del suolo (\S\ref{sec:lim-limiti}). (ii) Una sosta, un'uscita o un'identità certa: dietro un ostacolo non si distinguono; l'identità al ritorno è «presunta» o «ambigua». (iii) Un livello numerico oltre $3\unita{s}$ senza osservazione, con il predittore di ripiego, per chi entra nell'ombra da fermo o con il bordo dell'ombra non riconosciuto. (iv) Massa o accelerazione stimate (\S\ref{sec:sti}). (v) Una confidenza in percentuale o un intervallo «$\pm$ metri» attorno al punto: il punto previsto non ha una precisione dichiarata, la ha la regione. (vi) Una classe senza la sua origine: nelle prove l'etichetta è «assegnata a mano» e una riga fissa dichiara che nessun modello di classificazione è in esecuzione.

\paragraph{Due profili d'uscita.} Le regole sono le stesse; cambia il canale. Nel profilo di \emph{monitoraggio di un varco o di un'area di lavoro} l'uscita è lo schermo con il registro degli eventi. Nel profilo di \emph{avviso di presenza per una persona con disabilità visiva ferma} accanto al sensore l'uscita è sonora: un segnale di $0{,}3\unita{s}$ per ogni cambio di stato e la frase «direzione, distanza, stato» solo a stato stabile, con precedenza (1) affidabilità persa, ripetuta finché dura, (2) elemento visto in avvicinamento, (3) scomparso dietro l'ostacolo, (4) previsto ampio e poi non ricomparso, (5) ricomparso, (6) identità ambigua, (7) la frase. All'avvio il canale dichiara la sorgente dei dati e che il silenzio non significa assenza. Il canale sonoro è una specifica: che i segnali si distinguano, che diano una posizione e con quale ritardo arrivino non è stato provato con persone \prova{N}.

\subsection{Registro degli eventi}\label{sec:lim-registro}

Per il monitoraggio di attività il sistema scrive una riga per evento: ora, identità, evento, \emph{come lo sa} («osservato» oppure «dedotto»), età del dato, settore, configurazione. Gli eventi sono: ingresso; scomparsa dietro l'ostacolo; «non ricomparso da $n\unita{s}$» (dedotto); ricomparsa con identità presunta o ambigua; affidabilità persa e tornata; con le classi, etichetta assegnata, cambiata, smentita dalla cinematica, classe fuori scenario. Nessun evento afferma una sosta o un'uscita, e nessun evento nasce da un peso del predittore. La logica del registro è stata controllata in simulazione su scene sintetiche con una catena ridotta: 0 violazioni delle regole del registro su 600 scene; eventi mancati nel 16--52\,\% delle scene e scomparse spurie nel $12{,}5\,\%$ \prova{S}. Non è stata eseguita sulla catena completa \prova{N}.

\subsection{Salute del sensore e perdita di affidabilità}\label{sec:lim-salute}

\paragraph{Controllo all'avvio.} A scena vuota e ferma: (1) i filtri della camera che ricordano o riempiono i dati mancanti (persistenza temporale, riempimento dei buchi \cite{realsense-filtri-post}) sono spenti; (2) autoverifica: coprendo il sensore i pixel validi devono scendere sotto il $10\,\%$ entro 3 fotogrammi; (3) la posa viene dalla taratura, non è mai stimata dai primi fotogrammi (su una registrazione pubblica \cite{tum-rgbd} quella stima sbaglia di $5^\circ$ \prova{P}) e si confronta con una misura diretta: uno scarto di $0{,}5^\circ$ o più impedisce l'avvio; (4) un oggetto alto $10\unita{cm}$ a $2{,}5\unita{m}$ deve risultare «occupato»; (5) sul suolo vuoto, in $30\unita{s}$, non devono nascere gruppi né tracce confermate. Se cade una delle prime quattro prove il sistema non parte; se cade la quinta restano spenti l'avviso senza traccia e l'identità al ritorno, e l'uscita lo dichiara. Il motivo della (1): un filtro che riempie i buchi trasforma un dato mancante in un dato valido, cioè in un possibile «libero»; su profondità simulata dalle schede tecniche, con i filtri accesi e il sensore coperto, restano «libero» 989--995 celle per altri 3 fotogrammi \prova{S}. Il motivo della (3) e della (4) è nel limite~\ref{lim:posa}.

\paragraph{Indicatori in esercizio.} Il sensore è dichiarato non in salute se: non arriva un fotogramma nuovo per $0{,}13\unita{s}$ o ne arrivano tre identici; i pixel validi restano sotto la metà del valore di taratura per $0{,}2\unita{s}$; l'unità inerziale segnala più di $0{,}5^\circ$ o di $5^\circ/\mathrm{s}$; il piano del suolo stimato da fermo si scosta di $0{,}5^\circ$ dalla taratura. La quota dei validi si conta sui pixel e non sulle celle: con il $10\,\%$ di pixel validi 1\,444 celle su 1\,460 risultavano ancora note \prova{S}.

\paragraph{Che cosa accade.} «Affidabilità persa» è uno stato del sistema, non di una traccia. La mappa diventa tutta «non osservato»; ogni traccia confermata passa a «previsto ampio · dedotto, sensore non affidabile» e continua a uscire; il ciclo emette comunque, con il segnale dedicato. Al ritorno dopo un tempo $T$ le regioni di previsione e l'area d'avviso sono dilatate di $\vmax T+w$ e intersecate con ciò che non è visto libero nel fotogramma; una traccia la cui intersezione è vuota si chiude come «vista libera», e ogni riaggancio è marcato «identità ambigua». La scelta è quella del guasto verso lo stato più restrittivo, come nel segnalamento ferroviario \cite{us-49cfr-236-205}.

\emph{Livello di prova.} Segnale entro $0{,}13\unita{s}$ e 0 falsi segnali in $60\unita{s}$, in simulazione, con rumore dell'unità inerziale ipotizzato \prova{S}. Il ramo del ritorno è verificato per calcolo in pianta: intersecando con «non osservato» la persona passata dietro l'ostacolo era fuori dalla regione in 13 istanti su 13, intersecando con «non visto libero» era dentro in 13 su 13 \prova{C}; non è stato riprovato sulla catena \prova{N}.

\subsection{Profilo di chi porta il sensore}\label{sec:lim-portatore}

\paragraph{Configurazione.} Chi porta il sensore non è classificato: è descritto da una configurazione deterministica, fissata all'avvio e immutabile in esercizio. I campi sono: identificativo; velocità propria tipica e massima $v_e$, $v_{e,\max}$ (m/s); tempo di reazione $T_r$ (s), dall'avviso emesso alla reazione iniziata, più l'eventuale ritardo di collegamento; decelerazione di frenata $a_f$ e accelerazione massima (m/s$^2$); gradi di libertà e olonomia; canale d'avviso; indicatore «sensore in moto»; elenco delle classi ammesse nel sito (lo \emph{scenario}, \S\ref{sec:sti}).

\paragraph{Formule.} Sia $t_{\mathrm{conf}}=0{,}267\unita{s}$ il tempo di conferma di una traccia (\S\ref{sec:sti}), $\tau=0{,}155\unita{s}$ la latenza stimata (\S\ref{sec:lim-latenza}) e $C=0{,}95\unita{m}$ un margine costante. Per un oggetto che si avvicina frontalmente a velocità $v$ senza frenare, il tempo d'anticipo necessario e la distanza a cui l'oggetto deve essere rilevato sono
\begin{align}
T_{\mathrm{ant}} &= t_{\mathrm{conf}}+\tau+T_r+\frac{v_e}{a_f}, \label{eq:lim-tant}\\
d_{\mathrm{serve}}(v) &= (v+v_e)\,(t_{\mathrm{conf}}+\tau+T_r)+\frac{v_e^2}{2a_f}+\frac{v\,v_e}{a_f}+C. \label{eq:lim-dserve}
\end{align}
Il primo termine è lo spazio di chiusura fino all'inizio della reazione, il secondo lo spazio d'arresto del portatore, il terzo lo spazio percorso dall'oggetto durante la frenata. La forma «velocità per tempo più costante» è quella della distanza di protezione nella sorveglianza di velocità e separazione; $C$ somma $0{,}85\unita{m}$, valore che la stessa fonte riporta dalla norma sul posizionamento dei ripari per l'avvicinamento perpendicolare con capacità di rilevamento fra 40 e $70\unita{mm}$ \cite{marvel2017} (la norma non è stata consultata nel testo originale, e che la capacità di rilevamento del sistema stia in quell'intervallo non è verificato), e $0{,}10\unita{m}$ d'incertezza del sensore: il confronto dà un ordine di grandezza e non è una dichiarazione di conformità. La \emph{velocità coperta} è la $v$ per cui $d_{\mathrm{serve}}(v)$ uguaglia la portata; a sensore fermo ($v_e=0$) e con $3\unita{m}$ di portata vale $(3-C)/(t_{\mathrm{conf}}+\tau+T_r)$.

\paragraph{Uso.} All'avvio il sistema calcola $T_{\mathrm{ant}}$ e $d_{\mathrm{serve}}(\vmax)$ per ogni classe dello scenario; se $d_{\mathrm{serve}}(\vmax)$ supera la portata dichiara in una riga fissa «avviso tempestivo non dichiarato per» quelle classi; se la portata divisa per $\vmax$ è minore di $t_{\mathrm{conf}}+\tau$ dichiara che per quelle classi l'avviso esce dopo il contatto; se $v_e>0$ dichiara «soglie non provate a sensore in moto». In esercizio, a ogni «attenzione» si calcola il tempo disponibile come distanza divisa per la velocità d'avvicinamento osservata aumentata di $2\sigma_v$: se è minore di $T_{\mathrm{ant}}$ si aggiunge la parola «tardivo». Il profilo non abbassa mai un livello e nessuna etichetta di classe è argomento di queste regole.

\begin{table}[H]\centering\small
\caption{Profili d'esempio e grandezze derivate dalle \eqref{eq:lim-tant}--\eqref{eq:lim-dserve}, portata $3\unita{m}$ \prova{C}. Origine dei valori: (i) ipotesi di progetto; (r) valore riportato da una fonte non verificata. Le ultime due colonne sono analisi fuori dal dominio provato \prova{N}.}\label{tab:lim-portatore}
\begin{tabularx}{\linewidth}{@{}Lp{1.9cm}p{1.0cm}p{2.2cm}p{1.7cm}p{2.5cm}@{}}
\toprule
Profilo & $T_r$ (s) & $T_{\mathrm{ant}}$ (s) & Velocità coperta, sensore fermo (m/s) & In moto: $v_e$; $a_f$ & $d_{\mathrm{serve}}(0)$; velocità coperta\\
\midrule
stazione fissa, avviso a una persona vicina & 1,0 (i) & 1,42 & 1,44 & --- & ---\\
veicolo terrestre senza equipaggio, arresto automatico & 0,3 (i) & 0,72 & 2,84 & 1,0 (i); 1,5 (i) & $2{,}01\unita{m}$; $0{,}72\unita{m/s}$\\
lo stesso, con operatore remoto & 1,5 (r) + 0,3 (i) & 2,22 & 0,92 & --- & ---\\
multirotore senza equipaggio in volo stazionario & 0,3 (i) & 0,72 & 2,84 (l'oscillazione non è provata) & 3,0 (r); 3,0 (i) & $4{,}62\unita{m}$: oltre la portata\\
persona con disabilità visiva & 1,5 (i), campo 0,5--1,5 & 1,92 & 1,07 & 1,0 (i); 1,0 (i) & $3{,}37\unita{m}$: oltre la portata\\
\bottomrule
\end{tabularx}
\end{table}

\paragraph{Lettura e limiti.} (a) A sensore fermo del profilo conta solo $T_r$. Con $T_r=1{,}0\unita{s}$ la velocità coperta è $1{,}44\unita{m/s}$: una persona a passo normale rientra nel dominio con $6\unita{cm}$ di margine sulla distanza, il passo svelto no. (b) Il tempo di reazione di una persona con disabilità visiva a un avviso sonoro non ha una fonte: $1{,}0\unita{s}$ e $1{,}5\unita{s}$ sono due ipotesi di progetto, e con la seconda la velocità coperta scende a $1{,}07\unita{m/s}$, sotto il passo normale; il valore si decide solo con una misura su persone \prova{N}. Il valore di $1{,}5\unita{s}$ per l'operatore remoto è la reazione di frenata a sorpresa di un conducente riportata in letteratura, non verificata sulla fonte. (c) \emph{Ogni cifra con $v_e>0$ è un calcolo fuori dal dominio provato}: con il sensore in moto l'ombra si sposta, cade l'ipotesi di ostacolo fermo ed entra l'errore di posa (\S\ref{sec:lim-moto}); nessuna soglia del sistema è stata provata in quelle condizioni, nemmeno in simulazione. La fascia di guardia diventerebbe $w=\varepsilon_\Omega+(v_{\mathrm{scen}}+v_e)\cdot0{,}355\unita{s}$, con $\varepsilon_\Omega$ l'errore del bordo d'ombra e $v_{\mathrm{scen}}$ la velocità massima dello scenario. Chi si muove paga il proprio arresto: a $1{,}0\unita{m/s}$ con $a_f=1{,}0\unita{m/s^2}$ e $T_r=1{,}5\unita{s}$ la sola distanza per un ostacolo fermo ($3{,}37\unita{m}$) supera la portata. (d) Le classi veloci sono fuori portata comunque: la tabella del dominio per classe è nel capitolo~\ref{sec:occ}.

\subsection{Bilancio di latenza}\label{sec:lim-latenza}

\begin{table}[H]\centering\small
\caption{Voci note della latenza $\tau$, dall'evento alla scheda audio, a 30 fotogrammi al secondo \prova{C}. Nessuna voce è misurata sulla catena del prototipo \prova{N}.}\label{tab:lim-latenza}
\begin{tabularx}{\linewidth}{@{}p{5.2cm}p{2.0cm}L@{}}
\toprule
Tratto & Valore medio & Natura del dato\\
\midrule
evento $\to$ inizio dell'esposizione & $16{,}7\unita{ms}$ & calcolo: uniforme fra 0 e $33{,}3\unita{ms}$\\
esposizione $\to$ fotogramma disponibile & $95\unita{ms}$ & misura pubblicata dal costruttore sull'intero banco di prova, dal movimento allo schermo, nel modo $848\times100$ a 30 fotogrammi al secondo \cite{realsense-alta-velocita-d435}: stima per eccesso del solo sensore, non misurata alla risoluzione usata qui; media, senza percentili\\
fotogramma $\to$ comando d'uscita & $22\unita{ms}$ & misura su scene sintetiche con codice non ottimizzato\\
campioni $\to$ scheda audio & $21{,}3\unita{ms}$ & un quanto di 1\,024 campioni a $48\unita{kHz}$\\
\midrule
somma delle medie & $155\unita{ms}$ & \\
\bottomrule
\end{tabularx}
\end{table}

La somma si legge in due modi. Come \emph{media} vale $155\unita{ms}$. Come \emph{percentile}: la sola prima voce la fa variare fra 138 e $172\unita{ms}$, e se le altre fossero costanti il 95° percentile sarebbe $170\unita{ms}$; un requisito «95° percentile non oltre $150\unita{ms}$» non è quindi soddisfatto dalle sole voci note. Un secondo dato del costruttore per il tratto del sensore, un unico esempio di circa $71\unita{ms}$, porterebbe la somma a circa $131\unita{ms}$: le due cifre vengono da fonti diverse e nessuna decide. Mancano, perché senza dato, la sintesi vocale (nulla con suoni preregistrati) e il tratto dalla scheda al suono. La latenza entra nella fascia di guardia (\S\ref{sec:occ}) e nelle \eqref{eq:lim-tant}--\eqref{eq:lim-dserve}: ogni cifra che la usa è condizionata alla stima.

\subsection{Invarianti del sistema}\label{sec:lim-invarianti}

\begin{table}[H]\centering\footnotesize
\caption{Invarianti, livello di prova e capitolo che li tratta. «Per costruzione»: la proprietà segue dalla forma della regola e si controlla con un'asserzione a ogni ciclo; non è una misura.}\label{tab:lim-invarianti}
\begin{tabularx}{\linewidth}{@{}rLLp{1.1cm}@{}}
\toprule
 & Enunciato & Livello di prova & Cap.\\
\midrule
1 & «Libero» esiste solo nel fotogramma corrente, dove metà dell'impronta della cella e delle otto vicine legge il suolo; non si ricorda e non si deduce & \prova{S} 0 celle libere su un ostacolo in 300 condizioni (9\,000 fotogrammi), modello ideale del sensore, posa entro $0{,}5^\circ$ (limite superiore per condizione $1{,}3\,\%$); \prova{N} su sensore & \ref{sec:map}\\
2 & Fuori campo, oltre portata, suolo non visibile vicino, ombra, impronta letta per meno di metà (pixel non validi compresi), sensore non in salute: mai «libero» & \prova{D} per costruzione, senza ipotesi sul sensore. Vale per ciò che non è visto; non afferma che ciò che è «libero» sia libero & \ref{sec:map}\\
3 & Segno pieno solo con una misura non tagliata in porta; una previsione non esce come osservazione & per costruzione; soglia 0 non superata: 8 passaggi su 90 con la configurazione di riferimento su profondità simulata dalle schede tecniche, 1 su 140 con i gruppi a celle forti \prova{S} & \ref{sec:sti}\\
4 & $\Reg_\alpha(t)$ è contenuta nell'area d'avviso e non copre suolo «libero» oltre $0{,}20\unita{m}$ dall'ombra & \prova{D} l'annidamento; \prova{S} la fascia di $0{,}20\unita{m}$ (0 casi misurati, non garantiti dalla sola regola del peso) & \ref{sec:occ}\\
5 & Nessun timer, modello appreso, etichetta di classe, peso, identità o conferma dell'operatore produce un «libero», toglie o restringe un avviso, abbassa un livello & per costruzione; per l'etichetta \prova{S} sulla sola fascia: 0 violazioni su 20\,000 sequenze di etichette e 0 su altre 20\,000; \prova{N} sulla catena & \ref{sec:occ}\\
6 & L'evidenza negativa vale dove il suolo è letto ora; restringe $\Reg_\alpha$, mai l'avviso & per costruzione & \ref{sec:occ}\\
7 & L'avviso finisce solo guardando, cella per cella, oltre $w$ dall'ombra corrente & per costruzione; \prova{S} 0 fotogrammi con il nascosto fuori dall'avviso su 336, 169 e 2\,360, su profondità simulata dalle schede tecniche; che l'avviso termini quando il suolo è rivisto dopo la rimozione dell'ostacolo non è sostenuto: resta acceso dopo $2\unita{s}$ in 10 prove su 10 (72--100\,\%) & \ref{sec:occ}\\
8 & Ogni uscita porta contatore, età, livello, configurazione e «osservato» o «dedotto» & per costruzione & \ref{sec:lim}\\
9 & Massa e accelerazione non sono stimate & \prova{D} nel modello del punto materiale & \ref{sec:sti}\\
10 & Continuità: per ogni traccia con 4 fotogrammi di misura entrata in un'ombra, a ogni ciclo fino alla chiusura escono regione, età e stato & per costruzione, data la premessa; \prova{S} 0 passaggi con un vuoto su 23\,402 (6\,000 camminatori distinti: meno di 6,4 su 10\,000), modello ideale del sensore; la premessa manca in 24 passaggi su 23\,402 con l'avviso senza traccia, in 928 senza & \ref{sec:occ}, \ref{sec:val}\\
11 & Configurazione, scenario, profilo del portatore e interruttori sono fissati all'avvio & regola di progetto & \ref{sec:lim}\\
12 & Il livello più basso è «suolo visto libero da $1{,}3\unita{m}$ a $d\unita{m}$», $d\le2{,}6\unita{m}$; altrimenti «sconosciuto» & regola di progetto & \ref{sec:lim}\\
13 & Senza autoverifica, posa e prova della posa il sistema non parte; non in salute, il ciclo esce comunque & \prova{N} & \ref{sec:lim}\\
14 & Ogni evento dice se è osservato o dedotto; nessuno afferma sosta, uscita o identità certa & \prova{S} su 600 scene sintetiche, catena ridotta & \ref{sec:lim}\\
15 & L'etichetta di classe può solo allargare la fascia; senza etichetta, o con etichetta smentita dalla cinematica, vale il profilo prudente dello scenario, senza livello numerico & per costruzione; \prova{N} sulla catena & \ref{sec:sti}, \ref{sec:occ}\\
\bottomrule
\end{tabularx}
\end{table}

Tre forme più forti di questi enunciati non reggono, e per questo non sono fra gli invarianti: che il «libero» ricordato per tre fotogrammi sposti l'oggetto di meno della fascia (con la latenza lo spostamento è $0{,}72\unita{m}$ a $2{,}5\unita{m/s}$ \prova{C}); che la previsione non influenzi l'avviso (la regola d'ingresso nell'ombra usa la previsione: vale solo la monotonia, cioè la previsione accende e non restringe); che «il 100\,\% degli istanti dentro l'avviso» sia una misura (è una conseguenza della costruzione finché l'errore dell'ombra non supera $w$ più mezzo corpo).

\subsection{Condizioni di validità}\label{sec:lim-validita}

\begin{riquadro}{Condizioni di validità degli enunciati e delle cifre}
\begin{enumerate}[label=\arabic*.,itemsep=1pt]
\item \textbf{Camera ferma, posa nota.} Quota $h=1{,}2\unita{m}$, inclinazione $\theta=15^\circ$, errore d'inclinazione minore di $0{,}5^\circ$; suolo piano.
\item \textbf{Al chiuso, fino a $3\unita{m}$.} Oltre la portata non c'è mappa né avviso; la frase del suolo arriva al più a $2{,}6\unita{m}$.
\item \textbf{Superfici che rimandano segnale.} Ostacoli opachi che restituiscono una profondità valida e corretta; filtri della camera spenti.
\item \textbf{Un soggetto per ombra,} visto per almeno 4 fotogrammi prima di entrarvi; ostacolo fermo.
\item \textbf{Velocità entro il limite coperto dalla fascia.} Con $w=0{,}40\unita{m}$, tempo di rilevamento assunto $t_a=0{,}133\unita{s}$, latenza stimata $0{,}155\unita{s}$ e tenuta di un ciclo $0{,}067\unita{s}$: $0{,}40/0{,}355=1{,}13\unita{m/s}$; $0{,}56\unita{m/s}$ se il bordo dell'ombra sbaglia di $0{,}20\unita{m}$ \prova{C}.
\item \textbf{Ostacolo che nasconde per intero.} Per una persona alta $1{,}70\unita{m}$ servono $1{,}61\unita{m}$ d'ostacolo nella geometria provata.
\item \textbf{Per il livello numerico della regione:} ostacolo largo almeno $2\unita{m}$, persona che resta sul proprio percorso, bordo dell'ombra noto esattamente, non oltre $3\unita{s}$, classe «persona».
\end{enumerate}
Fuori da queste condizioni ogni cifra perde il suo dominio; restano le proprietà che valgono per costruzione sul singolo fotogramma: ciò che non è osservato non è mai «libero», e nessuna memoria, etichetta o timer produce un «libero».
\end{riquadro}

\subsection{Limiti}\label{sec:lim-limiti}

Ogni limite è dato con la cifra che lo mostra. Salvo diversa indicazione le frequenze sono simulazioni a sensore fermo; gli intervalli sono nel capitolo~\ref{sec:val}.
\begin{enumerate}[label=\arabic*.,ref=\arabic*,leftmargin=1.8em]
\item \textbf{Vetro e ostacoli che non rimandano segnale.} Il raggio attraversa la lastra e legge il suolo dietro: 654 punti della lastra su 800 sono dati per liberi a $2{,}0\unita{m}$ e 600 su 800 a $1{,}5\unita{m}$ \prova{S}. Un palo o una lastra scuri senza ritorno risultano «libero» su profondità simulata dalle schede tecniche; la regola dell'impronta regge un corpo esteso senza ritorno (0 fotogrammi con una cella libera su 900 \prova{S}), ma il palo di $3\unita{cm}$ non è stato provato con essa \prova{N}.
\item \textbf{Suolo non visibile vicino al sensore.} Il suolo è visto da $1{,}24\unita{m}$ sull'asse e da $1{,}71\unita{m}$ ai bordi del campo ($1{,}22$ e $1{,}70\unita{m}$ con la focale di 424 pixel, \S\ref{sec:geo-suolo}) \prova{C}. Non sono visti un ostacolo di $15\unita{cm}$ a meno di $1{,}09\unita{m}$ né la testa di una persona di $1{,}70\unita{m}$ a meno di $2{,}01\unita{m}$ \prova{C}; un oggetto alto $0{,}35\unita{m}$ fra $0{,}5$ e $0{,}8\unita{m}$ non riceve alcun pixel e il settore dichiara comunque il suolo letto più lontano \prova{S}.
\item \textbf{Ostacoli sospesi, sottili o più bassi del sensore.} Sotto un ostacolo sospeso il suolo resta «libero» (374 celle nella scena di prova \prova{S}); con una barra di $1\unita{cm}$ a quota $0{,}8\unita{m}$ la distanza dichiarata va oltre la barra nel 4,7--7,3\,\% dei settori per fotogramma \prova{S}. Dietro un ostacolo più basso di circa $1{,}65\unita{m}$ il sistema tratta come nascosta una persona ancora in parte visibile; con un pannello di $1{,}6\unita{m}$ la testa spunta in 10 passaggi su 90 \prova{S}.
\item \textbf{Secondo soggetto nella stessa ombra.} Un secondo soggetto mai visto che esce dalla stessa ombra riceve l'identità del primo: 2400 casi su 2400 nelle prove di base con modello ideale del sensore, fra $0{,}91$ e $0{,}995$ nelle prove d'insieme con modello ideale (3 disposizioni, oltre 3000 casi ciascuna), 19 su 20 (76--99\,\%) su profondità simulata dalle schede tecniche \prova{S}. Nessuna delle regole provate lo ripara: per questo l'identità è «presunta» e l'avviso non dipende da essa.
\item \textbf{Chi svolta, chi esce dal lato lontano, chi parte da fermo.} Per chi svolta la quota dentro $\Reg_\alpha$ nelle tre geometrie è $0{,}86/0{,}67/0{,}79$ a $2\unita{s}$ e $0{,}84/0{,}58/0{,}67$ a $3\unita{s}$; per chi esce dal lato lontano $0{,}82/0{,}61/0{,}79$ e $0{,}76/0{,}57/0{,}70$ \prova{S}, con modello ideale del sensore (\S\ref{sec:val}); con soggetto puntiforme e altri semi, chi esce dal lato lontano scende a $0{,}32$--$0{,}67$ a $2\unita{s}$ e a $0{,}13$--$0{,}47$ a $3\unita{s}$; per chi parte da fermo accanto all'ostacolo, $0{,}11/0{,}16/0{,}26$ a 1, 2 e $3\unita{s}$ nelle prove di base ($N=141$, 134, 54) e $0{,}82$--$0{,}83$ nella prova d'insieme, sugli istanti con il centro nascosto: in questi casi non si dichiara un livello e copre solo l'avviso. Con un ostacolo largo $1\unita{m}$ il livello è «almeno 87 su 100» fino a $1\unita{s}$, poi nessuna cifra.
\item \textbf{Errore del bordo d'ombra.} Con il bordo stimato e non riconosciuto la quota a $1\unita{s}$ scende a $0{,}30$--$0{,}46$ (estremo basso: chi prosegue); a $0{,}74$--$0{,}80$ se il bordo è riconosciuto in 8 cicli su 10 \prova{S}. Con l'ombra stimata più stretta di $0{,}10/0{,}20/0{,}45/0{,}60\unita{m}$ i passaggi in cui il nascosto esce dall'avviso sono $2/15/1\,704/1\,979$ su 2\,000 \prova{S}. Su profondità simulata dalle schede tecniche l'errore del bordo va da $+2$ a $+15\unita{cm}$ al rumore nominale e da $-65$ a $+73\unita{cm}$ al rumore limite.
\item\label{lim:posa} \textbf{Errore di posa.} Un errore d'inclinazione $\delta\theta$ sposta la quota ricostruita di circa $\rho\,\delta\theta$ (\S\ref{sec:geo}): un ostacolo di $10\unita{cm}$ scende sotto la soglia di $7\unita{cm}$ con $0{,}67^\circ$ a $2{,}55\unita{m}$ e $0{,}57^\circ$ a $3\unita{m}$ \prova{C}. Con $+1^\circ$ lo stesso ostacolo a $2{,}55\unita{m}$ risulta «libero» in 7 celle su 7, e una scatola di $10\unita{cm}$ fra $2{,}6$ e $2{,}9\unita{m}$ in 145--150 fotogrammi su 300 (tre esecuzioni) a $-1^\circ$ \prova{S}; le due prove usano convenzioni di segno opposte, e in entrambe l'errore abbassa la quota ricostruita. A $+0{,}5^\circ$ l'identità al ritorno scende a 6 passaggi su 10: la tolleranza è minore di $0{,}5^\circ$.
\item \textbf{Avviso che non si spegne.} Dietro un ostacolo fermo nessuna cella dell'avviso è rilasciabile (0 su 808; 140 passaggi su 140) \prova{S}: l'avviso finisce quando l'ombra si sposta. Nella realizzazione provata, rimosso l'ostacolo, l'avviso residuo resta acceso dopo $2\unita{s}$ in 10 prove su 10, (72--100\,\%), per una causa non determinata: che l'avviso termini quando il suolo è rivisto dopo la rimozione dell'ostacolo non è sostenuto \prova{S} (\S\ref{sec:occ}).
\item \textbf{Profondità realistica.} Su profondità simulata dalle schede tecniche, con la configurazione di riferimento, l'identità al ritorno tiene in 0 passaggi su 90 e il pannello fermo è etichettato mobile in 79 su 90; con i gruppi a celle forti (pannello mobile in 0 passaggi su 140 nel dominio e in 7 su 10 fuori) e l'avviso che resta, l'identità tiene in 176 passaggi su 180 [94--99\,\%] a 424 pixel e in 174 su 180 [93--98\,\%] a 848, in un dominio stretto (un adulto, $0{,}5$--$1{,}6\unita{m/s}$, entro $2{,}6\unita{m}$) e cade fuori (0--1 su 10 a $0{,}3\unita{m}$ dal pannello; 2 su 10 per un soggetto alto $1{,}0\unita{m}$) \prova{S}. Il simulatore sottostima i pixel senza dato ($7{,}8\,\%$ al livello limite contro 24--34\,\% in registrazioni pubbliche \prova{P}): le cifre al rumore nominale sono ottimistiche. Sulla stessa profondità simulata la regola della regione persistente (\S\ref{sec:occ}) non è stata eseguita: nella catena provata, che conserva uno stato «perso» senza regione, a soggetto nascosto la traccia è «previsto» in 1\,568 fotogrammi su 4\,552 e «perso» in 2\,983 a 424 pixel di larghezza (1\,584 e 2\,968 su 4\,553 a 848; a entrambe le risoluzioni le due voci sommano a un fotogramma meno del totale); la traccia esiste in tutti i fotogrammi, ma solo 44 e 43 passaggi su 140 hanno una previsione in ogni fotogramma nascosto \prova{S}. La continuità è quindi verificata solo con modello ideale del sensore.
\item \textbf{Classificazione.} Il classificatore su più classi è un componente richiesto dal sistema; in questo lavoro non è realizzato né valutato: le etichette sono assegnate a mano, eventualmente con errori iniettati per misurare la sensibilità. Nessuna classe diversa da «persona» è stata simulata nella catena completa: i percorsi del predittore e tutte le prove riguardano pedoni. Con $3\unita{m}$ di portata e un tempo di reazione ipotizzato $T_r=1{,}0\unita{s}$ l'avviso è tempestivo, per calcolo, fino a $1{,}44\unita{m/s}$: rientrano a velocità tipica 3 classi su 15 e a velocità massima nessuna \prova{C} (\S\ref{sec:occ}). La sicurezza delle indicazioni non dipende dall'etichetta (invarianti 5 e 15).
\item \textbf{Nessuna misura su sensore fisico, nessuna prova con persone.} Rumore, buchi, pixel misti, latenza e costo sulla scheda di bordo non sono misurati; il tempo per ciclo, $6{,}35\unita{ms}$ di media su profondità simulata, sale a $81\unita{ms}$ su una registrazione pubblica affollata \prova{P} (\S\ref{sec:cos}). Lo schermo e i segnali sonori non sono stati provati con persone, vedenti o no; la stessa carenza è documentata nella letteratura sugli ausili, dove su 70 lavori il $51{,}4\,\%$ è valutato in simulazione e il $22{,}9\,\%$ con persone con disabilità visiva \cite{romlay2021}.
\end{enumerate}

\subsection{Estensione a un sensore in moto: analisi dei requisiti}\label{sec:lim-moto}

Con il sensore in moto un residuo fra due fotogrammi può venire dal moto dell'oggetto, dall'errore del moto proprio, dalla parallasse, da un'occlusione che cambia o dal rumore. La regola «mobile» (\S\ref{sec:sti}) confronta due medie di 3 fotogrammi i cui capi distano $T_w=13\Delta=0{,}867\unita{s}$ con la soglia $5\,\sigma\sqrt{2/3}$: uno spostamento apparente $\delta$ dovuto alla posa non compensata produce un'etichetta falsa con probabilità $Q(5-\delta/(\sigma\sqrt{2/3}))+Q(5+\delta/(\sigma\sqrt{2/3}))$, che vale $10\,\%$ per $\delta=3{,}036\,\sigma$ e $1\,\%$ per $\delta=2{,}183\,\sigma$ \prova{D} (\S\ref{sec:sti-mobile}). Ne seguono i requisiti della tabella~\ref{tab:lim-moto}, che renderebbero valide le stesse regole; sono calcoli in forma chiusa nel modello a errori gaussiani indipendenti fra fotogrammi \prova{C}. Il solo controllo numerico è un Monte Carlo della formula nello stesso modello d'errore ($N=2\cdot10^6$ valutazioni: $0{,}0999$ e $0{,}0100$). Nessun requisito è verificato su una piattaforma in moto né con una simulazione del sensore in moto, e gli errori reali delle sorgenti di posa non sono noti \prova{N}: la tabella è un'analisi, non un risultato.

\begin{table}[H]\centering\small
\caption{Requisiti sulla posa per un sensore in moto (soglia: $10\,\%$ di falsi, salvo dove indicato) \prova{C}.}\label{tab:lim-moto}
\begin{tabularx}{\linewidth}{@{}Lp{3.4cm}p{3.4cm}@{}}
\toprule
Requisito & a $2\unita{m}$ & a $3\unita{m}$\\
\midrule
Imbardata relativa sulla finestra $T_w$ & \multicolumn{2}{@{}l}{$0{,}88^\circ$ ($1{,}01^\circ/\mathrm{s}$) a ogni distanza; $0{,}63^\circ$ ($0{,}73^\circ/\mathrm{s}$) per l'$1\,\%$}\\
Errore di velocità propria, trasversale / radiale & $3{,}5$ / $5{,}3\unita{cm/s}$ & $5{,}3$ / $11{,}9\unita{cm/s}$\\
Ritardo fra posa e profondità & \multicolumn{2}{@{}l}{$29\unita{ms}$ con un gradino di $30^\circ/\mathrm{s}$; $74\unita{ms}$ con un'oscillazione di $\pm3^\circ$ a $1\unita{Hz}$}\\
Inclinazione rispetto al suolo ($10\,\%$ di celle di suolo oltre la soglia di quota) & $1{,}70^\circ$; $1{,}27^\circ$ con $\pm2{,}5\unita{cm}$ d'oscillazione verticale & $1{,}03^\circ$; $0{,}79^\circ$\\
Scostamento non modellato perché la regione copra ancora $0{,}90$ & \multicolumn{2}{@{}l}{$9{,}7\unita{cm}$ a $0{,}67\unita{s}$; $16\unita{cm}$ a $1\unita{s}$; $41\unita{cm}$ a $2\unita{s}$}\\
Quota di fascia lasciata alla posa ($0{,}10\unita{m}$) & $2{,}9^\circ$; ritardo $95\unita{ms}$ a $30^\circ/\mathrm{s}$ & $1{,}9^\circ$; $64\unita{ms}$\\
\bottomrule
\end{tabularx}
\end{table}

\paragraph{Che cosa resta invariato e che cosa no.} La regola «mai libero dove non è osservato» è calcolata sul singolo fotogramma: non usa la posa relativa fra fotogrammi e resta valida in moto. Lo stesso vale per il divieto di produrre un «libero» da memoria, etichette o timer, e per la marca «osservato» o «dedotto». Dipendono invece dalla posa: la correttezza di ciò che è dichiarato «libero», che richiede fotogramma per fotogramma l'inclinazione rispetto al suolo entro la tolleranza del limite~\ref{lim:posa}; la regola «mobile», che senza compensazione del moto proprio scambia il moto apparente per moto dell'oggetto (nel seguimento di più oggetti la compensazione del moto della camera è usata per questo \cite{aharon2022}); l'ombra, che non è più quella di un ostacolo fermo; la fascia di guardia, che cresce con $v_e$. Una memoria dello spazio osservato tornerebbe necessaria, perché i fotogrammi vedono raggi diversi, e la forma compatibile con gli invarianti è una griglia locale che segue il sensore \cite{usenko2017}, usata a senso unico: può aggiungere un «occupato ricordato», non può produrre un «libero» né sciogliere una possibile presenza (\S\ref{sec:cos}); la conferma su tre fotogrammi vale solo sotto $0{,}375\unita{m/s}$ di traslazione propria \prova{C}.

\subsection{Misure che deciderebbero i punti non verificati sperimentalmente}\label{sec:lim-misure}

La tabella~\ref{tab:lim-misure} descrive, come metodo di verifica, le misure su un sensore fisico che deciderebbero gli enunciati marcati \prova{N}. Le soglie sono fissate prima della misura; dove le fonti non fissano una numerosità la casella lo dice.

\begin{table}[H]\centering\footnotesize
\caption{Protocollo di verifica: grandezza, numerosità, soglia di smentita.}\label{tab:lim-misure}
\begin{tabularx}{\linewidth}{@{}rLp{4.0cm}p{4.6cm}@{}}
\toprule
 & Grandezza misurata & Numerosità & Soglia\\
\midrule
1 & Gruppi e tracce sul suolo vuoto, a posa tarata; stato di un oggetto di $10\unita{cm}$ a $2{,}5\unita{m}$ & $30\unita{s}$ & 0 gruppi, 0 tracce confermate; oggetto «occupato»\\
2 & Celle «libero» dentro l'ombra vera di un pannello, tracciata sul suolo & 150 fotogrammi & 0 celle\\
3 & Quota dei pixel dell'impronta entro il margine di lettura del suolo, a scena ferma & $60\unita{s}$ & sotto l'$80\,\%$ oltre $2\unita{m}$: tarare il margine prima di dichiarare una distanza\\
4 & Quota dei cicli in cui il bordo dell'ombra è riconosciuto; errore del bordo sui due lati del pannello & non fissata & condizione per dichiarare un livello: errore entro $\pm20\unita{cm}$\\
5 & Persona dentro $\Reg_\alpha$ a 1 e $2\unita{s}$, con occlusione virtuale (una copia della catena riceve le misure con un cuneo cancellato; la persona resta visibile e fa da verità) & 100 passaggi nel dominio dichiarato, camminatori diversi da quelli della taratura & almeno 87 su 100: «più di 8 su 10» (esclude $0{,}80$ con potenza $0{,}88$ se il vero è $0{,}90$); almeno 96 su 100: «almeno 9 su 10»\\
6 & Segni pieni a soggetto nascosto, nella stessa prova & gli stessi 100 passaggi & 0\\
7 & Avvisi rilasciati senza aver rivisto il suolo, con un secondo soggetto; avviso dopo la rimozione dell'ostacolo & non fissata & 0 rilasci; avviso spento entro 5 fotogrammi e 0 nascosti fuori dall'avviso\\
8 & Latenza dall'evento all'uscita sonora & non fissata & 95° percentile entro $150\unita{ms}$\\
9 & Segnale di affidabilità persa coprendo il sensore per 1 e $2\unita{s}$; avvio con i filtri accesi o con la posa spostata di $1^\circ$ & $60\unita{s}$ per i falsi segnali & segnale entro $1\unita{s}$; 0 falsi; il sistema non parte\\
10 & Uscite di due esecuzioni sugli stessi fotogrammi che differiscono solo nelle etichette di classe (permutate, assenti, cambiate in occlusione) & non fissata & celle «libero» e distanze identiche; avviso con etichette non più stretto; nessun livello più basso\\
\bottomrule
\end{tabularx}
\end{table}

Una prova con dieci passaggi non distingue $0{,}80$ da $0{,}90$ (\S\ref{sec:val}); trenta passaggi senza errori non bastano a dichiarare «9 su 10» (limite inferiore $0{,}886$). Il tempo di reazione di una persona con disabilità visiva a un avviso sonoro, da cui dipende la riga corrispondente della tabella~\ref{tab:lim-portatore}, richiede una misura con persone, fuori da questo protocollo.

\section{Veicoli veloci: sensori a lunga portata, droni e aerei}\label{sec:lp}

I capitoli precedenti descrivono il sistema con una camera di profondità ferma, fino a $3\unita{m}$. Questo capitolo porta gli stessi principi a un'altra scala: oggetti che si muovono da qualche metro al secondo fino a $210\unita{m/s}$, osservati da sensori a lunga portata. Il caso di questo capitolo è un \emph{sito aerospaziale sorvegliato}: un piazzale, un'area di prova o il varco di un hangar, con droni in volo e aerei in rullaggio. Per il sito gli aerei in volo sono un limite dichiarato (\S\ref{sec:lp-volo}); il volo ad alta velocità, caso centrale del documento, è trattato nel capitolo~\ref{sec:vo}.

Tutto ciò che segue è calcolo o simulazione. I sensori sono simulati dai valori delle schede tecniche; il moto viene da scenari scritti a mano o da traiettorie pubbliche reali, con rumore di misura e occlusioni simulati. Nessuna prova usa un sensore fisico, e nessuna rilevazione reale di sensore è stata elaborata. Le etichette di classe sono assegnate a mano: il classificatore resta un componente richiesto e non realizzato (\S\ref{sec:sti}). Il sistema osserva e avvisa un operatore: nessuna uscita è un comando, e non vi sono funzioni d'ingaggio, di puntamento o di contrasto.

I principi sono gli stessi dei capitoli~\ref{sec:map} e~\ref{sec:occ}: ciò che non è osservato non è mai «libero»; l'avviso di possibile presenza si toglie solo guardando, o rivedendo l'oggetto; nessuna etichetta di classe riduce un avviso. Cambia ciò che serve per tenerli veri a questa scala, e il capitolo mostra anche dove, in simulazione, non lo sono.

\subsection{Il problema di scala}\label{sec:lp-scala}

Un avviso serve se arriva prima che chi lo riceve abbia bisogno di agire. Per un ricevente fermo e un oggetto che si avvicina con velocità relativa $v_{\mathrm{rel}}$, la distanza a cui l'avviso deve partire è
\begin{equation}\label{eq:lp-davv}
d_{\mathrm{avv}} = v_{\mathrm{rel}}\,\bigl(\tau_{\mathrm{lat}}+T_r+T_{\mathrm{man}}\bigr)+d_{\mathrm{man}}+m ,
\end{equation}
con $\tau_{\mathrm{lat}}$ la latenza del sistema, $T_r$ il tempo di reazione del ricevente, $T_{\mathrm{man}}$ il tempo della manovra o della messa in sicurezza, $d_{\mathrm{man}}$ lo spazio di manovra (per esempio la frenata $v^2/2a_f$) e $m$ un margine. Per un ricevente in moto vale la forma della \eqref{eq:lim-dserve}, con la velocità massima dell'inviluppo al posto di quella del soggetto. L'avviso è \emph{tempestivo} se la traccia è confermata, o la zona è dichiarata «libera», a una distanza non minore di $d_{\mathrm{avv}}$; altrimenti l'uscita dice «avviso non tempestivo per $\langle$classe$\rangle$».

\begin{table}[H]\centering\footnotesize
\caption{Distanza d'avviso richiesta in sei scenari nominali, dalla \eqref{eq:lp-davv} \prova{C}. Latenze, tempi di reazione e margini sono in gran parte ipotesi di progetto; il tempo di reazione su strada è preso da una fonte secondaria \cite{wiki-stopping-sight}. Nel primo scenario il $78\,\%$ della distanza viene dai $30\unita{s}$ ipotizzati per mettere in sicurezza l'area.}\label{tab:lp-scala}
\begin{tabularx}{\linewidth}{@{}>{\RaggedRight\arraybackslash}p{3.3cm}>{\RaggedRight\arraybackslash}p{3.2cm}Lp{1.3cm}>{\RaggedRight\arraybackslash}p{1.9cm}@{}}
\toprule
Scenario & Ricevente & Tempi e margini & $v_{\mathrm{rel}}$ & $d_{\mathrm{avv}}$\\
\midrule
drone (15--$21\unita{m/s}$) dietro un edificio del sito & operatore di sito & latenza $1\unita{s}$, reazione $5\unita{s}$, $30\unita{s}$ di messa in sicurezza, margine $50\unita{m}$ (ipotesi) & $21\unita{m/s}$ & $806\unita{m}$\\
aereo in rullaggio ($9\unita{m/s}$) dietro un hangar & conducente di un mezzo che attraversa la via & $1\unita{s}$ + $2{,}5\unita{s}$; frenate $1{,}0$ e $1{,}5\unita{m/s^2}$; margine $23\unita{m}$ & $17{,}3\unita{m/s}$ & $147\unita{m}$; tramite la torre $277\unita{m}$\\
aereo in finale e drone sotto la copertura radar & operatore di torre & $4+10+12{,}5\unita{s}$, più $10\unita{s}$ di riattaccata; margine $610\unita{m}$ & $91\unita{m/s}$ & $3{,}9\unita{km}$\\
drone con i sensori a bordo di un altro drone (solo analisi) & pilota remoto; sistema di bordo & $0{,}5+3\unita{s}$; di bordo $0{,}4\unita{s}$ & $21\unita{m/s}$ & $126\unita{m}$; $60\unita{m}$\\
auto a un incrocio dietro un camion & conducente; sistema di bordo & $0{,}1+2{,}5\unita{s}$, $3{,}4\unita{m/s^2}$; di bordo $0{,}5\unita{s}$, $6{,}9\unita{m/s^2}$ & $13{,}9\unita{m/s}$ & $67\unita{m}$; $24\unita{m}$\\
treno in stazione dietro un convoglio fermo; treno in linea & personale e macchinista & $1+3\unita{s}$, frenata $0{,}8\unita{m/s^2}$; in linea $600\unita{m}$ di frenata d'emergenza & $16{,}7$; $44{,}4\unita{m/s}$ & $261\unita{m}$; $798\unita{m}$\\
\bottomrule
\end{tabularx}
\end{table}

Dalla tabella~\ref{tab:lp-scala} seguono due fatti \prova{C}. Il primo: la portata di $3\unita{m}$ non entra in nessuno scenario, perché la distanza minima richiesta è $24\unita{m}$; serve un altro ordine di grandezza, da $10^2\unita{m}$ sul piazzale a $10^3\unita{m}$ e oltre in avvicinamento. Il secondo: nelle geometrie scelte l'occlusione cade dentro la distanza d'avviso. Il drone riappare a $375\unita{m}$ dall'operatore, che ne chiede $806$; l'aereo esce dall'ombra dell'hangar a $63\unita{m}$ dall'incrocio, contro $147$ richiesti. Un sistema che avvisa solo quando rivede l'oggetto arriva tardi: l'avviso deve partire prima dell'occlusione e restare acceso per tutta la sua durata. Il secondo fatto vale per le geometrie della tabella, non in generale.

\subsection{Sensori simulati dalle schede}\label{sec:lp-sensori}

La tabella~\ref{tab:lp-sensori} riassume i sensori simulati. I valori vengono dalle schede dei costruttori \cite{robin-iris,terma-scanter5000,hesai-at128,livox-tele15}, dalle norme statunitensi sull'ADS-B e sull'identificazione remota dei droni \cite{cfr-91-227,cfr-89-310} e da conti; ogni campo senza fonte è un'ipotesi. In tutto i parametri sono 124, di cui 57 ($46\,\%$) ipotesi: tutte le probabilità di rilevamento, i tassi di falso allarme, le latenze e i costi. Le schede non dichiarano a quale probabilità di rilevamento si riferisce la portata.

\begin{table}[H]\centering\footnotesize
\caption{Sensori simulati e valori usati. (i): ipotesi, senza una fonte aperta. Le portate valgono per la sezione o la riflettività indicate dalla scheda, non per qualunque oggetto.}\label{tab:lp-sensori}
\begin{tabularx}{\linewidth}{@{}>{\RaggedRight\arraybackslash}p{3.3cm}>{\RaggedRight\arraybackslash}p{3.0cm}>{\RaggedRight\arraybackslash}p{2.0cm}L>{\RaggedRight\arraybackslash}p{1.6cm}@{}}
\toprule
Sensore (esemplare) & Portata & Campo & Errore di misura & Rivisita\\
\midrule
radar per droni (Robin IRIS) \cite{robin-iris} & $1{,}1$--$7{,}6\unita{km}$ secondo il drone; strumentata 5 o $12\unita{km}$ & $360^\circ\times60^\circ$ & $0{,}6\unita{m}$; $0{,}6$--$0{,}75^\circ$; $1^\circ$ in elevazione & $1\unita{s}$\\
radar a pannello (tipo EchoGuard) & (i) & (i) & $<3{,}25\unita{m}$; $<1^\circ$; $<1{,}5^\circ$ in elevazione \cite{echodyne-echoguard} & (i)\\
radar del movimento a terra (Terma SCANTER 5000) \cite{terma-scanter5000} & $5\unita{km}$ per $1\unita{m^2}$ (i); $9{,}6\unita{km}$ in avvicinamento & $360^\circ$, senza quota & cella di $3\unita{m}$; $0{,}29^\circ$ (conto) & $1\unita{s}$\\
LiDAR a campo largo (Hesai AT128) \cite{hesai-at128} & $210\unita{m}$ al $10\,\%$ di riflettività & $120^\circ\times25{,}4^\circ$ & $0{,}1^\circ\times0{,}2^\circ$; $3\unita{cm}$ (i) & $0{,}1\unita{s}$ (i)\\
LiDAR a campo stretto (Livox Tele-15) \cite{livox-tele15} & $320\unita{m}$ al $10\,\%$, $500\unita{m}$ al $50\,\%$ & $14{,}5^\circ\times16{,}2^\circ$ & $2\unita{cm}$; $<0{,}03^\circ$ & $0{,}1\unita{s}$\\
radar 4D per veicoli & $300\unita{m}$ (i: scheda non aperta) & $120^\circ$ (i) & $0{,}22\unita{m}$; $1{,}2^\circ$; $0{,}1\unita{m/s}$ (i) & $0{,}05\unita{s}$ (i)\\
ADS-B \cite{cfr-91-227} & orizzonte del ricevitore & solo chi trasmette & $<92{,}6\unita{m}$; $<10\unita{m/s}$ & $\le1\unita{s}$; latenza $\le2\unita{s}$\\
identificazione remota dei droni \cite{cfr-89-310} & circa $1\unita{km}$ (i) & solo chi trasmette & $30{,}5\unita{m}$ al $95\,\%$; quota $45{,}7\unita{m}$ & $\le1\unita{s}$; latenza $\le1\unita{s}$\\
camera nel visibile, 12 Mpixel, $300\unita{mm}$ (i) & drone di $0{,}3\unita{m}$ a $13\unita{km}$, limite geometrico \cite{wiki-johnson} & $2{,}7^\circ$ & $11{,}5\unita{\mu rad}$ per pixel; nessuna distanza & $0{,}1\unita{s}$\\
camera termica $640\times480$, $100\unita{mm}$ (i) & drone a $1{,}2\unita{km}$, limite geometrico & $4{,}4^\circ$ & $120\unita{\mu rad}$; nessuna distanza & $0{,}1\unita{s}$\\
\bottomrule
\end{tabularx}
\end{table}

\paragraph{Portata del radar e sezione radar.} La scheda del radar per droni dà portate per classe di drone, senza la sezione radar. Attribuendo $0{,}01\unita{m^2}$ (i) al drone da $0{,}9\unita{kg}$ visto a $3100\unita{m}$, la portata scala come $R(\sigma)=3100\unita{m}\cdot(\sigma/0{,}01\unita{m^2})^{1/4}$: $3{,}1\unita{km}$ per un multirotore piccolo, $9{,}8\unita{km}$ per un velivolo leggero di $1\unita{m^2}$, $12\unita{km}$ (limite strumentato) per un aereo di linea \prova{C}. Un drone da $0{,}249\unita{kg}$ si vede solo entro $1{,}1\unita{km}$.

\paragraph{Probabilità di rilevamento (modello, i).} Per un oggetto fluttuante si usa il modello di Swerling~1 per un impulso, $P_d=P_{fa}^{1/(1+\mathrm{SNR})}$, con $\mathrm{SNR}\propto\sigma/R^4$. Con $P_d=0{,}5$ alla portata dichiarata e $P_{fa}=10^{-6}$ il rapporto segnale-rumore di riferimento è $18{,}9$ ($12{,}8\unita{dB}$). Se la portata di scheda valesse invece a $P_d=0{,}9$, la portata utile sarebbe il $62\,\%$ di quella usata: è l'ipotesi che pesa di più sull'intero modello.

\paragraph{Che cosa vuol dire «non osservato» a questa scala.} Una zona è non osservata per cinque cause, che l'uscita distingue: fuori campo; senza linea di vista (dietro un hangar, sotto l'orizzonte radar); \emph{sotto la soglia}, cioè in campo e in vista ma con una probabilità di rilevamento insufficiente per \emph{quella} classe; fra due rivisite (a $70\unita{m/s}$ un aereo si sposta di $70\unita{m}$ in un secondo, più di cento volte l'errore in distanza); oggetto che non coopera o non emette. La terza causa è nuova rispetto alla camera di profondità: la stessa zona può essere osservata per un aereo e non osservata per un drone piccolo. I sensori cooperativi (ADS-B, identificazione remota) aggiungono presenze e identità, ma il loro silenzio non osserva niente: un trasponditore spento non è un'assenza.

\subsection{Rappresentazione per sensore e «libero» per classe}\label{sec:lp-libero}

Ogni sensore ha una propria griglia polare in distanza, azimut ed elevazione, con i quattro stati del capitolo~\ref{sec:map}: libero, occupato, in ombra, non osservato. L'ombra è calcolata dagli ostacoli in pianta e dalla maschera del terreno. Il «libero» non si scrive più da un solo fotogramma, perché un radar che ruota illumina una cella una volta per giro e con una probabilità di rilevamento minore di uno: si scrive per classe, cumulando le passate.

\begin{definizione}[«Libero» per classe]\label{def:lp-libero}
Per una classe $c$ di sezione $\sigma_c$ e velocità massima $v_{\max,c}$, sia
\[W_c=\min\bigl(t_{\max},\,\ell/v_{\max,c}\bigr)\]
la finestra oltre la quale un oggetto può aver cambiato cella di lato $\ell$, e sia $n$ il numero di passate che hanno illuminato la cella in $[t-W_c,t]$ dopo l'ultimo ritorno. La probabilità cumulata di averlo visto è
\[
P_{c}=1-\bigl(1-P_d(R,\sigma_c)\bigr)^{n_{\mathrm{eff}}},\qquad n_{\mathrm{eff}}=1+(n-1)(1-\rho_c),
\]
con $\rho_c$ la correlazione fra passate, posta a 1 in mancanza di una misura ($n_{\mathrm{eff}}=1$), e presa come minimo sulla cella e sulle otto vicine. La cella è «libera per $c$» se non è in ombra, la linea di vista è accertata e $P_c\ge0{,}99$. È «libera» se lo è per ogni classe ammessa dallo scenario; altrimenti l'uscita dice «zona non verificata per $\langle$classe$\rangle$».
\end{definizione}

\begin{proposizione}[La regola con $n_{\mathrm{eff}}=1$ è prudente per ogni correlazione]\label{prop:lp-correlate}
Siano $D_1,\dots,D_n$ gli eventi di rilevamento nelle $n$ passate, con la stessa probabilità marginale $P_d$ e dipendenza qualunque. Allora
\[\Pr(\text{nessun rilevamento})\le 1-P_d :\]
la probabilità di non vedere un oggetto presente non supera quella che la regola con $n_{\mathrm{eff}}=1$ dichiara.
\end{proposizione}
\begin{proof}
L'evento «nessun rilevamento nelle $n$ passate» è contenuto nell'evento «nessun rilevamento nella prima», che ha probabilità $1-P_d$. Nella forma del modello a sezione fluttuante, con $P_d$ condizionata a $\sigma$ variabile, vale $\mathbb{E}[(1-P_d)^n]\le\mathbb{E}[1-P_d]$ perché $0\le1-P_d\le1$.
\end{proof}
La proposizione è \prova{D} nel modello delle passate. Il prezzo è la portata: per un multirotore di $0{,}01\unita{m^2}$, con un radar di riferimento di $P_d=0{,}9$ a $5\unita{km}$ su $1\unita{m^2}$, rivisita di $1\unita{s}$ e celle di $100\unita{m}$, il «libero» arriva a $2360\unita{m}$ se si contano 5 passate indipendenti e a $877\unita{m}$ con $n_{\mathrm{eff}}=1$ \prova{C}. Contare le passate come indipendenti è sbagliato quando l'aspetto dell'oggetto non cambia, per esempio un drone fermo in volo: con la sezione tenuta per 5 passate la probabilità di mancarle tutte alla portata di $2360\unita{m}$ è $0{,}263$, contro $0{,}010$ dichiarato (Monte Carlo di $400\,000$ casi) \prova{S}.

\paragraph{Più sensori.} Lo stato fuso segue l'ordine occupato, libero, in ombra, non osservato. Fra sensori di tipo diverso la probabilità di mancare l'oggetto è il prodotto delle singole; fra sensori dello stesso tipo si prende la migliore, perché l'aspetto può essere correlato (ipotesi). Un sensore con filtro Doppler non produce mai «libero» e il suo silenzio non conta per un oggetto con velocità radiale sotto la soglia del filtro; ogni ritorno toglie il «libero» alle celle entro $3\sigma$ dal punto misurato, più una cella. Un sensore in perdita di affidabilità rende tutta la sua griglia «non osservata». Un ritorno recente di un radar vieta il «libero» attorno al punto misurato per $5\unita{s}$: una sola passata mancata non basta ad aprirlo. Il prezzo di questa ritenuta, cioè le chiusure tardive dell'avviso, non è misurato.

\paragraph{La mappa come ingresso di sicurezza.} Una passata conta come sguardo solo se la linea di vista è accertata: o la mappa degli ostacoli è dichiarata completa, o sul raggio era atteso un ritorno dello sfondo oltre la cella ed è stato misurato alla distanza attesa entro $2\unita{m}$. Se lo sfondo manca il raggio non è accertato; se è più vicino dell'atteso c'è un ostacolo non in mappa, e le celle oltre vanno in ombra. In mancanza di un'indicazione esplicita la mappa non è dichiarata completa, e non c'è nessun «libero». Il cielo non ha sfondo: lì vale solo la mappa dichiarata. Su casi costruiti la regola regge: con un edificio assente dalla mappa e la mappa non dichiarata completa la cella dietro resta non osservata, e con lo sfondo misurato a $400$ invece che a $1500\unita{m}$ va in ombra e l'ostacolo è annotato \prova{S}. Se invece la mappa è dichiarata completa ed è sbagliata, un edificio che non vi compare apre un «libero» falso: è un controesempio (5 casi costruiti su 6). Per questo nel prototipo, per impostazione, la mappa non è dichiarata completa, e nei sei casi la cella dietro l'edificio resta «non osservata» \prova{S}. La convalida della mappa con lo sfondo misurato non è ancora usata nel simulatore (\S\ref{sec:lp-smentiti}).

\paragraph{Classi escluse.} Lo scenario dichiara le classi ammesse e quelle escluse; per le escluse l'uscita dice sempre «zona non verificata». Nel sito sorvegliato è escluso il drone sotto i $250\unita{g}$, che il radar non vede oltre circa $1{,}1\unita{km}$, e sono esclusi gli aerei in volo. Il «libero», quindi, vale solo per le classi dichiarate, con la mappa data e con passate trattate come correlate.

\subsection{Inseguimento multisensore}\label{sec:lp-inseguimento}

Lo stato di una traccia è posizione e velocità in tre dimensioni nel riferimento fisso Est--Nord--Alto. Il filtro è una miscela di modelli interagenti \cite{barshalom2001} (velocità quasi costante, virata coordinata a sinistra e a destra, manovra), con parametri per classe; una traccia senza classe usa tutti i modelli con i parametri dell'inviluppo. Le misure di tutti i sensori si fondono una alla volta, nell'ordine d'arrivo: previsione all'istante della misura, porta sulla distanza di Mahalanobis della posizione ($d^2\le16{,}27$, quantile $0{,}999$ di un chi quadro con 3 gradi di libertà), assegnazione, aggiornamento in forma di Joseph. Una misura appena fuori porta ($16{,}27<d^2\le30{,}66$) non apre una traccia: diventa una «misura dubbia», e solo la terza in $5\unita{s}$ apre una traccia nuova. La conferma chiede 2 misure su 3. Le misure cooperative portano un identificativo: lo stesso identificativo è la stessa traccia, identificativi diversi non si associano mai, e una traccia vista solo dal sensore cooperativo è dichiarata «non verificata da un sensore». Una traccia si chiude solo come «non più rilevata»: in vista, con la probabilità cumulata dei mancati rilevamenti, per la classe più piccola ammessa, almeno $0{,}99$. In ombra, fuori copertura o con il sensore non affidabile non si accumula nulla, e la traccia resta con l'avviso.

In simulazione \prova{S}, con un aereo a $70\unita{m/s}$ in virata per $15\unita{s}$, un radar a $1\unita{Hz}$, una camera a $5\unita{Hz}$ e un sensore cooperativo a $2\unita{Hz}$ spento a metà prova, l'errore quadratico medio di posizione è $39{,}9\unita{m}$ con il solo radar e $4{,}4\unita{m}$ con la fusione, che è migliore in 40 semi su 40, con 0 cambi d'identità. Con due oggetti senza classe che si incrociano a $100\unita{m}$ di quota di separazione, gli scambi d'identità sono 0 su 40 semi $[0;\,0{,}088]$. Sono casi costruiti senza falsi ritorni: con il disturbo e con più di due oggetti l'identità non è provata.

\subsection{Insieme di possibile presenza e spegnimento dell'avviso}\label{sec:lp-presenza}

L'avviso è l'insieme di possibile presenza del capitolo~\ref{sec:occ}, calcolato con l'\emph{inviluppo dello scenario}: la velocità e l'accelerazione massime fra tutte le classi ammesse, mai quelle dell'etichetta. Con centro nell'ultima posizione stimata, $\tau$ il tempo dall'ultima misura e $v_0=\min(|\hat v|+3\sigma_v,\,v_{\max})$, il raggio è
\begin{equation}\label{eq:lp-raggio}
r(\tau)=R(\tau;v_0,v_{\max},a_{\max})+3\sigma_p+w,\qquad
R=\begin{cases} v_0\tau+\tfrac12 a_{\max}\tau^2 & \tau\le\tau_1,\\ R(\tau_1)+v_{\max}(\tau-\tau_1) & \tau>\tau_1,\end{cases}
\end{equation}
con $\tau_1=(v_{\max}-v_0)/a_{\max}$ e $w=50\unita{m}$ (ipotesi). L'insieme è una sfera, salvo che lo scenario ammetta solo classi vincolate a un binario: allora è un tratto di binario, che a $30\unita{s}$ riduce una sfera di $2{,}5\unita{km}$ di raggio a un tratto di $675\unita{m}$ per il treno ad alta velocità. Il vincolo vale perché il treno non può lasciare il binario, non per l'etichetta.

\begin{proposizione}[Contenimento nell'inviluppo]\label{prop:lp-contenimento}
Se l'oggetto ha velocità al più $v_{\max}$ e accelerazione al più $a_{\max}$, e la sua posizione e velocità all'ultima misura stanno entro $3\sigma_p$ e $3\sigma_v$ dalla stima, allora a ogni $\tau\ge0$ l'oggetto sta nella sfera di raggio $r(\tau)$ della \eqref{eq:lp-raggio}.
\end{proposizione}
\begin{proof}
La velocità al tempo $s$ è al più $\min(v_0+a_{\max}s,\,v_{\max})$; integrando fra $0$ e $\tau$ si ottiene $R(\tau)$, a cui si aggiunge l'errore di posizione iniziale.
\end{proof}
La proposizione è \prova{D} nel modello e vale quanto valgono $v_{\max}$ e $a_{\max}$, che sono ipotesi o valori presi dai dati. Su $60\,000$ traiettorie simulate entro i limiti dichiarati (10 classi, 2000 traiettorie, 3 orizzonti) nessuna è uscita dall'insieme \prova{S}; è un controllo del conto, non una frequenza. L'inviluppo usato nel sito sorvegliato comprende i multirotori, i droni ad ala fissa e gli aerei in rullaggio; per i multirotori veloci assume $45\unita{m/s}$ e $25\unita{m/s^2}$, valori di ipotesi senza fonte. Per gli aerei in volo l'inviluppo è portato a $210\unita{m/s}$, perché i dati pubblici mostrano decolli a più di $200\unita{m/s}$ (\S\ref{sec:lp-traiettorie}). Su casi costruiti, un drone che scatta da fermo con $20$--$25\unita{m/s^2}$ fino a $45\unita{m/s}$ e un aereo a $205\unita{m/s}$ restano dentro l'avviso a ogni $\tau$ \prova{S}; oltre l'ipotesi, a $60\unita{m/s}$, il drone esce dall'avviso dopo $8\unita{s}$.

\paragraph{Spegnimento prudente.} L'avviso si spegne solo in due casi: l'oggetto è rivisto, oppure ogni cella che tocca l'insieme di possibile presenza è «libera» per ogni classe ammessa, con passate successive alla perdita. Il controllo è esatto sulle celle: una griglia affidabile deve contenere tutta la sfera, e ogni cella con il centro entro il raggio più mezza diagonale deve essere libera. Non c'è alcuno spegnimento a tempo. Poiché l'inviluppo non dipende dalla classe, nessuna etichetta riduce l'avviso, per costruzione \prova{D}; nel prototipo, su un passaggio, l'avviso è identico con e senza etichette in 228 cicli su 228, con l'inviluppo precedente degli aerei (cicli correlati: è un controllo, non una frequenza) \prova{S}.

\paragraph{Perché la regola ingenua cade.} Una regola più semplice conta come prova d'assenza una passata senza ritorno sul punto previsto dell'oggetto, o su un campione di punti dell'insieme, e chiude la traccia dopo abbastanza passate. Ma l'oggetto nascosto non è nel punto previsto: è da qualche parte nell'insieme, che a $10\unita{s}$ ha per un multirotore un raggio di circa $200\unita{m}$. In simulazione, con quella regola, l'avviso si spegne senza che si sia guardata tutta la zona in 564 passaggi su 2400 $[0{,}218;\,0{,}252]$; con la regola prudente, sugli stessi semi, in 0 su 600 $[0;\,0{,}006]$ \prova{S}. Con la regola prudente nel sistema, 0 passaggi su 200 per configurazione restano senza regione.

\paragraph{La zona prevista.} Oltre all'avviso, per una traccia con classe il sistema emette una zona di previsione al livello $0{,}90$: una sfera centrata sulla previsione a velocità costante, con raggio dato da un quantile conforme dell'errore \cite{angelopoulos2021}, tarato per classe a 1, 2, 5, 10 e $30\unita{s}$, più $3(\sigma_p+\sigma_v\tau)$. La zona usa l'etichetta, e per questo non è l'avviso. Diventa «previsto ampio» oltre l'ultimo orizzonte tarato, quando il suo raggio supera $0{,}8$ di quello dell'avviso, per una traccia senza classe o con un sensore non affidabile. Per un multirotore l'insieme raggiungibile ha raggio $105\unita{m}$ a $5\unita{s}$ e $210\unita{m}$ a $10\unita{s}$ \prova{C}: oltre pochi secondi la zona non dice nulla di più dell'avviso, e l'avviso si toglie solo tornando a guardare il volume, con un altro sensore o da un altro punto di vista.

\subsection{Previsione su traiettorie pubbliche reali}\label{sec:lp-traiettorie}

Il moto reale viene da sette archivi pubblici (tabella~\ref{tab:lp-dati}). Il rumore di misura è simulato dalle schede, le occlusioni sono inserite a mano (10 istanti a caso per tratto e per durata), e c'è un solo oggetto per scena, senza associazione. Su metà dei tratti si tarano il rumore di processo e il raggio della zona; le cifre sono sull'altra metà. Poiché più tratti appartengono allo stesso mezzo, la divisione è stata rifatta \emph{per mezzo}, e gli intervalli sono un bootstrap a grappoli per mezzo (2000 ricampionamenti); un Wilson con $n$ uguale al numero di mezzi dà un limite ancora più prudente. La verità è la posizione di bordo (GNSS o ADS-B), non una misura esterna.

\begin{table}[H]\centering\footnotesize
\caption{Archivi pubblici di traiettorie usati per il moto. I tratti sono quelli convertiti; il passo è quello della fonte.}\label{tab:lp-dati}
\begin{tabularx}{\linewidth}{@{}>{\RaggedRight\arraybackslash}p{2.2cm}L>{\RaggedRight\arraybackslash}p{1.6cm}>{\RaggedRight\arraybackslash}p{1.6cm}>{\RaggedRight\arraybackslash}p{3.3cm}@{}}
\toprule
Classe & Fonte e licenza & Tratti & Passo & Verità\\
\midrule
aerei & ricezione ADS-B attorno ad Atlanta, 13 minuti \cite{adsblol}; ODbL 1.0 & 53 & 3,1--$5{,}1\unita{s}$ & ADS-B, non elaborato\\
aerei di linea & traiettorie ADS-B Londra--Zurigo 2025 \cite{zenodo-lhr-zrh}; CC BY 4.0 & 80 & $1\unita{s}$ & ADS-B filtrato e interpolato dalla fonte\\
aviazione generale & TrajAir \cite{trajair}; CC BY 4.0 & 60 & $1\unita{s}$ & ADS-B elaborato dalla fonte\\
multirotore & voli di un quadricottero per consegne \cite{cmu-drone}; CC BY 4.0 & 196 & $0{,}12\unita{s}$ & GNSS e inerziale di bordo\\
ala fissa senza pilota & ALFA \cite{alfa}; CC BY 4.0 & 47 & $0{,}24\unita{s}$ & GNSS dell'autopilota; guasti indotti\\
auto & NGSIM, US-101 \cite{ngsim}; CC BY-SA 3.0 & 157 & $0{,}1\unita{s}$ & telecamere dall'alto\\
treni & posizioni dei treni finlandesi, 12 novembre 2026 \cite{digitraffic}; CC BY 4.0 & 111 & $3{,}0\unita{s}$ & GPS di bordo, senza quota\\
\bottomrule
\end{tabularx}
\end{table}

\begin{table}[H]\centering\footnotesize
\caption{Copertura della zona dichiarata al livello $0{,}90$ sulla metà di prova, con la divisione per mezzo \prova{P}. Il rumore è simulato dalle schede; le occlusioni sono virtuali.}\label{tab:lp-traiettorie}
\begin{tabularx}{\linewidth}{@{}>{\RaggedRight\arraybackslash}p{3.6cm}>{\RaggedRight\arraybackslash}p{1.6cm}>{\RaggedRight\arraybackslash}p{2.1cm}L>{\RaggedRight\arraybackslash}p{2.0cm}@{}}
\toprule
Classe; rumore simulato & $\tocc$ & Copertura & Bootstrap per mezzo & Wilson, $n$ = mezzi\\
\midrule
aerei, 43 aerei; $\sigma=7{,}5\unita{m}$ a $1\unita{s}$ & 2 / 10 / $30\unita{s}$ & 0,95 / 0,93 / 0,93 & $[0{,}92;0{,}98]$ / $[0{,}90;0{,}96]$ / $[0{,}88;0{,}99]$ & $\ge0{,}83$\\
multirotore, \emph{un} esemplare, 98 voli; radar per droni a $1{,}5\unita{km}$ & 5 / $10\unita{s}$ & 0,99 / 0,96 & per volo: $[0{,}99;1{,}00]$ / $[0{,}95;0{,}97]$ & per volo: $\ge0{,}95$ / $\ge0{,}90$\\
auto, 30 veicoli; radar 4D a $100\unita{m}$ & 2 / 10 / $30\unita{s}$ & 0,93 / 0,87 / 0,86 & $[0{,}90;0{,}95]$ / $[0{,}84;0{,}90]$ / $[0{,}82;0{,}90]$ & $\ge0{,}68$\\
treni, 19 treni; LiDAR a campo stretto a $300\unita{m}$ & 2 / 10 / $30\unita{s}$ & 0,91 / 0,88 / 0,90 & $[0{,}87;0{,}95]$ / $[0{,}84;0{,}92]$ / $[0{,}87;0{,}93]$ & $\ge0{,}67$\\
\bottomrule
\end{tabularx}
\end{table}

La forma più forte che resta vera per gli aerei è: su 43 aerei reali, dopo 2--$30\unita{s}$ senza vederli, la zona dichiarata conteneva l'aereo in circa 9 casi su 10 \prova{P}. Per i droni i dati reali sono due esemplari soltanto, un multirotore e un ala fissa: le cifre valgono per quel multirotore, e per l'ala fissa con guasto indotto la copertura a $10\unita{s}$ scende a $0{,}656$ $[0{,}584;\,0{,}721]$ (21 tratti, divisione per tratto). Nelle virate la zona copre meno: $0{,}74$ per gli aerei e $0{,}72$ per i droni fino a $10\unita{s}$ (divisione per tratto, casi correlati). Per i treni la verità GPS è sporca: con essa la zona tarata diventa larga quanto l'avviso, e non si dichiara una precisione.

\paragraph{L'avviso sulle traiettorie.} Con un inviluppo di $90\unita{m/s}$ per gli aerei l'insieme di possibile presenza non conteneva l'aereo vero in 105 casi su 722 a $2\unita{s}$, e in 167 su 730 a $10\unita{s}$, in 8--11 aerei su 43 \prova{P}: erano decolli e avvicinamenti reali fino a $205\unita{m/s}$. Il principio regge, perché l'avviso non usa l'etichetta; era sbagliato il valore del limite. Per questo l'inviluppo degli aerei in volo è $210\unita{m/s}$; con questo valore la prova sulle traiettorie non è stata rifatta \prova{N}. Anche per i treni l'avviso non conteneva il mezzo in 61 casi su 405 a $2\unita{s}$, in 18 treni su 19, con la verità GPS sporca: la causa non è stata accertata, e per i treni l'avviso non è garantito.

\subsection{Eventi improvvisi e visibilità degradata}\label{sec:lp-eventi}

Gli eventi trattati sono una collisione fra due droni con detriti, l'avaria di un motore con fiammata e fumo durante la corsa di un aereo, la nebbia e la pioggia, il fumo e i detriti in pista, e, come controllo, un passaggio ravvicinato senza collisione. Restano situazioni da osservare e segnalare. Le regole aggiungono alla cella non osservata una \emph{causa} (ostacolo, nube, abbagliamento, meteo, fuori portata, sensore non affidabile, mai guardata), che è un'etichetta in più e non uno stato nuovo. I tre principi non cambiano.

\paragraph{Frammentazione.} Un evento nasce \emph{dedotto} da un innesco di «separazione» (almeno 3 ritorni nuovi confermati entro $2\unita{s}$ attorno a una traccia) o di «contatto» (due tracce più vicine della somma dei mezzi ingombri più $3\sigma$, con almeno un indizio fra ritorni nuovi, residuo fuori porta per 2 passate o una traccia non più rilevata in vista). Diventa \emph{osservato} quando un sensore in salute risolve almeno 3 ritorni distinti per 2 passate; l'avviso è lo stesso nei due casi. Attorno all'ultimo stato comune osservato $p_0$, all'istante $t_a$, i frammenti sono contenuti in un volume $V(\tau)$: in pianta l'unione dei segmenti $[p_0,\,p_0+v_{c,h,i}\tau^*]$ percorsi alle ultime velocità delle tracce d'origine, ingrossata di
\[
r(\tau)=W\tau+(\Delta v+W)\,\tau^*+r_0 ,
\]
con $\Delta v$ la velocità relativa massima dei frammenti, $W$ il vento orizzontale massimo, $\tau^*$ il tempo di volo massimo e $r_0$ l'incertezza iniziale; in quota fra due estremi dati dalla gravità e dal vento verticale.

\begin{proposizione}[Volume dei detriti]\label{prop:lp-detriti}
Per un punto materiale con resistenza isotropa opposta alla velocità relativa all'aria, con coefficiente costante qualunque, vento costante entro $W$ e $W_z$, e velocità iniziale relativa entro $\Delta v$, il frammento sta in $V(\tau)$ per ogni $\tau\le\tau^*$.
\end{proposizione}
\begin{proof}[Traccia della dimostrazione]
Sia $r=v-w$ la velocità relativa all'aria. In orizzontale $\dot r_h=-k|r|\,r_h$: la direzione non cambia e il modulo non cresce, quindi lo spostamento rispetto all'aria è $\lambda\,r_h(0)\,\tau$ con $\lambda\in[0,1]$, e la posizione sta nel segmento ingrossato di $W\tau+(\Delta v+W)\tau$. In verticale, dove $r_z\le0$ vale $\dot r_z\ge-g$ e dove $r_z\ge0$ vale $\dot r_z\le-g$; integrando si hanno i due estremi di quota.
\end{proof}
Il risultato è \prova{D} nel modello per $\tau\le\tau^*$; il taglio a $\tau^*$, il tempo di caduta e il raggio dell'impronta a terra sono sostenuti dal calcolo, non dimostrati. Con $10\,000$ frammenti simulati per caso nessun campione è uscito dal volume (0 su $1\,131\,692$ in 8 casi conformi, 0 su $646\,957$ con raffiche entro i limiti) \prova{S}; il simulatore condivide il modello fisico con la prova, e vede gli errori d'algebra e di codice, non quelli del modello. Se la velocità vera dei frammenti è $1{,}5$ volte quella dichiarata, fra 7 e 88 campioni escono. Restano fuori dal modello la portanza dei pezzi piatti, la corrente ascendente di un incendio, il rotolamento a terra e i frammenti più veloci di $\Delta v$, come la punta di un'elica. Dentro il volume nessuna cella diventa «libera» se non con la probabilità cumulata $0{,}99$ per il frammento minimo ammesso ($5\unita{cm}$); a terra l'impronta resta «zona non verificata: possibili detriti» finché non è ispezionata, e con i sensori della tabella~\ref{tab:lp-sensori} non si scioglie mai da sola.

\paragraph{Nube, abbagliamento, meteo.} Per ogni sensore la nube è l'insieme delle celle fra il primo ritorno anomalo e lo sfondo atteso dalla mappa, stimato su 3 cicli, più $10\unita{m}$; dentro e dietro la nube la cella è non osservata per quel sensore, e solo i radar vedono dietro il fumo (ipotesi prudente). La nube si toglie quando lo sfondo torna visibile per $5\unita{s}$, e le celle tornano libere solo con la probabilità cumulata. L'abbagliamento si riconosce da indizi del sensore stesso (pixel saturi oltre il $5\,\%$ nella camera nel visibile e il $2\,\%$ nella termica, caduta dei ritorni LiDAR oltre il $30\,\%$, fondo radar più alto di $6\unita{dB}$): entro 3 cicli il sensore è dichiarato «affidabilità persa», con la causa, e le sue passate non contano dall'inizio dell'indizio. Con il meteo la portata dichiarata si accorcia per un fattore prudente e non si allunga mai da sola; se il «libero» scende sotto la distanza d'avviso l'uscita dice «avviso non tempestivo», e se poggia su un solo sensore «ridondanza assente». La tabella~\ref{tab:lp-meteo} mostra perché, con nebbia e pioggia forte, tutto poggia sul radar \prova{C}.

\begin{table}[H]\centering\footnotesize
\caption{Distanza del «libero» in metri per sensore e condizione, con i fattori prudenti assunti \prova{C}.}\label{tab:lp-meteo}
\begin{tabularx}{\linewidth}{@{}Lrrrrl@{}}
\toprule
Oggetto, condizione & radar per droni & camera nel visibile & termica & senza radar & avviso a $1490\unita{m}$ senza radar\\
\midrule
multirotore, sereno & 2858 & 2000 & 1200 & 2000 & sì\\
multirotore, pioggia forte & 2000 & 526 & 353 & 526 & no\\
multirotore, nebbia fitta & 2801 & 23 & 15 & 23 & no\\
aereo leggero, nebbia fitta & 4900 & 23 & 15 & 23 & no\\
\bottomrule
\end{tabularx}
\end{table}

\paragraph{Risultati in simulazione.} I conteggi sono per passaggio, con intervalli di Wilson al $95\,\%$; le due serie sono indipendenti nei semi, non nel codice.
\begin{itemize}
\item \emph{Collisione fra droni.} L'evento è dedotto da $1{,}2$ a $2{,}3\unita{s}$ dopo l'urto (mediana $1{,}8\unita{s}$); entro $2\unita{s}$ in 51 passaggi su 60 $[0{,}739;\,0{,}919]$, e con il tipo giusto («contatto») in 50 su 60 $[0{,}720;\,0{,}907]$: in due passaggi è dedotto come «separazione» \prova{S}. Queste cifre vengono da un simulatore dei frammenti che integrava la resistenza con il metodo di Eulero esplicito, instabile (\S\ref{sec:vo-scenari}). Con l'integratore stabile l'evento è dichiarato entro $2\unita{s}$ in 48 passaggi su 60 $[0{,}682;\,0{,}882]$, con 3 eventi su 60 $[0{,}017;\,0{,}137]$ aperti circa $11\unita{s}$ prima dell'urto, e i frammenti restano nel volume in 60 su 60 $[0{,}940;\,1]$ \prova{S}: per l'evento è un limite dichiarato. La soglia di progetto, entro $2\unita{s}$ in ogni passaggio, non è raggiunta. Nessuna cella è dichiarata «libera» su un frammento: 0 passaggi su 60 $[0;\,0{,}06]$, per costruzione. L'origine dei detriti è presa solo dalle tracce viste negli ultimi $3\unita{s}$: il raggio residuo della zona dei detriti va da $3{,}1$ a $5{,}3\unita{km}$ (mediana $4{,}5\unita{km}$), mai oltre $10\unita{km}$ (0 su 60 $[0;\,0{,}06]$). La zona non si riduce guardando: l'ispezione non ha tolto punti in nessun passaggio (0 su 60), e l'avviso resta su tutte le tracce dentro la zona. Nessun evento è dichiarato prima dell'urto (0 su 60 $[0;\,0{,}06]$), ma la regola delle tracce recenti è stata scritta dopo aver esaminato un passaggio con un evento falso prima dell'urto: questo esito non è una previsione. Con la versione precedente delle regole i frammenti stavano nel volume in una quota di $0{,}9973$, sotto il $0{,}999$ richiesto; con le regole attuali la ripetizione con semi nuovi non è rifatta.
\item \emph{Avaria con fiammata.} L'evento è dedotto in $0{,}2$--$0{,}3\unita{s}$ in 60 passaggi su 60 $[0{,}94;\,1]$, ma solo con un innesco ottimista, perché la saturazione è presa dal punto vero \prova{S}. Le camere sono dichiarate non affidabili entro 3 cicli in 90 passaggi su 90 $[0{,}959;\,1]$. Con le regole attuali le camere sono dichiarate non affidabili entro 3 cicli anche in 30 passaggi su 30. Il mezzo in pista è avvisato a non meno di $140\unita{m}$ dai frammenti veri in 60 passaggi su 60, e con le regole attuali entra davvero nell'impronta in 30 su 30; ma l'impronta di circa $3{,}5\unita{km}$ contiene già il mezzo, e un avviso che parta dal bordo dell'impronta non c'è in nessun passaggio (0 su 60): l'avviso è giusto e non distingue chi è in pericolo: tocca una mediana di 96 tracce.
\item \emph{Nebbia e pioggia.} Nessuna cella «libera» oltre la portata utile, su $4{,}2\cdot10^8$ celle\,$\cdot$\,ciclo, e 0 tracce perse senza avviso su 1080 cicli, con l'inviluppo precedente \prova{S}. Il conto della portata e il sistema usano la stessa formula: la prova non è indipendente.
\item \emph{Fumo e detriti in pista; passaggio ravvicinato.} Nessun falso evento: 0 su 120 e 0 su 60 passaggi col fumo; a $5\unita{m}$ di distanza senza collisione, 0 su 300 $[0;\,0{,}013]$ con l'inviluppo precedente degli aerei e 0 su 30 con quello di $210\unita{m/s}$, anche con le regole attuali (0 su 30 $[0;\,0{,}114]$) \prova{S}.
\end{itemize}
La perdita di un sensore è dichiarata in 460 casi simulati su 460 $[0{,}992;\,1]$, con un ritardo fino a $5{,}3\unita{s}$ (5,3 rivisite del radar del movimento a terra, contro le 3 della soglia di progetto), con l'inviluppo precedente degli aerei \prova{S}; la soglia di 3 rivisite non è quindi rispettata per quel sensore.

\subsection{Prove simulate del sito}\label{sec:lp-prove}

Il sito sorvegliato è descritto da tre scene, ciascuna simulata per passaggio con semi diversi, e con la regola di spegnimento prudente (tabella~\ref{tab:lp-prove}).
\begin{itemize}
\item \emph{Drone dietro un edificio.} Un multirotore a 15--$21\unita{m/s}$ parte a $990\unita{m}$, scende a $12\unita{m}$ di quota e passa per $8\unita{s}$ dietro un edificio di $80\times40\times15\unita{m}$ a $300\unita{m}$ dal sensore. Due configurazioni: radar per droni, radar a pannello, camera nel visibile e identificazione remota; oppure il solo radar per droni.
\item \emph{Aereo in rullaggio dietro un hangar.} Un aereo a $9\unita{m/s}$ sulla via dietro un hangar di $60\times60\times18\unita{m}$, per $17{,}1\unita{s}$; un mezzo a $8{,}3\unita{m/s}$ attraversa la via. Radar del movimento a terra, LiDAR a campo largo e ADS-B, oppure il solo radar.
\item \emph{Drone sotto la copertura del radar d'aeroporto.} Un multirotore scende da $120$ a $30\unita{m}$ sull'asse di una pista; sotto circa $43\unita{m}$ a $2{,}5\unita{km}$ una fascia di alberi ed edifici lo nasconde al radar d'aeroporto, mentre un aereo è in finale.
\end{itemize}

\begin{table}[H]\centering\footnotesize
\caption{Prove simulate del sito \prova{S}. Conteggi per passaggio; intervalli di Wilson al $95\,\%$. «Fuori avviso»: passaggi con almeno un istante in cui il mezzo nascosto è fuori dall'insieme di possibile presenza. Le probabilità di rilevamento sono ipotesi.}\label{tab:lp-prove}
\begin{tabularx}{\linewidth}{@{}>{\RaggedRight\arraybackslash}p{3.0cm}>{\RaggedRight\arraybackslash}p{2.7cm}L@{}}
\toprule
Scena & Configurazione & Esito\\
\midrule
drone dietro un edificio & quattro sensori & avviso tempestivo (servono $806\unita{m}$, traccia confermata a circa $950\unita{m}$) in 200 passaggi su 200 $[0{,}981;\,1]$, e in 100 su 100 $[0{,}963;\,1]$ con semi nuovi e la versione precedente delle regole; fuori avviso 0 su 200, ma solo 22 passaggi su 200 hanno un intervallo nascosto\\
 & solo radar per droni & tempestivo in 200 su 200 $[0{,}981;\,1]$; fuori avviso in 12 su 200 $[0{,}035;\,0{,}102]$ (6 in cui il drone rallenta, 6 in cui si ferma) e in 9 su 100 $[0{,}048;\,0{,}162]$ con semi nuovi e la versione precedente delle regole\\
aereo in rullaggio dietro un hangar & tre sensori & tempestivo in 200 su 200 $[0{,}981;\,1]$; 0 su 200 fuori avviso, 0 senza regione, 0 con «libero» sull'aereo; 34 passaggi su 200 hanno un intervallo nascosto\\
 & solo radar del movimento a terra & tempestivo in 50 su 200 $[0{,}195;\,0{,}314]$, con l'inviluppo precedente: un sensore solo non basta\\
drone sotto la copertura del radar d'aeroporto & radar d'aeroporto e ADS-B & il radar non conferma mai il drone (0 su 200); il drone non è mai dichiarato «libero» (0 su 200 $[0;\,0{,}019]$), e l'uscita dice «zona non verificata per i multirotori». La tempestività non è misurata: lo scenario parte più vicino della distanza richiesta, e non si dà alcuna cifra\\
\bottomrule
\end{tabularx}
\end{table}

Tre precisazioni. L'avviso perde il drone che rallenta o si ferma in circa il $6\,\%$ dei passaggi con il solo radar: la causa non è stata accertata. Il «libero» sull'oggetto con il radar a pannello regge solo con la mappa del sito dichiarata completa (\S\ref{sec:lp-smentiti}): la tempestività di questa tabella non dice che il «libero» sia corretto in ogni caso. Con le regole attuali è rimisurato solo il drone dietro un edificio (200 passaggi per configurazione); il rullaggio e il drone sotto la copertura hanno le cifre della versione precedente.

\paragraph{Altri mezzi, come confronto.} Con i sensori a bordo di un drone (radar 4D, LiDAR a campo largo, camera) l'avviso è tempestivo in 95 passaggi su 100 $[0{,}888;\,0{,}978]$, ma il mezzo esce dall'avviso in 9 su 100, perché il filtro Doppler perde il drone lento. Con l'inviluppo precedente, un'auto a un incrocio dietro un camion è avvisata in tempo in 200 passaggi su 200 $[0{,}981;\,1]$, e un treno in stazione in 186 su 200 $[0{,}886;\,0{,}958]$: il treno che si ferma o torna indietro prima della conferma non è annunciato in tempo \prova{S}.

\paragraph{La scala, riassunta.} In simulazione dalle schede la portata utile passa da $3\unita{m}$ a circa $1\unita{km}$, e la velocità per cui l'avviso è tempestivo da $1{,}44$ a circa $21\unita{m/s}$, per droni, mezzi di piazzale e aerei in rullaggio \prova{S}. Non è una misura su sensori.

\subsection{Aerei in volo: un limite dichiarato}\label{sec:lp-volo}

Con l'inviluppo di $210\unita{m/s}$, l'operatore di sito ha bisogno di $3410\unita{m}$ di preavviso, mentre il «libero» per un multirotore, con passate correlate, arriva a $877\unita{m}$. L'avviso per gli aerei in volo non è tempestivo, e il sistema lo calcola e lo scrive in ogni uscita: «avviso non tempestivo» \prova{C}\prova{S} (caso costruito). Per un treno ad alta velocità il macchinista avrebbe bisogno di $7{,}57\unita{km}$: anche lì l'avviso non è tempestivo, per calcolo \prova{C}. Con un treno sul binario il sistema può avvisare, non garantire l'arresto. Servirebbe un sensore a portata maggiore; nessuna cifra di tempestività è data per gli aerei in volo.

\subsection{Risultati smentiti, come limiti}\label{sec:lp-smentiti}

Le simulazioni hanno trovato casi in cui il sistema non fa ciò che dichiara. Li riportiamo come limiti.
\begin{enumerate}[label=\arabic*.]
\item \textbf{Il «libero» con il radar a pannello dipende dalla mappa.} Nella configurazione del sito, con la regola precedente, una sola passata mancata apriva un «libero» sulla cella del drone in 163 passaggi su 171 $[0{,}910;\,0{,}976]$: il ritorno cadeva nella cella vicina, e legare la finestra alla velocità massima, da solo, non lo toglieva. Con la ritenuta del ritorno per $5\unita{s}$ e con la mappa del sito dichiarata completa, che è un'ipotesi del simulatore, nessun passaggio su 171 ha istanti con «libero» sul drone $[0;\,0{,}022]$ \prova{S}; la ripetizione con semi nuovi non è rifatta. In una configurazione d'aeroporto il «libero» sull'oggetto non compare (0 su 171 $[0;\,0{,}022]$), ma con 8 falsi confermati di mediana per passaggio, e senza avviso tempestivo per gli aerei.
\item \textbf{Mappa dichiarata completa ma sbagliata.} Un edificio che non compare in una mappa dichiarata completa apre un «libero» falso dietro di sé, in 5 casi costruiti su 6 \prova{S}; per questo il prototipo, per impostazione, non dichiara la mappa completa, e lì la cella dietro l'edificio è «non osservata» (6 su 6).
\item \textbf{L'avviso perde il drone che rallenta o si ferma} in 12 passaggi su 200 con il solo radar (\S\ref{sec:lp-prove}); con il filtro Doppler a bordo, in 9 su 100, con la versione precedente delle regole.
\item \textbf{Classi non dichiarate.} Nel sito sorvegliato l'elicottero, l'aereo leggero, l'auto e il pedone non sono né ammessi né esclusi: per loro il «libero» non è definito.
\item \textbf{Eventi improvvisi.} La collisione non è dedotta entro $2\unita{s}$ in ogni passaggio: 51 su 60 con il primo simulatore dei frammenti (integratore di Eulero), 48 su 60 $[0{,}682;\,0{,}882]$ con l'integratore stabile, più 3 eventi aperti prima dell'urto; la zona dei detriti è di alcuni chilometri e non si riduce guardando; l'avviso non distingue chi è in pericolo; la fiammata è rilevata solo con un innesco ottimista (\S\ref{sec:lp-eventi}).
\item \textbf{Inviluppo.} Con $90\unita{m/s}$ l'avviso perdeva aerei reali (8--11 su 43); il valore di $210\unita{m/s}$ non è stato provato di nuovo sulle traiettorie. Il drone veloce oltre i $45\unita{m/s}$ ipotizzati esce dall'avviso.
\item \textbf{Falsi avvisi.} Nelle prove con l'inviluppo precedente si va da $3{,}8$ falsi avvisi all'ora con il solo radar per droni a 1594 all'ora con radar 4D e LiDAR \prova{S}; i falsi ritorni sono un'ipotesi del simulatore.
\end{enumerate}

\subsection{Condizioni di validità a lunga portata}\label{sec:lp-validita}

\begin{riquadro}{Condizioni di validità degli enunciati e delle cifre a lunga portata}
\begin{enumerate}[label=\arabic*.,itemsep=1pt]
\item \textbf{Sensori simulati dalle schede,} con le ipotesi della tabella~\ref{tab:lp-sensori}: probabilità di rilevamento, falsi allarmi, latenze e la sezione radar del multirotore ($0{,}01\unita{m^2}$) non hanno una fonte.
\item \textbf{Sensori fermi, posa nota;} il caso con i sensori a bordo di un drone è un'analisi con posa supposta nota.
\item \textbf{Classi dichiarate dallo scenario} e inviluppo dato: multirotori fino a $45\unita{m/s}$ e $25\unita{m/s^2}$ (ipotesi), aerei in volo fino a $210\unita{m/s}$.
\item \textbf{Mappa degli ostacoli corretta,} oppure sfondo misurato; senza l'una o l'altro non c'è «libero». Le cifre del «libero» sul drone nel sito valgono con la mappa dichiarata completa, che è un'ipotesi.
\item \textbf{Geometrie nominali,} con ostacoli estrusi e oggetti puntiformi; niente riflessioni multiple, uccelli o disturbo strutturato.
\item \textbf{Etichette di classe a mano;} il classificatore è richiesto e supposto esistente, ma non è realizzato: è lasciato ai lavori futuri.
\end{enumerate}
\end{riquadro}

\section{Volo ad alta velocità: collisioni, detriti ed eventi impulsivi}\label{sec:vo}
\sloppy

Il capitolo~\ref{sec:lp} ha portato i principi del sistema da $3\unita{m}$ a qualche chilometro, su un sito sorvegliato con droni e aerei in rullaggio. Questo capitolo tratta il caso più difficile: aerei e droni \emph{in volo}, da $20$ a $210\unita{m/s}$, con eventi improvvisi e fortemente impulsivi (una collisione fra droni veloci, l'avaria di un motore con fiammata e detriti, una cella di pioggia, l'accecamento di tutti i sensori ottici). Le domande sono tre: se le regole di sicurezza restano vere a queste velocità; quanto in fretta si perde la confidenza sulla posizione di un oggetto che non si vede più; che cosa possono dare sensori e calcolo allo stato dell'arte.

\begin{riquadro}{Natura delle prove di questo capitolo}
Tutto ciò che segue è \emph{simulazione dalle schede tecniche}, \emph{calcolo} o \emph{traiettorie pubbliche con un sensore simulato}. Nessuna prova usa un sensore fisico o un'unità di calcolo dedicata; le latenze dell'elaborazione dedicata sono stime per analogia, non misure. Le etichette di classe sono assegnate a mano: il classificatore è un componente richiesto e supposto esistente (\S\ref{sec:sti}). Il sistema osserva e avvisa: nessuna uscita è un comando, e non vi sono funzioni d'ingaggio, di puntamento o di contrasto. Le cifre sono per \emph{passaggio} (una ripetizione completa di uno scenario con semi diversi), con l'intervallo di Wilson al $95\,\%$ \cite{wilson1927}; dove si contano \emph{campioni} (cicli dentro un passaggio) i campioni sono correlati e l'intervallo è ottimista, e lo si scrive. Livelli di prova: \prova{D} dimostrato nel modello dichiarato, \prova{C} calcolato, \prova{S} simulato, \prova{P} su dati pubblici con sensore simulato, \prova{N} non verificato; «ipotesi» è scritto per esteso.
\end{riquadro}

\subsection{Gli scenari}\label{sec:vo-scenari}

Sei scenari nominali (tabella~\ref{tab:vo-scenari}) coprono l'avvicinamento, la salita in meteo avverso, la collisione, l'avaria, la cecità degli ottici e un controllo senza eventi. Ogni scenario è \emph{una} geometria: i passaggi variano soltanto i semi del rumore, dei mancati rilevamenti e dei frammenti. Le procedure di volo usate per fissare velocità e quote (sentiero di $3^\circ$, $250\unita{kt}$ sotto $10\,000\unita{ft}$) sono riprese da conoscenza generale e non sono state riaperte \prova{N}.

\begin{table}[H]\centering\footnotesize
\caption{Scenari di volo. Geometria: (C) calcolata, (i) ipotesi.}\label{tab:vo-scenari}
\begin{tabularx}{\linewidth}{@{}p{0.6cm}>{\RaggedRight\arraybackslash}p{6.4cm}>{\RaggedRight\arraybackslash}p{3.2cm}L@{}}
\toprule
& Geometria & Eventi & Che cosa mette alla prova\\
\midrule
F1 & aereo sul sentiero di $3^\circ$ a $75\unita{m/s}$; drone non cooperante a $20\unita{m/s}$ che lo attraversa; distanza minima $40{,}8\unita{m}$ a $t=80{,}0\unita{s}$ (C) & nessuno: quasi collisione & avviso di prossimità con anticipo; nessun falso evento\\
F2 & aereo in salita da $90$ a $125\unita{m/s}$ dentro una cella di pioggia in quota (raggio $1{,}5\unita{km}$, da $0$ a $1600\unita{m}$); aereo leggero non cooperante & nube in quota; portata ridotta (radar $\times0{,}70$, visibile $\times0{,}26$, termica $\times0{,}29$, fattori prudenti) & portata dichiarata, «ridondanza assente», regione che cresce fra due giri\\
F3 & due droni a $40$ e $35\unita{m/s}$ si urtano a $120\unita{m}$ di quota a $t=20\unita{s}$; uno cade, l'altro prosegue danneggiato a $25\unita{m/s}$ & collisione con frammenti (sezione radar $2\cdot10^{-4}\unita{m^2}$, i) & evento entro $2\unita{s}$; volume dei detriti; nessun «libero» dentro\\
F4 & aereo a $135\unita{m/s}$ e $450\unita{m}$; avaria a $t=30\unita{s}$, poi virata e discesa & fiammata (visibile cieco $5\unita{s}$, termica $13\unita{s}$), fumo, detriti a \emph{bassa energia} ($\Delta v\le100\unita{m/s}$, i) & affidabilità persa entro 3 cicli; evento entro $1\unita{s}$; impronta a terra\\
F5 & aereo a $80\unita{m/s}$ e drone a $45\unita{m/s}$ che si incrociano a $300\unita{m}$; distanza minima $39{,}2\unita{m}$ (C) & perdita totale degli ottici per 2, 3 o $4\unita{s}$, più una nube & solo radar; continuità e identità\\
F6 & controllo: due aerei a $120$ e $100\unita{m/s}$, separati di $300\unita{m}$ in quota e $1{,}5\unita{km}$ in pianta & nessuno & falsi eventi\\
\bottomrule
\end{tabularx}
\end{table}

\paragraph{Configurazioni.} Ogni scenario è simulato in due configurazioni. Nella configurazione~A il radar è il radar per droni del capitolo~\ref{sec:lp} (rotante, $1\unita{Hz}$), con camere nel visibile e termiche, e l'elaborazione è quella del prototipo su CPU ($80\unita{ms}$ aggiunti alla latenza di scheda). Nella configurazione~B il radar è a scansione elettronica a $10\unita{Hz}$ in traccia, di una classe commerciale: le tracce a $10\unita{Hz}$ sono riportate dalla stampa specializzata \citep{dronedj2026-echodyne}, la banda Ku e il campo di $120^\circ\times80^\circ$ dal produttore \citep{echodyne-tradegov-2024,echodyne-echoshield-news}; la scheda con le portate per bersaglio non è stata aperta. Quattro pannelli e la portata di riferimento di $1400\unita{m}$ per $0{,}01\unita{m^2}$ sono ipotesi; le camere vanno a $10\unita{Hz}$ e l'elaborazione dedicata aggiunge $1{,}2\unita{ms}$ (stima, \S\ref{sec:vo-hw}). Nelle due configurazioni cambiano \emph{insieme} il sensore e il calcolo: una differenza fra A e B non si può attribuire al calcolo. L'ADS-B è spento in tutte le prove (ipotesi prudente: dopo un evento l'oggetto colpito smette di trasmettere o trasmette un dato vecchio). In ciascuna cella si simulano 60 passaggi (F6: 200; F5: 20 per ciascuna durata di cecità).

\paragraph{Frammenti.} I frammenti partono con la velocità del mezzo più un incremento $u$ con $|u|\le\Delta v$, e cadono con resistenza aerodinamica e vento. Una prima versione del simulatore integrava la resistenza con il metodo di Eulero esplicito, che è instabile quando $g|v|\,\Delta t/v_t^2>2$ ($v_t$ velocità limite): con $|v|\approx195\unita{m/s}$, $v_t\approx6\unita{m/s}$ e $\Delta t=0{,}05\unita{s}$ il prodotto vale circa $2{,}6$, e alcuni frammenti finivano a centinaia di chilometri già dopo $0{,}5\unita{s}$. Le prove di questo capitolo usano un'integrazione semi-implicita della resistenza, stabile per ogni passo \prova{S}. Lo stesso difetto era presente nel simulatore degli eventi a bassa velocità del capitolo~\ref{sec:lp}, dove spostava il frammento più lontano della collisione fra multirotori da $130$ a $154\unita{m}$.

\paragraph{La collisione fra multirotori rifatta con l'integratore stabile.} La collisione fra droni lenti ($8$ e $6{,}7\unita{m/s}$) del capitolo~\ref{sec:lp} è stata ripetuta su 60 passaggi con l'integratore semi-implicito \prova{S}:
\begin{itemize}
\item i frammenti restano tutti nel volume dichiarato in $60/60$ passaggi $[0{,}940;\,1]$, contro $48/60$ con l'integratore di Eulero;
\item l'evento è dichiarato entro $2\unita{s}$ dall'urto in $48/60$ passaggi $[0{,}682;\,0{,}882]$; contando solo i ritardi non negativi, entro $2{,}5\unita{s}$ sono $57/60$ $[0{,}863;\,0{,}983]$ (soglia notata dopo la prova);
\item in $3/60$ passaggi $[0{,}017;\,0{,}137]$ un evento di tipo «separazione» è aperto da $11{,}0$ a $11{,}5\unita{s}$ \emph{prima} dell'urto: va contato come falso evento. La causa (tracce di camera al primo aggancio) non è verificata.
\end{itemize}
Il conteggio di 51 su 60 riportato nel capitolo~\ref{sec:lp} vale quindi solo con l'integratore di Eulero. La ripetizione cambia anche l'estrazione del rumore delle camere, e per questo non isola l'effetto dell'integratore. Per il volume dei detriti il risultato è solido; per l'evento è un limite dichiarato.

\subsection{Il nucleo: l'inviluppo a 210 m/s anche nell'insieme di presenza}\label{sec:vo-nucleo}

Il nucleo è quello del capitolo~\ref{sec:lp}: «libero» per classe con passate correlate, inseguimento a più modelli, insieme di possibile presenza dall'inviluppo dello scenario, spegnimento dell'avviso solo guardando tutta la zona, eventi dedotti da separazione o contatto, volume dei detriti. Il cielo non ha sfondo: il «libero» vale solo con la mappa dichiarata (\S\ref{sec:lp-libero}).

\paragraph{Una miglioria trovata dalle prove.} Nelle prime prove di volo l'insieme di possibile presenza del nucleo era calcolato con la velocità massima \emph{delle classi} ($90\unita{m/s}$), mentre l'inviluppo di $210\unita{m/s}$ per gli aerei in volo, deciso dopo le traiettorie pubbliche (\S\ref{sec:lp-traiettorie}), era applicato solo allo strato degli eventi. Con aerei a $120$--$135\unita{m/s}$ l'oggetto vero usciva dalla regione dichiarata, in violazione dell'ipotesi della proposizione~\ref{prop:lp-contenimento} ($|v|\le\vmax$). La correzione è un cambiamento di un solo parametro, ma il suo effetto è grande \prova{S}:
\begin{itemize}
\item avaria (F4), configurazione A: passaggi senza alcun ciclo privo di regione da $13/60$ $[0{,}13;\,0{,}34]$ a $60/60$ $[0{,}940;\,1]$; quota dei cicli di perdita con il vero dentro la regione da $0{,}144$ a $1{,}000$ ($1072/1072$ campioni dopo l'avaria);
\item controllo (F6), configurazione A: da $126/200$ a $186/200$ $[0{,}886;\,0{,}958]$ passaggi senza buchi; vero dentro la regione da $0{,}272$ a $0{,}984$ $[0{,}966;\,0{,}993]$ (campioni, intervallo ottimista).
\end{itemize}
La versione corretta del nucleo è quella usata in tutto il capitolo; la precedente è conservata per confronto. Il caso illustra un fatto formale: una dimostrazione di contenimento (\S\ref{sec:vo-presenza}) vale solo se il parametro usato nel codice è quello dell'enunciato. La violazione non era nella prova ma nel collegamento fra la prova e il programma.

\subsection{Il volume dei detriti in volo}\label{sec:vo-detriti}

Il volume $V(\tau)$ è quello della proposizione~\ref{prop:lp-detriti}: l'unione dei segmenti percorsi alle velocità delle tracce d'origine, ingrossata di $W\tau+(\Delta v+W)\tau^*+r_0$, fra due estremi di quota. In volo la velocità d'origine $v_c$ è grande (fino a $210\unita{m/s}$), e la pianta si allunga lungo la rotta. La proposizione è \prova{D} nel modello del punto materiale con resistenza isotropa, vento costante entro $W$ e $W_z$, incremento iniziale entro $\Delta v$; vale per $\tau\le\tau^*$.

\paragraph{Prove.} Con l'integratore stabile \prova{S}:
\begin{itemize}
\item F3: $600/600$ frammenti dentro $V(\tau)$ per tutta la caduta $[0{,}9936;\,1]$; anche con $|u|=120\unita{m/s}$ sul bordo dei parametri e vento a $15\unita{m/s}$, e con una velocità d'origine dichiarata sbagliata di $8\unita{m/s}$ ($600/600$, 10 passaggi per variante);
\item F4 nel nucleo completo: passaggi con tutti i frammenti dentro $V$: $58/60$ $[0{,}886;\,0{,}991]$ in ciascuna configurazione; nei due passaggi peggiori i frammenti dentro sono il $96\,\%$ (A) e il $74\,\%$ (B); al bordo dei parametri (aereo a $210\unita{m/s}$, $\Delta v=100\unita{m/s}$, vento $15\unita{m/s}$): $1000/1000$ $[0{,}9962;\,1]$, con un'impronta di $8{,}65\unita{km}$;
\item impronta a terra: mediana $6{,}9\unita{km}$ (A) e $6{,}8\unita{km}$ (B) in F4; $5{,}1\unita{km}$ in F3 \prova{C}.
\end{itemize}
Il simulatore dei frammenti condivide il modello fisico con la proposizione: la prova vede gli errori di conto e di codice, non quelli del modello. L'impronta è prudente e grande: per renderla utile serve misurare la velocità dei frammenti (Doppler) invece di prenderla a priori (\S\ref{sec:ps-tabella}).

\paragraph{Che cosa il volume non contiene.} Un'avaria \emph{non contenuta} libera pezzi di disco a centinaia di metri al secondo, fuori dal $\Delta v$ dichiarato. Con il $5\,\%$ dei frammenti come pezzi di disco a $250\unita{m/s}$ e velocità limite di $80\unita{m/s}$ (ipotesi), $965/1000$ frammenti restano dentro $V$, ma i pezzi di disco solo $15/50$, e i passaggi con tutti i frammenti dentro sono $0/10$ $[0;\,0{,}278]$ \prova{S}. A terra nessuno cade fuori dall'impronta di $6{,}8\unita{km}$; in aria il volume non li contiene. Per questo F4 è un'\emph{avaria con detriti a bassa energia}, e l'avaria non contenuta è un caso non provato.

\subsection{Risultati}\label{sec:vo-risultati}

La tabella~\ref{tab:vo-risultati} raccoglie i risultati delle prove simulate, con quattro classi:
\begin{itemize}
\item \textbf{solido}: limite inferiore di Wilson $\ge0{,}90$, oppure 0 su $N$ con limite superiore $\le0{,}06$;
\item \textbf{incoraggiante}: stima $\ge0{,}80$, limite inferiore $\ge0{,}70$, e una base matematica dichiarata;
\item \textbf{limite}: il sistema non fa ciò che dovrebbe, e lo si scrive;
\item \textbf{non provato}: nessuna prova sostiene l'enunciato.
\end{itemize}
Le classi si riferiscono alla simulazione dalle schede, con etichette a mano e nessun sensore fisico.

\begin{table}[H]\centering\footnotesize
\caption{Volo ad alta velocità: risultati simulati per passaggio, con l'intervallo di Wilson al $95\,\%$ \prova{S}. A: radar rotante a $1\unita{Hz}$ ed elaborazione su CPU; B: radar a scansione elettronica a $10\unita{Hz}$ ed elaborazione dedicata stimata. Nucleo adottato: quello con l'inviluppo a $210\unita{m/s}$ (\S\ref{sec:vo-nucleo}), senza la regola di contatto in prova. Simulazione dalle schede, etichette a mano, nessun sensore fisico.}\label{tab:vo-risultati}
\begin{tabularx}{\linewidth}{@{}>{\RaggedRight\arraybackslash}p{5.0cm}>{\RaggedRight\arraybackslash}p{2.1cm}>{\RaggedRight\arraybackslash}p{4.4cm}L@{}}
\toprule
Enunciato & Classe & Conteggio e intervallo & Base\\
\midrule
Nei sei scenari, dall'avaria alla collisione di droni, nessuna cella è dichiarata «libera» sul punto vero dell'oggetto o su un frammento (con la mappa dichiarata) & solido & F6: $0/200$ $[0;\,0{,}019]$; le altre celle: $0/60$ $[0;\,0{,}060]$; ancora 0 in tutte le 10 celle con la regola di contatto in prova & «libero» solo ciò che si osserva; insieme di presenza a $210\unita{m/s}$\\
Avaria di un aereo a $135\unita{m/s}$ con fiammata: evento dedotto entro $1\unita{s}$, camere dichiarate non affidabili & incoraggiante (innesco ottimista) & $60/60$ $[0{,}940;\,1]$ in A e in B & separazione di tracce e salute del sensore; la saturazione è presa dalla finestra vera dell'evento\\
Dopo l'avaria i detriti a bassa energia restano nel volume dichiarato & incoraggiante & $58/60$ $[0{,}886;\,0{,}991]$; nei peggiori $96\,\%$ (A) e $74\,\%$ (B) dei frammenti & proposizione~\ref{prop:lp-detriti}, $\Delta v\le100\unita{m/s}$\\
Avviso prima del primo frammento a terra; margine minimo per un operatore $0{,}25$--$0{,}45\unita{s}$ & incoraggiante & $60/60$ $[0{,}940;\,1]$; risposta umana di $5{,}4\unita{s}$ (ipotesi; \cite{faa-ac-90-48e}) & $d=v\,(L+t_{\mathrm{risp}}+t_{\mathrm{mezzo}})$\\
Due droni veloci in rotta di collisione: avviso di prossimità almeno $10\unita{s}$ prima dell'urto & incoraggiante (soglia di $10\unita{s}$ non scritta prima) & $60/60$ $[0{,}940;\,1]$ & massimo avvicinamento a velocità costante\\
Dopo la collisione veloce, punto vero e frammenti dentro la regione dichiarata & solido come copertura & $60/60$ passaggi $[0{,}940;\,1]$; $939/939$ campioni & proposizione~\ref{prop:vo-presenza}, $|v|\le210\unita{m/s}$\\
Quasi collisioni a $40\unita{m}$ e incroci veloci: nessun falso evento & solido (nucleo adottato) & F6: $0/200$ $[0;\,0{,}019]$; F1, F5: $0/60$ $[0;\,0{,}060]$ & conferma su più cicli\\
Con il radar a $10\unita{Hz}$ (B) la regione non si perde in nessun ciclo & solido in B & $60/60$ e $200/200$ $[0{,}981;\,1]$; i $10\unita{Hz}$ vengono da una fonte di stampa & rivisita più breve della crescita dell'inviluppo\\
Con il radar rotante a $1\unita{Hz}$ (A) l'aereo veloce perde continuità, e il sistema lo dichiara & limite & F1: $39/60$ $[0{,}524;\,0{,}758]$; F6: $186/200$ $[0{,}886;\,0{,}958]$ & --\\
Ottici ciechi per $2$--$4\unita{s}$: ogni traccia conserva una regione; l'identità resta «presunta» & solido (regione); limite (identità) & $60/60$ $[0{,}940;\,1]$; scambio $28/60=0{,}47$ $[0{,}35;\,0{,}59]$ (A), $38/60=0{,}63$ $[0{,}51;\,0{,}74]$ (B) & --\\
Collisione \emph{veloce} dichiarata come evento & limite & nucleo adottato: $0/60$ $[0;\,0{,}060]$ in A e in B. Regola di contatto in prova, non adottata: B $57/60$ $[0{,}863;\,0{,}983]$, A $34/60$ $[0{,}441;\,0{,}684]$; falsi eventi F5 $2/60$ $[0{,}009;\,0{,}114]$; prima dell'urto $1/60$; varianti di geometria 16--19 su 20 in B & conferma su tracce risolte e Doppler simulato ($\sigma_{v_r}=0{,}5\unita{m/s}$, ipotesi)\\
Avaria non contenuta (pezzi di disco a $250\unita{m/s}$) & non provato & $0/10$ $[0;\,0{,}278]$ passaggi con tutti i frammenti dentro & --\\
Collisione fra droni \emph{lenti} ($8$ e $6{,}7\unita{m/s}$), integratore stabile: evento entro $2\unita{s}$; frammenti nel volume & limite (evento); solido (volume) & evento $48/60$ $[0{,}682;\,0{,}882]$, con $3/60$ $[0{,}017;\,0{,}137]$ aperti prima dell'urto; volume $60/60$ $[0{,}940;\,1]$ & fusione e separazione di tracce\\
\bottomrule
\end{tabularx}
\end{table}

\paragraph{Perché la collisione veloce non è dichiarata.} In F3 i ritorni dei frammenti, con una sezione radar di $2\cdot10^{-4}\unita{m^2}$ (ipotesi) a circa $2{,}1\unita{km}$, non superano la soglia di rilevamento; l'innesco di «separazione» non scatta. L'innesco di «contatto» richiede un indizio fra ritorni nuovi, residui fuori porta o una traccia non più rilevata in vista, e nelle due configurazioni non si combina entro $2\unita{s}$. Il sistema mostra invece l'avviso di prossimità prima dell'urto, la perdita delle tracce, le regioni che crescono e le tracce nuove non associate: la collisione non è riconosciuta come evento, ma la regione è mantenuta. Il risultato della collisione fra multirotori lenti ($48/60$ con l'integratore stabile) non si trasferisce al volo veloce.

\paragraph{La regola di contatto provata, e non adottata.} Per la collisione veloce è stata provata una regola di contatto: due tracce \emph{risolte e mature} (non doppie, \S\ref{sec:vo-impulso}) che convergono a più di $10\unita{m/s}$, seguite da una fusione o dalla sparizione di una delle due, e da tracce nuove o da un Doppler fuori norma (Doppler simulato con $\sigma_{v_r}=0{,}5\unita{m/s}$, ipotesi). La soglia fissata prima chiedeva F3 entro $2\unita{s}$ in almeno $0{,}85$ dei passaggi, senza falsi eventi nel controllo. Gli esiti, su semi nuovi \prova{S}:
\begin{itemize}
\item in B la collisione è dichiarata entro $2\unita{s}$ (ritardo fra $0$ e $2\unita{s}$) in $57/60$ passaggi $[0{,}863;\,0{,}983]$; in A in $34/60$ $[0{,}441;\,0{,}684]$, sotto la soglia;
\item nel controllo F6 nessun falso evento, $0/200$ $[0;\,0{,}019]$ in A e in B; ma nel quasi incrocio di F5 a $39\unita{m}$ la regola apre falsi eventi in $2/60$ $[0{,}009;\,0{,}114]$, in A e in B, e in F1 (A) in $1/60$ $[0{,}003;\,0{,}089]$; in F3 un evento è aperto prima dell'urto in $1/60$;
\item l'avaria (F4) non peggiora: $60/60$ $[0{,}940;\,1]$; «libero» sul punto vero resta 0 in tutte le celle.
\end{itemize}
La tabella~\ref{tab:vo-contatto} dà la regola fuori dalla geometria nominale, su 20 passaggi per variante. Con 20 passaggi il limite inferiore di Wilson non arriva a $0{,}85$ nemmeno con 20 su 20 ($0{,}839$): le varianti servono a cercare il caso che rompe la regola, non a sostenerla.

\begin{table}[H]\centering\footnotesize
\caption{Regola di contatto in prova (non adottata) sulla collisione veloce fuori dalla geometria nominale \prova{S}. «Prima»: evento aperto prima dell'urto, da contare come falso evento. Le varianti di distanza spostano tutti i sensori.}\label{tab:vo-contatto}
\begin{tabularx}{\linewidth}{@{}>{\RaggedRight\arraybackslash}p{3.6cm}LLLL@{}}
\toprule
Variante & A, entro $2\unita{s}$ & A, prima & B, entro $2\unita{s}$ & B, prima\\
\midrule
nominale (60 passaggi) & $34/60$ $[0{,}441;\,0{,}684]$ & $1/60$ & $57/60$ $[0{,}863;\,0{,}983]$ & $1/60$\\
incrocio a $30^\circ$ & $10/20$ $[0{,}299;\,0{,}701]$ & $4/20$ $[0{,}081;\,0{,}416]$ & $19/20$ $[0{,}764;\,0{,}991]$ & $1/20$\\
incrocio a $150^\circ$ & $16/20$ $[0{,}584;\,0{,}919]$ & $0/20$ & $17/20$ $[0{,}640;\,0{,}948]$ & $2/20$ $[0{,}028;\,0{,}301]$\\
quota $40\unita{m}$ & $14/20$ $[0{,}481;\,0{,}855]$ & $0/20$ & $19/20$ $[0{,}764;\,0{,}991]$ & $0/20$\\
quota $400\unita{m}$ & $9/20$ $[0{,}258;\,0{,}658]$ & $0/20$ & $16/20$ $[0{,}584;\,0{,}919]$ & $0/20$\\
sensori a circa $570\unita{m}$ dall'urto & $17/20$ $[0{,}640;\,0{,}948]$ & $0/20$ & $7/7$, prova parziale & $0/7$\\
\bottomrule
\end{tabularx}
\end{table}
In ogni variante il «libero» sul punto vero è $0/20$ $[0;\,0{,}161]$. Tre varianti del controllo con l'incrocio portato più vicino danno 0 falsi eventi su 20 in A e in B, ma nessuna scende sotto i $300\unita{m}$ e quindi non tocca il punto debole, il quasi incrocio a $39\unita{m}$ di F5. La regola non è adottata: la collisione veloce resta un limite dichiarato, e la regola si presenta solo come miglioria in prova (\S\ref{sec:ps-tabella}). In breve: \emph{la collisione veloce non è ancora dichiarata come evento; una regola di contatto in prova la riconosce entro $2\unita{s}$ in 57 passaggi simulati su 60 con il radar a scansione elettronica, ma dà ancora falsi eventi in un quasi incrocio a $39\unita{m}$ (2 su 60), e per questo non è adottata.}

\paragraph{Lettura per il caso centrale.} Ciò che regge nelle simulazioni di volo è il principio di sicurezza (nessun «libero» sull'oggetto o sui frammenti), il contenimento della regione con l'inviluppo corretto, la gestione della perdita dei sensori, l'avaria con detriti a bassa energia e l'avviso di prossimità. Ciò che non regge ancora è il riconoscimento della collisione veloce come evento, l'identità dopo la cecità e la continuità con un radar rotante. Sono risultati ottenuti su una geometria per scenario: non coprono la varietà dei casi reali.

\subsection{Un controesempio istruttivo: la regola «impulso»}\label{sec:vo-impulso}

Per riconoscere la collisione veloce è stata provata una regola più semplice: si dichiara un evento quando due tracce vicine mostrano un cambio brusco di velocità relativa, senza chiedere i ritorni dei frammenti.

\begin{controesempio}[Una regola sensibile non è una regola corretta]\label{ce:vo-impulso}
Con la regola «impulso» la collisione veloce è dedotta entro $2\unita{s}$ in $59/60$ passaggi $[0{,}911;\,0{,}997]$. Nello scenario di controllo, dove non c'è alcuna collisione, la stessa regola apre un falso evento in $198/200$ passaggi (A) e $158/200$ (B), e in F1 in $57/60$; in F3 l'evento si apre in mediana $10\unita{s}$ \emph{prima} dell'urto \prova{S}.
\end{controesempio}
Il meccanismo è chiaro. Un radar che vede lo stesso aereo da due pannelli, o una traccia vecchia non ancora chiusa, produce due tracce dello stesso oggetto. Quando una delle due è corretta o abbandonata, la velocità relativa della coppia cambia di colpo: il predicato «impulso fra due tracce vicine» è vero senza che vi sia contatto. La regola confonde una proprietà dell'\emph{inseguitore} (le tracce doppie) con una proprietà del \emph{mondo} (l'urto). Il controesempio fissa due criteri per ogni regola d'evento: si giudica sulla coppia (sensibilità, falsi eventi nel controllo), e il predicato deve essere falso per due tracce dello stesso oggetto. La regola proposta al capitolo~\ref{sec:ps} chiede tracce \emph{risolte e non doppie} per almeno 3 aggiornamenti prima del contatto proprio per questo. La regola «impulso» è respinta.

\subsection{Perdita di confidenza nel tempo}\label{sec:vo-confidenza}

Quando un oggetto non è più osservato, il sistema deve dire quanto sa ancora della sua posizione, in funzione del tempo $\tau$ trascorso dall'ultima misura. Si danno due risposte di natura diversa: una \emph{probabilità} nel modello di moto (\S\ref{sec:vo-c-modello}--\ref{sec:vo-c-imm}) e un \emph{contenimento} fuori dal modello (\S\ref{sec:vo-presenza}). La prima è utile ma vale solo se il modello è giusto; il secondo vale sotto ipotesi più deboli e dà una regione molto più grande.

\subsubsection{Varianza di posizione nel moto a velocità quasi costante}\label{sec:vo-c-modello}

\begin{ipotesi}[Moto e stima all'ultimo aggiornamento]\label{ip:vo-cv}
Su ciascun asse lo stato $x=(p,v)$ segue $\dot p=v$, $\dot v=w(t)$, con $w$ rumore bianco gaussiano di densità spettrale $q$ ($\mathrm{m^2/s^3}$), indipendente dallo stato. All'istante dell'ultimo aggiornamento ($\tau=0$) l'errore di stima è gaussiano, a media nulla, con covarianza
\[P_0=\begin{pmatrix}\sigma_0^2 & \sigma_{0v}\\ \sigma_{0v} & \sigma_v^2\end{pmatrix},\]
dove $\sigma_{0v}$ è la \emph{covarianza} fra errore di posizione e di velocità (può essere negativa). Dopo $\tau=0$ non arriva alcuna misura.
\end{ipotesi}

\begin{proposizione}[Varianza di posizione]\label{prop:vo-sigma}
Sotto l'ipotesi~\ref{ip:vo-cv} l'errore di posizione al tempo $\tau\ge0$ è gaussiano a media nulla con varianza
\begin{equation}\label{eq:vo-sigma}
\sigma^2(\tau)=\sigma_0^2+2\sigma_{0v}\,\tau+\sigma_v^2\,\tau^2+\frac{q\,\tau^3}{3}.
\end{equation}
\end{proposizione}
\begin{proof}
Il sistema è lineare: $\dot x=Ax+Bw$ con $A=\bigl(\begin{smallmatrix}0&1\\0&0\end{smallmatrix}\bigr)$, $B=(0,1)^\top$. La matrice di transizione è $\Phi(\tau)=e^{A\tau}=\bigl(\begin{smallmatrix}1&\tau\\0&1\end{smallmatrix}\bigr)$, perché $A^2=0$. La soluzione è $x(\tau)=\Phi(\tau)x(0)+\int_0^\tau\Phi(\tau-s)Bw(s)\,ds$. L'errore di previsione $e(\tau)=x(\tau)-\Phi(\tau)\hat x(0)$ è la somma di $\Phi(\tau)e(0)$ e dell'integrale stocastico, indipendenti per ipotesi; è quindi gaussiano a media nulla, con covarianza
\[P(\tau)=\Phi(\tau)P_0\Phi(\tau)^\top+Q(\tau),\qquad Q(\tau)=q\int_0^\tau\Phi(s)BB^\top\Phi(s)^\top ds=q\begin{pmatrix}\tau^3/3&\tau^2/2\\ \tau^2/2&\tau\end{pmatrix},\]
poiché $\Phi(s)B=(s,1)^\top$ e $\int_0^\tau s^2ds=\tau^3/3$. L'elemento $(1,1)$ di $\Phi P_0\Phi^\top$ è $(1,\tau)P_0(1,\tau)^\top=\sigma_0^2+2\sigma_{0v}\tau+\sigma_v^2\tau^2$; sommando $q\tau^3/3$ si ha la \eqref{eq:vo-sigma}.
\end{proof}
È la forma continua del modello a velocità quasi costante \cite{barshalom2001}. Nel filtro discreto la stessa $Q(\tau)$ è quella della discretizzazione esatta, quindi la \eqref{eq:vo-sigma} vale anche ai passi del filtro.

\begin{lemma}[Convessità]\label{lem:vo-convessa}
$\sigma^2(\tau)$ è convessa su $[0,\infty)$; se $\sigma_{0v}\ge0$ è anche non decrescente.
\end{lemma}
\begin{proof}
$\tfrac{d^2}{d\tau^2}\sigma^2=2\sigma_v^2+2q\tau\ge0$ per $\tau\ge0$. Inoltre $\tfrac{d}{d\tau}\sigma^2=2\sigma_{0v}+2\sigma_v^2\tau+q\tau^2\ge0$ se $\sigma_{0v}\ge0$.
\end{proof}

\subsubsection{Confidenza su un disco e su una sfera}\label{sec:vo-c-chiusa}

\begin{ipotesi}[Isotropia]\label{ip:vo-iso}
In $d$ dimensioni ($d=2$ o $3$) gli assi sono indipendenti e hanno la stessa $\sigma^2(\tau)$: $P_0$ e $q$ uguali per asse, nessuna covarianza fra assi. La regione è il disco ($d=2$) o la sfera ($d=3$) di raggio $r$ centrata sulla previsione $m(\tau)$.
\end{ipotesi}

\begin{proposizione}[Confidenza in forma chiusa]\label{prop:vo-c}
Sotto le ipotesi~\ref{ip:vo-cv} e~\ref{ip:vo-iso}, la probabilità che l'oggetto stia nella regione al tempo $\tau$ è, con $k(\tau)=r/\sigma(\tau)$,
\begin{align}
C_2(\tau)&=1-\exp\!\Bigl(-\frac{r^2}{2\sigma^2(\tau)}\Bigr),\label{eq:vo-c2}\\
C_3(\tau)&=\operatorname{erf}\!\Bigl(\frac{k}{\sqrt2}\Bigr)-\sqrt{\frac{2}{\pi}}\;k\,e^{-k^2/2}.\label{eq:vo-c3}
\end{align}
\end{proposizione}
\begin{proof}
Per la proposizione~\ref{prop:vo-sigma} e l'isotropia, $e(\tau)/\sigma(\tau)$ ha componenti normali standard indipendenti, e $S=|e|^2/\sigma^2$ segue un chi quadro con $d$ gradi di libertà. Allora $C=\Pr(S\le k^2)$. Per $d=2$ la densità di $S$ è $\tfrac12e^{-s/2}$, e $\int_0^{k^2}\tfrac12e^{-s/2}ds=1-e^{-k^2/2}$. Per $d=3$, passando al raggio $\rho=\sqrt S$, la densità è $\sqrt{2/\pi}\,\rho^2e^{-\rho^2/2}$; integrando per parti,
$\int_0^k\sqrt{2/\pi}\,\rho^2e^{-\rho^2/2}d\rho=\sqrt{2/\pi}\bigl[-\rho e^{-\rho^2/2}\bigr]_0^k+\sqrt{2/\pi}\int_0^ke^{-\rho^2/2}d\rho$, che è la \eqref{eq:vo-c3}.
\end{proof}

\paragraph{La regione che cresce.} Sia $k_\alpha$ il quantile della distribuzione del raggio al livello $\alpha$: $k_\alpha=\sqrt{-2\ln(1-\alpha)}$ in due dimensioni ($k_{0,9}=2{,}146$, $k_{0,5}=1{,}177$); in tre dimensioni la radice di $C_3=\alpha$, trovata per bisezione ($k_{0,9}=2{,}500$) \prova{C}. La regione di raggio
\begin{equation}\label{eq:vo-ralpha}
r_\alpha(\tau)=k_\alpha\,\sigma(\tau)
\end{equation}
tiene la confidenza al livello $\alpha$ per ogni $\tau$, per la proposizione~\ref{prop:vo-c}; per $\tau$ grande cresce come $k_\alpha\sqrt{q/3}\,\tau^{3/2}$. È la forma da mostrare quando l'oggetto non è più visto: invece di una cifra che cala, una regione che si allarga a livello fissato.

\subsubsection{Il tempo di soglia}\label{sec:vo-tc}

Per una regione di raggio $r$ fissato, il tempo di soglia $t_C$ è l'ultimo istante in cui la confidenza è almeno $\alpha$.

\begin{proposizione}[Esistenza e unicità del tempo di soglia]\label{prop:vo-tc}
Sotto le ipotesi~\ref{ip:vo-cv} e~\ref{ip:vo-iso}, con $q>0$, sia
\[f(t)=\frac q3t^3+\sigma_v^2t^2+2\sigma_{0v}t+\sigma_0^2-\Bigl(\frac r{k_\alpha}\Bigr)^2 .\]
\begin{enumerate}[label=(\roman*)]
\item $C(\tau)\ge\alpha$ se e solo se $f(\tau)\le0$.
\item Se $\sigma_0<r/k_\alpha$, la cubica $f$ ha una e una sola radice positiva $t_C$, e $C(\tau)\ge\alpha$ per $\tau\in[0,t_C]$, $C(\tau)<\alpha$ per $\tau>t_C$.
\item Se $\sigma_0\ge r/k_\alpha$ e $\sigma_{0v}\ge0$, allora $C(\tau)<\alpha$ per ogni $\tau>0$ e si pone $t_C=0$. Se invece $\sigma_{0v}<0$, l'insieme $\{\tau>0:C(\tau)\ge\alpha\}$ è vuoto oppure un intervallo $[t_-,t_+]$ delimitato dalle due radici positive di $f$.
\end{enumerate}
\end{proposizione}
\begin{proof}
(i) $C$ è funzione strettamente crescente di $k^2=r^2/\sigma^2$ nelle \eqref{eq:vo-c2} e \eqref{eq:vo-c3} (è una funzione di ripartizione con densità positiva), quindi $C\ge\alpha$ equivale a $k\ge k_\alpha$, cioè a $\sigma^2(\tau)\le(r/k_\alpha)^2$, cioè a $f(\tau)\le0$.
(ii) $f$ è continua, $f(0)=\sigma_0^2-(r/k_\alpha)^2<0$ e $f(t)\to+\infty$ perché $q>0$: per il teorema dei valori intermedi esiste una radice positiva. Per il lemma~\ref{lem:vo-convessa} $f$ è convessa. Se vi fossero due radici $0<t_1<t_2$, scrivendo $t_1=\lambda\cdot0+(1-\lambda)t_2$ con $\lambda\in(0,1)$ la convessità darebbe $0=f(t_1)\le\lambda f(0)+(1-\lambda)f(t_2)=\lambda f(0)<0$, assurdo. Dunque la radice è unica; $f<0$ su $[0,t_C)$ e, per continuità e $f\to+\infty$, $f>0$ su $(t_C,\infty)$; si conclude con (i).
(iii) Se $\sigma_{0v}\ge0$ la $f$ è non decrescente (lemma~\ref{lem:vo-convessa}) e $f'(t)\ge qt^2>0$ per $t>0$, quindi $f(\tau)>f(0)\ge0$. Se $\sigma_{0v}<0$, l'insieme di sottolivello $\{f\le0\}$ di una funzione convessa è convesso, quindi un intervallo; con $f(0)\ge0$ i suoi estremi positivi, se esistono, sono radici di $f$, e una cubica ha al più tre radici reali.
\end{proof}

La radice si calcola con la formula di Cardano o, più stabilmente, con il metodo di Newton partendo da un $t$ con $f(t)>0$: per la convessità le iterate decrescono in modo monotono verso $t_C$. Nel prototipo la radice dalla cubica coincide con quella per bisezione entro $10^{-14}\unita{s}$ \prova{C}. Un controllo a campione su $40\,000$ casi casuali ($\sigma_0$ da $1$ a $50\unita{m}$, correlazione fra $-0{,}99$ e $0{,}99$, $q$ da $10^{-2}$ a $10^2$, $d=2$ e $3$) conferma la (ii) in $26\,870$ casi su $26\,870$ con $f(0)<0$; nei casi con $f(0)\ge0$ si hanno zero radici positive ($11\,169$) oppure due ($1961$, tutti con $\sigma_{0v}<0$), come dice la (iii) \prova{C}. La (iii) corregge una semplificazione: «$t_C=0$ se $\sigma_0\ge r/k_\alpha$» è vero solo con $\sigma_{0v}\ge0$.

\paragraph{Un esempio calcolato.} Per un aereo con $\sigma_0=7{,}5\unita{m}$, $\sigma_v=2\unita{m/s}$, $\sigma_{0v}=0$ e $q=4\unita{m^2/s^3}$ (valori di ipotesi, a scopo illustrativo), e una regione di $100\unita{m}$, il tempo di soglia al $90\,\%$ è $10{,}7\unita{s}$ sul disco e $9{,}6\unita{s}$ sulla sfera; la regione che tiene il $90\,\%$ ha raggio $17$, $39$, $91$, $238$ e $427\unita{m}$ a $1$, $5$, $10$, $20$ e $30\unita{s}$ (disco) \prova{C}. Lo stesso aereo, fuori dal modello, sta nell'insieme di presenza di raggio $r_0+210\,\tau$, cioè $2{,}1\unita{km}$ a $10\unita{s}$: la differenza di un ordine di grandezza è il prezzo della garanzia (\S\ref{sec:vo-presenza}).

\subsubsection{Il caso anisotropo}\label{sec:vo-c-aniso}

Per un aereo la covarianza è lunga lungo la rotta e stretta di traverso, e l'ipotesi~\ref{ip:vo-iso} non vale. Sia la posizione $x\sim\mathcal N(m,\Sigma)$ e la regione l'ellissoide $E=\{x:(x-c)^\top A^{-1}(x-c)\le1\}$, con $A=LL^\top$. Con $z=L^{-1}(x-c)$ la regione diventa la palla unitaria e $z\sim\mathcal N(L^{-1}(m-c),\,L^{-1}\Sigma L^{-\top})$. Diagonalizzando $L^{-1}\Sigma L^{-\top}=U\operatorname{diag}(\lambda)U^\top$ e ponendo $\mu=U^\top L^{-1}(m-c)$,
\begin{equation}\label{eq:vo-aniso}
C=\Pr\Bigl(\sum_{i=1}^d\lambda_i\bigl(y_i+\mu_i/\sqrt{\lambda_i}\bigr)^2\le1\Bigr),\qquad y_i\sim\mathcal N(0,1)\ \text{indipendenti},
\end{equation}
perché la rotazione $U$ lascia invariata la palla. È la funzione di ripartizione di un chi quadro generalizzato non centrato. Nel prototipo la si calcola per riduzione esatta a integrali annidati: in due dimensioni
\[C=\int_{-1}^{1}\frac{1}{\sqrt{\lambda_1}}\,\varphi\!\Bigl(\frac{z_1-\mu_1}{\sqrt{\lambda_1}}\Bigr)\Bigl[\Phi\Bigl(\frac{h-\mu_2}{\sqrt{\lambda_2}}\Bigr)-\Phi\Bigl(\frac{-h-\mu_2}{\sqrt{\lambda_2}}\Bigr)\Bigr]dz_1,\qquad h=\sqrt{1-z_1^2},\]
con $\mu_i$ qui nelle coordinate $z$ ruotate; in tre dimensioni con un secondo integrale sulla corda $\sqrt{h^2-z_2^2}$. Con la sostituzione $z_1=\sin\theta$ l'integranda diventa liscia, e la quadratura di Gauss--Legendre a 64 nodi dà un errore dominato da quello della funzione d'errore ($<1{,}2\cdot10^{-7}$ relativo). La regione può essere fissa nel mondo (la zona dichiarata alla perdita, $c$ costante) o seguire la previsione ($c=m(\tau)$). L'approssimazione con un $\sigma$ equivalente $\sqrt{\lambda_1\lambda_2}$ non è usata: sbaglia proprio quando $\lambda_1/\lambda_2$ è grande.

\paragraph{Verifica numerica.} Contro le forme chiuse \eqref{eq:vo-c2}--\eqref{eq:vo-c3} lo scarto massimo è $1{,}4\cdot10^{-8}$ in due dimensioni e $7{,}6\cdot10^{-8}$ in tre, su 20 coppie $(\sigma,r)$ per dimensione; contro un Monte Carlo di $10^6$ campioni su 20 casi anisotropi e fuori centro, tutti i casi cadono nell'intervallo al $99{,}9\,\%$ ($|z|\le2{,}11$) \prova{C}. Il controllo vede gli errori di codice e di algebra; condivide con il calcolo il modello gaussiano.

\subsubsection{Più regimi di moto}\label{sec:vo-c-imm}

Il filtro del sistema è una miscela di modelli interagenti \cite{blom1988,mazor1998,barshalom2001} con tre regimi: velocità quasi costante, virata coordinata con tasso stimato (linearizzata), accelerazione di Wiener \cite{li2003}. Durante la perdita ogni regime $j$ evolve con la sua previsione gaussiana $\mathcal N(m_j(\tau),\Sigma_j(\tau))$, e i pesi con la catena di Markov dei regimi, $\mu(\tau+h)=\Pi(h)^\top\mu(\tau)$.

\begin{proposizione}[Confidenza della miscela]\label{prop:vo-imm}
Se la densità della posizione è la miscela $\sum_j\mu_j(\tau)\,\mathcal N(m_j(\tau),\Sigma_j(\tau))$, allora per ogni regione misurabile $E$
\[C(\tau)=\sum_j\mu_j(\tau)\,P_j(\tau),\qquad P_j(\tau)=\Pr\nolimits_{\mathcal N(m_j,\Sigma_j)}(x\in E),\]
e ogni $P_j$ si calcola con la \eqref{eq:vo-aniso}.
\end{proposizione}
\begin{proof}
È la formula della probabilità totale, condizionando sul regime.
\end{proof}
La formula è esatta \emph{per la miscela}; la miscela approssima la densità vera, per la fusione dei momenti nell'interazione e per la linearizzazione della virata \prova{D}. Due conseguenze. Per $\tau$ molto maggiore del tempo di soggiorno nei regimi i pesi tendono alla distribuzione stazionaria della catena e la regione si allarga per effetto del regime più disperso: è la causa della prudenza misurata oltre $10\unita{s}$. E $C$ della miscela non è necessariamente monotona: se il peso passa verso un regime meno disperso, $C$ può crescere senza alcuna misura. Nei registri simulati degli scenari F2, F3 ed F5 si osservano 76 piccole violazioni (il denominatore, in passi, non è riportato nei registri; è un insieme diverso dai 1440 passi della taratura del \S\ref{sec:vo-taratura}) della regola «la confidenza non cresce durante la perdita» \prova{S}. Una correzione che si dimostra in una riga è dichiarare $\tilde C(\tau)=\min_{0\le s\le\tau}C(s)$: è non crescente per costruzione e non supera mai $C$, quindi non è mai più ottimista del modello; non corregge la taratura, che resta da provare (\S\ref{sec:ps-tabella}).

\subsubsection{Fuori dal modello: l'insieme di presenza}\label{sec:vo-presenza}

Dopo un urto, un'avaria o una manovra non modellata la covarianza del filtro non descrive più l'errore: nel simulato, dopo la collisione veloce, la confidenza prevista era $0{,}976$ e il vero stava nella regione in $14/24$ passaggi, $0{,}583$ $[0{,}388;\,0{,}755]$ \prova{S}. Lì il sistema mostra un'altra grandezza, che non è una probabilità.

\begin{proposizione}[Contenimento garantito]\label{prop:vo-presenza}
Siano $\hat p_0$ la posizione stimata all'ultima misura e $r_0\ge0$. Se
\begin{enumerate}[label=(\alph*)]
\item la posizione vera all'ultima misura soddisfa $|p(0)-\hat p_0|\le r_0$, e
\item la velocità vera rispetto al riferimento fisso soddisfa $|v(s)|\le\vmax$ per ogni $s\in[0,\tau]$,
\end{enumerate}
allora $|p(\tau)-\hat p_0|\le r(\tau)=r_0+\vmax\,\tau$. Se inoltre $|a|\le a_{\max}$ e $|v(0)-\hat v_0|\le3\sigma_v$, vale la stima più stretta della proposizione~\ref{prop:lp-contenimento}.
\end{proposizione}
\begin{proof}
$|p(\tau)-\hat p_0|\le|p(0)-\hat p_0|+\bigl|\int_0^\tau v(s)\,ds\bigr|\le r_0+\int_0^\tau|v(s)|\,ds\le r_0+\vmax\tau$, per la disuguaglianza triangolare e quella dell'integrale.
\end{proof}
Il contenimento vale per ogni traiettoria che rispetti (a) e (b), comunque manovri: non serve alcun modello di moto. Le sue ipotesi vanno però dichiarate con precisione, perché sono esse a portare la garanzia:
\begin{itemize}
\item (b) è un'ipotesi sulla classe, non una misura: $\vmax=210\unita{m/s}$ per gli aerei in volo, dai dati pubblici (\S\ref{sec:lp-traiettorie}), e $45\unita{m/s}$ per i droni (ipotesi senza fonte). Un drone da corsa più veloce di $45\unita{m/s}$ esce dalla regione;
\item (a) non è garantita: l'errore iniziale è limitato solo in probabilità. Nel sistema $r_0=3{,}4\,\sigma_0$ più un margine; la garanzia è quindi condizionata all'evento (a), la cui probabilità sotto l'ipotesi~\ref{ip:vo-cv} si calcola con la proposizione~\ref{prop:vo-c};
\item (b) usa il riferimento fisso: per i frammenti, trascinati dal vento, vale invece il volume $V(\tau)$ della proposizione~\ref{prop:lp-detriti};
\item come in ogni dimostrazione di questo tipo, il parametro nel codice deve essere quello dell'enunciato (\S\ref{sec:vo-nucleo}).
\end{itemize}
Nel simulato l'insieme ha contenuto il vero in tutti i campioni dopo l'urto ($939/939$, F3, B) e dopo l'avaria ($1072/1072$, F4, A) \prova{S}: è un controllo del codice e dei parametri, non una frequenza, perché le traiettorie simulate rispettano l'inviluppo per costruzione.

\subsubsection{Verifica su traiettorie reali}\label{sec:vo-adsb}

La formula è stata messa alla prova su traiettorie pubbliche (gli archivi della tabella~\ref{tab:lp-dati}, fra cui la rete ADS-B OpenSky \cite{schafer2014}): aerei con velocità mediana oltre $60\unita{m/s}$ (104 sequenze, 31 mezzi; nella metà di prova 12 aerei identificati dal codice ICAO e 4 numeri di volo di una sola rotta, che su più giorni possono essere aerei diversi), droni ad ala fissa (47 sequenze) e un multirotore (196 voli). Il rumore del sensore è simulato (aerei: $7{,}5\unita{m}$ isotropi a $1\unita{Hz}$; droni: il radar per droni a $1{,}5\unita{km}$). Le occlusioni sono virtuali, di 1--$30\unita{s}$, 8 istanti a caso per sequenza e per durata. Taratura e prova sono divise per mezzo; il fattore di scala di $q$ è scelto sulla taratura. La prova usa la trasformata integrale di probabilità: per ogni caso $u=C$ della sfera centrata sulla previsione con raggio pari all'errore vero; la verità sta nella regione al livello $\alpha$ se e solo se $u\le\alpha$. La tabella~\ref{tab:vo-adsb} dà la copertura della regione al $90\,\%$ del filtro a più regimi \emph{prima} della taratura descritta più avanti.

\begin{table}[H]\centering\footnotesize
\caption{Copertura osservata della regione al $90\,\%$ del filtro a più regimi sulla metà di prova, per durata dell'occlusione, prima della taratura \prova{P}. Wilson per caso: i casi di uno stesso mezzo sono correlati e l'intervallo è ottimista. Sono limiti le celle con il limite inferiore sotto $0{,}70$ o la stima sotto $0{,}80$.}\label{tab:vo-adsb}
\begin{tabularx}{\linewidth}{@{}>{\RaggedRight\arraybackslash}p{4.0cm}LLLLL@{}}
\toprule
Mezzi & 1--$2\unita{s}$ & 3--$5\unita{s}$ & 6--$10\unita{s}$ & 11--$20\unita{s}$ & 21--$30\unita{s}$\\
\midrule
aerei, prima della taratura & $0{,}71$ $[0{,}67;\,0{,}74]$ (limite) & $0{,}92$ $[0{,}89;\,0{,}94]$ & $0{,}95$ $[0{,}92;\,0{,}97]$ & $0{,}98$ $[0{,}96;\,0{,}99]$ & $0{,}96$ $[0{,}94;\,0{,}98]$\\
aerei in virata & -- & -- & $29/37=0{,}78$ $[0{,}63;\,0{,}89]$ (limite) & $68/71=0{,}96$ & $96/105=0{,}91$ $[0{,}85;\,0{,}95]$\\
aerei in accelerazione & $2/6$ (limite) & $14/17$ & $15/18$ & $22/25$ & $35/38$\\
droni ad ala fissa (forse un solo velivolo) & $0{,}91$ $[0{,}88;\,0{,}93]$ & $0{,}92$ & $0{,}96$ & $1{,}00$ & $1{,}00$\\
multirotore (un esemplare) & $0{,}95$ $[0{,}94;\,0{,}96]$ & $0{,}99$ & $1{,}00$ & $1{,}00$ & $1{,}00$\\
\bottomrule
\end{tabularx}
\end{table}

\paragraph{Lettura.} Prima della taratura il filtro era prudente fra $3$ e $30\unita{s}$, con una regione più larga del necessario, e ottimista a 1--$2\unita{s}$ \prova{P}. I limiti erano tre.
\begin{enumerate}[label=\arabic*.]
\item \emph{A 1--$2\unita{s}$ la regione è ottimista} ($0{,}71$ invece di $0{,}90$), e lo è anche il filtro a velocità costante ($0{,}72$). La causa più probabile è la verità stessa: le posizioni ADS-B hanno un rumore e un'età variabili, con un campionamento irregolare di 1--$4\unita{s}$, che la \eqref{eq:vo-sigma} non contiene. La taratura che segue aggiunge per questo un termine di riferimento alla varianza.
\item \emph{In virata a 6--$10\unita{s}$} la copertura è $0{,}78$, con il limite superiore sotto $0{,}90$; il filtro a velocità costante, per confronto, scende a $0{,}24$ $[0{,}13;\,0{,}40]$.
\item \emph{In accelerazione longitudinale forte} (corsa di decollo) il filtro non ha il regime adatto: in un caso a $10\unita{s}$ dichiara un raggio di $144\unita{m}$ contro un errore vero di $801\unita{m}$. Al livello del $50\,\%$, a 1--$2\unita{s}$, la copertura è $0{,}36$ $[0{,}33;\,0{,}40]$: la curva sbaglia la forma, non solo un livello.
\end{enumerate}
Nel nucleo simulato la confidenza della covarianza è vicina all'osservato nelle fasce alte nel solo scenario in salita (F2, A: prevista $0{,}95$--$0{,}99$, osservata $56/60$ passaggi $[0{,}841;\,0{,}974]$) e con il radar a $10\unita{Hz}$; nel controllo con il radar a $1\unita{Hz}$ una confidenza dichiarata di almeno $0{,}9$ è stata osservata in $1/78$ campioni, e nelle fasce medie la confidenza è ottimista di $0{,}1$--$0{,}3$ \prova{S}. Per questo la cifra della confidenza del nucleo non è mostrata come una probabilità tarata; e la zona costruita con il quantile $0{,}99$ della sola covarianza, che ha contenuto il vero in una frazione da $0{,}003$ a $0{,}91$ dei cicli, non è usata. Si mostrano la regione $r_{0,9}(\tau)$ con la sua copertura osservata per tipo di moto e, fuori dal modello, l'insieme di presenza.
\subsubsection{Taratura della confidenza}\label{sec:vo-taratura}

La taratura modifica la varianza e la lettura della confidenza in quattro punti, tutti fissati sulla metà di taratura:
\begin{equation}\label{eq:vo-sigma-rif}
\sigma^2(\tau)=\sigma_0^2+2\sigma_{0v}\,\tau+\sigma_v^2\,\tau^2+\frac{q\,\tau^3}{3}+\sigma_{\mathrm{rif}}^2,\qquad C(\tau)=1-\exp\!\Bigl(-\frac{r^2}{2\sigma^2(\tau)}\Bigr),\qquad r_{0,9}(\tau)=2{,}146\,\sigma(\tau).
\end{equation}
\begin{enumerate}[label=\arabic*.]
\item un termine $\sigma_{\mathrm{rif}}^2$, costante, per il rumore e l'età della verità di riferimento ($\sigma_{\mathrm{rif}}=20\unita{m}$ per gli aerei); il gonfiamento che serve sta in questo termine, non in un fattore su $q$. La copertura è quindi misurata rispetto alla posizione ADS-B, non alla posizione vera: con un'altra verità di riferimento $\sigma_{\mathrm{rif}}$ va rifissato (tappa 2 del \S\ref{sec:ps-piano});
\item un regime di accelerazione longitudinale nella miscela, con i pesi $\mu_j$ fissati all'istante della perdita: $C=\sum_j\mu_jP_j$ (proposizione~\ref{prop:vo-imm});
\item una ricalibrazione isotonica per fascia di $\tau$ e per regime dominante alla perdita;
\item la lettura con il minimo progressivo $\tilde C(\tau)=\min_{s\le\tau}C(s)$, non crescente per costruzione: la $C$ grezza ha 45 violazioni della monotonia su 1440 passi \prova{C}.
\end{enumerate}
La \eqref{eq:vo-sigma-rif} è dimostrata nel modello a velocità quasi costante con rumore gaussiano (proposizioni~\ref{prop:vo-sigma} e~\ref{prop:vo-c}, con $\sigma_{\mathrm{rif}}$ indipendente) e verificata numericamente. La \eqref{eq:vo-sigma-rif} è scritta per un solo regime nel piano; la copertura della tabella~\ref{tab:vo-taratura} è misurata sulla sfera ($k_{0,9}=2{,}500$), con la miscela della proposizione~\ref{prop:vo-imm} e la ricalibrazione isotonica: la cifra viene dalla catena intera, non dalla sola \eqref{eq:vo-sigma-rif}. La ricalibrazione è empirica. La tabella~\ref{tab:vo-taratura} confronta la copertura prima e dopo, sulla metà di prova \prova{P}.

\begin{table}[H]\centering\footnotesize
\caption{Copertura della regione al $90\,\%$ per gli aerei, metà di prova, prima e dopo la taratura \prova{P}. Wilson per caso: i casi di uno stesso mezzo sono correlati e l'intervallo è ottimista; fra graffe il bootstrap per mezzo. Sono limiti le celle in grassetto corsivo.}\label{tab:vo-taratura}
\setlength{\tabcolsep}{3pt}
\begin{tabularx}{\linewidth}{@{}>{\RaggedRight\arraybackslash}p{2.0cm}LLLLLLL@{}}
\toprule
$\tau$ & 1--$2\unita{s}$ & 3--$5\unita{s}$ & 6--$10\unita{s}$ & 11--$20\unita{s}$ & 21--$30\unita{s}$ & in virata & in accelerazione\\
\midrule
prima & $0{,}71$ $[0{,}67;\,0{,}74]$ & $0{,}92$ & $0{,}95$ & $0{,}98$ & $0{,}96$ & $0{,}905$ & $0{,}841$\\
dopo & $655/733=0{,}89$ $[0{,}87;\,0{,}91]$ $\{0{,}88;\,0{,}92\}$ & $372/416=0{,}89$ $[0{,}86;\,0{,}92]$ & $344/385=0{,}89$ $[0{,}86;\,0{,}92]$ & $334/374=0{,}89$ $[0{,}86;\,0{,}92]$ & $328/388=0{,}85$ $[0{,}81;\,0{,}88]$ $\{0{,}78;\,0{,}89\}$ & \textbf{\emph{161/221 = 0,73}} $[0{,}67;\,0{,}78]$ & \textbf{\emph{81/107 = 0,76}} $[0{,}67;\,0{,}83]$\\
$r_{0,9}$ mediano & $39\unita{m}$ (era 27) & -- & $138\unita{m}$ (era 264) & -- & $1106\unita{m}$ (era 2065) & -- & --\\
\bottomrule
\end{tabularx}
\end{table}

\paragraph{Lettura e condizioni.} Lo scarto medio dal livello dichiarato scende da $0{,}081$ a $0{,}016$, e la fascia 1--$2\unita{s}$ passa da $0{,}71$ a $0{,}89$ $[0{,}87;\,0{,}91]$: un risultato incoraggiante, con tre condizioni che vanno lette insieme alla cifra.
\begin{enumerate}[label=\arabic*.]
\item \emph{Concentrazione.} L'$89{,}5\,\%$ dei casi di prova degli aerei viene da 4 numeri di volo di una sola rotta, ripresa in più giorni; la metà di taratura contiene altri voli della stessa rotta. Non c'è fuga di mezzi fra le due metà, ma c'è correlazione di rotta. Senza i due mezzi più pesanti la fascia 21--$30\unita{s}$ vale $0{,}81$ $[0{,}75;\,0{,}86]$.
\item \emph{Virata.} In virata la copertura scende da $0{,}905$ a $0{,}73$ $[0{,}67;\,0{,}78]$: la taratura la peggiora, ed è un limite dichiarato. In accelerazione resta a $0{,}76$ $[0{,}67;\,0{,}83]$, limite.
\item \emph{Pochi droni.} Per i droni ad ala fissa e i multirotori la copertura dopo la taratura è $0{,}89$--$0{,}93$, ma da uno o pochi esemplari.
\end{enumerate}
In una frase: \emph{dopo la perdita, nella simulazione su traiettorie di volo pubbliche, la regione al $90\,\%$ ha contenuto l'aereo da 85 a 89 volte su 100 fra $1$ e $30\unita{s}$; i casi vengono per nove decimi da quattro voli della stessa rotta, e in virata la copertura scende a 73 su 100, un limite dichiarato.} Della confidenza non si mostrano: la cifra del nucleo con il radar a $1\unita{Hz}$, la zona al $99\,\%$ della sola covarianza e la cifra numerica dei registri di F2, F3 ed F5, ai quali il minimo progressivo non è applicato.

\paragraph{Fuori dal modello.} Dopo un urto o un'avaria si mostra l'insieme di presenza $r_0+\vmax\tau$ della proposizione~\ref{prop:vo-presenza}, con $\vmax=210\unita{m/s}$ per gli aerei e $45\unita{m/s}$ per i droni: un contenimento condizionato a $|p(0)-\hat p_0|\le r_0$, non una probabilità.

\subsection{Sensore a bordo di un drone in moto: lo scenario «scorta» (F7)}\label{sec:vo-scorta}

Negli scenari F1--F6 i sensori sono fermi. Nello scenario F7 il sistema è a bordo di un drone in volo (il \emph{portatore}) e mantiene la traccia di un drone veloce che lo precede; accanto, due altri droni collidono, con fiammata, fumo e detriti. Il caso mette insieme le tre difficoltà del capitolo: il moto proprio del sensore, un evento impulsivo vicino e la perdita degli ottici. Il sistema resta di sola osservazione e avviso: non comanda il portatore, non insegue e non intercetta; «mantenere la traccia» vuol dire tenere per il drone seguito una traccia o una regione dichiarata.

\paragraph{Modello (tutto ipotesi).}
\begin{itemize}
\item \emph{Portatore:} volo livellato a $120\unita{m}$, a $25$--$35\unita{m/s}$, in rotta retta (il caso peggiore per il campo visivo, perché il portatore non orienta il sensore); accelerazione massima $4\unita{m/s^2}$, tempo di reazione $T_r=0{,}25\unita{s}$, profilo deterministico come al \S\ref{sec:lim-portatore}.
\item \emph{Navigazione inerziale} di classe tattica con GNSS: posizione $0{,}5\unita{m}$, deriva $0{,}02\unita{m/s}$, passeggiata aleatoria $0{,}05\unita{m/\sqrt{s}}$; assetto $0{,}1^\circ$, deriva $0{,}005^\circ/\mathrm{s}$. Il moto proprio è compensato con la posa \emph{stimata}; l'errore di navigazione, correlato e che il filtro non riduce, entra nella regione come margine $3(\sigma_p+b\,t+\sigma_{rw}\sqrt t)+\rho\cdot3(\sigma_{att}+\omega t)$, con $\rho$ la distanza.
\item \emph{Drone seguito:} sezione radar $0{,}02\unita{m^2}$, parte a $150$--$300\unita{m}$ davanti, a $35$--$45\unita{m/s}$, con uno zig-zag laterale di $3\unita{m/s^2}$ e periodo $8\unita{s}$; nella variante «virata» vira a $8\unita{m/s^2}$ durante l'accecamento e per altri $2\unita{s}$ (circa $45$--$60^\circ$ di cambio di rotta). È designato all'avvio con un'etichetta a mano.
\item \emph{Evento} a $t=6\unita{s}$: due droni a $30$--$40\unita{m/s}$, con angolo d'incontro di $60$--$180^\circ$, collidono a $30$--$100\unita{m}$ di lato, davanti (fra il portatore e il drone seguito) o dietro. La fiammata satura le camere per $0{,}5$--$3\unita{s}$; il fumo è una sfera di $5$--$30\unita{m}$ che deriva, dura $5$--$20\unita{s}$ e blocca la vista ottica; i frammenti sono 30 più i 2 corpi, con velocità pari alla media dei mezzi più un incremento $|u|\le50\unita{m/s}$, velocità limite da $5$ a $30\unita{m/s}$ e sezione radar da $10^{-4}$ a $10^{-2}\unita{m^2}$. Il radar riceve gli echi dei frammenti, echi spuri nella nube per $3\unita{s}$ (in media 3 per scansione), un fondo di $0{,}5$ falsi per scansione e il mascheramento del drone seguito ($P_D\times0{,}3$) quando un frammento gli è vicino in distanza e in angolo.
\item \emph{Controllo F7c:} una quasi collisione con $30$--$60\unita{m}$ di separazione verticale e nessun urto; in metà dei passaggi un riflesso solare satura la camera per $0{,}5\unita{s}$.
\end{itemize}

\paragraph{Configurazioni} (classi di sensore dichiarate come ipotesi; nessuna scheda aperta per questo scenario). A: radar FMCW piccolo di bordo, $5\unita{Hz}$, campo di $\pm60^\circ$ in avanti, con camera fissa di $\pm30^\circ$. B: radar a scansione elettronica a $20\unita{Hz}$ con quattro pannelli ($360^\circ$), con camera nel visibile e nell'infrarosso su brandeggio. Le latenze di elaborazione non sono simulate in F7, e non è simulata una camera a eventi.

\paragraph{Il tracciatore.} È un componente nuovo, scritto per il sensore in moto: il nucleo degli scenari F1--F6 lavora su griglie di sensori fermi. Ne riprende le regole, non il codice: mai «libero» sul non osservato; insieme di presenza che cresce con il tempo; confidenza che non cresce durante la perdita; un evento non sparisce; un sensore perso non torna «affidabile» da solo.
\begin{itemize}
\item filtro di Kalman a velocità costante ($Q=80\unita{m^2/s^3}$); associazione con porta di Mahalanobis al $99{,}9\,\%$ e porta Doppler di $3\sigma+2\unita{m/s}$, assegnazione globale; conferma 3 su 5; il drone seguito non si cancella mai: senza misure diventa una regione;
\item regione del drone seguito: il massimo fra il raggio al $99{,}99\,\%$ della covarianza più il margine della navigazione e l'insieme di presenza $r+3\sigma_v\tau+a\tau^2/2$;
\item evento: un grappolo di almeno 5 misure nuove entro $40\unita{m}$ in $1\unita{s}$, insieme alla saturazione delle camere o a due tracce che convergono entro $20\unita{m}$;
\item una misura che cade nel volume dei detriti rende l'identità «presunta, dubbia»;
\item «libero» solo dentro il campo del radar, entro la portata alla quale un oggetto di $10^{-4}\unita{m^2}$ è visto con $P_D\ge0{,}99$ in 3 scansioni ($193\unita{m}$ in A, $339\unita{m}$ in B), e fuori da ogni regione e dal volume dei detriti;
\item avviso al portatore se il suo percorso stimato nei prossimi $5\unita{s}$ entra nel volume dei detriti allo stesso istante futuro.
\end{itemize}

\paragraph{Il volume dei detriti con il cammino massimo.} Il volume di F7 usa la resistenza per stringere la proposizione~\ref{prop:lp-detriti}.

\begin{proposizione}[Cono e cammino massimo dei frammenti]\label{prop:vo-cammino}
Un frammento parta da $p_0$ con velocità $v(0)=v_c+u$, $|u|\le\Delta v$, e si muova con $\dot v=\mathbf g-k|v|\,v$, dove $\mathbf g$ è la gravità ($|\mathbf g|=g$), $k=g/v_t^2>0$ e il vento è nullo. Sia $v_0=|v(0)|$. Allora per ogni $\tau\ge0$:
\begin{enumerate}[label=(\roman*)]
\item la proiezione orizzontale dello spostamento è $S\,(v_c+u)_h$ per un $S\in[0,\tau]$;
\item $|p(\tau)-p_0|\le L(\tau)=\dfrac{v_t^2}{g}\ln\!\Bigl(\cosh\dfrac{g\tau}{v_t}+\dfrac{v_0}{v_t}\sinh\dfrac{g\tau}{v_t}\Bigr)$;
\item $L(\tau)$ cresce con $v_t$ e con $v_0$: basta dichiarare $v_t\le v_{t,\max}$ e $v_0\le|v_c|+\Delta v$.
\end{enumerate}
\end{proposizione}
\begin{proof}
(i) La gravità è verticale, quindi la parte orizzontale soddisfa $\dot v_h=-k|v|\,v_h$: $v_h(s)=\lambda(s)\,v_h(0)$ con $\lambda(s)=\exp(-\int_0^sk|v|)\in(0,1]$. Lo spostamento orizzontale è $\bigl(\int_0^\tau\lambda\bigr)v_h(0)$, e $S=\int_0^\tau\lambda\in[0,\tau]$.
(ii) Sia $w=|v|$. Dove $w>0$, $\dot w=v\cdot\dot v/w=(v\cdot\mathbf g)/w-kw^2\le g-kw^2$. Sia $U$ la soluzione di $\dot U=g-kU^2$, $U(0)=v_0$: con $x=g\tau/v_t$, $c=v_0/v_t$, $N=\sinh x+c\cosh x$ e $D=\cosh x+c\sinh x>0$, si verifica che $U=v_tN/D$ (poiché $dN/dx=D$, $dD/dx=N$, si ha $dU/d\tau=g(1-N^2/D^2)=g-kU^2$, e $U(0)=v_0$). Il secondo membro è lipschitziano su ogni intervallo limitato, e per il lemma di confronto $w(\tau)\le U(\tau)$. Allora $|p(\tau)-p_0|\le\int_0^\tau w\le\int_0^\tau U=\frac{v_t^2}{g}\ln D(\tau)=L(\tau)$, perché $\frac{d}{d\tau}\frac{v_t^2}{g}\ln D=\frac{v_t^2}{g}\cdot\frac{g}{v_t}\cdot\frac ND=U$ e $D(0)=1$.
(iii) Un $k$ più piccolo (cioè $v_t$ più grande) o un $v_0$ più grande aumentano il secondo membro o il dato iniziale dell'equazione di confronto, e quindi $U$ e il suo integrale.
\end{proof}
Il volume dichiarato è il cono orizzontale $\{p_0+S(v_c+u)_h:\ S\in[0,\tau],\ |u|\le\Delta v\}$ fra due estremi di quota, intersecato con la sfera di raggio $L(\tau)$ e ingrossato di $r_0$. Nel sistema $\Delta v=80+10\unita{m/s}$ (più $60\unita{m/s}$ se le velocità dei mezzi non sono note), $v_{t,\max}=40\unita{m/s}$, $r_0=15\unita{m}$ più l'estensione del grappolo, e $v_c$ è presa dalle tracce dei mezzi. La proposizione è \prova{D} nel modello dichiarato (resistenza lungo $-v$, vento nullo); il simulatore dei frammenti condivide lo stesso modello, quindi la verifica in simulazione vede gli errori di codice, non quelli del modello.

\paragraph{Risultati.} La tabella~\ref{tab:vo-scorta} riporta la valutazione, fatta una sola volta su semi nuovi dopo aver fissato i parametri su semi di sviluppo: 60 passaggi per cella e 200 per il controllo, con l'intervallo di Wilson al $95\,\%$ per passaggio. Le celle sono le quattro geometrie (lato, davanti, dietro, virata) per le due configurazioni.

\begin{table}[H]\centering\footnotesize
\caption{Scenario F7 «scorta»: risultati simulati per passaggio \prova{S}. A: radar di bordo frontale ($\pm60^\circ$); B: radar a $360^\circ$. Simulazione dalle schede e da ipotesi, etichette a mano, nessun sensore fisico.}\label{tab:vo-scorta}
\begin{tabularx}{\linewidth}{@{}>{\RaggedRight\arraybackslash}p{5.0cm}>{\RaggedRight\arraybackslash}p{2.0cm}>{\RaggedRight\arraybackslash}p{4.6cm}L@{}}
\toprule
Enunciato & Classe & Conteggio e intervallo & Base\\
\midrule
Durante una collisione a $30$--$100\unita{m}$ con fiammata, fumo e detriti, la regione che contiene il drone seguito non si perde in nessun ciclo & solido & B: $60/60$ $[0{,}940;\,1]$ in ognuna delle 4 geometrie; A: $59/60$ $[0{,}911;\,0{,}997]$ davanti, $60/60$ nelle altre & regione con il margine della navigazione e l'insieme di presenza\\
Mai «libero» sul drone seguito né su un frammento & solido (in parte per costruzione) & $0/60$ $[0;\,0{,}060]$ in ogni cella & «libero» solo se osservato, entro la portata per $10^{-4}\unita{m^2}$; il drone seguito sta spesso oltre quella portata ($193$ e $339\unita{m}$)\\
Nessun frammento attribuito al drone seguito senza dirlo & solido & $0/60$ $[0;\,0{,}060]$ in ogni cella; le misure nel volume sono dichiarate «presunta, dubbia» (davanti 20/60, virata 23--25/60) & porta di Mahalanobis e Doppler; marcatura nel volume\\
Il volume dei detriti contiene tutti i frammenti simulati & solido nel modello & B: $60/60$ $[0{,}940;\,1]$ in ogni geometria; campioni $44\,051$, $43\,681$, $43\,600$ e $42\,950$, tutti dentro & proposizione~\ref{prop:vo-cammino}\\
Camere dichiarate non affidabili in ogni istante di saturazione; nessun falso evento nel controllo simulato, senza disturbo a grappoli & solido (innesco ottimista; controllo debole) & 0 cicli saturi dichiarati affidabili; F7c: $0/200$ $[0;\,0{,}019]$ in A e in B & --\\
Con il radar non disturbato, se il percorso del portatore entra nel volume, l'avviso arriva prima del frammento entro $10\unita{m}$ & incoraggiante (radar indisturbato; soglia dei $10\unita{m}$ scelta dagli autori) & davanti $56/56$ $[0{,}936;\,1]$, virata $54/54$ $[0{,}934;\,1]$; margine minimo $0{,}4\unita{s}$ & percorso stimato a $5\unita{s}$ intersecato con il volume\\
Il volume è prudente: l'avviso scatta anche quando nessun frammento passa vicino & limite & B: lato $60/60$, dietro $60/60$ senza intersezione & --\\
Radar solo frontale (A): un urto di lato o dietro non è visto (il sistema dice «non so», mai «libero») & limite & evento di lato $1/60$ $[0{,}003;\,0{,}089]$, dietro $0/60$; in B $60/60$ in tutte le geometrie & --\\
Radar solo frontale (A): se il drone seguito vira fuori dal campo, l'identità non è più garantita & limite & virata: identità $43/60$ $[0{,}592;\,0{,}815]$ (B: $60/60$); su semi nuovi $15/20$ $[0{,}531;\,0{,}888]$ & la regione lo contiene sempre: $60/60$\\
\bottomrule
\end{tabularx}
\end{table}

Il fallimento in virata con il radar frontale ha una causa precisa, non un errore d'associazione: il drone seguito vira, il portatore resta in rotta retta, e il drone esce dal campo di $\pm60^\circ$. La regione lo contiene in ogni ciclo, la traccia resta «non vista» con la regione che cresce, e l'errore finale supera i $30\unita{m}$ (da $30$ a $83\unita{m}$); le associazioni spurie sono 0. Con il radar a $360^\circ$ (B) l'identità in virata e l'evento di lato e dietro passano a $60/60$: la miglioria è applicata in simulazione.

\paragraph{Robustezza su semi e geometrie nuovi.} Su 20 passaggi per cella con semi nuovi, la continuità è $20/20$ $[0{,}839;\,1]$ in tutte le celle e il «libero» sul drone seguito o su un frammento è $0/20$ in tutte le 16 celle; il controllo dà $0/60$ falsi eventi in A e in B \prova{S}. Con l'urto portato a $10$--$30\unita{m}$ davanti, continuità e «libero» reggono ($20/20$). Con frammenti a $100\unita{m/s}$, più veloci dei $90$ dichiarati, il volume li contiene ancora: $15\,019/15\,019$ e $14\,557/14\,557$ campioni in B, $20/20$ passaggi in due varianti. Una quasi collisione a $5$--$20\unita{m}$ senza urto non apre falsi eventi ($0/60$ in A e in B), ma è un controllo debole: il controllo non contiene fenomeni che generano grappoli di echi, come stormi di uccelli, pioggia o contromisure.

\paragraph{Il radar accecato dall'esplosione (variante F7r).} In F7 il radar non è accecato: fumo e luce non lo degradano, e la riacquisizione del drone seguito è facile per ipotesi. La tabella~\ref{tab:vo-fisica} riassume, come ordini di grandezza, che cosa degrada ciascun sensore.

\begin{table}[H]\centering\footnotesize
\caption{Effetto dei fenomeni di una collisione vicina sui sensori: ordini di grandezza, quasi tutti ipotesi (i).}\label{tab:vo-fisica}
\begin{tabularx}{\linewidth}{@{}>{\RaggedRight\arraybackslash}p{3.0cm}L>{\RaggedRight\arraybackslash}p{3.6cm}>{\RaggedRight\arraybackslash}p{2.0cm}@{}}
\toprule
Fenomeno & Radar (banda X, Ku, millimetrica) & Camere nel visibile e nell'infrarosso & Durata (i)\\
\midrule
fumo (particelle da $0{,}1$ a $100\unita{\mu m}$) & non degradato: particelle molto più piccole della lunghezza d'onda ($3\unita{cm}$, $1{,}9\unita{cm}$, $3{,}9\unita{mm}$), diffusione trascurabile & visibile accecato lungo la linea di vista; infrarosso degradato & $5$--$20\unita{s}$\\
fiammata e luce & nessun effetto della luce & saturate; l'infrarosso medio e lungo satura vicino al picco di emissione & $0{,}5$--$3\unita{s}$\\
fiamma ionizzata & attenuazione solo se la linea di vista attraversa la fiamma: assorbimento misurato in banda X su una fiamma vegetale a $1050\unita{K}$ \cite{mphale2007}; trasferimento ad altre fiamme ipotesi & come sopra & $0{,}5$--$3\unita{s}$\\
nube densa di detriti vicina & degradato: eco forte (va come $\sigma/R^4$), possibile saturazione del ricevitore sotto i 20--$30\unita{m}$, controllo automatico del guadagno che abbassa la $P_D$, mascheramento nella stessa cella di distanza e Doppler, cammini multipli & occlusioni sparse e falsi bersagli caldi & $1$--$3\unita{s}$\\
impulso elettromagnetico & trascurabile per due droni senza esplosivo (ipotesi) & trascurabile & $<1\unita{ms}$\\
\bottomrule
\end{tabularx}
\end{table}

La variante F7r acceca il radar con un modello \emph{pessimista} (ipotesi): nessuna misura radar per $0{,}2$--$2\unita{s}$ secondo la distanza dell'urto ($2\unita{s}$ a $30\unita{m}$, $0{,}2\unita{s}$ a $100\unita{m}$), occultamento del drone seguito ($P_D\times0{,}3$) quando la linea di vista passa nella nube, e una cella «totale» con l'urto a $30$--$50\unita{m}$ davanti, radar cieco $2\unita{s}$ e camere sature $2$--$3\unita{s}$. Un controllo automatico del guadagno reale recupera di norma in tempi molto più brevi (da millisecondi a $100\unita{ms}$, ipotesi). Il radar è dichiarato «non affidabile» in ogni ciclo saturo, e non torna affidabile da solo. Esiti, 60 passaggi per cella e 200 per il controllo \prova{S}:
\begin{itemize}
\item \emph{reggono:} «libero» sul drone seguito o su un frammento $0/60$ $[0;\,0{,}060]$ in ogni cella; continuità $59/60$ $[0{,}911;\,0{,}997]$ nella cella totale con il radar frontale, $60/60$ in tutte le altre; falsi eventi nel controllo $0/200$ $[0;\,0{,}019]$; tutti i campioni dei frammenti dentro il volume, quando il volume esiste;
\item \emph{crolla l'avviso al portatore} prima di un frammento entro $10\unita{m}$: davanti $27/52$ $[0{,}387;\,0{,}649]$ (A) e $28/52$ $[0{,}405;\,0{,}667]$ (B); virata $19/55$ $[0{,}234;\,0{,}477]$ e $21/55$ $[0{,}265;\,0{,}514]$; cella totale $0/59$ $[0;\,0{,}061]$. A $30$--$50\unita{m}$, con frammenti a circa $50\unita{m/s}$, i frammenti arrivano prima che il radar torni: nessun avviso basato sul volume localizzato può precederli;
\item \emph{scendono con il radar frontale (A):} la riacquisizione entro $2\unita{s}$ dal ritorno del radar, davanti $41/60$ $[0{,}558;\,0{,}787]$ e virata $43/60$ $[0{,}592;\,0{,}815]$, perché il drone seguito resta dentro il volume dei detriti e il sistema, per prudenza, non lo riaggancia lì; l'evento dichiarato, cella totale $7/60$ $[0{,}058;\,0{,}222]$ e virata $33/60$ $[0{,}425;\,0{,}669]$. In un solo passaggio (cella totale, A) una riacquisizione presunta viene da una traccia che non era del drone seguito.
\end{itemize}
In breve: \emph{con il radar accecato fino a $2\unita{s}$ insieme alle camere, il sistema non ha mai detto «libero» né sul drone seguito né sui frammenti, e il drone è rimasto nella regione in 59--60 passaggi simulati su 60; l'avviso dei detriti al portatore, invece, con il radar cieco non è affidabile.}

\paragraph{Correzione provata (c1).} Alla saturazione congiunta di radar e camere il sistema dichiara subito un «evento vicino non localizzato», con un volume prudente attorno al portatore dentro il quale non dice mai «libero». Su semi nuovi, con confronto appaiato \prova{S}:
\begin{itemize}
\item l'avviso prima del frammento entro $10\unita{m}$ passa, nella cella totale, da $0/59$ $[0;\,0{,}061]$ a $59/59$ $[0{,}939;\,1]$ (McNemar appaiato, $p=3{,}5\cdot10^{-18}$); davanti a $51/51$ $[0{,}930;\,1]$ e in virata a $58/58$ $[0{,}938;\,1]$;
\item continuità $60/60$ in ogni cella, «libero» $0$; identità spurie non dichiarate da $1/60$ a $0/60$; falsi avvisi nel controllo $0/200$ $[0;\,0{,}019]$, una prova debole perché nel controllo simulato il radar non satura mai.
\end{itemize}
Limiti: l'avviso non localizza l'evento (una sfera fino a $565\unita{m}$ di raggio, accesa per una mediana di $8\unita{s}$); se il radar satura per un'altra causa, con il bagliore naturale della camera, i falsi avvisi sono $11/200$ $[0{,}031;\,0{,}096]$ (e $200/200$ con il bagliore forzato a $6\unita{s}$); la riacquisizione entro $2\unita{s}$ resta sotto la soglia di $51/60$ fissata prima: davanti $49/60$ $[0{,}701;\,0{,}894]$, in virata $40/60$ $[0{,}541;\,0{,}773]$. È una strada promettente, non adottata.

\paragraph{Confidenza nel tempo del drone seguito.} Con la \eqref{eq:vo-sigma} nel piano e i parametri dello stato stazionario del filtro di F7 (misura a $250\unita{m}$; A: $\sigma_0=3{,}05\unita{m}$, $\sigma_{0v}=13{,}6\unita{m^2/s}$, $\sigma_v=6{,}84\unita{m/s}$; B: $\sigma_0=1{,}26\unita{m}$, $\sigma_{0v}=4{,}44\unita{m^2/s}$, $\sigma_v=5{,}17\unita{m/s}$; per entrambi $q=80\unita{m^2/s^3}$ per asse; insieme di presenza $r(\tau)=4{,}59\,\sigma_0+3\sigma_v\tau+a\tau^2/2$ con $a=10\unita{m/s^2}$), per una regione fissa di $30\unita{m}$ la tabella~\ref{tab:vo-curva-f7} dà la confidenza dopo una perdita \emph{totale}, di radar e ottici insieme \prova{C}, calcolata da ipotesi. Con il radar frontale la confidenza scende sotto $0{,}9$ a $1{,}35\unita{s}$.

\begin{table}[H]\centering\footnotesize
\caption{Drone seguito in F7: confidenza $C(\tau)$ per una regione di $30\unita{m}$, raggio $r_{0,9}(\tau)$ e insieme di presenza, dopo una perdita totale \prova{C}.}\label{tab:vo-curva-f7}
\begin{tabularx}{\linewidth}{@{}>{\RaggedRight\arraybackslash}p{4.2cm}LLLLLLL@{}}
\toprule
$\tau$ (s) & 0 & 0,5 & 1,0 & 1,5 & 2,0 & 3,0 & 5,0\\
\midrule
A: $C$ & 1,00 & 1,00 & 0,98 & 0,84 & 0,62 & 0,31 & 0,09\\
A: $r_{0,9}$ (m) & 6,5 & 13,2 & 22,5 & 33,6 & 46,2 & 75,3 & 146\\
A: insieme di presenza (m) & 14 & 26 & 40 & 56 & 75 & 121 & 242\\
B: $C$ & 1,00 & 1,00 & 1,00 & 0,93 & 0,73 & 0,37 & 0,11\\
B: $r_{0,9}$ (m) & 2,7 & 8,6 & 17,1 & 27,6 & 39,5 & 67,5 & 137\\
\bottomrule
\end{tabularx}
\end{table}
Durante l'accecamento \emph{ottico} simulato, quando il radar continua a vedere, la quota di cicli con il drone dentro $r_{0,9}$ è $0{,}86$--$0{,}93$ in A e $0{,}88$--$0{,}92$ in B nei primi $5\unita{s}$ (cicli correlati, senza intervallo): in A la regione al $90\,\%$ del filtro è leggermente ottimista, ed è un limite, non un «$90\,\%$». Dentro la regione dichiarata ci sono tutti i cicli \prova{S}.

\paragraph{Limiti dello scenario F7.}
\begin{enumerate}[label=\arabic*.]
\item \emph{Ritardo dell'evento non presentabile.} L'evento è dichiarato alla prima scansione in tutti i passaggi: è un artefatto ottimista del simulatore, perché i frammenti nascono tutti nel punto d'urto e a $30$--$100\unita{m}$ anche sezioni di $10^{-4}\unita{m^2}$ ritornano subito, e perché la saturazione delle camere è presa dalla finestra vera dell'evento. Il ritardo dell'evento in F7 non è un risultato.
\item \emph{Velocità vera del portatore.} Il tracciatore usa la velocità vera del portatore nella porta Doppler e nel percorso dell'avviso, e la sua posizione vera nel margine della navigazione; le misure usano invece la posa stimata. A velocità costante l'errore di velocità della navigazione vale circa $0{,}02\unita{m/s}$ (ipotesi). La miglioria è usare la velocità stimata dalla navigazione inerziale.
\item \emph{Margine dell'avviso.} Con l'urto a $10$--$30\unita{m}$ davanti, il margine dell'avviso al portatore è sotto $T_r=0{,}25\unita{s}$ in 7 casi su 20 $[0{,}181;\,0{,}567]$ (minimo $0{,}1\unita{s}$), in A e in B: è fisica, perché a $10\unita{m}$ e $30\unita{m/s}$ restano circa $0{,}3\unita{s}$. Serve un avviso di prossimità dalle tracce \emph{prima} dell'urto, come in F3.
\item \emph{Avviso non discriminante.} Il volume cresce di circa $90\unita{m/s}$ e l'avviso scatta anche quando nessun frammento passa vicino.
\item \emph{Controllo debole e prove limitate.} Il controllo non contiene disturbo a grappoli; il portatore non manovra; il vento è nullo; le varianti di geometria sono tre, con 20 passaggi ciascuna; il tracciatore non è integrato nel nucleo.
\end{enumerate}

\subsection{Sensori allo stato dell'arte e calcolo dedicato}\label{sec:vo-hw}

\paragraph{Sensori.} Per il volo veloce contano la rivisita e la portata sull'oggetto più piccolo. La tabella~\ref{tab:vo-sensori} riassume le classi di sensore considerate; per il rilevamento radar dei droni con reti di sensori si veda la rassegna \cite{coluccia2020}. Le schede commerciali non sono state aperte, salvo dove è detto: i valori sono da leggere come ordini di grandezza dichiarati dai produttori, da verificare prima di ogni uso.

\begin{table}[H]\centering\footnotesize
\caption{Classi di sensore per il volo veloce. (i): ipotesi; (n): fonte non aperta.}\label{tab:vo-sensori}
\begin{tabularx}{\linewidth}{@{}>{\RaggedRight\arraybackslash}p{4.0cm}>{\RaggedRight\arraybackslash}p{1.8cm}>{\RaggedRight\arraybackslash}p{3.8cm}L@{}}
\toprule
Sensore & Rivisita & Portata & Che cosa non vede\\
\midrule
radar a scansione elettronica in banda Ku \citep{echodyne-tradegov-2024,echodyne-echoshield-news,dronedj2026-echodyne} & $0{,}1\unita{s}$ in traccia (stampa) & portate per bersaglio non verificate: $2{,}5$, $1{,}4$ e $1{,}0\unita{km}$ per velivolo leggero, drone medio e drone piccolo sono forse di un altro modello dello stesso produttore; le distanze che ne dipendono sono ipotesi & fuori dal pannello ($120^\circ\times80^\circ$, produttore); i primi $0{,}3$--$0{,}5\unita{s}$ di una traccia nuova (i)\\
radar a schiera attiva multifunzione (i: nessun esemplare con scheda) & $0{,}025\unita{s}$ sull'oggetto & $5\unita{km}$ su un velivolo leggero (i) & i settori in ricerca; con molti frammenti la rivisita di ciascuno cala\\
radar di avvicinamento in banda S \citep{stone1989,wikipedia-asr11} & $4{,}8\unita{s}$ (requisito della classe; calcolato da $12{,}5$ giri al minuto) & $111\unita{km}$ ($60$ miglia nautiche) & $1{,}2\unita{km}$ di percorso fra due giri a $250\unita{m/s}$; bassa quota; droni piccoli\\
ADS-B \cite{cfr-91-227} & almeno un messaggio al secondo; latenza $\le2{,}0\unita{s}$ & portata radio & chi non trasmette; i detriti; l'avionica danneggiata\\
identificazione remota \cite{cfr-89-310} & $1\unita{s}$ & portata radio & droni non conformi; detriti\\
torretta nel visibile e nell'infrarosso medio (n: scheda non aperta) & $30\unita{Hz}$ (n) & oltre $25\unita{km}$ per un obiettivo di questa classe: dato del produttore, su un bersaglio grande \citep{ophir-supir-50-1350} & fuori dal campo stretto ($1{,}7^\circ$, n); fumo, nebbia, fiammata\\
LiDAR a campo stretto \cite{livox-tele15} & $0{,}1\unita{s}$ & $500\unita{m}$ al $50\,\%$ di riflettività & oltre $500\unita{m}$; fumo\\
\bottomrule
\end{tabularx}
\end{table}

In un evento impulsivo i sensori cooperativi perdono proprio l'oggetto colpito, gli ottici si accecano nella fiammata, e resta il radar, che deve seguire molti più oggetti. Con una schiera attiva il costo è la rivisita: il tempo di fascio dato a ciascun oggetto cala con il numero dei frammenti.

\paragraph{Bilancio della latenza.} La tabella~\ref{tab:vo-latenza} dà il ritardo dalla cattura all'avviso, per stadio, nel caso peggiore. L'unica misura è quella del prototipo in Python su CPU; tutto il resto è stimato da ipotesi \prova{C}.

\begin{table}[H]\centering\footnotesize
\caption{Latenza per stadio in millisecondi, caso peggiore, traccia a $10\unita{Hz}$ (fra parentesi a $40\unita{Hz}$, ipotesi). Nessuna cifra della colonna dedicata è misurata.}\label{tab:vo-latenza}
\begin{tabularx}{\linewidth}{@{}>{\RaggedRight\arraybackslash}p{5.4cm}>{\RaggedRight\arraybackslash}p{2.4cm}>{\RaggedRight\arraybackslash}p{1.9cm}>{\RaggedRight\arraybackslash}p{2.6cm}L@{}}
\toprule
Stadio & CPU & GPU & FPGA o ASIC & Origine\\
\midrule
età della misura (attesa della rivisita) & 0--100 (0--25) & uguale & uguale & sensore\\
integrazione coerente & 5--20 & 5--20 & 5--20 & sensore (i)\\
trasferimento dei dati grezzi & 0,5--2 & 0,6--2,5 & 0,001--0,01 & i\\
elaborazione del segnale (fasci, FFT, CFAR) & 5--20 & 1--3 & 0,05--0,5 & i\\
estrazione dei punti & 0,5--2 & 0,2--1 & 0,01--0,1 & i\\
associazione, filtro, regione, regole & 0,09--55 & 0,05--5 & 0,01--0,5 & CPU: misura del prototipo; resto i\\
uscita verso una macchina & 0,1--1 & uguale & 0,01--0,1 & i\\
\midrule
per aggiornamento, verso una macchina & 11--200 (11--125) & 7--133 (7--58) & 5--121 (5--46) & C\\
conferma di una traccia nuova (3 rivisite su 5) & +300--500 (+75--125) & uguale & uguale & sensore (i)\\
primo avviso, verso una macchina & 311--700 (86--250) & 307--633 (82--183) & 305--621 (80--171) & C\\
\bottomrule
\end{tabularx}
\end{table}

La somma dei soli stadi di calcolo della colonna dedicata va da $0{,}08$ a $1{,}2\unita{ms}$ \prova{C}, da ipotesi: non è inferiore al millisecondo in ogni caso, e non è misurata. Con la traccia a $10\unita{Hz}$, l'hardware dedicato porta il primo avviso verso una macchina da $0{,}70$ a $0{,}62\unita{s}$; portare la rivisita da $10$ a $40\unita{Hz}$ lo porta da $0{,}62$ a $0{,}17\unita{s}$ \prova{C}. I circa $80\unita{ms}$ guadagnati sul calcolo vengono in gran parte dal passaggio dal Python a codice compilato, che darebbe anche una CPU generica. \emph{La leva vera è la rivisita del sensore, non il calcolo.} Fa eccezione la modalità prudente del prototipo (95° percentile di $340\unita{ms}$, misurato): lì il codice compilato o un FPGA servono a tenere il calcolo entro un massimo garantito.

\paragraph{L'analogia con gli acceleratori per modelli linguistici (ipotesi).} Gli acceleratori dedicati alla generazione di testo dichiarano guadagni grandi rispetto alle GPU generiche: circa 17\,000 contro circa 230 \emph{token} al secondo per utente su un modello da 8 miliardi di parametri, cioè circa $74$ volte, per un chip con i pesi cablati \citep{heise2026-taalas}: è un dato dell'azienda da prove interne, e a febbraio 2026 non c'erano misure indipendenti; circa $8{,}4$ volte per un processore a scala di fetta, con circa 1936 \emph{token} al secondo misurati da terzi contro i 1800 dichiarati dal produttore \citep{heise2026-taalas,cerebras2024-inference}. Il guadagno nasce da un fatto preciso: la generazione è \emph{limitata dalla memoria}, e portare i pesi sul chip elimina il trasferimento dalla memoria esterna. L'analogia si trasferisce quindi solo ai carichi con la stessa struttura:
\begin{itemize}
\item \emph{vale}, come ipotesi, per l'elaborazione del segnale radar, un flusso di dati grezzi limitato dalla banda: su un FPGA in linea con il convertitore i dati non passano dalla memoria del calcolatore. Il fattore usato è $40$ volte nel caso peggiore rispetto alla CPU (da $20$ a $0{,}5\unita{ms}$), con un intervallo plausibile di $10$--$100$ volte. Due lavori pubblici sostengono l'ordine di grandezza: un CFAR su FPGA che opera in $0{,}6\unita{\mu s}$ \citep{sim2023}, e un'elaborazione distanza-Doppler su FPGA 85 volte più rapida che su un processore ARM, per un radar ad apertura sintetica \citep{choi2021}. Restano un'analogia, non una misura del caso in esame;
\item \emph{non vale} per il filtro e l'associazione: pochi oggetti, stato piccolo, calcolo dominato dai tempi di contorno. Una misura pubblica su FPGA dà un passo del filtro di Kalman più lento in latenza, circa $725\unita{ns}$ contro circa $250\unita{ns}$ su CPU, e circa $4{,}5$ volte più rapido in portata, perché il circuito accetta un nuovo stato a ogni ciclo di clock; nel caso misurato la banda del bus PCIe e $0{,}5\unita{ms}$ di latenza iniziale annullano il guadagno \citep{summers2018}. Qui l'hardware dedicato può dare il \emph{determinismo} (il 99° percentile vicino al massimo, ipotesi), non la velocità.
\end{itemize}
Un circuito integrato dedicato a tutta la catena, analogo al chip con i pesi cablati, è un'ipotesi di progetto: non esiste e non è stato provato. Anche se esistesse, toglierebbe al più i circa $75\unita{ms}$ di calcolo; la rivisita e la conferma resterebbero dove sono.

\paragraph{Quando l'avviso è tempestivo.} La distanza minima è $d=v\,(L+t_{\mathrm{risp}}+t_{\mathrm{mezzo}})$, con $L$ il primo avviso, $t_{\mathrm{risp}}$ la risposta di chi riceve e $t_{\mathrm{mezzo}}$ il ritardo del mezzo. Per un umano senza ausilio si usa il totale di $12{,}5\unita{s}$ di una tabella della FAA sul «vedere ed evitare» (AC 90-48E, Tabella 1 \citep{faa-ac-90-48e}), che \emph{comprende già} i $2{,}0\unita{s}$ di ritardo dell'aereo; per un umano avvisato con la direzione $5{,}4\unita{s}$, dalla stessa tabella, più $2{,}0\unita{s}$ di ritardo dell'aereo ($0{,}5\unita{s}$ per un drone, ipotesi). Ne segue \prova{C}:
\begin{itemize}
\item aereo a $250\unita{m/s}$: umano senza ausilio $250\cdot(0{,}75+12{,}5)=3312\unita{m}$ (con $0{,}75\unita{s}$, il primo avviso su CPU della stima dei sensori; con i $0{,}70\unita{s}$ della tabella~\ref{tab:vo-latenza} sarebbero $3300\unita{m}$); umano avvisato con la direzione $250\cdot(0{,}62+5{,}4+2{,}0)=2005\unita{m}$; verso una macchina, con elaborazione dedicata e rivisita a $40\unita{Hz}$, $250\cdot(0{,}17+2{,}0)=543\unita{m}$;
\item drone a $90\unita{m/s}$: umano senza ausilio $90\cdot(0{,}75+12{,}5-2{,}0+0{,}5)=1058\unita{m}$.
\end{itemize}
A $250\unita{m/s}$ la differenza fra un umano senza ausilio e una macchina vale circa $2{,}8\unita{km}$; l'hardware dedicato, nella stessa colonna, ne toglie qualche decina di metri. Le distanze non contano la geometria della manovra: sono un limite inferiore del bisogno. Con un radar a scansione elettronica che traccia un velivolo leggero a $2{,}5\unita{km}$ (portata non verificata, ipotesi), l'avviso sarebbe tempestivo per una macchina e al limite per un umano avvisato con la direzione; per un umano senza ausilio no.

\subsection{Il sito con droni e aerei in rullaggio, e gli altri usi}\label{sec:vo-altri}

Il sito sorvegliato del capitolo~\ref{sec:lp} resta il caso in cui l'avviso è tempestivo in simulazione \prova{S}: il drone a 15--$21\unita{m/s}$ dietro un edificio è annunciato in tempo in $200/200$ passaggi $[0{,}981;\,1]$ per ciascuna delle due serie, e in $100/100$ $[0{,}963;\,1]$ con semi nuovi; il drone non è dichiarato «libero», con la mappa dichiarata, in $0/171$ $[0;\,0{,}022]$; il drone che rallenta o si ferma resta nell'avviso in $188/200$ $[0{,}898;\,0{,}965]$ (incoraggiante); l'aereo in rullaggio è avvisato in tempo con tre sensori fusi in $200/200$ $[0{,}981;\,1]$, e con il solo radar in $50/200$ $[0{,}195;\,0{,}314]$ (limite); l'avviso si spegne solo guardando tutta la zona, con $0/600$ $[0;\,0{,}006]$ spegnimenti senza aver guardato. Per gli aerei in volo l'avviso \emph{al sito} non è tempestivo: servono $3410\unita{m}$ \prova{C}, e il sistema lo scrive. Sui 43 aerei reali la zona dichiarata li conteneva in $0{,}95$ / $0{,}93$ / $0{,}93$ dei casi a 2 / 10 / $30\unita{s}$ (bootstrap per mezzo $[0{,}92;\,0{,}98]$ / $[0{,}90;\,0{,}96]$ / $[0{,}88;\,0{,}99]$; Wilson con $n=43$ mezzi, estremo inferiore $\ge0{,}83$; tabella~\ref{tab:lp-traiettorie}) \prova{P}, con l'inviluppo di $90\unita{m/s}$, non rimisurato a $210$: è un esperimento diverso dalla taratura del \S\ref{sec:vo-taratura} (zona prevista con quantile conforme del \S\ref{sec:lp-traiettorie}, non la regione al $90\,\%$ della formula chiusa). A corto raggio, con la camera di profondità ferma, valgono i risultati dei capitoli~\ref{sec:map}--\ref{sec:val}: la proposizione di sicurezza del «libero», dimostrata sotto le sue condizioni (posa entro $0{,}5^\circ$, ostacolo che rimanda il segnale) \prova{D}; una persona nascosta senza zona sullo schermo in $0$ passaggi su $23\,402$ $[0;\,0{,}0002]$, con sensore ideale \prova{S}; su profondità simulata la persona nascosta fuori dall'avviso in $0$ fotogrammi su $4324$ $[0;\,0{,}001]$ \prova{S}. Gli altri usi (un veicolo senza equipaggio, una persona che usa il sistema come ausilio) condividono lo stesso nucleo, con un profilo deterministico di chi porta il sensore. Per una persona non vedente il sistema è un aiuto d'avviso, non sostituisce il bastone né il cane guida, e non ha cifre provate con persone.

\subsection{Condizioni di validità}\label{sec:vo-validita}

\begin{riquadro}{Condizioni di validità degli enunciati e delle cifre in volo}
\begin{enumerate}[label=\arabic*.,itemsep=1pt]
\item \textbf{Simulazione dalle schede:} probabilità di rilevamento, sezioni radar dei frammenti ($2\cdot10^{-4}\unita{m^2}$), falsi ritorni, latenze e la classe del radar a scansione elettronica sono ipotesi o valori di schede non aperte.
\item \textbf{Una geometria nominale per scenario,} senza disturbo strutturato, uccelli, riflessioni multiple; ADS-B spento.
\item \textbf{Inviluppo dichiarato:} aerei fino a $210\unita{m/s}$, droni fino a $45\unita{m/s}$ (ipotesi); oltre, l'insieme di presenza non contiene l'oggetto.
\item \textbf{Frammenti a bassa energia} ($\Delta v\le100$--$120\unita{m/s}$), punto materiale con resistenza isotropa; l'avaria non contenuta non è coperta.
\item \textbf{Confidenza} dimostrata nel modello a velocità quasi costante gaussiano, approssimata per la miscela, verificata su traiettorie pubbliche con un sensore simulato e con i limiti a 1--$2\unita{s}$ e in manovra; fuori dal modello vale solo l'insieme di presenza.
\item \textbf{Calcolo dedicato:} tutte le latenze dell'FPGA o dell'ASIC sono stime per analogia; nessuna misura su hardware.
\item \textbf{Sensore in moto (scorta):} navigazione inerziale, portatore, bersaglio, frammenti e accecamento del radar sono ipotesi; il tracciatore usa la velocità vera del portatore in due punti; vento nullo, portatore in rotta retta, controllo senza disturbo a grappoli (\S\ref{sec:vo-scorta}).
\item \textbf{Etichette a mano;} il classificatore è supposto esistente e non è realizzato.
\end{enumerate}
\end{riquadro}

\section{Limiti attuali e aree di miglioramento: da proposta a sistema attuato}\label{sec:ps}

Questo capitolo dice, per ogni limite trovato, che cosa manca per passare dalla proposta a un sistema attuato: la miglioria (che cosa e come) e la prova che la decide, con la soglia fissata prima della prova e il risultato che la smentirebbe. Nessuna delle migliorie è ancora provata su un sensore fisico o su hardware dedicato.

\paragraph{Che cosa si può dire dei risultati.} Il prototipo è stato costruito e provato in un giorno. Le prove svolte non coprono tutti i casi, e non potrebbero: per lo più una geometria per scenario, sensori simulati dalle schede, etichette a mano, nessun sensore fisico. Dentro questi limiti i risultati sono incoraggianti, e ciascuno ha una base matematica e un intervallo: nel volo veloce il principio di sicurezza (nessun «libero» sull'oggetto o sui frammenti, con la mappa dichiarata) regge in tutte le celle simulate, $0/200$ $[0;\,0{,}019]$ nel controllo e $0/60$ $[0;\,0{,}060]$ nelle altre; l'insieme di presenza ha contenuto il vero dopo l'urto e l'avaria; a bordo di un drone in moto la regione del drone seguito non si è persa in $60/60$ passaggi $[0{,}940;\,1]$ per geometria con il radar a $360^\circ$, anche durante una collisione vicina; dopo la taratura, su traiettorie di volo pubbliche con un sensore simulato, la regione al $90\,\%$ ha contenuto l'aereo da 85 a 89 volte su 100 fra $1$ e $30\unita{s}$, con i casi concentrati su una rotta e la virata come limite ($0{,}73$ $[0{,}67;\,0{,}78]$). Le cifre sotto soglia sono riportate come limiti, e la tabella che segue dice, per ciascuno, quale miglioria serve e come la si decide.

\subsection{Tabella dei limiti e delle migliorie}\label{sec:ps-tabella}

Le latenze sono stime da ipotesi; nessuna è misurata. «Applicata in simulazione» vuol dire che la miglioria è stata scritta e provata nel simulatore, non su un sensore.

{\footnotesize\setlength{\tabcolsep}{4pt}
\begin{longtable}{@{}>{\RaggedRight\arraybackslash}p{3.3cm}>{\RaggedRight\arraybackslash}p{5.6cm}>{\RaggedRight\arraybackslash}p{5.8cm}@{}}
\caption{Limiti, migliorie e prove che le decidono.}\label{tab:ps-migliorie}\\
\toprule
Limite & Miglioria: che cosa e come & Stato, esito e prova che la decide\\
\midrule
\endfirsthead
\toprule
Limite & Miglioria: che cosa e come & Stato, esito e prova che la decide\\
\midrule
\endhead
\bottomrule
\endfoot
collisione veloce (F3) non dichiarata come evento & regola di contatto: tracce risolte e mature, convergenza $\ge10\unita{m/s}$, fusione o sparizione, tracce nuove o Doppler fuori norma & \emph{applicata in simulazione, non adottata}: B $57/60$ $[0{,}863;\,0{,}983]$, ma falsi eventi in F5 $2/60$ ed eventi prima dell'urto; A $34/60$. Prossimi passi: Doppler sulla cella unita, riagganci esclusi, probabilità di collisione raggiunta, radar vicino su una griglia di geometrie fissata prima. Soglia: F3 in B $\ge54/60$ con F5, F1 e F6 a 0, su semi nuovi e su almeno 10 geometrie\\
insieme di presenza a $90\unita{m/s}$ nel nucleo & inviluppo a $210\unita{m/s}$ anche nell'insieme di presenza & \emph{applicata in simulazione}: F4 in A da $13/60$ a $60/60$; F6 in A da $126/200$ a $186/200$\\
simulatore dei frammenti con l'integratore di Eulero, instabile & integratore semi-implicito & \emph{applicata in simulazione}: collisione fra multirotori, volume da $48/60$ a $60/60$ $[0{,}940;\,1]$; evento $48/60$ e 3 eventi precoci, ora limite\\
confidenza ottimista a 1--$2\unita{s}$ & $\sigma_{\mathrm{rif}}^2$ ($20\unita{m}$ per gli aerei), regime di accelerazione, pesi fissati, ricalibrazione isotonica, minimo progressivo & \emph{applicata su dati pubblici}: da $0{,}71$ a $0{,}89$ $[0{,}87;\,0{,}91]$, scarto medio $0{,}016$. Peggiora la virata (da $0{,}91$ a $0{,}73$), l'accelerazione resta a $0{,}76$. Prossimo: ricalibrazione condizionata a una manovra rilevata in modo causale (soglia: virata $\ge0{,}85$; smentita sotto $0{,}80$), cinque divisioni, divisione per rotta\\
confidenza del nucleo mal tarata con il radar a $1\unita{Hz}$ & minimo progressivo e covarianza gonfiata sui registri di F1 ed F6 & non fatto; soglia: con $C\ge0{,}9$ copertura osservata $\ge0{,}80$\\
avaria non contenuta: $0/10$ $[0;\,0{,}278]$ & profilo dei frammenti con $\Delta v\ge300\unita{m/s}$, impronta $\ge10\unita{km}$, Doppler sui frammenti & non fatto; soglia: $60/60$ passaggi con tutti i frammenti dentro $V$\\
innesco della fiammata ottimista (anche in F7) & saturazione rilevata dall'istogramma della camera & non fatto; soglia: ritardo $\le1$ ciclo in almeno $0{,}9$ dei casi, su registrazioni vere\\
identità dopo la cecità (scambio $0{,}47$--$0{,}63$) & Doppler, sezione radar e firma termica per riassociare & non fatto; soglia: scambio $\le0{,}2$\\
continuità con il radar rotante a $1\unita{Hz}$ & radar a scansione elettronica a $10\unita{Hz}$ (tracce a $10\unita{Hz}$ da stampa, campo di $120^\circ\times80^\circ$ dal produttore; portate per bersaglio non verificate) & simulato in B: $60/60$; manca la misura sul sensore\\
F7: radar solo frontale ($\pm60^\circ$) & radar a $360^\circ$ (quattro pannelli) o camera su brandeggio puntata dalla traccia & \emph{applicata in simulazione}: in B identità in virata $60/60$, evento di lato e dietro $60/60$ (contro $43/60$, $1/60$, $0/60$)\\
F7: radar non accecato dall'esplosione & ricevitore saturo per $0{,}2$--$2\unita{s}$ e occultamento (variante F7r, modello pessimista) & \emph{applicata in simulazione}: «libero» $0/60$ e continuità $59$--$60/60$ reggono; l'avviso prima dei frammenti crolla (cella totale $0/59$); la riacquisizione con il radar frontale scende a $41$--$43/60$. Correzione c1 (avviso immediato «evento vicino non localizzato» alla saturazione congiunta di radar e camere) \emph{applicata in simulazione}: avviso $59/59$ $[0{,}939;\,1]$ nella cella totale; falsi $11/200$ $[0{,}031;\,0{,}096]$ con saturazione del radar d'altra causa; riacquisizione davanti $49/60$, sotto la soglia di $51/60$; non adottata. Prossimo: localizzare l'avviso (direzione del lampo, Doppler); riacquisizione presunta dentro il volume con coerenza Doppler; modello del controllo del guadagno dalla scheda. Soglie: falsi $\le0{,}05$ con avviso $\ge0{,}9$; riacquisizione davanti, radar frontale, $\ge51/60$\\
F7: volume prudente, avviso non discriminante ($60/60$ senza frammenti vicini) & velocità per frammento dal Doppler; volume per grappoli tracciati & non fatto; soglia: avviso falso $\le10\,\%$ con contenimento al $100\,\%$\\
F7: margine dell'avviso sotto $T_r$ con l'urto a 10--$30\unita{m}$ (7 su 20) & previsione dell'urto dalle tracce prima del contatto (avviso di prossimità, come in F3) & non fatto; soglia: margine $\ge0{,}25\unita{s}$ in $20/20$ a 10--$30\unita{m}$\\
F7: velocità vera del portatore nel tracciatore & velocità dalla navigazione inerziale; covarianza della navigazione correlata nello stato & non fatto; stesse soglie della tabella~\ref{tab:vo-scorta}\\
F7: controllo senza disturbo a grappoli & controllo con stormi, pioggia e contromisure simulati & non fatto; soglia: falsi eventi $0/200$\\
F7: tracciatore nuovo, non il nucleo & integrare F7 nel nucleo con il sensore in moto & non fatto\\
una geometria per scenario & almeno 10 geometrie per scenario & in parte: F3 su 5 varianti (16--19 su 20 in B); F7 su 3 varianti\\
tempestività verso l'umano & avviso con la direzione: $5{,}4\unita{s}$ invece di $12{,}5\unita{s}$ (AC 90-48E, Tabella 1 \cite{faa-ac-90-48e}); $2005\unita{m}$ invece di $3312\unita{m}$ a $250\unita{m/s}$ & calcolato; manca la prova con operatori\\
latenza del calcolo & FPGA o ASIC per segnale e CFAR (CFAR su FPGA in $0{,}6\unita{\mu s}$ \cite{sim2023}; distanza-Doppler 85 volte più rapida che su ARM \cite{choi2021}); calcolo dedicato totale $0{,}08$--$1{,}2\unita{ms}$ (ipotesi). Analogia con i chip per modelli linguistici (circa 17\,000 contro circa 230 \emph{token} al secondo, dato aziendale da prove interne \cite{heise2026-taalas}), valida solo per i carichi limitati dalla memoria. Passo di Kalman su FPGA: più lento in latenza, circa $4{,}5$ volte in portata \cite{summers2018} & misura su FPGA reale; soglia: $\le1{,}2\unita{ms}$ per aggiornamento\\
primo avviso verso una macchina & la leva è la rivisita: a $10\unita{Hz}$ $0{,}62\unita{s}$, a $40\unita{Hz}$ $0{,}17\unita{s}$ (calcolato da ipotesi) & misura dalla cattura all'avviso\\
mappa, classificatore, sensore fisico & LiDAR per la mappa, classificatore per classe, prove sul campo & tutte le soglie precedenti\\
\end{longtable}}

\subsection{Piano a tappe}\label{sec:ps-piano}

Ogni tappa riusa le soglie della tabella~\ref{tab:ps-migliorie} e aggiunge una sola fonte di realtà alla volta, in modo che un fallimento si possa attribuire.
\begin{enumerate}[label=\arabic*.]
\item \textbf{Simulazione estesa.} Geometrie variate (almeno 10 per scenario), disturbo a grappoli e uccelli simulati, i passi successivi della regola di contatto e della taratura della confidenza, i frammenti di un'avaria non contenuta, il portatore in manovra e il vento in F7, il tracciatore di F7 nel nucleo. Uscita: le stesse classi della tabella~\ref{tab:vo-risultati}, con le soglie fissate prima. Non richiede hardware.
\item \textbf{Banco con registrazioni.} Il nucleo, compilato, gira su registrazioni vere di radar e camere (archivi pubblici o di terzi, con licenza), senza modificarle. Si misurano mediana, 99° percentile e massimo della latenza di calcolo, l'innesco della fiammata dal segnale della camera, la confidenza contro una verità di bordo più pulita dell'ADS-B, i falsi eventi con il disturbo reale.
\item \textbf{Sensore su banco.} Un radar a scansione elettronica e una camera, fermi, su oggetti noti (droni in volo controllato in un'area autorizzata): rivisita, portata sull'oggetto piccolo, sezione radar dei frammenti simulati con oggetti di prova, latenza dalla cattura all'avviso. Qui si misura per la prima volta la colonna «dedicato» della tabella~\ref{tab:vo-latenza}, su FPGA.
\item \textbf{Campo.} Prove in un sito autorizzato con traffico vero, con operatori: tempestività verso l'umano, falsi avvisi all'ora, mappa verificata con il LiDAR, classificatore per classe. Solo dopo questa tappa una cifra del sistema si potrebbe dire misurata su un sensore.
\end{enumerate}

\subsection{Che cosa resta fuori}\label{sec:ps-fuori}

Il sistema resta un sistema di osservazione e di avviso: non decide manovre, non comanda mezzi, non ha funzioni d'ingaggio o di contrasto. Non sostituisce le procedure di separazione del traffico aereo né i sistemi anticollisione di bordo. Le prove descritte in questo documento non danno alcuna base per un uso operativo senza le tappe 2--4 del piano. Nello scenario F7 il sistema mantiene la traccia di un drone e avvisa il portatore: non lo guida e non agisce sul drone seguito.

\section{Conclusioni}\label{sec:con}

Il lavoro risponde a una domanda: che cosa può dichiarare un sistema di percezione nell'intervallo in cui non osserva un oggetto, perché un ostacolo lo nasconde o perché un evento improvviso acceca i sensori. La risposta adottata è un'uscita che tiene distinti l'osservato, il previsto con la sua regione e il suo livello, e il non noto, e che non dichiara mai «libero» ciò che non è stato osservato. I risultati hanno tre livelli di prova diversi, che qui si tengono separati.

\paragraph{Che cosa è dimostrato, nel modello dichiarato \prova{D}.}
\begin{itemize}
\item Nessuna cella in ombra, fuori campo, oltre portata o letta per meno di metà dell'impronta è dichiarata «libero» (per costruzione); e, se la profondità valida è corretta, la posa è esatta e il suolo è piano, i raggi dei pixel che leggono il suolo di una cella «libero» sono privi di superfici fino al suolo, a meno del margine di lettura, all'istante d'acquisizione (Proposizione~\ref{prop:map-sicurezza}). L'enunciato non copre la cella intera né la colonna sopra di essa: la sicurezza della cella intera e la tenuta con un errore d'inclinazione entro $0{,}5^\circ$ sono verificate in simulazione su 300 condizioni dichiarate, non dimostrate \prova{S} (\S\ref{sec:map-banco}). Fuori da queste ipotesi l'enunciato cade: i controesempi sono il vetro, l'ostacolo che non rimanda segnale e l'errore d'inclinazione di $1^\circ$ (\S\ref{sec:map-controesempi}).
\item Un soggetto nascosto, con velocità non superiore a $\vmax$, resta nell'insieme di possibile presenza $\Pres(t)$, che si riduce solo per osservazione e mai per il tempo trascorso (Proposizioni~\ref{prop:occ-contenimento} e~\ref{prop:occ-monotonia}): a posa esatta, o con il libero eroso del margine di posa, su suolo piano e con la lettura per punto esatta, senza margine di lettura; con il margine e con la mappa a celle il contenimento poggia sulla verifica in simulazione \prova{S}.
\item La fascia di guardia $w$ contiene il centro del soggetto se vale la condizione~\eqref{eq:occ-fascia}, l'ostacolo è fermo, il soggetto resta nel campo e in portata ed è riportato come osservato entro $t_a$ da quando il suo centro è fuori dall'ombra; con $w=0{,}40\unita{m}$ la velocità coperta, per calcolo, è $1{,}13\unita{m/s}$, e $0{,}56\unita{m/s}$ se il bordo dell'ombra sbaglia di $0{,}20\unita{m}$; per l'intero corpo, di raggio $r_c$, la condizione chiede $r_c$ in più (Proposizione~\ref{prop:occ-fascia}) \prova{C}.
\item La regione di previsione è contenuta nell'area d'avviso e, con l'ostacolo fermo e l'ombra stimata costante, non si svuota finché la traccia è nell'ombra; per ogni ingresso riconosciuto, l'area d'avviso con la previsione contiene quella senza (Proposizione~\ref{prop:occ-annidamento}). Un'etichetta di classe non produce un «libero» e, per ogni ingresso riconosciuto, non riduce l'avviso (Proposizione~\ref{prop:occ-etichetta}).
\item La massa non è identificabile da misure di posizione, e l'accelerazione non è stimabile con precisione utile alla frequenza e al rumore assunti (Proposizione~\ref{prop:sti-massa} per la massa \prova{D}; conto del \S\ref{sec:sti-massa} per l'accelerazione \prova{C}).
\item La geometria d'osservazione è in forma chiusa: visibilità del suolo, ombra di un ostacolo, durata dell'occlusione, rumore di profondità ed effetto di un errore di posa (\S\ref{sec:geo}).
\end{itemize}

\paragraph{Che cosa è verificato in simulazione \prova{S}.}
\begin{itemize}
\item \emph{Con modello ideale del sensore.} La regione è emessa a ogni ciclo dell'occlusione dopo l'ingresso: 0 passaggi con un'interruzione su $23\,402$ passaggi nascosti, da 6000 soggetti distinti; la premessa, cioè l'ingresso nell'ombra riconosciuto, manca in 928 passaggi ($3{,}97\,\%$; 76 con $P_D=0{,}9$ e 852 con $P_D=0{,}7$), nei quali non c'è regione; con l'avviso acceso anche per le tracce chiuse entro $0{,}75\unita{m}$ da un'ombra, regola attivata solo dopo la prova d'ingresso sul suolo vuoto (\S\ref{sec:lim-salute}), restano senza avviso 24 passaggi ($0{,}10\,\%$; 2 e 22). Con un ostacolo largo $2\unita{m}$, il bordo dell'ombra noto e il soggetto sul proprio percorso, la regione contiene il soggetto con estremo inferiore dell'intervallo sopra $0{,}90$ fino a $2\unita{s}$ e sopra $0{,}85$ a $3\unita{s}$, in 30 combinazioni su 30. Con un pannello largo $1\unita{m}$ la soglia di $0{,}90$ non è raggiunta a $1\unita{s}$ ($942/1045$, $\ell^-=0{,}882$) e oltre non si dichiara alcuna cifra (\S\ref{sec:val}).
\item La regola del «libero» della cella non dichiara libera alcuna cella sopra un ostacolo in 300 condizioni dichiarate, con errore d'inclinazione entro $0{,}5^\circ$; lo zero vale per quell'insieme di condizioni e non è una frequenza (\S\ref{sec:map}).
\item \emph{Su profondità simulata dalle schede tecniche.} Con i gruppi a celle forti e l'avviso che resta dopo il riaggancio, un soggetto nascosto entro $3\unita{m}$ non resta senza avviso (0 fotogrammi su 4324 a 424 pixel di larghezza e 0 su 4325 a 848, e 0 su 3608 per risoluzione in una replica); con la configurazione di riferimento, sugli stessi passaggi, resta senza avviso in 180 e in 271 fotogrammi. L'identità alla ricomparsa, che con la configurazione di riferimento tiene in 0 passaggi su 90, con i gruppi a celle forti e le condizioni aggiuntive del riaggancio tiene in 176 passaggi su 180 a 424 pixel di larghezza e in 174 su 180 a 848, in un dominio ristretto: un adulto solo, a $0{,}5\unita{m}$ dietro il pannello, entro $2{,}6\unita{m}$, rumore nominale (\S\ref{sec:sti}, \S\ref{sec:val}). Queste regole non fanno parte della configurazione di riferimento: si attivano all'avvio solo se le prove d'ingresso sono superate (\S\ref{sec:lim-salute}). Una cella «libero» dentro l'ombra vera di un ostacolo: 0 su $2\,907\,032$ celle\,$\cdot$\,fotogramma in una prima serie di 135 sequenze, 5 su $3\,992\,186$ in una seconda di 381 (una al livello di rumore nominale); entrambe con la variante che conserva il «libero» per tre fotogrammi, non con la regola della cella adottata. Non è risolto se le 5 dipendano dalla regola o dal bordo della maschera di riferimento.
\item \emph{Su traiettorie pubbliche di pedoni} \prova{P}. Il modello di moto a velocità quasi costante, con $q=0{,}2\unita{m^2/s^3}$, copre il soggetto in $0{,}977$, $1{,}000$ e $1{,}000$ dei casi a 1, 2 e $3\unita{s}$ (219 tratti di 9 persone; numerosità efficace fra 9 e 219; posizioni note al millimetro fino all'ultimo istante, senza profondità): riguarda il modello di moto, non la catena.
\item Il costo di calcolo è in media $6{,}35\unita{ms}$ per ciclo su profondità simulata, su un filo di un processore da portatile, e $81\unita{ms}$ su una registrazione pubblica affollata (\S\ref{sec:cos}).
\end{itemize}

\paragraph{Dove le simulazioni mostrano che il sistema cede.} Chi svolta o esce dal lato lontano dell'ombra è nella regione, a 2 e $3\unita{s}$, in una quota di casi fra $0{,}57$ e $0{,}86$ nelle prove con il corpo esteso e fra $0{,}13$ e $0{,}70$ in quelle con il soggetto puntiforme; chi parte da fermo accanto all'ostacolo, in $0{,}11$--$0{,}26$ dei casi fra 1 e $3\unita{s}$ nelle prove di base; con il bordo dell'ombra stimato e non riconosciuto la quota a $1\unita{s}$ scende a $0{,}30$--$0{,}46$ (estremo basso: chi prosegue). In questi casi non si dichiara un livello e resta la sola area d'avviso. Un secondo soggetto mai visto che esce dalla stessa ombra riceve l'identità del primo: 2400 casi su 2400 nelle prove di base con modello ideale del sensore, fra $0{,}91$ e $0{,}995$ nelle prove d'insieme con modello ideale, 19 su 20 (76--99\,\%) su profondità simulata dalle schede tecniche; nessuna delle regole provate lo ripara, e per questo l'identità è presunta e l'avviso non ne dipende. Su profondità simulata dalle schede tecniche un pannello fermo è dichiarato «in movimento» in 79 passaggi su 90 con i gruppi per componenti connesse; con i gruppi a celle forti in 0 su 140 nel dominio dell'identità al ritorno e in 7 su 10 fuori da esso. Dietro un ostacolo fermo l'avviso non si spegne.

\paragraph{Che cosa resta non verificato sperimentalmente \prova{N}.}
\begin{itemize}
\item \emph{Ogni grandezza su un sensore fisico.} Nessuna prova è stata eseguita su una camera: rumore, pixel senza dato, errore del bordo dell'ombra, copertura della regione, identità alla ricomparsa, latenza e costo sulla scheda di bordo non sono misurati. La latenza non è misurata: la somma delle voci note dà $155\unita{ms}$ in media, o circa $131\unita{ms}$ con un altro dato del costruttore per il tratto del sensore; le due fonti non concordano e nessuna voce è misurata sulla catena.
\item \emph{La continuità su profondità simulata.} La regola della regione persistente è stata eseguita solo con modello ideale del sensore (\S\ref{sec:lim-limiti}).
\item \emph{Il termine dell'avviso per osservazione.} Che l'avviso non termini senza osservazione vale per costruzione; che termini quando il suolo è rivisto dopo la rimozione dell'ostacolo non è sostenuto: nella realizzazione provata resta acceso dopo $2\unita{s}$ in 10 prove su 10 (72--100\,\%), per una causa non determinata.
\item \emph{Il classificatore.} È un componente richiesto dal sistema, supposto esistente e lasciato ai lavori futuri; in questo lavoro non è realizzato né valutato: è simulato con etichette assegnate a mano. Nessuna classe diversa da «persona» è stata simulata nella catena completa, e la regola che lega la classe alla fascia di guardia è stata provata sulla sola fascia (0 riduzioni dell'avviso su $20\,000$ sequenze di etichette, due volte). Un'etichetta non produce un «libero» né riduce un avviso acceso; può cambiare l'esito di un riaggancio, e quindi gli ingressi successivi (\S\ref{sec:occ-classe}).
\item \emph{L'avviso tempestivo per classe e per portatore.} Con $3\unita{m}$ di portata, sensore fermo e un tempo di reazione ipotizzato di $1{,}0\unita{s}$, l'avviso è tempestivo per calcolo fino a $1{,}44\unita{m/s}$: 3 classi su 15 alla velocità tipica, nessuna alla massima. A corto raggio il sensore in moto è trattato solo come analisi dei requisiti (\S\ref{sec:lim-moto}); il sensore in moto a lunga portata è simulato nel \S\ref{sec:vo-scorta}.
\item \emph{L'uso da parte di persone.} Lo schermo e i segnali sonori non sono stati provati con persone, vedenti o no.
\end{itemize}

Le misure che deciderebbero ciascuno di questi punti, con la soglia che li smentirebbe, sono elencate in \S\ref{sec:lim-misure}.

\paragraph{Nel volo ad alta velocità, il caso centrale.} Con aerei e droni da $20$ a $210\unita{m/s}$, collisioni, detriti ed eventi impulsivi, tutto in simulazione dalle schede o su traiettorie pubbliche con un sensore simulato:
\begin{itemize}
\item \emph{dimostrato nel modello} \prova{D}: la varianza di posizione dopo l'ultima misura, la confidenza in forma chiusa su disco e sfera, il raggio $r_\alpha(\tau)$ che tiene il livello, l'esistenza e l'unicità del tempo di soglia quando $\sigma_0<r/k_\alpha$ (e $t_C=0$ solo con $\sigma_{0v}\ge0$), la confidenza della miscela di regimi con il minimo progressivo non crescente, il contenimento fuori modello condizionato a $|p(0)-\hat p_0|\le r_0$, il cono e il cammino massimo dei frammenti con la resistenza;
\item \emph{solido in simulazione} \prova{S}: nessun «libero» sull'oggetto o su un frammento nei sei scenari ($0/200$ $[0;\,0{,}019]$ e $0/60$ $[0;\,0{,}060]$); regione mantenuta dopo la collisione veloce ($60/60$ $[0{,}940;\,1]$); continuità con il radar a $10\unita{Hz}$; a bordo di un drone in moto, regione del drone seguito mantenuta in $60/60$ passaggi per geometria con il radar a $360^\circ$, anche con il radar accecato fino a $2\unita{s}$ ($59$--$60/60$);
\item \emph{incoraggiante}: l'avaria con fiammata ($60/60$, innesco ottimista), l'avviso di prossimità ($60/60$), e, su traiettorie pubbliche, la copertura della regione al $90\,\%$ da $0{,}85$ a $0{,}89$ fra $1$ e $30\unita{s}$ dopo la taratura \prova{P}, con i casi concentrati su una rotta;
\item \emph{limiti dichiarati}: la collisione veloce non è ancora dichiarata come evento (la regola in prova, $57/60$, dà falsi eventi e non è adottata); la confidenza in virata ($0{,}73$); la continuità con il radar a $1\unita{Hz}$; l'avviso dei detriti con il radar accecato (fino a $0/59$ $[0;\,0{,}061]$; con la correzione c1 $59/59$ $[0{,}939;\,1]$, ma non localizzato e con $11/200$ falsi avvisi se il radar satura per un'altra causa); l'avaria non contenuta, non provata; le latenze del calcolo dedicato, solo stimate.
\end{itemize}
Sono risultati incoraggianti, ottenuti con un prototipo costruito e provato in un giorno: ciascuno ha una base matematica e un intervallo, e i limiti sono scritti con la miglioria che li affronterebbe e la prova che la deciderebbe (\S\ref{sec:ps}). Nessuno è misurato su un sensore fisico o su hardware dedicato.

\paragraph{Al sito con droni e aerei in rullaggio.} Portati a una scala di chilometri, con sensori simulati dalle schede tecniche, i tre principi restano gli stessi, ma per tenerli veri servono regole nuove: il «libero» per classe con la probabilità di rilevamento cumulata, con le passate trattate come correlate, e prudente per qualunque correlazione \prova{D}; l'insieme di possibile presenza con l'inviluppo dello scenario, che contiene ogni moto entro i limiti dichiarati \prova{D}; lo spegnimento dell'avviso solo quando tutta la zona è stata guardata. In un sito aerospaziale sorvegliato l'avviso per un drone nascosto da un edificio e per un aereo in rullaggio nascosto da un hangar arriva in tempo in 200 passaggi simulati su 200 \prova{S}; la regola di spegnimento più semplice perdeva l'avviso in 564 passaggi su 2400. Su traiettorie pubbliche di 43 aerei reali la zona dichiarata conteneva l'aereo in $0{,}95$ / $0{,}93$ / $0{,}93$ dei casi a 2 / 10 / $30\unita{s}$ (bootstrap per mezzo $[0{,}92;\,0{,}98]$ / $[0{,}90;\,0{,}96]$ / $[0{,}88;\,0{,}99]$) \prova{P}, con l'inviluppo di $90\unita{m/s}$; è un esperimento diverso dalla taratura della confidenza del volo veloce (zona prevista, non regione al $90\,\%$).
Restano limiti dichiarati: per gli aerei in volo l'avviso non è tempestivo; il «libero» sul drone nel sito regge solo con la mappa degli ostacoli dichiarata completa, ed è falso con una mappa dichiarata e sbagliata; l'avviso perde per qualche ciclo il drone che rallenta o si ferma (12 passaggi su 200); dopo una collisione fra droni lenti l'evento è dichiarato entro $2\unita{s}$ in 48 passaggi su 60 con l'integratore stabile, con 3 eventi aperti prima dell'urto, e la zona dei detriti è di alcuni chilometri e non si riduce guardando; e le probabilità di rilevamento dei sensori sono ipotesi, da misurare su sensori fisici (\S\ref{sec:lp-smentiti}, \S\ref{sec:lim-lunga}).

Il sistema svolge solo osservazione, monitoraggio e avviso: non è un dispositivo di protezione.

\clearpage
\appendix
\section{Parametri del sistema}\label{sec:par}

La tabella~\ref{tab:par-parametri} raccoglie i parametri del sistema descritto nei capitoli precedenti. Nell'ultima colonna «scelto» indica un valore fissato a tavolino, «tarato» un valore ricavato da dati simulati, «stimato» un valore preso da una fonte esterna o da un conto e non misurato. Nessun parametro è tarato su un sensore fisico; i valori marcati «da misurare» vanno rifissati sulla camera usata. Le soglie di tempo valgono in secondi: il numero di fotogrammi è dato per il ciclo di $15\unita{Hz}$.

\begin{table}[H]\centering\footnotesize
\caption{Parametri del sistema (prima parte): sorgente, mappa, gruppi, filtro e ciclo di vita della traccia.}\label{tab:par-parametri}
\begin{tabularx}{\linewidth}{@{}p{1.25cm}Lp{3.5cm}Lp{1.2cm}@{}}
\toprule
Simbolo & Significato & Valore & Origine & Tipo\\
\midrule
$h$, $\theta$ & quota e inclinazione della camera & $1{,}2\unita{m}$; $15^\circ$ verso il basso & geometria d'osservazione (\S\ref{sec:geo}) & scelto\\
$\Delta$ & periodo del ciclo d'elaborazione & $1/15\unita{s}$ su un flusso a 30 fotogrammi al secondo & valore di progetto, non una prestazione misurata & scelto\\
--- & portata della mappa & $3\unita{m}$ & intervallo utile della camera; da misurare & scelto\\
--- & tolleranza sull'inclinazione & meno di $0{,}5^\circ$ & controesempio a $1^\circ$ (\S\ref{sec:lim-limiti}) & scelto\\
--- & autoverifica all'avvio & pixel validi sotto il $10\,\%$ entro 3 fotogrammi a sensore coperto & regola di progetto & scelto\\
--- & prove d'ingresso & oggetto di $10\unita{cm}$ a $2{,}5\unita{m}$ «occupato»; suolo vuoto per $30\unita{s}$ senza gruppi né tracce & regola di progetto & scelto\\
--- & cella della mappa polare & $1^\circ\times0{,}1\unita{m}$ & \S\ref{sec:map} & scelto\\
--- & quota d'ostacolo & da $0{,}07$ a $2{,}0\unita{m}$ sopra il suolo & \S\ref{sec:map} & scelto\\
--- & pixel minimi per «occupato» & 3 nel fotogramma corrente & \S\ref{sec:map} & scelto\\
--- & margine di lettura del suolo alla profondità attesa $z_0$ & $0{,}03+3\,k\,z_0^2$ (m); $k=0{,}004$ nominale, $0{,}0100$ al limite & modello del rumore del sensore; da misurare sul suolo & stimato\\
--- & condizione di «libero» & metà dei pixel dell'impronta della cella e di ognuna delle 8 vicine entro il margine & fissata dopo una prima simulazione & scelto\\
--- & memoria del «libero» & 0 fotogrammi & controesempi alle regole con memoria (\S\ref{sec:map}) & scelto\\
--- & inizio della frase sul suolo & $1{,}3\unita{m}$ & suolo visibile da $1{,}24\unita{m}$ sull'asse & scelto\\
--- & gruppi a celle forti (attivi solo dopo le prove d'ingresso) & celle con almeno $0{,}30$ del massimo della componente; $0{,}05\unita{m}$ e $0{,}004\unita{m^2}$ per le soglie attese & $0{,}30$ ricavato dopo i dati simulati; le altre fissate prima & tarato; scelto\\
$\sigma_z(d)$ & rumore di profondità & $1{,}5$ / $3{,}4$ / $6{,}0\unita{cm}$ a 2 / 3 / $4\unita{m}$ & dato del costruttore; da misurare & stimato\\
--- & soglia della regola «mobile» & $5\,\sigma\sqrt{2/m}$, $m=3$ fotogrammi per media, misure distanti al più $1\unita{s}$ & \S\ref{sec:sti} & scelto\\
$q$ & densità del rumore di processo & $0{,}2\unita{m^2/s^3}$ & prudente: $0{,}05$--$0{,}1$ su traiettorie pubbliche di pedoni & scelto\\
$\gamma$ & soglia di porta & quantile al $99\,\%$ & \S\ref{sec:sti} & scelto\\
--- & conferma della traccia & 3 misure su 5 & \S\ref{sec:sti} & scelto\\
--- & misure minime per la garanzia d'ingresso in ombra & 4 fotogrammi & con 3 la traccia manca nel $27{,}1\,\%$ dei casi, con 4 nel $5{,}2\,\%$ & tarato\\
--- & ingresso in ombra senza misura & dopo 2 cicli, con ultima misura o punto previsto entro $w$ dall'ombra & fissato dopo una prova sugli stessi semi & scelto\\
--- & durata di «non confermato» & $0{,}2\unita{s}$ & regola di progetto & scelto\\
--- & avviso senza traccia (attivo solo dopo le prove d'ingresso); distanza massima del riaggancio dall'ombra d'ingresso & $0{,}75\unita{m}$ & regola di progetto & scelto\\
--- & covarianza di misura al riaggancio & $R+(\text{larghezza}/2)^2$ & \S\ref{sec:sti} & scelto\\
\bottomrule
\end{tabularx}
\end{table}

\begin{table}[H]\centering\footnotesize
\caption{Parametri del sistema (seconda parte): occlusione, avviso, salute del sensore, uscita, profilo di chi porta il sensore.}\label{tab:par-parametri-due}
\begin{tabularx}{\linewidth}{@{}p{1.25cm}Lp{3.5cm}Lp{1.2cm}@{}}
\toprule
Simbolo & Significato & Valore & Origine & Tipo\\
\midrule
$w$ & fascia di guardia attorno all'ombra & $0{,}40\unita{m}$ & copre $1{,}13\unita{m/s}$ con età $0{,}133\unita{s}$, latenza $0{,}155\unita{s}$, tenuta $0{,}067\unita{s}$ & scelto\\
--- & fascia di ritaglio della regione sul suolo «libero» & $0{,}20\unita{m}$ & \S\ref{sec:occ} & scelto\\
$\alpha$ & livello della regione di previsione & $0{,}90$ & \S\ref{sec:occ} & scelto\\
--- & insieme di percorsi del predittore & 400, di cui 200 proseguono e 200 si fermano a un istante uniforme & \S\ref{sec:occ} & scelto\\
--- & peso di un percorso su suolo «libero» oltre la fascia di ritaglio & moltiplicato per $0{,}1$ a ogni ciclo & \S\ref{sec:occ} & scelto\\
--- & lisciature della densità & 3 & \S\ref{sec:occ} & tarato\\
--- & soglia e orizzonte di «previsto ampio» & regione oltre l'$80\,\%$ delle celle dell'ombra; $3\unita{s}$ & fissati dopo aver visto i dati simulati & scelto\\
--- & ingresso da fermo & velocità sotto $0{,}3\unita{m/s}$ all'ingresso o negli ultimi $2\unita{s}$ & \S\ref{sec:occ} & scelto\\
--- & punto previsto fuori dal non visto libero & fermo 8 cicli, poi al più $0{,}167\unita{m}$ per ciclo & \S\ref{sec:occ} & scelto\\
--- & libreria di taratura dei percorsi & almeno 200 percorsi, due persone, tre andature & \S\ref{sec:val} & scelto\\
--- & rilascio di una cella dell'avviso & «libero» per 3 fotogrammi di fila e oltre $w$ dall'ombra corrente & \S\ref{sec:occ} & scelto\\
--- & fotogramma mancante & $0{,}13\unita{s}$ senza fotogramma nuovo, o tre identici & da rifissare con $60\unita{s}$ di conteggio a vuoto e sotto carico & scelto\\
--- & pixel validi & sotto metà della taratura per $0{,}2\unita{s}$ & \S\ref{sec:lim-salute} & scelto\\
--- & unità inerziale; piano del suolo & $0{,}5^\circ$ o $5^\circ/\mathrm{s}$; $0{,}5^\circ$ & \S\ref{sec:lim-salute} & scelto\\
--- & segnale sonoro & $0{,}3\unita{s}$ per cambio di stato & \S\ref{sec:lim-uscita} & scelto\\
$\tau$ & latenza dall'evento all'uscita & $0{,}155\unita{s}$ (circa $0{,}131\unita{s}$ con un altro dato del costruttore) & somma di voci note (\S\ref{sec:lim-latenza}); non misurata & stimato\\
$t_{\mathrm{conf}}$ & tempo di conferma & $0{,}267\unita{s}$ & 4 periodi a $15\unita{Hz}$: valore usato nei conti per la conferma «3 misure su 5»; il legame fra i due numeri non è documentato & scelto\\
$C$ & margine della distanza d'avviso & $0{,}95\unita{m}$ & $0{,}85\unita{m}$ attribuiti a una norma, non riscontrati, più $0{,}10\unita{m}$ & stimato\\
$T_r$, $v_e$, $a_f$ & tempo di reazione, velocità propria, frenata del portatore & per profilo (tabella~\ref{tab:lim-portatore}) & ipotesi di progetto o valori riportati & scelto\\
--- & soglia di fiducia del classificatore; coerenza cinematica & $0{,}7$; 4 cicli su 6 & ipotesi; il classificatore non è realizzato (\S\ref{sec:sti}) & scelto\\
\bottomrule
\end{tabularx}
\end{table}

\section{Riproducibilità}\label{sec:rip}

\paragraph{Tipi di prova.} Ogni cifra del documento appartiene a uno di quattro tipi: dimostrazione in un modello dichiarato; calcolo in forma chiusa; simulazione; elaborazione di dati pubblici. Le simulazioni sono di due specie: con modello ideale del sensore (misure di posizione con rumore gaussiano e ombra geometrica, senza immagini di profondità) e su profondità simulata dalle schede tecniche di tre camere a stereo attivo (RealSense D435i e D455, Orbbec Gemini 335), a due livelli di rumore, nominale e limite. I dati pubblici sono traiettorie di pedoni al chiuso con verità di riferimento \cite{thor}, che non contengono profondità, sequenze di profondità registrate \cite{tum-rgbd,depthtrack} e una registrazione di esempio distribuita dal costruttore di una delle camere, usata solo per il confronto dei pixel senza dato (\S\ref{sec:val}). Nessuna prova usa un sensore fisico.

\paragraph{Linguaggi.} Le simulazioni con modello ideale del sensore e il predittore sono in C (compilatore gcc 16.2.1, opzione \texttt{-O2}); la catena di riferimento su profondità simulata, la regola del «libero» della cella e i calcoli in forma chiusa sono in Python con una libreria numerica. I tempi di calcolo sono misurati su un calcolatore portatile con processore AMD Ryzen AI 9 HX 370 e sistema Linux, con un solo filo d'esecuzione; non sono tempi della scheda di bordo, e le misure rifatte a macchina carica non sono stabili (\S\ref{sec:cos}).

\paragraph{Generatore di numeri casuali e semi.} Le simulazioni in C usano un generatore a incremento costante: stato $s\in\mathbb{Z}_{2^{64}}$, a ogni estrazione $s\leftarrow s+G \pmod{2^{64}}$ con $G$ dispari, uscita $\mu(s)$ con $\mu$ funzione di mescolamento biiettiva. Lo stato iniziale del seme $\sigma$ è $s_0(\sigma)=\mu(\sigma G+c)$. La forma senza mescolamento, $s_0(\sigma)=\sigma G+c$, rende il flusso del seme $\sigma+1$ uguale a quello del seme $\sigma$ avanzato di un'estrazione (1\,000 valori uguali su 1\,000; con il mescolamento 0 su 1\,000), e gli intervalli binomiali calcolati su semi contigui risultano troppo stretti di un fattore fra $1{,}3$ e $2{,}5$ (\S\ref{sec:val}). Le cifre di continuità e di copertura sono calcolate con semi inizializzati nella forma corretta; taratura e verifica usano insiemi di semi disgiunti. In nessuna prova la libreria di taratura viene da un generatore di camminate e i passaggi di verifica da un altro.

\paragraph{Numerosità tipiche.} Continuità e copertura con modello ideale del sensore: 23\,402 passaggi nascosti di 6\,000 camminatori distinti; sull'insieme di tutti i lanci, 217\,861 passaggi confermati su 226\,400 lanciati, di al più 70\,400 camminatori distinti; da 608 a 2\,359 casi per cella della tavola di copertura. Regola del «libero» della cella: 10 scene, 300 condizioni, 9\,000 fotogrammi, in due lanci a semi diversi. Profondità simulata dalle schede tecniche: 90 passaggi (3 disposizioni per 30 semi) per le prove di base, 180 passaggi per risoluzione per l'identità al ritorno. Traiettorie pubbliche di pedoni: 219 tratti di 9 persone, con numerosità efficace fra 9 e 219. Regola delle etichette: 20\,000 sequenze, due volte. Durata delle esecuzioni sul calcolatore indicato: circa 4 minuti per la simulazione d'insieme, 2 minuti per la regola del «libero», da 12 a 15 minuti per le prove d'identità.

\paragraph{Come si rifanno i conti principali.} Dal solo testo si ricalcolano:
\begin{enumerate}[label=(\alph*)]
\item l'intervallo di Wilson al $95\,\%$ \cite{wilson1927,brown2001} di una frequenza $k/N$: per $k=826$, $N=879$ si ottiene $[0{,}9220;\,0{,}9536]$; per zero eventi su $N$ il limite superiore è $z^2/(N+z^2)$ con $z^2=3{,}8415$, cioè $6{,}4$ su $10\,000$ per $N=6\,000$;
\item la velocità coperta dalla fascia di guardia, $w/(0{,}133+0{,}155+0{,}067)\unita{s}=0{,}40/0{,}355=1{,}13\unita{m/s}$, e con $0{,}20\unita{m}$ d'errore del bordo $(0{,}40-0{,}20)/0{,}355=0{,}56\unita{m/s}$;
\item la velocità coperta dell'avviso a sensore fermo, $(3-0{,}95)/(0{,}267+0{,}155+T_r)$: $1{,}44\unita{m/s}$ con $T_r=1{,}0\unita{s}$, $2{,}84$ con $0{,}3$, $1{,}07$ con $1{,}5$, $0{,}92$ con $1{,}8$ (\S\ref{sec:lim-portatore});
\item la somma di latenza $16{,}7+95+22+21{,}3=155\unita{ms}$ (\S\ref{sec:lim-latenza});
\item la soglia d'inclinazione oltre la quale un ostacolo alto $a$ scende sotto la quota $y_s=0{,}07\unita{m}$, $\delta\theta=(a-y_s)/\rho$: $0{,}67^\circ$ a $2{,}55\unita{m}$ e $0{,}57^\circ$ a $3\unita{m}$ per $a=0{,}10\unita{m}$;
\item la probabilità di un falso «mobile» per deriva di posa (\S\ref{sec:lim-moto}), controllata con un Monte Carlo indipendente di $2\cdot10^6$ valutazioni: $0{,}0999$ e $0{,}0100$ alle due soglie.
\end{enumerate}
I conti (b)--(e) richiedono solo le quattro operazioni. Le simulazioni non si rifanno dal solo testo: i capitoli~\ref{sec:map}--\ref{sec:occ} e l'appendice~\ref{sec:par} danno regole e parametri, ma non la geometria completa delle disposizioni $G_1$--$G_3$, i parametri del generatore di camminate e dei sette comportamenti, le dieci scene della regola del «libero», il modello del simulatore da scheda (pixel misti, pixel senza dato), l'elenco delle registrazioni pubbliche usate e i semi: per rifarle servono i programmi e i dati del progetto.

\clearpage
% Gli indirizzi lunghi della bibliografia possono andare a capo anche ai trattini.
\begingroup\sloppy\def\UrlBreaks{\do\/\do\-\do\_\do\.}
\bibliographystyle{unsrtnat}
\bibliography{allegato}
\endgroup
\end{document}
